% !TeX program = xelatex
% Standalone English volume. Use XeLaTeX until the references stabilize.
\documentclass[a4paper,11pt]{article}
\usepackage[T1]{fontenc}
\usepackage{lmodern}
\usepackage[english]{babel}
\usepackage[left=26mm,right=26mm,top=25mm,bottom=25mm,headheight=15pt,headsep=7mm,footskip=11mm]{geometry}
\usepackage{amsmath,amssymb,mathtools,amsthm}
\newtheorem{proposition}{Proposition}[section]
\newtheorem{lemma}[proposition]{Lemma}
\newtheorem{theorem}[proposition]{Theorem}
\newtheorem{corollary}[proposition]{Corollary}
\newtheorem{maintheorem}[proposition]{Main Theorem}
\theoremstyle{definition}
\newtheorem{definition}[proposition]{Definition}
\theoremstyle{remark}
\newtheorem{remark}[proposition]{Remark}
\makeatletter
\renewcommand*\l@section{\@dottedtocline{1}{0em}{2.4em}}
\renewcommand*\l@subsection{\@dottedtocline{2}{1.45em}{3.45em}}
\renewcommand*\l@subsubsection{\@dottedtocline{3}{3.8em}{4.55em}}
\makeatother
\usepackage{xcolor,fancyhdr,titlesec,needspace,xurl}
\usepackage{longtable,multicol}
\usepackage{hyperref}
\definecolor{ink}{HTML}{1A3542}
\definecolor{muted}{HTML}{4F606A}
\hypersetup{unicode=true,colorlinks=true,linkcolor=ink,filecolor=ink,urlcolor=ink,citecolor=ink,bookmarksnumbered=true,
 pdfdisplaydoctitle=true,pdflang={en},
 pdftitle={The generalized Riemann hypothesis for Dirichlet L-functions. Edition 930, Volume XXI},
 pdfsubject={A natural-language reconstruction of the proposed proof},
 pdfauthor={A mathematical reconstruction}}
\titleformat{\section}{\Large\bfseries\color{ink}}{\thesection}{.65em}{}
\titlespacing*{\section}{0pt}{.8em}{.65em}
\setlength{\parindent}{1.2em}
\linespread{1.045}
\setlength{\parskip}{.22em plus .05em minus .03em}
\setlength{\emergencystretch}{2em}
\setcounter{topnumber}{2}
\setcounter{tocdepth}{3}\setcounter{secnumdepth}{3}
\widowpenalty=4500
\clubpenalty=4500
\pagestyle{fancy}\fancyhf{}
\fancyhead[L]{\small\color{muted}Dirichlet GRH \textperiodcentered\ Volume I}
\fancyhead[R]{\small\color{muted}Dirichlet zeros and power intervals}
\fancyfoot[C]{\small\thepage}
\renewcommand{\headrulewidth}{.3pt}
\newcommand{\lead}[1]{\par\smallskip\noindent{\bfseries\color{ink}#1.}\enspace}
\newcommand{\sour}[2]{\unskip\penalty150\hspace{.35em}{\small\hyperlink{ref#1}{[#2]}}}
\titleformat{\subsection}{\large\bfseries}{\thesubsection}{.6em}{}
\titlespacing*{\subsection}{0pt}{1.15em plus .15em}{.4em}
\setlength{\abovedisplayskip}{9pt plus 2pt minus 3pt}
\setlength{\belowdisplayskip}{9pt plus 2pt minus 3pt}
\newcommand{\RR}{\mathbb R}\newcommand{\CC}{\mathbb C}\newcommand{\QQ}{\mathbb Q}
\newcommand{\ZZ}{\mathbb Z}\newcommand{\ord}{\operatorname{ord}}
\newcommand{\Frob}{\operatorname{Frob}}\newcommand{\Gal}{\operatorname{Gal}}
\newcommand{\len}{\operatorname{length}}\newcommand{\Res}{\operatorname*{Res}}
\newcommand{\wmean}{\operatorname{av}}\newcommand{\Ker}{\operatorname{ker}}
\newcommand{\Lh}{\mathcal L}\newcommand{\Gg}{\mathcal G}
\newcommand{\real}{\operatorname{Re}}\newcommand{\imag}{\operatorname{Im}}
\newcommand{\abs}[1]{\left|#1\right|}\newcommand{\eps}{\varepsilon}
\newcommand{\anchoritem}[2]{\par\Needspace{2.8\baselineskip}\hypertarget{ref#1}{}\noindent{\bfseries\color{ink}#2.}\enspace}
\newcommand{\srcloc}[2]{{\small\nolinkurl{#1.lean}}, line~#2}
\numberwithin{equation}{section}
\raggedbottom

\newcommand{\Sel}{\mathcal S}
\newcommand{\Cl}{\operatorname{Cl}}
\newcommand{\Aut}{\operatorname{Aut}}
\newcommand{\GL}{\operatorname{GL}}
\newcommand{\OO}{\mathcal O}
\hypersetup{bookmarksdepth=3}


\addto\captionsenglish{\renewcommand{\partname}{Volume}}


\makeatletter
\expandafter\def\csname r@cor:content125\endcsname{{92.19}{15}{}{proposition.92.19}{Dirichlet-GRH-v930-Volume-XVIII.pdf}}
\expandafter\def\csname r@cor:cost128\endcsname{{100.17}{79}{}{proposition.100.17}{Dirichlet-GRH-v930-Volume-XIX.pdf}}
\expandafter\def\csname r@cor:mom156\endcsname{{87.25}{49}{}{proposition.87.25}{Dirichlet-GRH-v930-Volume-XVII.pdf}}
\expandafter\def\csname r@cor:original200\endcsname{{83.7}{57}{}{proposition.83.7}{Dirichlet-GRH-v930-Volume-XVI.pdf}}
\expandafter\def\csname r@cor:rank111\endcsname{{9.14}{119}{}{proposition.9.14}{Dirichlet-GRH-v930-Volume-I.pdf}}
\expandafter\def\csname r@cor:square215\endcsname{{98.33}{31}{}{proposition.98.33}{Dirichlet-GRH-v930-Volume-XIX.pdf}}
\expandafter\def\csname r@cor:square216\endcsname{{98.45}{39}{}{proposition.98.45}{Dirichlet-GRH-v930-Volume-XIX.pdf}}
\expandafter\def\csname r@cor:transportarray533\endcsname{{96.29}{94}{}{proposition.96.29}{Dirichlet-GRH-v930-Volume-XVIII.pdf}}
\expandafter\def\csname r@cor:whittakerterminal527\endcsname{{96.8}{78}{}{proposition.96.8}{Dirichlet-GRH-v930-Volume-XVIII.pdf}}
\expandafter\def\csname r@cor:zero135\endcsname{{86.8}{8}{}{proposition.86.8}{Dirichlet-GRH-v930-Volume-XVII.pdf}}
\expandafter\def\csname r@eq:cross132\endcsname{{86.2}{5}{}{equation.86.2}{Dirichlet-GRH-v930-Volume-XVII.pdf}}
\expandafter\def\csname r@eq:fprod215\endcsname{{98.22}{31}{}{equation.98.22}{Dirichlet-GRH-v930-Volume-XIX.pdf}}
\expandafter\def\csname r@eq:fratio215\endcsname{{98.23}{31}{}{equation.98.23}{Dirichlet-GRH-v930-Volume-XIX.pdf}}
\expandafter\def\csname r@eq:gram158\endcsname{{87.11}{50}{}{equation.87.11}{Dirichlet-GRH-v930-Volume-XVII.pdf}}
\expandafter\def\csname r@eq:label209\endcsname{{93.6}{28}{}{equation.93.6}{Dirichlet-GRH-v930-Volume-XVIII.pdf}}
\expandafter\def\csname r@eq:letters111\endcsname{{9.1}{117}{}{equation.9.1}{Dirichlet-GRH-v930-Volume-I.pdf}}
\expandafter\def\csname r@eq:marked215\endcsname{{98.20}{29}{}{equation.98.20}{Dirichlet-GRH-v930-Volume-XIX.pdf}}
\expandafter\def\csname r@eq:match209\endcsname{{93.7}{30}{}{equation.93.7}{Dirichlet-GRH-v930-Volume-XVIII.pdf}}
\expandafter\def\csname r@eq:orb133\endcsname{{98.1}{3}{}{equation.98.1}{Dirichlet-GRH-v930-Volume-XIX.pdf}}
\expandafter\def\csname r@eq:real127\endcsname{{92.15}{19}{}{equation.92.15}{Dirichlet-GRH-v930-Volume-XVIII.pdf}}
\expandafter\def\csname r@eq:region198\endcsname{{20.12}{33}{}{equation.20.12}{Dirichlet-GRH-v930-Volume-IV.pdf}}
\expandafter\def\csname r@eq:seed127\endcsname{{92.16}{20}{}{equation.92.16}{Dirichlet-GRH-v930-Volume-XVIII.pdf}}
\expandafter\def\csname r@eq:success135\endcsname{{86.4}{8}{}{equation.86.4}{Dirichlet-GRH-v930-Volume-XVII.pdf}}
\expandafter\def\csname r@eq:tail114\endcsname{{19.15}{17}{}{equation.19.15}{Dirichlet-GRH-v930-Volume-IV.pdf}}
\expandafter\def\csname r@eq:tr128\endcsname{{100.5}{79}{}{equation.100.5}{Dirichlet-GRH-v930-Volume-XIX.pdf}}
\expandafter\def\csname r@eq:window\endcsname{{{W}}{2}{}{AMS.2}{Dirichlet-GRH-v930-Volume-I.pdf}}
\expandafter\def\csname r@lem:archdet303\endcsname{{126.6}{15}{}{proposition.126.6}{Dirichlet-GRH-v930-Volume-XXIV.pdf}}
\expandafter\def\csname r@lem:binary230\endcsname{{86.5}{6}{}{proposition.86.5}{Dirichlet-GRH-v930-Volume-XVII.pdf}}
\expandafter\def\csname r@lem:blockphysical295\endcsname{{100.12}{73}{}{proposition.100.12}{Dirichlet-GRH-v930-Volume-XIX.pdf}}
\expandafter\def\csname r@lem:classes304\endcsname{{119.2}{4}{}{proposition.119.2}{Dirichlet-GRH-v930-Volume-XXIII.pdf}}
\expandafter\def\csname r@lem:cosets163\endcsname{{86.21}{18}{}{proposition.86.21}{Dirichlet-GRH-v930-Volume-XVII.pdf}}
\expandafter\def\csname r@lem:covol127\endcsname{{92.23}{19}{}{proposition.92.23}{Dirichlet-GRH-v930-Volume-XVIII.pdf}}
\expandafter\def\csname r@lem:deep163\endcsname{{86.23}{19}{}{proposition.86.23}{Dirichlet-GRH-v930-Volume-XVII.pdf}}
\expandafter\def\csname r@lem:density207\endcsname{{98.15}{16}{}{proposition.98.15}{Dirichlet-GRH-v930-Volume-XIX.pdf}}
\expandafter\def\csname r@lem:dimension239\endcsname{{94.22}{55}{}{proposition.94.22}{Dirichlet-GRH-v930-Volume-XVIII.pdf}}
\expandafter\def\csname r@lem:factorial291\endcsname{{105.9}{40}{}{proposition.105.9}{Dirichlet-GRH-v930-Volume-XX.pdf}}
\expandafter\def\csname r@lem:flag215\endcsname{{98.30}{29}{}{proposition.98.30}{Dirichlet-GRH-v930-Volume-XIX.pdf}}
\expandafter\def\csname r@lem:flagcollision530\endcsname{{96.13}{82}{}{proposition.96.13}{Dirichlet-GRH-v930-Volume-XVIII.pdf}}
\expandafter\def\csname r@lem:flags211\endcsname{{98.19}{19}{}{proposition.98.19}{Dirichlet-GRH-v930-Volume-XIX.pdf}}
\expandafter\def\csname r@lem:markedtransport533\endcsname{{96.27}{92}{}{proposition.96.27}{Dirichlet-GRH-v930-Volume-XVIII.pdf}}
\expandafter\def\csname r@lem:mass209\endcsname{{93.7}{26}{}{proposition.93.7}{Dirichlet-GRH-v930-Volume-XVIII.pdf}}
\expandafter\def\csname r@lem:match209\endcsname{{93.10}{28}{}{proposition.93.10}{Dirichlet-GRH-v930-Volume-XVIII.pdf}}
\expandafter\def\csname r@lem:middlelattice531\endcsname{{96.18}{86}{}{proposition.96.18}{Dirichlet-GRH-v930-Volume-XVIII.pdf}}
\expandafter\def\csname r@lem:norm173\endcsname{{72.12}{56}{}{proposition.72.12}{Dirichlet-GRH-v930-Volume-XIV.pdf}}
\expandafter\def\csname r@lem:power291\endcsname{{107.11}{82}{}{proposition.107.11}{Dirichlet-GRH-v930-Volume-XX.pdf}}
\expandafter\def\csname r@lem:real287\endcsname{{105.4}{36}{}{proposition.105.4}{Dirichlet-GRH-v930-Volume-XX.pdf}}
\expandafter\def\csname r@lem:size97\endcsname{{87.1}{33}{}{proposition.87.1}{Dirichlet-GRH-v930-Volume-XVII.pdf}}
\expandafter\def\csname r@lem:stab205\endcsname{{98.6}{7}{}{proposition.98.6}{Dirichlet-GRH-v930-Volume-XIX.pdf}}
\expandafter\def\csname r@lem:translate163\endcsname{{86.19}{16}{}{proposition.86.19}{Dirichlet-GRH-v930-Volume-XVII.pdf}}
\expandafter\def\csname r@lem:transport129\endcsname{{100.14}{76}{}{proposition.100.14}{Dirichlet-GRH-v930-Volume-XIX.pdf}}
\expandafter\def\csname r@lem:whittakermatrix526\endcsname{{96.2}{74}{}{proposition.96.2}{Dirichlet-GRH-v930-Volume-XVIII.pdf}}
\expandafter\def\csname r@lem:zero133\endcsname{{98.1}{3}{}{proposition.98.1}{Dirichlet-GRH-v930-Volume-XIX.pdf}}
\expandafter\def\csname r@nav:VI-15\endcsname{{92.8}{18}{}{subsection.92.8}{Dirichlet-GRH-v930-Volume-XVIII.pdf}}
\expandafter\def\csname r@nav:VI-16\endcsname{{100.14}{76}{}{subsection.100.14}{Dirichlet-GRH-v930-Volume-XIX.pdf}}
\expandafter\def\csname r@nav:VII-02\endcsname{{86.2.1}{6}{}{subsubsection.86.2.1}{Dirichlet-GRH-v930-Volume-XVII.pdf}}
\expandafter\def\csname r@prop:average163\endcsname{{86.25}{21}{}{proposition.86.25}{Dirichlet-GRH-v930-Volume-XVII.pdf}}
\expandafter\def\csname r@prop:composition298\endcsname{{100.9}{71}{}{proposition.100.9}{Dirichlet-GRH-v930-Volume-XIX.pdf}}
\expandafter\def\csname r@prop:compositionfinite298\endcsname{{100.4}{65}{}{proposition.100.4}{Dirichlet-GRH-v930-Volume-XIX.pdf}}
\expandafter\def\csname r@prop:cost177\endcsname{{72.23}{65}{}{proposition.72.23}{Dirichlet-GRH-v930-Volume-XIV.pdf}}
\expandafter\def\csname r@prop:covol304\endcsname{{100.6}{68}{}{proposition.100.6}{Dirichlet-GRH-v930-Volume-XIX.pdf}}
\expandafter\def\csname r@prop:depth245\endcsname{{95.3}{67}{}{proposition.95.3}{Dirichlet-GRH-v930-Volume-XVIII.pdf}}
\expandafter\def\csname r@prop:dilation179\endcsname{{104.3}{4}{}{proposition.104.3}{Dirichlet-GRH-v930-Volume-XX.pdf}}
\expandafter\def\csname r@prop:fibres211\endcsname{{98.20}{20}{}{proposition.98.20}{Dirichlet-GRH-v930-Volume-XIX.pdf}}
\expandafter\def\csname r@prop:finite237\endcsname{{94.20}{53}{}{proposition.94.20}{Dirichlet-GRH-v930-Volume-XVIII.pdf}}
\expandafter\def\csname r@prop:finiteall294\endcsname{{101.5}{84}{}{proposition.101.5}{Dirichlet-GRH-v930-Volume-XIX.pdf}}
\expandafter\def\csname r@prop:index90\endcsname{{94.6}{41}{}{proposition.94.6}{Dirichlet-GRH-v930-Volume-XVIII.pdf}}
\expandafter\def\csname r@prop:lattice282\endcsname{{99.11}{51}{}{proposition.99.11}{Dirichlet-GRH-v930-Volume-XIX.pdf}}
\expandafter\def\csname r@prop:loops302\endcsname{{119.1}{3}{}{proposition.119.1}{Dirichlet-GRH-v930-Volume-XXIII.pdf}}
\expandafter\def\csname r@prop:marked215\endcsname{{98.31}{30}{}{proposition.98.31}{Dirichlet-GRH-v930-Volume-XIX.pdf}}
\expandafter\def\csname r@prop:marked291\endcsname{{99.3}{42}{}{proposition.99.3}{Dirichlet-GRH-v930-Volume-XIX.pdf}}
\expandafter\def\csname r@prop:markedpowers292\endcsname{{99.7}{46}{}{proposition.99.7}{Dirichlet-GRH-v930-Volume-XIX.pdf}}
\expandafter\def\csname r@prop:measure283\endcsname{{105.19}{50}{}{proposition.105.19}{Dirichlet-GRH-v930-Volume-XX.pdf}}
\expandafter\def\csname r@prop:partnerweights531\endcsname{{96.20}{87}{}{proposition.96.20}{Dirichlet-GRH-v930-Volume-XVIII.pdf}}
\expandafter\def\csname r@prop:phase301\endcsname{{126.5}{14}{}{proposition.126.5}{Dirichlet-GRH-v930-Volume-XXIV.pdf}}
\expandafter\def\csname r@prop:places288\endcsname{{105.13}{44}{}{proposition.105.13}{Dirichlet-GRH-v930-Volume-XX.pdf}}
\expandafter\def\csname r@prop:powers290\endcsname{{99.2}{41}{}{proposition.99.2}{Dirichlet-GRH-v930-Volume-XIX.pdf}}
\expandafter\def\csname r@prop:quot215\endcsname{{98.28}{27}{}{proposition.98.28}{Dirichlet-GRH-v930-Volume-XIX.pdf}}
\expandafter\def\csname r@prop:rank11\endcsname{{6.9}{30}{}{proposition.6.9}{Dirichlet-GRH-v930-Volume-I.pdf}}
\expandafter\def\csname r@prop:ray50\endcsname{{6.30}{63}{}{proposition.6.30}{Dirichlet-GRH-v930-Volume-I.pdf}}
\expandafter\def\csname r@prop:scale284\endcsname{{105.14}{45}{}{proposition.105.14}{Dirichlet-GRH-v930-Volume-XX.pdf}}
\expandafter\def\csname r@prop:schurframe259\endcsname{{75.7}{120}{}{proposition.75.7}{Dirichlet-GRH-v930-Volume-XIV.pdf}}
\expandafter\def\csname r@prop:sharp237\endcsname{{94.21}{54}{}{proposition.94.21}{Dirichlet-GRH-v930-Volume-XVIII.pdf}}
\expandafter\def\csname r@prop:shell263\endcsname{{75.5}{117}{}{proposition.75.5}{Dirichlet-GRH-v930-Volume-XIV.pdf}}
\expandafter\def\csname r@prop:split287\endcsname{{100.10}{72}{}{proposition.100.10}{Dirichlet-GRH-v930-Volume-XIX.pdf}}
\expandafter\def\csname r@prop:terminal293\endcsname{{100.11}{73}{}{proposition.100.11}{Dirichlet-GRH-v930-Volume-XIX.pdf}}
\expandafter\def\csname r@prop:terminalcore255\endcsname{{95.4}{68}{}{proposition.95.4}{Dirichlet-GRH-v930-Volume-XVIII.pdf}}
\expandafter\def\csname r@prop:weight163\endcsname{{86.20}{17}{}{proposition.86.20}{Dirichlet-GRH-v930-Volume-XVII.pdf}}
\expandafter\def\csname r@prop:word166\endcsname{{83.13}{62}{}{proposition.83.13}{Dirichlet-GRH-v930-Volume-XVI.pdf}}
\expandafter\def\csname r@sec:char94\endcsname{{121.16}{51}{}{subsection.121.16}{Dirichlet-GRH-v930-Volume-XXIII.pdf}}
\expandafter\def\csname r@sec:comparison\endcsname{{6}{24}{}{section.6}{Dirichlet-GRH-v930-Volume-I.pdf}}
\expandafter\def\csname r@sec:flags215\endcsname{{98.9}{29}{}{subsection.98.9}{Dirichlet-GRH-v930-Volume-XIX.pdf}}
\expandafter\def\csname r@sec:fullrankresidue533\endcsname{{96.8}{92}{}{subsection.96.8}{Dirichlet-GRH-v930-Volume-XVIII.pdf}}
\expandafter\def\csname r@sec:missing\endcsname{{6.6}{32}{}{subsection.6.6}{Dirichlet-GRH-v930-Volume-I.pdf}}
\expandafter\def\csname r@sec:partnerlifts531\endcsname{{96.6}{85}{}{subsection.96.6}{Dirichlet-GRH-v930-Volume-XVIII.pdf}}
\expandafter\def\csname r@sec:realtower169\endcsname{{72.3}{50}{}{subsection.72.3}{Dirichlet-GRH-v930-Volume-XIV.pdf}}
\expandafter\def\csname r@sec:terminal209\endcsname{{93.3}{26}{}{subsection.93.3}{Dirichlet-GRH-v930-Volume-XVIII.pdf}}
\expandafter\def\csname r@seriesguide\endcsname{{1.1}{4}{}{subsection.1.1}{Dirichlet-GRH-v930-Volume-I.pdf}}
\expandafter\def\csname r@thm:adjoint157\endcsname{{86.14}{12}{}{proposition.86.14}{Dirichlet-GRH-v930-Volume-XVII.pdf}}
\expandafter\def\csname r@thm:all202\endcsname{{20.15}{34}{}{proposition.20.15}{Dirichlet-GRH-v930-Volume-IV.pdf}}
\expandafter\def\csname r@thm:allblocks297\endcsname{{107.13}{84}{}{proposition.107.13}{Dirichlet-GRH-v930-Volume-XX.pdf}}
\expandafter\def\csname r@thm:allfinite215\endcsname{{98.32}{30}{}{proposition.98.32}{Dirichlet-GRH-v930-Volume-XIX.pdf}}
\expandafter\def\csname r@thm:balanced179\endcsname{{104.2}{4}{}{proposition.104.2}{Dirichlet-GRH-v930-Volume-XX.pdf}}
\expandafter\def\csname r@thm:blocktail295\endcsname{{107.1}{73}{}{proposition.107.1}{Dirichlet-GRH-v930-Volume-XX.pdf}}
\expandafter\def\csname r@thm:char125\endcsname{{92.18}{15}{}{proposition.92.18}{Dirichlet-GRH-v930-Volume-XVIII.pdf}}
\expandafter\def\csname r@thm:common111\endcsname{{9.13}{118}{}{proposition.9.13}{Dirichlet-GRH-v930-Volume-I.pdf}}
\expandafter\def\csname r@thm:common133\endcsname{{98.2}{4}{}{proposition.98.2}{Dirichlet-GRH-v930-Volume-XIX.pdf}}
\expandafter\def\csname r@thm:commonnorm128\endcsname{{100.16}{78}{}{proposition.100.16}{Dirichlet-GRH-v930-Volume-XIX.pdf}}
\expandafter\def\csname r@thm:commonorbit305\endcsname{{121.5}{34}{}{proposition.121.5}{Dirichlet-GRH-v930-Volume-XXIII.pdf}}
\expandafter\def\csname r@thm:compositionreal298\endcsname{{107.14}{85}{}{proposition.107.14}{Dirichlet-GRH-v930-Volume-XX.pdf}}
\expandafter\def\csname r@thm:conductor232\endcsname{{19.15}{12}{}{proposition.19.15}{Dirichlet-GRH-v930-Volume-IV.pdf}}
\expandafter\def\csname r@thm:cross132\endcsname{{86.3}{4}{}{proposition.86.3}{Dirichlet-GRH-v930-Volume-XVII.pdf}}
\expandafter\def\csname r@thm:deep163\endcsname{{86.24}{20}{}{proposition.86.24}{Dirichlet-GRH-v930-Volume-XVII.pdf}}
\expandafter\def\csname r@thm:disk230\endcsname{{86.6}{7}{}{proposition.86.6}{Dirichlet-GRH-v930-Volume-XVII.pdf}}
\expandafter\def\csname r@thm:doubleorbit532\endcsname{{96.26}{91}{}{proposition.96.26}{Dirichlet-GRH-v930-Volume-XVIII.pdf}}
\expandafter\def\csname r@thm:fiber283\endcsname{{105.16}{47}{}{proposition.105.16}{Dirichlet-GRH-v930-Volume-XX.pdf}}
\expandafter\def\csname r@thm:finite216\endcsname{{98.44}{39}{}{proposition.98.44}{Dirichlet-GRH-v930-Volume-XIX.pdf}}
\expandafter\def\csname r@thm:finite281\endcsname{{99.10}{50}{}{proposition.99.10}{Dirichlet-GRH-v930-Volume-XIX.pdf}}
\expandafter\def\csname r@thm:first250\endcsname{{71.10}{25}{}{proposition.71.10}{Dirichlet-GRH-v930-Volume-XIV.pdf}}
\expandafter\def\csname r@thm:fullrankseparation533\endcsname{{96.28}{93}{}{proposition.96.28}{Dirichlet-GRH-v930-Volume-XVIII.pdf}}
\expandafter\def\csname r@thm:gap132\endcsname{{86.4}{5}{}{proposition.86.4}{Dirichlet-GRH-v930-Volume-XVII.pdf}}
\expandafter\def\csname r@thm:global158\endcsname{{87.26}{50}{}{proposition.87.26}{Dirichlet-GRH-v930-Volume-XVII.pdf}}
\expandafter\def\csname r@thm:identity127\endcsname{{92.24}{20}{}{proposition.92.24}{Dirichlet-GRH-v930-Volume-XVIII.pdf}}
\expandafter\def\csname r@thm:invar148\endcsname{{20.11}{31}{}{proposition.20.11}{Dirichlet-GRH-v930-Volume-IV.pdf}}
\expandafter\def\csname r@thm:localmass156\endcsname{{87.24}{48}{}{proposition.87.24}{Dirichlet-GRH-v930-Volume-XVII.pdf}}
\expandafter\def\csname r@thm:long232\endcsname{{5.6}{23}{}{proposition.5.6}{Dirichlet-GRH-v930-Volume-I.pdf}}
\expandafter\def\csname r@thm:low285\endcsname{{105.18}{49}{}{proposition.105.18}{Dirichlet-GRH-v930-Volume-XX.pdf}}
\expandafter\def\csname r@thm:mass205\endcsname{{98.8}{9}{}{proposition.98.8}{Dirichlet-GRH-v930-Volume-XIX.pdf}}
\expandafter\def\csname r@thm:midlevel131\endcsname{{86.1}{3}{}{proposition.86.1}{Dirichlet-GRH-v930-Volume-XVII.pdf}}
\expandafter\def\csname r@thm:midnorm131\endcsname{{86.2}{4}{}{proposition.86.2}{Dirichlet-GRH-v930-Volume-XVII.pdf}}
\expandafter\def\csname r@thm:moment126\endcsname{{92.21}{17}{}{proposition.92.21}{Dirichlet-GRH-v930-Volume-XVIII.pdf}}
\expandafter\def\csname r@thm:moment283\endcsname{{105.10}{41}{}{proposition.105.10}{Dirichlet-GRH-v930-Volume-XX.pdf}}
\expandafter\def\csname r@thm:moment286\endcsname{{105.3}{34}{}{proposition.105.3}{Dirichlet-GRH-v930-Volume-XX.pdf}}
\expandafter\def\csname r@thm:norm177\endcsname{{72.21}{64}{}{proposition.72.21}{Dirichlet-GRH-v930-Volume-XIV.pdf}}
\expandafter\def\csname r@thm:norm211\endcsname{{98.21}{21}{}{proposition.98.21}{Dirichlet-GRH-v930-Volume-XIX.pdf}}
\expandafter\def\csname r@thm:norm212\endcsname{{98.41}{36}{}{proposition.98.41}{Dirichlet-GRH-v930-Volume-XIX.pdf}}
\expandafter\def\csname r@thm:norm284\endcsname{{105.17}{48}{}{proposition.105.17}{Dirichlet-GRH-v930-Volume-XX.pdf}}
\expandafter\def\csname r@thm:orbitaverage532\endcsname{{96.24}{90}{}{proposition.96.24}{Dirichlet-GRH-v930-Volume-XVIII.pdf}}
\expandafter\def\csname r@thm:orbitnoncancellation533\endcsname{{96.30}{94}{}{proposition.96.30}{Dirichlet-GRH-v930-Volume-XVIII.pdf}}
\expandafter\def\csname r@thm:ordinary163\endcsname{{86.22}{18}{}{proposition.86.22}{Dirichlet-GRH-v930-Volume-XVII.pdf}}
\expandafter\def\csname r@thm:prime209\endcsname{{93.9}{28}{}{proposition.93.9}{Dirichlet-GRH-v930-Volume-XVIII.pdf}}
\expandafter\def\csname r@thm:prime210\endcsname{{93.17}{34}{}{proposition.93.17}{Dirichlet-GRH-v930-Volume-XVIII.pdf}}
\expandafter\def\csname r@thm:product301\endcsname{{100.5}{66}{}{proposition.100.5}{Dirichlet-GRH-v930-Volume-XIX.pdf}}
\expandafter\def\csname r@thm:range239\endcsname{{88.1}{52}{}{proposition.88.1}{Dirichlet-GRH-v930-Volume-XVII.pdf}}
\expandafter\def\csname r@thm:rangenew50\endcsname{{6.29}{62}{}{proposition.6.29}{Dirichlet-GRH-v930-Volume-I.pdf}}
\expandafter\def\csname r@thm:ratio248\endcsname{{89.3}{76}{}{proposition.89.3}{Dirichlet-GRH-v930-Volume-XVII.pdf}}
\expandafter\def\csname r@thm:rationalrank302\endcsname{{102.4}{91}{}{proposition.102.4}{Dirichlet-GRH-v930-Volume-XIX.pdf}}
\expandafter\def\csname r@thm:real282\endcsname{{105.15}{46}{}{proposition.105.15}{Dirichlet-GRH-v930-Volume-XX.pdf}}
\expandafter\def\csname r@thm:realtrans300\endcsname{{107.3}{75}{}{proposition.107.3}{Dirichlet-GRH-v930-Volume-XX.pdf}}
\expandafter\def\csname r@thm:saturate215\endcsname{{98.29}{28}{}{proposition.98.29}{Dirichlet-GRH-v930-Volume-XIX.pdf}}
\expandafter\def\csname r@thm:skew289\endcsname{{105.6}{37}{}{proposition.105.6}{Dirichlet-GRH-v930-Volume-XX.pdf}}
\expandafter\def\csname r@thm:split281\endcsname{{101.1}{80}{}{proposition.101.1}{Dirichlet-GRH-v930-Volume-XIX.pdf}}
\expandafter\def\csname r@thm:success135\endcsname{{86.7}{8}{}{proposition.86.7}{Dirichlet-GRH-v930-Volume-XVII.pdf}}
\expandafter\def\csname r@thm:tensor249\endcsname{{106.9}{62}{}{proposition.106.9}{Dirichlet-GRH-v930-Volume-XX.pdf}}
\expandafter\def\csname r@thm:transport233\endcsname{{94.16}{50}{}{proposition.94.16}{Dirichlet-GRH-v930-Volume-XVIII.pdf}}
\makeatother
\begin{document}
\setcounter{part}{20}
\fancyhead[L]{\small\color{muted}Dirichlet GRH \textperiodcentered\ Volume XXI}
\fancyhead[R]{\small\color{muted}Rational flag masses and traces}
\thispagestyle{plain}\begin{center}
{\LARGE\bfseries The generalized Riemann hypothesis\\[2pt]for Dirichlet \(L\)-functions}\\[6pt]
{\small 27 November 2026 \quad\textperiodcentered\quad Edition 930}\end{center}
\part{Rational flag masses and ordinary traces at the marks}\label{part:XXI}
\begin{abstract}\noindent We count rational flags and their stabilizers, compute the mass of the terminal classes and use it in the ordinary trace. The last section identifies the mixed terms and the identity coefficient at the two marked primes.\end{abstract}
\setcounter{section}{108}
\noindent The \hyperref[seriesguide]{guide in Volume~I} lists all twenty-seven volumes.
\begingroup\small\setlength{\parskip}{0pt}\tableofcontents\endgroup
\clearpage

\section{Rational flags and the mass formula}\label{nav:g-XI-0}\label{nav:volume-XXI}

\subsection{The weight needed to sum the rational periods}\label{nav:IX-09}
\label{sec:moment217}

We keep \(h=k-1\), \(S_*=d(d-1)\), \(m=h-2d+1\), and the
normalizations of Section~\ref{sec:flags215}. For a supported
nondegenerate label write \(A_j=|N_{E/\QQ}D_j|\) at nonzero
endpoints, and use \(A_j^\#\) for the corresponding norms after
the integral change of variables in Lemma~\ref{lem:match209}.
Put \(G=|N_{F/\QQ}\mathcal G|\), \(G^\#=|N_{F/\QQ}\mathcal G^\#|\),
and \(B_\gamma^\#=G^\#\prod_\nu A_{j_\nu}^\#\).
The test pair and the two unit column tuples are fixed while the
label varies. If \(\mathcal A\) is any finite set of these labels,
write
\[
   \mathcal H_\theta(\mathcal A)
     =\sum_{\gamma\in\mathcal A} w(\beta_\gamma)(B_\gamma^\#)^\theta
               P_\infty^\#\|\widehat K_{\gamma^\#}\|_{\rm HS}^2 .
\]
Every term uses its own orbit and the same transported sections.
Here \(w(\beta)=\prod_{\nu,\tau}(1+|\tau D_{j_\nu}^\#|^2)^2\)
is the weight of the reindexed integral endpoints.

\begin{theorem}[An absolute period bound from a real weighted sum]
\label{thm:moment217}
For \(h\ge4(d-1)\), which holds on the common intervals for
large \(x\), one has
\begin{align}
 \sum_{\gamma\in\mathcal A,c}|J_{\gamma,c}|
 &\le
 \frac{A_*\,\mathfrak m_h\,8^{fd/2}}{\sqrt{c_{\rho_1}^*c_{\rho_2}^*}}\,
  e^{O(f2^{-m})} R^{-fS_*}
   Q^{f(S_*/2+d-2-d^2h)}\,
      \mathcal H_{h/2}(\mathcal A)^{1/2}.              \label{eq:absolute217}
\end{align}
There is no constant depending on the truncation \(\mathcal A\).
\end{theorem}
\begin{proof}
The proof of Theorem~\ref{thm:norm211} bounds the absolute cell
sum for a fixed label by its compatible mass times
\(A_*c_E^a\|K_{\infty,\gamma}^{\rm area}\|_{\rm op}
 /\sqrt{c_{\rho_1}^*c_{\rho_2}^*}\).
Insert \eqref{eq:match209} and \eqref{eq:fratio215}. Apart from
the common scalar
\(A_*\mathfrak m_h e^{O(f2^{-m})}R^{-fS_*}/
\sqrt{c_{\rho_1}^*c_{\rho_2}^*}\), the result is
\[
 Q^{f(a-d-2dh)} P_\infty^{1/2}W_\gamma
                 (P_\infty^\#)^{1/2}\|\widehat K_{\gamma^\#}\|_{\rm op}.
\]
We used \(d(d-2)-S_*=-d\). The powers of the endpoints in this
expression can be found exactly. By Lemma~\ref{lem:density207}
and the telescoping calculation in Corollary~\ref{cor:square215},
\begin{equation}
       P_\infty^{1/2}W_\gamma
             =G^{h/2}\prod_{\nu<l}A_{j_\nu}^{h_{\nu+1}} .
       \label{eq:linear217}
\end{equation}
Indeed the two exponents on a proper endpoint add to
\((h_\nu+h_{\nu+1})/2+(h_{\nu+1}-h_\nu)/2=h_{\nu+1}\).
At the final endpoint they add to zero.

Substitute \(G=Q^{-2fd(d-2)}G^\#\) and
\(A_j=Q^{-2fj(j-2)}A_j^\#\). The overall exponent of \(Q\) is
now
\[
 f\Bigl(a-d-d^2h-2\sum_{\nu<l}j_\nu(j_\nu-2)h_{\nu+1}\Bigr)
 \le f(S_*/2+d-2-d^2h).
\]
Only \(j_\nu=1\) can give a negative summand, with
\(h_{\nu+1}\le d-1\), and \(a\le S_*/2\).
At each real place
\(0<\tau\mathcal G^\#\le|\tau D_d^\#|^2\). Thus \(G^\#\le A_d^\#\)
and \((G^\#)^{h/2}\le(G^\# A_d^\#)^{h/4}\).
All the proper endpoint norms \(A_{j_\nu}^\#\) are positive
integers and their exponents in \eqref{eq:linear217} are at
most \(d-1\le h/4\). This bounds the whole reindexed norm
factor by \((B_\gamma^\#)^{h/4}\).

Replace the operator norm by the Hilbert--Schmidt norm, insert
\(w(\beta_\gamma)^{-1/2}w(\beta_\gamma)^{1/2}\), and apply
Cauchy--Schwarz in \(\gamma\). The reindexing is injective
by Lemma~\ref{lem:match209}, and all its endpoints are integral.
Theorem~\ref{thm:allorbits203} gives
\(\sum_\gamma w(\beta_\gamma)^{-1}\le8^{fd}\).
The second squared sum is precisely \(\mathcal H_{h/2}\).
This proves \eqref{eq:absolute217} for every finite truncation,
and hence for the full nonnegative sums by monotone convergence.
\end{proof}

\begin{proposition}[A family with growing finite orbital mass]
\label{prop:escape217}
Fix one of the fields and sets of parameters above. There are
distinct rational labels \(\gamma_t\), \(t=1,2,\ldots\), all
regular and supported at the original level, with
\begin{align*}
 \mathcal M_{\gamma_t}^{=}\asymp t^{2fh},\qquad&
          P_\infty(\gamma_t)\asymp t^{2fh},\\
 w(\beta_{\gamma_t})\asymp t^{4f},\qquad&
          B_{\gamma_t}^\#\asymp t^{4f}.
\end{align*}
The comparison constants can depend on the fixed field, ranks
and levels. They are independent of \(t\).
\end{proposition}
\begin{proof}
Take the signed unit basis of the original row lattice.
For \(1\le i<d\) write \(u_i=e_{+,i}+e_{-,i}\) as in
Theorem~\ref{thm:midlevel131} and set \(w_i=e_{+,i}-e_{-,i}\).
There is a further positive unit basis vector \(e_0\) in the
positive middle space. With \(e\) the negative terminal basis
vector put \(w_d=e_0+e\). These \(d\) vectors \(w_i\) are
mutually orthogonal and isotropic. Choose nonzero
\(\alpha,\lambda\in Q^2T\mathcal O_E\) of trace zero to \(F\).
Such elements are obtained by multiplying a nonzero difference
\(b-\bar b\), \(b\in\mathcal O_E\), by \(Q^2T\).
For an isotropic \(w\) and \(\bar a=-a\), the map
\[
                   x\longmapsto x+a w H(w,x)
\]
is a rational isometry. Expanding the pairing gives linear
terms with coefficient \(a+\bar a\) and a quadratic term
containing \(H(w,w)\). Its inverse uses \(-a\).
Define \(\gamma_t\) by taking the product of these maps for
\((w_i,\alpha)\), \(i<d\), and \((w_d,t\lambda)\). All factors
preserve the standard lattice and are the identity modulo
\(Q^2T\). At a marked prime the principal congruence subgroup
lies in the row level. In the first split component at \(T\)
the modified lattice lies between the standard lattice and
its \(T\)-multiple. Thus an integral map congruent to the
identity modulo \(T\), as well as its inverse, preserves the
modified lattice. The dual component is then preserved as
well. Hence \(\gamma_t\) itself belongs to the original finite
row level, giving finite support.

Since \(H(u_i,w_i)=2\) and \(H(e,w_d)=-1\), the cross
matrix for \(Y,\gamma_tY\) is diagonal with entries
\[
                4\alpha,\ldots,4\alpha,\quad T(-1+t\lambda).
\]
All blocks therefore have size one. There is no common
unipotent in Lemma~\ref{lem:stab205}, so the relevance
condition on characters holds. The last minor distinguishes
the labels, and
\[
 \mathcal G_t
    =(4\alpha\overline{4\alpha})^{d-1}
                        T^2t^2\lambda\bar\lambda
\]
is totally positive. Its norm is a fixed positive constant
times \(t^{2f}\). All proper endpoint minors are fixed.
The absolute value of the final minor at each real place
is bounded above and below by positive fixed multiples of
\(t\), since \(\lambda\) is nonzero and purely imaginary
there. The same assertions, with fixed scaling factors, hold
after \eqref{eq:label209}. These facts give the asserted
growth of the two arithmetic weights. They give the growth of
\(P_\infty\) by Lemma~\ref{lem:density207}, since the exponent
of the final endpoint is zero.

For a lower bound on the finite mass use \eqref{eq:fprod215}.
Here \(D\) is trivial, and the integral subgroup in the two
changes of frame has volume one. Its cross matrices are integral
everywhere. At a marked place the diagonal cross matrix is
\(H_0\bmod\pi^2\) by construction. Integral upper unitriangular
changes on either side preserve \(H_0=-T E_{dd}\), so the new
cross matrices satisfy both congruences of \eqref{eq:marked215}.
Consequently \(V_v^{=}\ge1\) for all \(v\). Formula
\eqref{eq:fprod215} bounds the mass below by a positive fixed
multiple of \(G_t^h\), where \(G_t=|N_{F/\QQ}\mathcal G_t|\).
For the reverse bound apply \eqref{eq:fratio215} with
\(W_{\gamma_t}=1\) and \(a=0\). This completes the proof.
\end{proof}

\begin{corollary}[A necessary exponent for separating the masses]
\leavevmode\label{cor:diverge217}
Consider the Cauchy--Schwarz argument in which the real squared
sum is \(\mathcal H_\theta\) and the masses are put entirely
in the other factor. Up to constants independent of the label,
the sum in that other squared factor is
\begin{equation}
 \sum_{\gamma}
  \frac{Q^{-2fa}(\mathcal M_\gamma^{=})^2}
       {w(\beta_\gamma)(B_\gamma^\#)^\theta P_\infty(\gamma)} .
 \label{eq:diverge217}
\end{equation}
It diverges if
\(\theta\le(h-2)/2+1/(4f)\). In particular on the selected
intervals a moment with \(\theta=A d^2\) for a fixed \(A\)
does not make this arithmetic factor finite for all large \(x\).
\end{corollary}
\begin{proof}
Before using the bound on the mass, the single-label estimate
in the first paragraph of the proof of Theorem~\ref{thm:moment217}
has label-dependent coefficient
\(\mathcal M_\gamma^{=}Q^{-fa}P_\infty(\gamma)^{-1/2}\)
on \((P_\infty^\#)^{1/2}\|\widehat K_{\gamma^\#}\|_{\rm HS}\).
Separating the prescribed real weight gives \eqref{eq:diverge217}.
On the sequence in Proposition~\ref{prop:escape217} each summand
is comparable to \(t^{2f(h-2-2\theta)}\). The sum over positive
integers diverges when this exponent is at least \(-1\). Finally
the ranks on the common intervals have \(d\asymp x^{1/2-\epsilon/16}\)
and \(k\asymp x^{9/2+10\epsilon}\), so \(h/d^2\to\infty\).
\end{proof}

The smaller squared weight occurs in
\srcloc{Formalism/SharpRestriction/CuspSpanNormBounds/RotatedWeightedBound}{26}.
Corollary~\ref{cor:diverge217} concerns the arithmetic factor
of the specified absolute-value argument. Its proof uses a sum
over all supported labels. It yields no lower bound on the
analytic kernel of any one of those labels. A different
quantity appears in the source's mean over row lattices. With
the initial flag and a single row lattice fixed, it counts
equivalence classes of two terminal flags together with a
comparison map, subject to a determinant bound. Its rank-one
count fixes one terminal flag as well and counts the admissible
partners in that same lattice. The definitions are in
\srcloc{Formalism/SharpRestriction/CuspSpanNormBounds/CntPartners}{190}
and \srcloc{Formalism/SharpRestriction/CuspSpanNormBounds/CntPartners}{199}.
The row-weighted inequality is stated in
\srcloc{Formalism/SharpRestriction/CuspSpanNormBounds/CongruenceAtSplitFrobeniusHolds}{556}.
Keeping the common lattice couples the two flags. Establishing
the necessary estimates for that joint count and transporting
them to the designated period would require a further argument.

\subsection{Saturation of two flags in the same row lattice}\label{nav:IX-10}
\label{sec:saturation219}

Fix a supported, relevant rational label of nonzero Gram
determinant. Work on the fields and intervals of
Section~\ref{sec:flags215}. Denote its decorated frames by
\(Y_0,Y_1\), the spans of their first \(d-1\) vectors by \(X_i\),
and the spans of all \(d\) vectors by \(W_i\). Put
\(W=W_0\oplus W_1\), \(C=W^\perp\), \(n=2d\), \(h=k-1\)
and \(m=h-2d+1\). The form on \(W\) is nondegenerate, and \(C\)
is totally positive definite. In a compatible full-level row
lattice \(\Lambda\) put
\[
 U_i=\Lambda\cap W_i,\qquad U=U_0\oplus U_1,\qquad
 B=\Lambda\cap W,\qquad C_\Lambda=\Lambda\cap C .
\]
For an integral nondegenerate lattice \(M\) write
\(\delta(M)=[M^\vee:M]^{1/2}\), with dual in its own span.
All indices in this section are cardinalities of quotients.

At an unmarked finite place put \(B_v^+=U_v\).
At \(v\mid Q=g_1g_2\) choose primitive integral representatives
\(Y_{0,v},Y_{1,v}\) with equal reductions, which exist by
Lemma~\ref{lem:flags211}. Let \(B_v^+\) be generated by the
columns of \(Y_{0,v}\) and \((Y_{1,v}-Y_{0,v})/\pi_v\).
The generators belong to \(B_v\). Changing the uniformizer
multiplies the second tuple by a unit. If the two representatives
are changed by integral upper unitriangular matrices \(V_i\),
equality of their new reductions forces \(V_0\equiv V_1\bmod\pi_v\)
by independence of the first reduction. The new difference tuple
is the old one times \(V_1\), plus \(Y_{0,v}(V_1-V_0)/\pi_v\).
This proves independence of choices. The local lattices define
\(B^+\) by taking their intersection in \(W\). Localization
recovers them since they differ from \(U_v\) at finitely many
places. Write \(I_v=[B_v:B_v^+]\) and \(I=[B:B^+]=\prod_v I_v\).
The saturation condition for this row is \(I=1\).

\begin{proposition}[The determinant and the congruence index]\label{prop:index219}
The following identities hold on every compatible row.
\[
 [B^+:U]=Q^{2df},\qquad
 \delta(C_\Lambda)=\delta(B)
    =\frac{|N_{F/\QQ}\mathcal G|}{Q^{2df}I}.
\]
Let \(\mathcal L\) be the module of \(c\in C\) for which there
are \(x_i\in X_i\) with \(c+x_i\in\Lambda\) and
\[
 x_0-x_1\in Q\Lambda+(\Lambda\cap X_0)+(\Lambda\cap X_1).
\]
Then \(C_\Lambda\subset\mathcal L\) and
\([\mathcal L:C_\Lambda]\mid I\). In particular the two
lattices agree on every row satisfying the saturation condition.
\end{proposition}
\begin{proof}
At each place a primitive integral representative of a flag is
a basis of \(U_{i,v}\), by Lemma~\ref{lem:flag215}. The
determinant of the concatenated Gram matrix has valuation
\(v(\mathcal G)\). A square Hermitian Gram matrix \(H\) gives
local discriminant index \(q_v^{2v(\det H)}\), by taking the
determinant of the map \(H\) on the integer lattice of the
quadratic algebra. Thus \(\delta(U)=|N_{F/\QQ}\mathcal G|\).
At a marked place the change to the divided differences has
determinant \(\pi_v^{-d}\), hence
\([B_v^+:U_v]=q_v^{2d}\). Both primes split completely.
An over-lattice of finite index \(a\) divides \(\delta\) by
\(a\), by the Gram determinant change of basis locally.
This gives the formula for \(B\).

The orthogonal projections of \(\Lambda\) onto \(W,C\) have
images \(B^\vee,C_\Lambda^\vee\). Locally, an integral
functional on \(B_v\) extends to one on \(\Lambda_v\) because
\(B_v\) is primitive. Self-duality represents it by a vector
in \(\Lambda_v\). Orthogonal projection gives the claimed image,
and the proof for \(C_{\Lambda,v}\) is the same. The quotient
\(\Lambda/(B\oplus C_\Lambda)\) therefore projects bijectively
onto both discriminant modules. For example a class whose
\(W\)-projection belongs to \(B\) has, after subtracting that
element, a representative in \(C_\Lambda\). Hence their orders
agree.

For \(c\in\mathcal L\) choose the indicated \(x_i\).
The difference \(x_0-x_1\) belongs to \(B\), and its class
modulo \(U\) is independent of the choices since two choices
of \(x_i\) differ by an element of \(\Lambda\cap X_i\).
The resulting homomorphism kills \(C_\Lambda\). Conversely if
\(x_0-x_1\in U\), the direct decomposition in \(W_0\oplus W_1\)
implies \(x_i\in U_i\), so \(c\in C_\Lambda\). The congruence
places the image in \(Q(B/U)\), because \(Q\Lambda\cap W=QB\).
Now \(QB^+\subset U\) at every place. Multiplication by \(Q\)
gives a surjection \(B/B^+\to Q(B/U)\). The asserted divisibility
follows. This gives the index in the compatibility module defined
in \srcloc{Formalism/SharpRestriction/CuspSpanNormBounds/IsPrimitiveSub}{178}
for the actual integral flags and congruence used here.
\end{proof}

\begin{theorem}[Saturation in the orbital measure]\label{thm:moment219}
Normalize the compatible finite orbital measure of the fixed
label to probability, using the masses \(L_v^{=}\) of
Proposition~\ref{prop:marked215}. For every
\(0\le t\le m-2\) one has, with an absolute constant \(C_1\),
\[
 \Pr(I>1)\le C_1 f\,2^{-m},\qquad
 \mathbb E[I^t]\le\exp(C_1f\,2^{-(m-t)}).
\]
The second bound also holds with \(I\) replaced by
\([\mathcal L:C_\Lambda]\).
\end{theorem}
\begin{proof}
In the disintegration of Proposition~\ref{prop:marked215}
fix the two changes of basis. At an unmarked place the ordered
tuple is the concatenation of the two frames. At a marked place
it uses the divided differences. Its generated lattice in the
ambient row is precisely \(B_v^+\). Pulling back the saturation
under the column map identifies \(I_v\) with \([M':M_0]\) in
Theorem~\ref{thm:saturate215}. For each integral Gram target
the compatible ordered-frame test contains every primitive
tuple and is contained in the set of all integral tuples.
A primitive tuple contributes exactly \(p(h,u)\) after division
by the determinant factor in that theorem.

The normalized contribution of tuples with \(I_v>1\), weighted
by \(I_v^t\), is bounded above by
\[
 e^{C_2q_v^{-m}}\sum_{M'\supsetneqq M_0}[M':M_0]^{-(h-t)}
                  \ \le\ C_3q_v^{-(m-t)} .
\]
The first inequality uses the over-lattice formula and
\(|\log p(h,u)|\le C_2q_v^{-m}\). The last sum is at most the
product of geometric series in the proof of
Theorem~\ref{thm:saturate215}, with the term one removed.
Its logarithm after adding one is \(O(q_v^{-(m-t)})\), uniformly
for \(m-t\ge2\), in both quadratic algebras. Dividing by the
primitive contribution costs at most \(e^{C_2q_v^{-m}}\).
The bound holds in every nonempty Gram fibre. Hence it holds
relative to the full compatible measure after integration over
the two changes of basis. The marked multiplicity \(q_v^{d(d-1)}\)
is common to the weighted and unweighted integrals.

Each \(I_v\) is intrinsic to the full-level cell, so inserting
its weight in the orbital mass formula multiplies the local
factors under the same disintegration as in
Theorem~\ref{thm:mass205}. The definite complement contributes
its usual constant mass. For a fixed label, \(\delta(B)\ge1\)
and its local Gram ideal is integral. The local version of
Proposition~\ref{prop:index219} therefore forces \(I_v=1\)
away from the divisors of \(\mathcal G\). The product is finite.
The union bound at \(t=0\) and the product of the local bounds
\(1+C_3q_v^{-(m-t)}\) give the assertions since
\(\sum_v q_v^{-b}\le f\sum_{j\ge2}j^{-b}\ll f2^{-b}\) for
\(b\ge2\). The index assertion follows from divisibility.
\end{proof}

\begin{corollary}[Contribution of the rows failing saturation]
\label{cor:bad219}
For unit column tuples at the original finite level, the sum
of \(|J_{\gamma,c}|\) over cells with \(I>1\) is at most
\[
 C_1 f2^{-m}\,
 \frac{A_*\mathcal M_\gamma^{=}}{\sqrt{c_{\rho_1}^*c_{\rho_2}^*}}\,
                  c_E^a\|K_{\infty,\gamma}^{\rm area}\|_{\rm op}.
\]
For \(I\ge b_0>1\) one may replace the factor \(C_1 f2^{-m}\) by
\(b_0^{-t}\exp(C_1f2^{-(m-t)})\).
\end{corollary}
\begin{proof}
Every refinement of a compatible coarse cell has the same row
lattice and the same two saturated flag modules. The definition
of \(B_v^+\) above is independent of the matching representatives.
Thus \(I\) is constant across all those refinements.
Apply Proposition~\ref{prop:fibres211} only to the specified
coarse cells and use the unit column norms as in
Theorem~\ref{thm:norm211}. The measure left in this inequality
is \(\mathcal M_\gamma^{=}\) times the probability of the
specified event. Theorem~\ref{thm:moment219} proves the first
bound and Markov's inequality proves the second.
\end{proof}

For generic rows of the original period the compatibility module
is consequently the integral perpendicular \(C_\Lambda\).
Controlling pairs of terminal flags with prescribed comparison
map also requires a count of the common definite subspaces they
can determine in the initial column.

\subsection{Mass with a prescribed discriminant module}\label{nav:IX-11}
\label{sec:discmass221}

Let \(E/F\) be one of the CM extensions in
Lemma~\ref{lem:mass209}, with \(f=[F:\QQ]\). It is unramified
at every finite place. For a positive definite integral
Hermitian lattice \(M\), let \(\mathcal D(M)=M^\vee/M\).
Its pairing has values in \(E/\OO_E\), given by the Hermitian
pairing of representatives. It is perfect by double duality.
Put \(\delta(M)=|\mathcal D(M)|^{1/2}\), and let
\(\mathfrak m_a\) be the mass of the standard self-dual
genus in rank \(a\). Mass is the sum of the reciprocals of
the isometry group orders over isometry classes in a genus.
All isometries are \(\OO_E\)-linear.

At a finite place \(v\), let \(A=E\otimes_F F_v\) with
integer ring \(O\), let \(q\) be the residue order of \(F_v\),
and let \(\pi\) be a uniformizer of \(F_v\). Put \(\eta=1\)
if \(A\) splits and \(\eta=-1\) otherwise. Write
\[
       P_b=\prod_{i=1}^{b}(1-\eta^i q^{-i}),\qquad P_0=1 .
\]
In a basis over \(O\) an integral nondegenerate Hermitian
lattice \(L\) of rank \(a\) has Gram matrix
\(H=\operatorname{diag}(\pi^{b_1},\ldots,\pi^{b_a})\),
with \(b_i\ge0\). In the split algebra this is Smith reduction
of the pairing between the two factors. In the unramified
field divide the form by its smallest entry scale. Some
residue pairing is nonzero. If every norm vanished, taking
the norm of \(x+c y\) for all residue scalars \(c\) and using
the perfect trace pairing would make every pairing zero.
A vector of unit norm splits off a line by orthogonal
projection. Induction and surjectivity of the norm on units
give the asserted diagonalization. The latter surjectivity
follows from the residue norm and correction of successive
digits by the residue trace. This argument includes primes
over two. The sorted exponents are unique by the elementary
divisor theorem. Write \(e=\sum_i b_i\) and
\(u=\#\{i:b_i>0\}\). Then \(\delta_v(L)=q^e\), and the
positive exponents determine the finite paired module.

\begin{lemma}[The stabilizer of a discriminant marking]
\label{lem:discstab221}
Let \(K_L\) be the stabilizer of \(L\) in its rational unitary
group, and \(K_L^\#\) the kernel of its action on
\(\mathcal D(L)\). The map \(K_L\to\Aut\mathcal D(L)\)
is onto. There is a Haar measure on each local unitary group,
invariant under transport by rational isometry, for which
\begin{equation}
                 \operatorname{vol}(K_L^\#)
                       =q^{-ae}\frac{P_{a-u}}{P_a}.
                       \label{eq:discvol221}
\end{equation}
\end{lemma}
\begin{proof}
In the diagonal basis \(L^\vee=H^{-1}O^a\). Write
\(g=1+XH\). The condition \((g-1)L^\vee\subset L\) is
\(X\in M_a(O)\), and the isometry equation becomes
\[
                      X+X^*+X^*HX=0 .
\]
Every such solution gives an integral isometry with integral
inverse. Its determinant has norm one and is integral in
both factors when \(A\) splits. Set \(t=a-u\). Modulo
\(\pi\), order the basis so that \(H=\operatorname{diag}(I_t,0)\)
and write the blocks of \(X\) as \(A_0,B,C,D\). The matrix
\(1+A_0\) is unitary, \(B\) is arbitrary,
\(C^*=-(1+A_0)^*B\), and \(D+D^*=-B^*B\).
The last equation has \(q^{u^2}\) solutions by the surjectivity
of the trace on diagonal coordinates. The unitary group of the
residue algebra in rank \(t\) has order \(q^{t^2}P_t\).
This follows by choosing successive orthonormal vectors as in
Proposition~\ref{prop:quot215}, or successive independent
columns for the split general linear group. Thus there are
\(q^{a^2}P_t\) residue solutions. At every solution the
linearized equation is onto the Hermitian matrices, because
it is \(Y^*(1+HX)+(1+HX)^*Y\), where \(1+HX\) is
invertible modulo \(\pi\) by the block description. Every
solution consequently has \(q^{a^2}\) lifts at each further
digit, and completeness gives an integral solution in every
residue class so obtained.

To fix the Haar normalization, at the identity use the
coordinates \(X\) on the tangent space of skew-Hermitian
matrices. Give its integral lattice additive volume one and
translate that density by left multiplication. In this
measure the solutions with \(X\equiv0\bmod\pi^j\) have volume
\(q^{-ja^2}\) for sufficiently large \(j\). Indeed
\(Y\mapsto Y+Y^*\) has an integral linear section (use the
trace on the diagonal). In the complementary coordinates the
implicit function theorem expresses the equation near zero
as a graph over the integral skew-Hermitian coordinates, with
identity differential on the latter. On sufficiently small
balls the absolute Jacobian of the invariant density is one.
Left translation by \(1+X_0H\) acts by
\(X\mapsto X_0+(1+X_0H)X\), preserving these congruences.
The digit counts therefore give volume \(P_t\) to \(K_L^\#\)
in this density.

Under a basis change \(S\), the new Gram matrix is \(S^*HS\)
and the tangent coordinate becomes \(S^{-1}X(S^*)^{-1}\).
The absolute Jacobian on skew-Hermitian matrices is
\(|N_{A/F_v}\det S|_v^{-a}\). Integral changes preserve the
integral tangent lattice. For a diagonal change with entries
\(s_i\), the \(i\)-th diagonal coordinate contributes
\(|N s_i|_v^{-1}\) and each independent off-diagonal pair
contributes \(|N s_i\,N s_j|_v^{-1}\). Elementary divisors
give the formula for general \(S\), also in the split algebra.
Multiplication of the density by \(|\det H|_v^a/P_a\)
makes it independent of the basis. This proves
\eqref{eq:discvol221}.

For surjectivity discard the unimodular orthogonal summand,
whose discriminant is zero. In a basis of the dual of the
remaining summand the Gram matrix is \(D_+^{-1}\) and the
lattice is \(D_+O^u\), where \(D_+\) is the diagonal matrix
of the positive powers of \(\pi\). Lift the images of the
basis under a given discriminant isometry to columns \(S\)
in \(O^u\). The reduction of \(S\) is invertible because
\(D_+O^u\subset\pi O^u\). The lift preserves \(D_+O^u\)
because the relations downstairs give inclusion and its unit
determinant gives equality. Pairing preservation gives
\(S^*D_+^{-1}S=D_+^{-1}+W\) with \(W\) integral and
Hermitian. If \(W\in\pi^j M_u(O)\), multiply \(S\) on the
right by \(1+\pi^jD_+C\). This fixes its action on the
quotient, and modulo \(\pi^{j+1}\) changes the error by
\(\pi^j(C+C^*)\), since \(D_+\equiv0\bmod\pi\). Solve by
the residue trace and continue starting at \(j=0\). The
convergent corrections give an exact isometry, which extends
by the identity on the discarded summand.
\end{proof}

\begin{theorem}[Genus mass with its discriminant marked]
\label{thm:discmass221}
For every totally positive definite integral Hermitian lattice
\(M\) of rank \(a\ge1\), with \(u_v\) positive elementary
divisors at \(v\),
\begin{equation}
 \operatorname{mass}(M)\,|\Aut\mathcal D(M)|
   =\mathfrak m_a\,\delta(M)^a
            \prod_{v:u_v>0}\frac{P_a(q_v,\eta_v)}
                                   {P_{a-u_v}(q_v,\eta_v)} .
 \label{eq:markedmass221}
\end{equation}
\end{theorem}
\begin{proof}
First choose a maximal integral lattice locally in the
rational space of \(M\). At a split place or at a nonsplit
place of even determinant valuation the space admits a
unimodular lattice. At a nonsplit place of odd determinant
valuation it admits the matrix
\(\operatorname{diag}(I_{a-1},\pi)\). These assertions use
the classification by dimension and determinant modulo norms
over a nonarchimedean local field. In the unramified case
units are norms. A lattice with the latter matrix is maximal
because an integral over-lattice would give an isotropic
submodule in its discriminant, an anisotropic line over the
quadratic residue field. For our unramified extension the
condition that all norms are integral is equivalent to
integrality of the pairing. Indeed, polarize after multiplying one
argument by each integer of \(E_v\) and use the unimodular
trace pairing. The local maximal lattices can be chosen equal
to \(M_v\) almost everywhere and patched by intersection in
the rational space. For such patching, clear all denominators
in one fixed global lattice. The prescribed changes are then
submodules of a finite quotient, which is the direct sum of
its primary parts. Taking their inverse image gives exactly
the prescribed completions.

In the measure of Lemma~\ref{lem:discstab221} a unimodular
stabilizer has volume one. For the odd-valuation maximal
lattice its finite discriminant has \(q+1\) isometries, the
norm-one residue scalars. Its full stabilizer therefore has
volume \((q+1)/(q^a-(-1)^a)\). The maximal definite
Hermitian mass formula of Gan, Hanke and Yu gives
\(\mathfrak m_a\) times the product of the reciprocals of
precisely these volumes. Its hypotheses hold since \(F\)
is totally real, \(E/F\) is CM, and the lattice is maximal
for the integral norm condition. The common archimedean and
\(L\)-value factors are those of Lemma~\ref{lem:mass209}.
See \cite{XXI-XVI-IV-ghy209}, Propositions 4.4 and 4.5.

For two compact open lattice levels in the same rational
unitary group pass to their intersection. A class with finite
stabilizer \(J\) acquires a finite set of level structures of
cardinality the index. Their \(J\)-orbits contribute total
weight equal to that index divided by \(|J|\), by the orbit
formula. The masses at the two levels are thus inversely
proportional to their compact volumes. Apply this first to
the maximal level, then to the kernel of the discriminant
action for \(M\). The finite paired module is the direct sum
of its local parts with orthogonality across distinct
places of \(F\), so its isometry group is the product of the
local isometry groups. Surjectivity in
Lemma~\ref{lem:discstab221} gives the index of the new level
in the full stabilizer. Substitution of
\eqref{eq:discvol221} now gives \eqref{eq:markedmass221}.
\end{proof}

Write \(N_j(L,D)\) for the number of primitive rank-\(j\)
sublattices \(M=L\cap EM\) with \(\delta(M)=D\).
Here \(EM\) denotes the rational span of \(M\). Write
\(\mathbb E_{\rm sd}\) for expectation in the self-dual
positive definite genus in rank \(k\), normalized by
\(\mathfrak m_k\). These are the primitive sublattices counted
in \srcloc{Formalism/SharpRestriction/CuspSpanNormBounds/CntPartners}{79}.

\begin{theorem}[The self-dual genus mean]
\label{thm:sdmean221}
Let \(j+s=k\) with \(j,s\ge1\), and \(D\ge1\) an integer.
At \(v\) choose a partition \(\lambda_v\) of \(e_v\ge0\)
into \(u_v\le\min(j,s)\) positive parts, allowing the empty
partition. Sum over choices with \(\prod_v q_v^{e_v}=D\)
and \(\sum_{\eta_v=-1} e_v\) even. Write \(\mathcal D_v\)
for the paired module of that partition and put
\(R_v(a)=P_a/P_{a-u_v}\). Then
\begin{equation}
 \mathbb E_{\rm sd} N_j(L,D)
 =\frac{\mathfrak m_j\mathfrak m_s}{\mathfrak m_k}D^k
      \sum_{(\lambda_v)}
      \prod_{v:e_v>0}\frac{R_v(j)R_v(s)}{|\Aut\mathcal D_v|}.
 \label{eq:sdmean221}
\end{equation}
In particular the mean is at most
\((\mathfrak m_j\mathfrak m_s/\mathfrak m_k)D^{k+\nu}\),
where \(\nu=4+\log_2(2f)\).
\end{theorem}
\begin{proof}
For a primitive \(M\subset L\), let \(N=L\cap(EM)^\perp\).
The projection argument of Proposition~\ref{prop:index219}
identifies \(L/(M\oplus N)\) with both discriminant modules.
In their direct sum it is the graph of an anti-isometry, since
integrality of \(L\) says that the two linking pairings on
the graph sum to zero. Conversely the inverse image of the
graph of any such anti-isometry in \(M^\vee\oplus N^\vee\)
is integral and self-dual, since the graph is its own
annihilator by perfectness. Its intersections with the two
rational summands are \(M\) and \(N\). These constructions
are inverse up to isometry with \(M\) distinguished.

The two local partitions must agree by the elementary divisor
description. Conversely this condition supplies local
anti-isometries, since the negative of each form is
isometric to it by a unit scalar of norm \(-1\). Gluing the
local anti-isometries on the finite module supplies a global
one. For a proposed list of partitions, the rank-\(j\) and
rank-\(s\) local rational spaces each have determinant norm
symbol \((-1)^{e_v}\) at nonsplit places and \(+1\) at all
other finite places. Take positive forms at infinity, with
symbol \(+1\). Landherr's local-to-global theorem gives
existence precisely under the displayed parity condition
and uniqueness of each rational space then. The local
lattice classes give a single genus in each space by
patching as above. Every glued self-dual \(k\)-lattice lies
in the standard genus by the same classification and the
uniqueness of the local unimodular lattice. For the
local and global classification see
\href{https://people.math.harvard.edu/~gross/preprints/incoherent.pdf}
{Gross, \emph{Incoherent definite spaces and Shimura varieties},
Section 3 and Theorem 3.1}.

For fixed isometry classes \(M,N\) the number of graph
identifications is \(|\Aut\mathcal D(M)|\).
The stabilizer of a graph under \(\Aut M\times\Aut N\)
is exactly the stabilizer of the distinguished sublattice in
the glued lattice, by restriction to the two rational
summands. The sum of the reciprocal orders of the graph
stabilizers over orbits is therefore
\(|\Aut\mathcal D(M)|/(|\Aut M||\Aut N|)\).
This also holds for \(j=s\) since \(M\) is distinguished.
Alternatively summing over the orbits of distinguished
sublattices in each \(L\) gives \(N_j(L,D)/|\Aut L|\).
Sum over the two genera, use \eqref{eq:markedmass221}
twice and divide by \(\mathfrak m_k\). This gives
\eqref{eq:sdmean221}, including finiteness.

Each \(R_v(a)\le\exp(\sum_{i\ge1}q_v^{-i})\le\exp(1)\).
There are at most \(\log_2 D\) nonempty places. Above a
rational prime \(p\) with \(p^b\Vert D\), encode the ordered
local partitions by concatenating their parts in a fixed
order of places and multiplying each part by its residue
degree. This is a composition of \(b\) with each summand
labelled by one of at most \(f\) places. It recovers the
partitions and gives at most \((2f)^b\) possibilities.
Multiplication over \(p\), followed by
\(\sum_p v_p(D)\le\log_2 D\) and
\(2/\log 2<4\), gives the asserted upper bound. The empty
collection gives it at \(D=1\).
\end{proof}

\subsection{Primitive subspaces in the marked terminal column}\label{nav:IX-12}
\label{sec:termmean221}

Use the same field, prime \(T\) and original terminal flag as
in Theorem~\ref{thm:midlevel131}. Thus \(T\) splits completely
in \(E\). If \(X\) is the isotropic part of this flag and \(z\)
its terminal vector, the positive column is the rank-\(k\)
Hermitian space \(P=(X+Ez)^\perp/X\). In a compatible
row lattice \(\Lambda\) its integral lattice \(L\) is the
image in \(P\) of \(\Lambda\cap(X+Ez)^\perp\). Indeed an
integral lift from \(X^\perp\) of a class perpendicular to
\(z\) is already perpendicular to \(z\), since \(X\perp z\).
We average over the genus of these column lattices with
reciprocal automorphism weights. This is the measure on
initial columns before a second rational flag is specified.

\begin{lemma}[Recovering the self-dual parent]
\label{lem:termparent221}
Fix, for each \(v\mid T\), the split factor on which the
chosen \(\varpi\) has valuation one. A column in the terminal
genus is equivalent to a self-dual lattice \(L_*\) of rank
\(k\) together with one hyperplane in \(L_{*,v}^+/\pi L_{*,v}^+\)
at each such \(v\). Its genus mass equals
\[
           \mathfrak m_k\left(\frac{T^k-1}{T-1}\right)^f .
\]
Its normalized average can be obtained by sampling \(L_*\)
with its normalized genus mass and then sampling the
hyperplanes independently and uniformly.
\end{lemma}
\begin{proof}
Away from \(T\), split off integral hyperbolic partners to
the primitive isotropic columns as in
Lemma~\ref{lem:flag215}. In their unimodular middle quotient
the terminal vector has unit norm, so its perpendicular
is self-dual. At \(v\mid T\) write \(O=\ZZ_T\) in the chosen
factor. The middle lattice in the original integral
frame has first factor
\[
          \{(w,z_0)\in O^k\oplus O:
                      z_0\equiv \ell_v(w)\pmod T\},
\]
and second factor its dual. The terminal perpendicular has
\(z_0=0\) in each factor. Its first lattice is the inverse
image of \(\ker\bar\ell_v\) in \(O^k\). Its second is
\((O^k)^\vee\), since projection of the middle lattice on
\(w\) is all of \(O^k\). Transporting by isometry gives this
description for every member of the genus. The self-dual
parent is intrinsic. At each \(v\mid T\) replace the first
lattice by the dual of the second, leaving all other local
lattices as given. Patching produces \(L_*\). Conversely
every hyperplane gives the prescribed local genus because
\(\GL_k(O)\) acts transitively on hyperplanes modulo \(T\).
These constructions commute with isometries since the split
factors have been fixed.

For fixed \(L_*\), its finite isometry group acts on the
product of hyperplane sets. Its orbits give the isometry
classes over \(L_*\), and the stabilizers are their isometry
groups. The mass of each orbit equals its size divided by
\(|\Aut L_*|\). There are \((T^k-1)/(T-1)\) hyperplanes
at each of the \(f\) places. Summing proves both assertions.
Retaining a nonzero covector above each hyperplane multiplies
the mass by \((T-1)^f\). Fixed compact level refinements at
the other primes multiply it by their indices, by the same
orbit argument. Normalized averages of functions of \(L\)
are consequently the same at these finer column levels.
\end{proof}

\begin{theorem}[The primitive count at the actual terminal prime]
\label{thm:termmean221}
Let \(\mathbb E_T\) be the normalized terminal genus mean.
For \(j+s=k,\ j,s\ge1\), put \(p=(T^s-1)/(T^k-1)\). Then
\begin{align}
  \mathbb E_T N_j(L,D)
   &=\sum_{e=0}^f\binom fe p^{f-e}(1-p)^e\,
                      \mathbb E_{\rm sd}N_j(L_*,D/T^e),
                   \label{eq:termmean221}\\
  \mathbb E_T N_j(L,D)
   &\le \frac{\mathfrak m_j\mathfrak m_s}{\mathfrak m_k}
          (D/T^f)^k (T^f)^s D^\nu . \label{eq:termupper221}
\end{align}
Here \(D\ge1\) is integral and the count is read as zero at
any argument which is not a positive integer.
\end{theorem}
\begin{proof}
A rational \(j\)-plane determines exactly one primitive
sublattice in \(L_*\) and one in \(L\), by intersection.
If \(A\subset L_*\) is the first, the second is \(A\cap L\).
At \(v\mid T\) its index is \(1\) when the hyperplane
contains the reduction of \(A_v^+\), and \(T\) otherwise.
Indeed \(A_v^+\) is a direct summand of the parent in this
factor and the second factor is unchanged. A uniform
nonzero residue covector vanishes on the resulting
\(j\)-space with probability \(p\), by counting covectors
on the quotient of dimension \(s\). The determinant in the
convention of Section~\ref{sec:discmass221} changes by the
index, so \(\delta(A\cap L)=T^e\delta(A)\) when precisely
\(e\) hyperplanes fail to contain the reduction. Independence
over the places gives the binomial factor for every \(A\).
Summation over \(A\) gives \eqref{eq:termmean221}.

At a nonzero summand use Theorem~\ref{thm:sdmean221}, bound
\((D/T^e)^\nu\le D^\nu\) and sum the factor \(T^{-ke}\).
The binomial average is
\((p+(1-p)T^{-k})^f=T^{-jf}\). This proves
\eqref{eq:termupper221}.
\end{proof}

\begin{proposition}[Size of the genus factor in the middle ranks]
\label{prop:massdecay221}
Put \(R=|\operatorname{disc}E|^{1/(2f)}
       =|\operatorname{disc}F|^{1/f}\).
There is an absolute \(C\) such that for \(2\le s\le k-2\),
\begin{equation}
        \frac{\mathfrak m_{k-s}\mathfrak m_s}{\mathfrak m_k}
          \le \mathfrak m_2
                   \left(\frac{C}{kR}\right)^{s(k-s)f}.
                   \label{eq:massdecay221}
\end{equation}
On the common intervals \(\mathfrak m_2\le b^{C f}\) for
all sufficiently large \(x\), where \(b=\exp(x^\eps)\).
\end{proposition}
\begin{proof}
Write \(D_F=|\operatorname{disc}F|\). The discriminant tower
formula gives \(D_E=D_F^2\) by finite unramifiedness.
For \(n\ge2\), put \(a_n=\mathfrak m_n/\mathfrak m_{n-1}\)
and let \(\eta\) denote the quadratic character of \(E/F\).
The mass formula in Lemma~\ref{lem:mass209} gives
\[
 a_n=D_F^{\,n-1/2}
       \left(\frac{(n-1)!}{(2\pi)^n}\right)^f
                       L_F(n,\eta^n).
\]
Indeed \(a_n=2^{-f}|L_F(1-n,\eta^n)|\). The latter \(L\)-value
is taken on the \(f\) even Dirichlet factors of \(F\) for
even \(n\), and on the \(f\) odd factors of
\(\zeta_E/\zeta_F\) for odd \(n\). For a primitive character
\(\chi\) of the parity of \(n\), its functional equation and
the gamma reflection and duplication formulas give
\[
 |L(1-n,\chi)|=
  \frac{2(n-1)!}{(2\pi)^n}
       \operatorname{cond}(\chi)^{n-1/2}|L(n,\bar\chi)|.
\]
By the abelian conductor--discriminant formula the conductors
in the even product multiply to \(D_F\), and those in the
odd product to \(D_E/D_F=D_F\). The factors are closed under
conjugation and their Euler product at \(n\) is positive.
It lies between \(\zeta(n)^{-f}\) and \(\zeta(n)^f\), since
each local factor is between \(1-q_v^{-n}\) and
\((1-q_v^{-n})^{-1}\), and \(\zeta_F(n)\le\zeta(n)^f\).

Set \(\ell=\min(s,k-s)\) and \(h=k-\ell\ge k/2\).
In the ratio on the left of \eqref{eq:massdecay221}, after
retaining \(\mathfrak m_2\), the factors left are
\[
       \frac1{a_{h+1}a_{h+2}}
                  \prod_{i=3}^{\ell}\frac{a_i}{a_{h+i}} .
\]
The factorial quotient \((h+i-1)!/(i-1)!\ge h!\),
and \(h!\ge(h/\exp(1))^h\). Thus each displayed quotient
\(a_i/a_{h+i}\) is at most \((C_0/(hR))^{hf}\) for an
absolute \(C_0\), by the Euler product bounds. Each of
\(1/a_{h+1},1/a_{h+2}\) satisfies the same bound by the
formula for \(a_n\), using \(R\ge1\). There are \(\ell\)
such factors and \(h\ell=s(k-s)\). This proves the assertion.

Finally \(\mathfrak m_1=2^{1-f}|L_F(0,\eta)|\).
For an odd primitive Dirichlet character of conductor \(c\),
the Bernoulli formula at zero gives \(|L(0,\chi)|\le c\).
Thus \(\mathfrak m_1\le 2^{1-f}D_F\), and the formula for
\(a_2\) gives \(\mathfrak m_2\le C_1^f D_F^{5/2}\).
All conductors divide \(N\) with \(\log(2N)=o(\log b)\)
by Section~\ref{sec:terminal209}, and \(D_F\le N^f\).
This proves the asserted interval bound.
\end{proof}

In particular \eqref{eq:termupper221} has the precise genus
ratio \(\Xi_s\) used in the middle ranks of the source. The
factor \((T^f)^sD^\nu\) is bounded by
\((D T^f)^{C(s+\log f)}\). These comparisons use the same
\(T\), \(E\), and interval as in the period. The genus ratio
and the primitive mean under comparison are specified in
\srcloc{Formalism/SharpRestriction/CuspSpanNormBounds/CntPartners}{127}
and
\srcloc{Formalism/SharpRestriction/CuspSpanNormBounds/CongruenceAtSplitFrobeniusHolds}{533}.

\begin{proposition}[Two reductions at one terminal place]
\label{prop:joint221}
Fix two primitive sublattices in the same self-dual parent,
and at a place above \(T\) let their first-factor reductions
have dimensions \(a,b\ge1\) and span of dimension \(c\).
Write \(\epsilon_A,\epsilon_B\) for the two indicators of
failing the hyperplane condition. For \(z=T^{-k}\),
\[
 \mathbb E\,z^{\epsilon_A+\epsilon_B}
   =(1-z)T^{-c}+z(T^{-a}+T^{-b}-1)\le T^{-c}.
\]
For fixed rational subspaces in a fixed parent the joint
moment over all \(v\mid T\) is the product of these expressions.
\end{proposition}
\begin{proof}
The probability that a uniform nonzero covector kills a
\(d_0\)-space is \(p_{d_0}=(T^{k-d_0}-1)/(T^k-1)\).
It kills both reductions exactly when it kills their span.
Since \(z^{\epsilon_A}=z+(1-z)1_{\{\epsilon_A=0\}}\), the
expectation is \(z^2+z(1-z)(p_a+p_b)+(1-z)^2p_c\).
Substitute \((1-z)p_{d_0}=T^{-d_0}-z\). The bound follows
from \(T^{-a}+T^{-b}\le1\), and the product assertion from
Lemma~\ref{lem:termparent221}.
\end{proof}

\begin{corollary}[Distinct subspaces occurring in flag comparisons]
\label{cor:subspaces221}
Fix an initial flag \(\mathcal F\) in the row and the
determinant cutoff \(D_{\max}>0\) for comparing two terminal
flags. For \(1\le j=k-s\le k-1\) and \(b_0\ge1\), consider
their distinct rational comparison subspaces of dimension
\(j\) in \(P\) for which the product of the two compatibility
indices is at least \(b_0\). Their number, averaged in the
initial column genus (or its indicated refinements), is at
most
\[
    \frac{\mathfrak m_j\mathfrak m_s}{\mathfrak m_k}\,
       (D_1/T^f)^k (T^f)^s D_1^{\nu+1},
          \qquad D_1=D_{\max}/b_0 ,
\]
and is zero if \(D_1<1\).
\end{corollary}
\begin{proof}
Write \(W_{\mathcal F}=X+Ez\) for the original frame span,
and \(W_i=X_i+Ez_i\) for the two terminal frame spans.
In the nondegenerate pair range put
\(C_i=(W_{\mathcal F}+W_i)^\perp\). This space embeds
isometrically in \(P\) by reduction modulo \(X\), since
its intersection with the nondegenerate frame sum is zero.
It embeds likewise in \(P_i=W_i^\perp/X_i\). Denote the
quotient maps by \(p,p_i\) on their respective domains.
A comparison isometry \(\beta\) carrying the first terminal
flag to the second induces \(\bar\beta:P_1\to P_2\). Its
comparison subspace \(D_\beta\subset P\) is the set of
\(p(c_1)=p(c_2)\) for which \(c_i\in C_i\) and
\(\bar\beta p_1(c_1)=p_2(c_2)\). This agrees with the two
linear congruences modulo \(X\) and \(X_2\) in the
definition of the rational comparison. Both \(c_i\) are
uniquely determined by their common image.

Let \(\mathcal L_i\) be the compatibility lattice of
Proposition~\ref{prop:index219} for the \(i\)-th pair and
\(i_i=[\mathcal L_i:\Lambda\cap C_i]\).
The comparison lattice
\(M_\beta=\{c_1\in\mathcal L_1:p(c_1)\in D_\beta\}\)
has full rank in the corresponding subspace of \(C_1\),
since \(\mathcal L_1\) contains \(\Lambda\cap C_1\).
Its isometric image belongs to \(L\cap D_\beta\), because
each \(c_1\in\mathcal L_1\) has an integral lift modulo
\(X\) perpendicular to the initial frame. The cutoff used
for comparing the flags is
\(\delta(M_\beta)i_1i_2\le D_{\max}\). Thus the primitive
saturation \(L\cap D_\beta\) has determinant at most \(D_1\),
by the index change of determinant. This is precisely the
lattice and cutoff in
\srcloc{Formalism/SharpRestriction/CuspSpanNormBounds/IsPrimitiveSub}{213}
and \srcloc{Formalism/SharpRestriction/CuspSpanNormBounds/CntPartners}{190}.

Distinct rational subspaces have distinct primitive
saturations in \(L\). In each column the count is therefore
at most \(\sum_{1\le D\le D_1}N_j(L,D)\), regardless of the
choice of compatible row above that column. Average and apply
\eqref{eq:termupper221} at the at most \(D_1\) positive
integers in this range. This also proves the zero assertion.
\end{proof}

\begin{theorem}[Uniqueness of the comparison for two fixed flags]
\label{thm:unique221}
Fix three flags in one original row \(\Lambda\), the initial
one and two admissible terminal flags of the same terminal
norm. At \(Q=g_1g_2\) suppose each terminal flag has the same
reduced isotropic submodule as the initial flag. This holds
in every compatible row of Section~\ref{sec:saturation219}.
Among integral row isometries carrying the first terminal
flag and its terminal vector to the second, the comparison
congruence is equivalent to inducing the identity on the common
residual middle quotient at each place above \(Q\). Among the
isometries satisfying it there is at most one induced
isometry \(X_1^\perp/X_1\to X_2^\perp/X_2\). Consequently a
fixed ordered pair of flags contributes at most one comparison
class in the equivalence relation used for the pair count.
\end{theorem}
\begin{proof}
Write \(X_0=X\) and, for \(a=0,1,2\), put
\(A_a=\Lambda\cap X_a\), \(B_a=\Lambda\cap X_a^\perp\)
and \(J_a=B_a/A_a\). The latter is an integral lattice in
the middle quotient. At each \(v\mid Q\) the three reductions
of \(A_a\) in the residue row are the same submodule \(U\).
Indeed admissibility gives equality of the reduced frames up
to an invertible upper triangular change. The reductions of
the \(B_a\) are all \(U^\perp\). To see this, restriction of
the row pairing from \(\Lambda_v\) to the \(O\)-dual of
\(A_{a,v}\) is onto by self-duality and primitivity. Its
kernel is \(B_{a,v}\). This exact sequence splits as a
sequence of modules, and hence remains exact on reduction.
The same reasoning in the two factors at a split place uses
the perfect pairing between the factors. In particular
the reductions of the \(J_a\) identify canonically with
\(U^\perp/U\) in the residue row.

For \(x_1\in B_1\), choose \(l\in B_0\) with
\(l\equiv x_1\bmod Q\Lambda\). Such an \(l\) exists since
the reductions of \(B_0\) and \(B_1\) agree at the places
above \(Q\), the ideal \(Q\OO_E\) is squarefree, and the
Chinese remainder theorem for the module \(B_0\) realizes
all its prescribed residues simultaneously. For an admissible
comparison \(\beta\) the congruence says that for every
\(l\in B_0,x_1\in B_1\),
\[
 l-x_1\in Q\Lambda+A_0+A_1
       \quad\Longleftrightarrow\quad
 l-\beta x_1\in Q\Lambda+A_0+A_2 .
\]
Our choice makes the left side true. In each residue row
the right side implies \(\beta x_1\equiv x_1\bmod U\).
It follows that every comparison induces the same map on
\(J_{1,v}/\pi J_{1,v}\to J_{2,v}/\pi J_{2,v}\) at each marked place, namely
the identity under the common identification. Here \(\pi\)
denotes the local uniformizer. The global elements of \(B_1\)
surject on its residue module. Conversely the two membership
conditions can be checked in the residue rows at \(v\mid Q\),
since each subgroup contains \(Q\Lambda\). There each tests
equality of the two images in \(U^\perp/U\). Identity on
the residue quotients therefore implies the equivalence for
every \(l,x_1\). This is the congruence of
\srcloc{Formalism/SharpRestriction/CuspSpanNormBounds/IsPrimitiveSub}{206}.

If two comparisons exist, their quotient acts on \(J_1\) as
an integral isometry fixing the class of \(z_1\). Its
rational action decomposes as the identity on the nonisotropic
line \(Ez_1\) and a unitary action on its totally positive
definite perpendicular \(P_1\). It has finite order on \(J_1\).
Indeed it preserves the positive lattice \(J_1\cap P_1\),
whose isometry group is finite. In all archimedean embeddings
the group is compact and the images of any fixed integral
basis of the underlying \(\ZZ\)-lattice form a discrete set.
Its action on the two rational summands determines the action
on \(J_1\).

At one marked place take a single factor over the odd
completely split prime \(p=g_1\). In a basis of the completion
of \(J_1\) the quotient has matrix \(1+H\) over \(\ZZ_p\)
with \(H\equiv0\bmod p\). This congruence kernel is
torsion-free. If the least valuation of an entry of
\(H\ne0\) is \(t\ge1\), the binomial expansion shows that
raising to a power prime to \(p\) preserves \(t\). Raising
to the \(p\)-th power changes it to \(t+1\), since \(pH\)
is the unique term of this order and every entry of the
remaining terms is divisible by \(p^{t+2}\) for odd \(p\).
Thus a finite-order element in the kernel is the identity.
The embedding of \(E\) into this completion is injective,
so the two rational maps coincide on the entire quotient.
For fixed terminal flags equality there gives the prescribed
equivalence by taking the two unipotent changes of frame to
be the identity in
\srcloc{Formalism/SharpRestriction/CuspSpanNormBounds/CntPartners}{186}.
\end{proof}

The comparison subspace is thus determined by the two flags whenever a comparison exists. To apply Corollary~\ref{cor:subspaces221} we need its determinant cutoff in terms of their Gram data and the multiplicity with which different flag pairs produce it. These are computed in Theorems~\ref{thm:cutoff223} and~\ref{thm:pairmean225}.

\subsection{The determinant cutoff for two terminal flags}\label{nav:IX-13}
\label{sec:cutoff223}

Fix the row, the initial flag \(0\), and two terminal flags
\(1,2\) of Section~\ref{sec:termmean221}. Both terminal norms
are \(-T'\), and the initial norm is \(-T\). In the common-prime
application \(T'=T\). Assume each of the two frame sums is
nondegenerate, the reduced frames agree up to the allowed
changes at \(Q\), and an admissible comparison \(\beta\) exists.
Let \(W_a=X_a+Ez_a\) be the frame span of flag \(a\),
\(P_a=W_a^\perp/X_a\), and \(L_a\) its integral column.
Write \(S_i=W_0\oplus W_i\) and \(C_i=S_i^\perp\) for \(i=1,2\).
Thus \(C_i\) is positive definite of rank \(k-1\).

Put \(P=P_0\), \(P'=P_2\). For \(c\in C_i\) let \(p_i(c)\)
be its image in \(P\), and \(q_i(c)\) its image in \(P'\),
using the induced comparison \(\bar\beta:P_1\to P_2\) for
\(i=1\). Both embeddings are isometric. Denote their joint
image in \(P\oplus P'\) by \(\Gamma_i\). It is the graph of
an isometry between two hyperplanes in definite spaces.
Their intersection is the graph of an isometry
\(D\to D'\), where \(D\subset P\) is the comparison subspace
of Corollary~\ref{cor:subspaces221}. Write \(j=\dim_E D=k-s\).
We prove the cutoff for every \(2\le s\le k\), allowing
\(D=0\). For each of the two initial--terminal pairs put
\[
     G_{0,i}=\frac{|N_{F/\QQ}\mathcal G_i|}{Q^{2df}},\qquad
     I_i=[B_i:B_i^+],\qquad i_i=[\mathcal L_i:\Lambda\cap C_i]
\]
as in Proposition~\ref{prop:index219}, with \(B_i=\Lambda\cap S_i\)
and \(U_i=(\Lambda\cap W_0)\oplus(\Lambda\cap W_i)\).
The notation \(U_i\) here denotes the sum for the pair.
We have \(i_i\mid I_i\) and
\(\delta(\mathcal L_i)i_i^2=G_{0,i} i_i/I_i\). Here
\(\mathcal G_i\) is the determinant of the Gram matrix on
the concatenated frames up to sign.

\begin{lemma}[The common lattice and the difference of the pairings]
\label{lem:graphs223}
Under either graph map, the compatibility lattice
\(\mathcal L_i\) is the intersection of \(\Gamma_i\) with
the same lattice in \(P\oplus P'\). In particular the two
lattices intersect the common graph in the same lattice
\(M\), measured with the metric of its first projection.
Write \(H_i\) for the orthogonal projection of
\(p_i(\mathcal L_i)\) on \(D^\perp\) in \(P\). Then
\(\operatorname{rank}H_i=s-1\) and
\[
                       \delta(\mathcal L_i)=\delta(M)\delta(H_i).
\]
The pairing
\[
       B(c_1,c_2)=
         \langle p_1(c_1),p_2(c_2)\rangle
            -\langle q_1(c_1),q_2(c_2)\rangle
\]
descends to the rational spans of \(H_1,H_2\), takes values
in \(Q\OO_E\) on these lattices, and has operator norm at
most two at every archimedean place with their projected metrics.
\end{lemma}
\begin{proof}
Use the lattice of pairs \((l,l')\in L_0\oplus L_2\) with
equal images at every \(v\mid Q\) in the residue middle
quotient \(U^\perp/U\) of Theorem~\ref{thm:unique221}.
A class from \(C_i\) has integral projections exactly when
it admits integral lifts in \(X_0^\perp,X_i^\perp\), respectively,
using \(\beta^{-1}\) in the second coordinate for \(i=1\).
They differ from \(c\) by vectors in \(X_0,X_i\). Their
difference lies in \(Q\Lambda+(\Lambda\cap X_0)+(\Lambda\cap X_i)\)
exactly when their residue classes agree. This follows by
reducing modulo \(Q\Lambda\) and the Chinese remainder
theorem, since both isotropic lattices reduce to \(U\).
The comparison identifies the residue middle quotients by
the identity. These conditions define \(\mathcal L_i\).
Intersecting the resulting lattice of pairs with the common
graph proves the assertion about \(M\). In particular \(M\)
is the lattice \(M_\beta\) of Corollary~\ref{cor:subspaces221},
under the first projection.

Here and below a possibly fractional lattice \(A\) in a
positive space has determinant defined as follows. In a
rational basis \(v_1,\ldots,v_a\), write its top exterior
power as \(\mathfrak a(v_1\wedge\cdots\wedge v_a)\).
Then \(\delta(A)=N\mathfrak a\prod_\tau\det
(\langle v_i,v_j\rangle_\tau)\), with one \(\tau\) above
each real place, positive absolute ideal norm, and empty
determinant one. Change of basis cancels between these
factors. For integral lattices this is the discriminant
convention of Section~\ref{sec:discmass221}, as seen on a
local basis. For projection on an orthogonal complement,
the determinant is the product of those of the kernel
lattice and the projected image. Indeed choose a rational
basis of the kernel and lifts of a basis of the image.
The exterior ideals multiply by the exact lattice sequence,
which splits locally. Subtracting kernel components from
the lifts preserves the exterior vector and makes the Gram
matrix block diagonal. Applying this to \(p_i(\mathcal L_i)\)
gives the asserted identity and the rank is \(k-1-j=s-1\).

On \(D\) the two graph isometries agree, so pairing a common
graph vector with either graph gives zero for the difference
of forms. Removing the \(D\)-component on the source removes
its common isometric image in \(D'\) on the target. This
proves descent. Each of the resulting inner products is
bounded by the product of the projected norms, by the
Cauchy--Schwarz inequality. Integrality in \(Q\OO_E\)
follows because the two original column pairings are
integral and reduce to the same Hermitian pairing in
\(U^\perp/U\) at every \(v\mid Q\).
\end{proof}

\begin{lemma}[Rank and determinant of the reduced pairing]
\label{lem:detpair223}
The descended pairing has rank \(s-1\) or \(s-2\). In the
second case let \(R_i\) denote the intersection of \(H_i\)
with its radical line on its own side. In the first case
put \(R_i=0\). For \(n=\operatorname{rank}B\),
\[
  \frac{\delta(H_1)\delta(H_2)}{\delta(R_1)\delta(R_2)}
                               \ \ge\ (Q/2)^{2nf}.
\]
\end{lemma}
\begin{proof}
After removing the common graph, the resulting graphs are
disjoint totally isotropic \((s-1)\)-planes in
\(D^\perp\oplus(D')^\perp\) with the difference of the two
forms. Indeed the original graphs each split orthogonally
in their positive source coordinate as the common graph
plus their projected graph. At any real place the ambient
difference form has signature \((s,s)\). The radical on
the first graph is its intersection with the perpendicular
of the second. It injects into the perpendicular of the
second modulo the second, a space of signature \((1,1)\).
This signature follows by splitting off hyperbolic partners
for a basis of the second isotropic graph. An isotropic
subspace of a plane of signature \((1,1)\) has dimension
at most one. The ranks of a matrix and its adjoint agree,
so the two radical dimensions are equal. Ranks are
preserved on extending \(E\) to \(\CC\).

Project \(H_i\) orthogonally off the rational radical to
give a lattice \(H_i'\) of rank \(n\). Pairings are unchanged
on projection, so the descended pairing still has norm at
most two and takes values in \(Q\OO_E\). In rational bases
let \(B_0\) be its matrix and \(G_1,G_2\) the Gram matrices.
If \(\mathfrak a_i\) are the two exterior ideals, the
nonzero ideal
\((\det B_0)\overline{\mathfrak a_1}\mathfrak a_2\)
lies in \(Q^n\OO_E\) by taking determinants on local bases
(conjugate the opposite ideal for the opposite pairing
convention). Its norm is at least \(Q^{2nf}\). At \(\tau\),
in orthonormal bases all singular values are at most two,
so \(|\det B_0|_\tau^2\le 2^{2n}\det G_{1,\tau}\det G_{2,\tau}\).
Multiply over \(\tau\) and insert the two ideal norms. The
result is \(Q^{2nf}\le 2^{2nf}\delta(H_1')\delta(H_2')\).
Orthogonal projection gives \(\delta(H_i')=\delta(H_i)/\delta(R_i)\)
by Lemma~\ref{lem:graphs223}. This argument also covers \(n=0\).
\end{proof}

We give the additional estimate needed when the descended
pairing has a radical. Let \(D_{i,d-1}\), \(D_{i,d}\) denote
the leading minors of the cross pairing of the two decorated
frames \(0,i\), ending with the isotropic and full frames,
respectively. For each archimedean embedding define
\[
 k_{i,\tau}=
 \begin{cases}
  (1+c_{i,\tau})/\sqrt{c_{i,\tau}^2-1},&
            D_{i,d-1}\ne0,\quad
     c_{i,\tau}=\dfrac{|\tau(D_{i,d}/D_{i,d-1})|}{\sqrt{TT'}},\\[4pt]
  1,&D_{i,d-1}=0 .
 \end{cases}
 \qquad K_i=\prod_\tau k_{i,\tau}.
\]
These factors are invariant under integral or rational
upper unitriangular changes of frame.

\begin{proposition}[The extra radical line]
\label{prop:radical223}
If the descended pairing is singular, then for
\(\{i,a\}=\{1,2\}\),
\begin{equation}
                \delta(R_i)\ \ge\
                  \frac{T^f i_a}{I_a K_a^2}.
                \label{eq:radical223}
\end{equation}
\end{proposition}
\begin{proof}
Start with an element of the rational left or right radical
on side \(i\), lifted to its graph before removing \(D\).
View its two coordinates in \(P_0,P_a\), undoing
\(\bar\beta\) in the second coordinate when \(a=1\).
Choose rational row lifts \(l\in W_0^\perp,\ l'\in W_a^\perp\).
The radical condition says that their projections onto
\(C_a\) coincide, since their pairings against every
\(c\in C_a\) coincide. Thus \(l-l'\in S_a\). Modulo
\(X_0+X_a\) it has unique coefficients
\(\alpha z_0+\alpha' z_a\), by the direct decomposition
\(S_a=W_0\oplus W_a\). Changing either lift does not change
these coefficients. An element of the common graph has
lifts equal to a single vector of \(C_a\), so adding such
an element also leaves them fixed. This defines \(E\)-linear
functionals \(\alpha,\alpha'\) on the projected radical line.

First bound the denominators of \(\alpha\). Each element of
\(R_i\) is the projection of an integral compatibility graph
vector. The two lifts may then be chosen integral, and
their difference belongs to
\(Q\Lambda+(\Lambda\cap X_0)+(\Lambda\cap X_a)\), since the
graph coordinates have equal marked reductions. Hence
\(l-l'\) belongs to \(B_a\) and its class modulo \(U_a\)
belongs to \(Q(B_a/U_a)\). The \(z_0\)-coefficient gives a
map \(B_a/U_a\to E/\OO_E\). Indeed at each place a
primitive representative of the initial frame is a basis
of \(\Lambda_v\cap W_0\), with last vector equal to \(z_0\)
modulo \(X_0\), so \(U_a\) has integral \(z_0\)-coefficient.
The subgroup of order \(i_a\) injected in \(Q(B_a/U_a)\)
by Proposition~\ref{prop:index219} is represented by
differences in \(X_0+X_a\), so its coefficient is zero.
That proposition also gives \(|Q(B_a/U_a)|\le I_a\).
The inverse image of the coefficient image of \(Q(B_a/U_a)\)
in \(E\) is therefore a fractional ideal \(J\supset\OO_E\)
with \([J:\OO_E]\le I_a/i_a\), and \(\alpha(R_i)\subset J\).

We now compute the archimedean norm of this functional.
Remove the common \(C_a\)-component of the two lifts. Their
images \(y,y'\) in the respective positive normal lines to
\(p_a(C_a),q_a(C_a)\) have equal norm, say \(t\), because
the graph on side \(i\) is an isometry and the removed
component was the same on both sides. Both normal components
are unchanged on first removing the common \(D\)-component.
Thus \(t\le\|r\|_\tau\) for the resulting vector \(r\)
in the projected radical line. If \(t=0\), the two original
coordinates are the graph of a vector of \(C_a\). They are
also on graph \(i\), so their vector vanishes after removing
the common graph. Thus \(t>0\) for nonzero \(r\).

The normal components have representatives in
\(S_a\cap W_0^\perp\) and \(S_a\cap W_a^\perp\). Write their
difference as \(w_0+w_a\), with
\(w_0=x_0+\alpha z_0\), \(w_a=x_a+\alpha'z_a\),
\(x_0\in X_0,\ x_a\in X_a\). Comparing the norms in
\(y-w_0=y'+w_a\) gives
\[
       \|y\|^2-\|y'\|^2=T|\alpha|^2-T'|\alpha'|^2.
\]
Here the norms on the left are those of the positive quotient
lines, equal to the row norms of their representatives.
The identity holds for arbitrary normal components \(y,y'\).
By polarization the map from their direct sum to the two
coefficients preserves a nondegenerate form of signature
\((1,1)\) and is injective. In particular for the nonzero
equal-norm pair above, \(\alpha\ne0\), since otherwise the
identity forces both coefficients to vanish.

When \(D_{a,d-1}\ne0\), the pairing of \(X_0\) and \(X_a\)
is nonsingular. Modify \(z_0,z_a\) by \(X_0,X_a\) to make
both perpendicular to \(X_0+X_a\). Denote the modified
vectors by \(v_0,v_a\). Their norms are \(-T,-T'\), and
their cross entry \(b_a\) has magnitude
\(|\tau(D_{a,d}/D_{a,d-1})|\) by the determinant of the
block matrix after these row and column operations. In
their span the positive normal to \(v_0\) has a generator
of the form \(v_a+\lambda v_0\), with \(|\lambda|=|b_a|/T\)
and squared norm \(T'(c_{a,\tau}^2-1)\). This is strictly
positive since its class is a nonzero vector in the positive
quotient. The normal to \(v_a\) has a generator
\(v_0+\lambda'v_a\) of squared norm \(T(c_{a,\tau}^2-1)\).
Reading the coefficient of \(v_0\) for normal components of
length \(t\), then using the triangle inequality on their
difference, gives
\[
                 \sqrt T\,|\tau\alpha(r)|
                          \le k_{a,\tau}t .
\]
If \(D_{a,d-1}=0\), the pairing of \(X_0,X_a\) has defect
one. Adding the column of \(z_a\) must give full row rank
on \(X_0\), for a residual kernel would be in the radical
of \(S_a\). The same holds with the flags exchanged.
Choose \(x_a\in X_a\) perpendicular to \(X_0\) and scale
it so that \(\langle z_0,x_a\rangle=T\), possible since
this kernel line pairs nontrivially with \(z_0\). Then
\(z_0+x_a\) is a positive normal representative on the
initial side, of squared norm \(T\). Exchanging flags gives
a normal \(z_a+x_0\) of squared norm \(T'\). Their coefficient
pairs modulo \(X_0+X_a\) are \((1,0)\) and \((0,1)\).
Thus \(\sqrt T\,|\tau\alpha(r)|=t\), giving the displayed
bound with \(k_{a,\tau}=1\).

Write \(R_i=\mathfrak a r_0\) for a fractional ideal and a
nonzero rational vector. Since \(\mathfrak a\alpha(r_0)
\subset J\) and \(\alpha(r_0)\ne0\),
\[
     N\mathfrak a\,|N_{E/\QQ}\alpha(r_0)|
                          \ge NJ\ge i_a/I_a .
\]
Multiplying the squared archimedean inequalities and using
\(\delta(R_i)=N\mathfrak a\prod_\tau\|r_0\|_\tau^2\)
proves \eqref{eq:radical223}.
\end{proof}

\begin{theorem}[Determinant bound in terms of the two Gram matrices]
\label{thm:cutoff223}
Let \(M\) be the comparison lattice. Its determinant and the
two compatibility indices satisfy
\begin{equation}
 \delta(M)i_1i_2\ \le\ \sqrt{G_{0,1}G_{0,2}}
 \begin{cases}
  (2/Q)^{(s-1)f}\sqrt{i_1i_2/(I_1I_2)},&
                    \operatorname{rank} B=s-1,\\[2pt]
  (2/Q)^{(s-2)f}K_1K_2/T^f,&
                    \operatorname{rank} B=s-2 .
 \end{cases}
 \label{eq:cutoff223}
\end{equation}
In particular, for fixed Gram data of the two pairs one may
use the respective right side, replacing the square root of
the index ratio by \(1\), as \(D_{\max}\) in
Corollary~\ref{cor:subspaces221} for \(2\le s\le k-1\).
\end{theorem}
\begin{proof}
Multiply the two determinant identities of
Lemma~\ref{lem:graphs223}. After multiplication by
\(i_1^2i_2^2\) the numerator is
\(G_{0,1}G_{0,2} i_1i_2/(I_1I_2)\). For the denominator
\(\delta(H_1)\delta(H_2)\), apply Lemma~\ref{lem:detpair223}.
When \(n=s-1\) the radical determinants are one. When
\(n=s-2\) their product is at least
\(T^{2f}i_1i_2/(I_1I_2K_1^2K_2^2)\) by
Proposition~\ref{prop:radical223}. Taking square roots gives
both bounds. All the quantities used the same \(\Lambda\)
and original finite levels. The determinant alternatives
being compared are stated in
\srcloc{Formalism/SharpRestriction/CuspSpanNormBounds/CongruenceAtSplitFrobeniusHolds}{88}.
\end{proof}

We next count the ordered pairs above each subspace satisfying \eqref{eq:cutoff223}. An integral Witt decomposition makes it possible to do this with the unipotent equivalence of the original row.

\subsection{Counting flags above one sublattice of the initial column}\label{nav:IX-14}
\label{sec:frcount225}

Fix the initial integral decorated flag in a row lattice
\(\Lambda\) at the original finite levels. Write \(t=d-1\)
for its isotropic rank, \(X\) for its isotropic space, \(z\)
for its terminal vector of norm \(-T\), and \(n_T=T^f\).
A terminal flag has isotropic rank \(t\) and last norm
\(-T'\), with \(T'=T\) in the period being studied.
We use integral representatives modulo upper unitriangular
changes of their \(d\) columns. An integral row isometry
which preserves the initial isotropic flag, acts trivially
on each of its graded lines, and acts trivially on
\(X^\perp/X\) is an allowed change of row representative.
Write \(N_\Lambda\) for this group. The equivalence of
terminal flags uses both changes. This is the relation in
\srcloc{Formalism/SharpRestriction/CuspSpanNormBounds/FlagEEqHyperbolic}{122}
and \srcloc{Formalism/SharpRestriction/CuspSpanNormBounds/CntPartners}{177}.

\begin{lemma}[Integral coordinates for frames]
\label{lem:coordinates225}
The initial frame can be chosen integral and primitive over
\(\OO_E\). With this choice the row has an integral Witt
decomposition \(A\oplus J\oplus A^*\), with
\(A=\OO_E^t\) the initial isotropic lattice, \(A^*=\OO_E^t\)
its hyperbolic partner, and \(J\) a self-dual middle lattice
containing the primitive vector \(z\). In these coordinates
\(L=J\cap z^\perp\) is the positive initial column. Each
integral upper unitriangular change of the initial frame is
realized by an element of \(N_\Lambda\).
\end{lemma}
\begin{proof}
At each finite place choose the primitive representative
of Lemma~\ref{lem:flag215}. Relative to the given rational
frame this is an upper unitriangular matrix. The identity
works away from finitely many places. Approximate the
entries of the finitely many changes simultaneously by
elements of \(E\) integral elsewhere. This is the Chinese
remainder theorem after clearing the finitely many
denominators in a fractional ideal. Integrality and
independence modulo the residue maximal ideals persist for
sufficiently close changes. Thus the new frame spans a
globally primitive free submodule \(A\oplus\OO_Ez\) of
\(\Lambda\). The same argument applies to any admissible
terminal frame. Two integral primitive representatives of
the same decoration differ by an integral unitriangular
matrix since their local change is integral at every place.

The quotient by \(A\oplus\OO_Ez\) is torsion-free, hence
projective over the Dedekind ring. Its integral functionals
on the frame summand extend to \(\Lambda\). By self-duality
choose partners to the ordered basis of \(A\) having the
identity pairings with that basis and zero pairings with
\(z\). Correct their mutual pairings by adding vectors of
\(A\). The equation for the coefficients is
\(S+S^*=-H\), for the Hermitian matrix \(H\) to be removed.
This has an integral solution by assigning one entry in
each off-diagonal pair and using the integral trace on the
diagonal. The trace ideal is the unit ideal, as its
localization is the unit ideal at every finite place in
our unramified extension. The resulting hyperbolic lattice
splits off orthogonally, leaving \(J\) and \(z\) as claimed.

For any integral map \(V:A^*\to J\), with adjoint \(V^*\)
in these coordinates, the transformation
\begin{equation}
 (x,b,y)\longmapsto
       (x-V^*b+Sy,\ b+Vy,\ y),\qquad S+S^*=-V^*V ,
 \label{eq:shear225}
\end{equation}
is an integral isometry in \(N_\Lambda\), as is seen by
expanding the pairing. The same trace argument produces
\(S\). It fixes \(z\) when the columns of \(V\) belong to
\(L\). More generally \(V^*z\) can be prescribed arbitrarily
in \(A\) since \(z\) is primitive in the self-dual \(J\).
One may also change the basis of \(A\) by an integral upper
unitriangular matrix and that of \(A^*\) by its inverse
adjoint. Together these isometries realize every initial
frame change in the statement.
\end{proof}

Let \(Z\) be an integral \(d\)-by-\(d\) cross Gram matrix
of a supported nondegenerate frame pair. Its class allows
an integral lower unitriangular change on the left and an
integral upper unitriangular change on the right. Let \(C\)
be its first \(t\) rows. Its rank is \(t\), since a kernel
vector on the row side would give an initial isotropic
vector in the radical of the full frame sum. Define
\[
  J_Z=[\,\OO_E^t:C\OO_E^d\,],\qquad
  G_Z=|N_{F/\QQ}\mathcal G|,\qquad
  A_s(Z)=9^{sf}\max(1,G_Z/n_T)^s J_Z^{s+t}.
\]
Here \(\mathcal G=\det Z\,\overline{\det Z}
-TT'D_t\bar D_t\), where \(D_t\) is the leading
\(t\)-by-\(t\) minor of \(Z\). This is the full Gram
determinant up to sign, by expansion in the two last
self-pairings. If \(\mathfrak c\) is the ideal generated by
the maximal minors of \(C\), then \(J_Z=N\mathfrak c\)
by local elementary divisors. Both \(D_t\) and \(\det Z\)
belong to \(\mathfrak c\), by expansion along the last row
for the latter. Thus \((\mathcal G)\subset\mathfrak c\bar{\mathfrak c}\)
and \(J_Z\le G_Z\). These quantities are invariant under
both allowed changes of \(Z\).

For a terminal flag let \(C_{\rm flag}\) denote the
perpendicular of the full frame sum, and \(p\) its embedding
in the initial positive column. Say that an integral
sublattice \(P\subset L\) lifts to this perpendicular if
\(P\subset p(\Lambda\cap C_{\rm flag})\). Write
\(a_P=[L\cap EP:P]\) and \(\operatorname{rank}P=k-s\).

\begin{theorem}[Flags with an integral lift of a prescribed sublattice]
\label{thm:frcount225}
For \(1\le s\le k\), the number of terminal flag classes
for the integral cross class of \(Z\) for which \(P\)
lifts to the perpendicular is at most
\begin{equation}
                 (n_T^2 a_P)^t\, A_s(Z) .\label{eq:frcount225}
\end{equation}
When \(P=0\) the prefactor can be replaced by \(1\).
\end{theorem}
\begin{proof}
Work in the decomposition of Lemma~\ref{lem:coordinates225}.
Every class under consideration has a representative of
exact cross matrix \(Z\). A right change is a basis change
of the terminal decoration. The left change is the adjoint
of an integral initial frame change and can be realized by
applying the inverse row isometry to the terminal frame.
These operations keep \(P\) fixed in the initial quotient.
An upper bound for exact frames modulo row isometries
fixing the initial frame gives an upper bound for the
requested classes.

For a frame admitting \(P\), each \(p_0\in P\) has a unique
lift to \(C_{\rm flag}\), since this perpendicular meets
\(X\) trivially. The integral lift is
\((\phi(p_0),p_0,0)\), for a linear map \(\phi:P\to A\).
The isometries \eqref{eq:shear225} with columns of \(V\)
in \(L\) change this map by \(-V^*|_P\). Put \(P_* =L\cap EP\).
Primitivity of \(P_*\) in \(L\) implies that restriction of
integral functionals from \(L\) to \(P_*\) is onto. Hence
orthogonal projection on \(EP\) maps \(L^\vee\) onto
\(P_*^\vee\). The image of \(L\) there has index at most
\([L^\vee:L]=n_T^2\), by Lemma~\ref{lem:termparent221}.
Moreover \([P^\vee:P_*^\vee]=a_P\) by duality on the same
span. It follows that at most \((n_T^2 a_P)^t\) classes of
\(\phi\) occur, by considering the \(t\) functionals
independently. For \(P=0\) there is one. Fix a representative
of such a class and move each frame to it by an integral shear.

Write the columns of the terminal frame in three blocks
\((U,B,C)\), so that \(C\) is the fixed block from \(Z\).
Let \(v\in\OO_E^d\) be the signed vector of maximal
cofactors of \(C\), normalized in sign to give the expansion
of \(\det Z\). It spans the rational kernel of \(C\),
and its last entry is \(D_t\) up to sign. Put \(w\) equal
to the projection of \(Bv\in J\) on the positive space
\(z^\perp\) in \(EJ\). Since \(Cv=0\),
\[
       \langle Bv,Bv\rangle=-T'|D_t|^2,\qquad
       \langle z,Bv\rangle=\det Z,\qquad
       \langle w,w\rangle=\mathcal G/T ,
\]
where \(|D_t|^2=D_t\bar D_t\) in \(F\). The first identity
uses the self Gram matrix of the terminal frame. The
fiber condition on \(\phi\) gives
\(\langle p_0,Bv\rangle=0\) for all \(p_0\in P\), by
pairing the lift with the linear combination of columns
of coefficients \(v\). Thus \(w\in (EP)^\perp\) in the
positive column. Its prescribed squared norm is positive
at every embedding by nondegeneracy. For two possible
values of \(w\), their difference belongs to
\(L_P=L\cap(EP)^\perp\), since the corresponding \(Bv\)'s
have the same pairing with \(z\).

At \(\tau\) let \(r_\tau=(\tau(\mathcal G)/T)^{1/2}\).
In the product of the complex \(s\)-spaces give a vector
the norm \(\max_\tau\|x_\tau\|/r_\tau\). All the candidates
lie in its unit ball. A nonzero difference \(x\in L_P\)
satisfies \(\prod_\tau\|x_\tau\|^2\ge1\), by the norm of the
nonzero totally positive integer \(\langle x,x\rangle\).
The separation is therefore at least
\(b_0^{-1}\), where \(b_0=(\prod_\tau r_\tau)^{1/f}\).
Disjoint open balls of half that separation lie in the
ball of radius \(1+1/(2b_0)\). Comparing real volumes in
dimension \(2sf\) bounds the number by
\((1+2b_0)^{2sf}\le9^{sf}\max(1,G_Z/n_T)^s\).

Fix one such \(w\). For two possible middle blocks \(B\)
their difference takes values in \(L_P\). In fact their
pairings with \(z\) are the same last row of \(Z\), and
the pairings with \(p_0\) are fixed by
\(\langle\phi(p_0),C\rangle\). The difference vanishes on
\(v\), since \(Bv\) is fixed by \(w\) and its pairing with
\(z\). Hence it factors through a homomorphism
\(C\OO_E^d\to L_P\). To check factorization, two integral
preimages of a vector differ in the rational kernel \(Ev\),
where the difference vanishes. The shears \eqref{eq:shear225}
alter \(B\) by \(VC\) for all integral \(V:\OO_E^t\to L_P\),
fixing \(z\) and \(\phi\). Restriction of such \(V\)'s has
index \(J_Z^s\) in the group of homomorphisms on \(C\OO_E^d\).
Indeed local diagonal reduction of the inclusion multiplies
the index by the \(s\)-th power of each elementary-divisor
index. This argument works for projective \(L_P\).

Finally fix \(B\) itself after these shears. The difference
\(\Delta U\) of two possible upper blocks satisfies
\((\Delta U)^* C+C^*\Delta U=0\) by the equality of the
two self Gram matrices. Evaluation at \(v\) gives
\(\Delta Uv=0\) since \(C^*\) is injective. Thus
\(\Delta U=S C\) for a rational matrix \(S\). Surjectivity
of \(C\) over \(E\) in the Gram identity gives \(S+S^*=0\).
The allowed integral skew matrices give shears with \(V=0\).
Modulo them, the possible \(S\)'s inject into
\[
 \operatorname{Hom}(C\OO_E^d,\OO_E^t)/
            \operatorname{res}\operatorname{Hom}(\OO_E^t,\OO_E^t).
\]
Indeed a difference which extends integrally is already
integral and skew. This quotient has order \(J_Z^t\)
by the same local calculation. Multiplying the four
counts proves \eqref{eq:frcount225}.
\end{proof}

\subsection{A mean for ordered pairs with prescribed integral Gram data}\label{nav:IX-15}
\label{sec:pairmean225}

Let \(Z_1,Z_2\) be prescribed integral cross classes, each
with nondegenerate full Gram matrix. Count ordered terminal
flag pairs together with admissible comparisons as in
Section~\ref{sec:cutoff223}. Equivalence allows independent
row changes from \(N_\Lambda\) on the two flags and identifies
the two comparisons when their induced middle maps correspond
under these changes. The initial flag remains fixed throughout.

\begin{lemma}[Forgetting the comparison and counting finite indices]
\label{lem:indices225}
For fixed \(\Lambda\), the map from pair classes to ordered
pairs of individual flag classes is injective. Also, in
every projective \(\OO_E\)-lattice of rank \(j\), the number
of submodules of absolute index \(n\) is at most
\(n^{j+\lambda}\), where \(\lambda=\log_2(2f)\).
\end{lemma}
\begin{proof}
Suppose two pairs have the same individual classes, via
two changes \(n_1,n_2\in N_\Lambda\) and two unitriangular
changes of the decorations. Transport the first comparison
on the middle quotients by \(n_2\) and \(n_1^{-1}\). The
transported map and the second comparison take the class
of the first terminal vector to that of the second. Indeed
a decoration change alters the terminal vector only by the
terminal isotropic subspace. Both maps induce the identity
on the common residue middle quotient at \(v\mid Q\):
each \(n_a\) acts trivially there because it induces the
identity on the initial middle quotient. Their difference
by composition is therefore an integral isometry fixing a
nonisotropic vector with positive definite perpendicular,
and congruent to the identity at the odd split prime.
The finite-order and torsion-free argument of
Theorem~\ref{thm:unique221} makes it the identity on the
entire middle quotient. This is precisely the equivalence
required of the comparison. The rational common subspace
in the initial column is also unchanged because the
\(n_a\) act identically on that column.

For the index bound, induct on \(n\), keeping the stronger
bound \(n^j(2f)^{\Omega(n)}\), where \(\Omega\) counts prime
factors with multiplicities. The case \(n=1\) is immediate.
For a prime \(p\mid n\), the \(p\)-primary part of the
finite quotient has a simple quotient with residue field of
order \(q\) at some prime of \(E\) above \(p\). Thus the
submodule lies in a maximal submodule of index \(q\mid n\).
At a specified prime the possibilities are hyperplanes
in a residue \(j\)-space, at most \(q^j\). Their kernels are
again projective of rank \(j\). For each of the at most
\(2f\) primes over \(p\), induction inside such a kernel
gives a total at most \(n^j(2f)^{\Omega(n)-1}\).
Summation, followed by \(\Omega(n)\le\log_2 n\), proves
the assertion.
\end{proof}

\begin{theorem}[Mean number of admissible ordered comparisons]
\label{thm:pairmean225}
Let \(1\le s\le k-1\), \(j=k-s\), \(\nu=4+\lambda\).
Impose \(\delta(M)i_1i_2\le Y\) on the comparison lattice
and indices of Section~\ref{sec:cutoff223}, with \(Y\ge1\).
In the normalized initial column genus the mean number of
comparison classes of dimension \(j\) and integral cross
classes \(Z_1,Z_2\) is at most
\begin{equation}
  8\,\frac{\mathfrak m_j\mathfrak m_s}{\mathfrak m_k}
       (Y/n_T)^k n_T^{s+4t} Y^{\nu+1+2t}
                    A_s(Z_1)A_s(Z_2).
                    \label{eq:pairmean225}
\end{equation}
The bound is uniform over compatible row representatives
above each initial column. It also holds for their initial
unipotent-covolume average at the original level. For
\(Y<1\) the count is zero. For \(s=k\) the number per row
is at most \(A_k(Z_1)A_k(Z_2)\).
\end{theorem}
\begin{proof}
Let \(D\subset P_0\) be the rational common subspace and
write \(S=L\cap D\). Identify \(M\) with its integral
image in \(L\). Put \(a=[S:M]\) and
\[
     P_i=M\cap p_i(\Lambda\cap C_i),\qquad b_i=[M:P_i].
\]
The entire \(M\) lifts to \(\mathcal L_i\) on either side
by Lemma~\ref{lem:graphs223}. Hence \(b_i\le i_i\), by
injection of \(M/P_i\) into \(\mathcal L_i/(\Lambda\cap C_i)\).
We have \(\delta(S)a=\delta(M)\) and so
\(\delta(S)a b_1 b_2\le Y\). Also \(L\cap EP_i=S\), hence
\(a_{P_i}=a b_i\). For fixed \(S,a,b_1,b_2\), there are
at most \((a b_1 b_2)^{j+\lambda}\) choices for \(M\) and
the two \(P_i\), by Lemma~\ref{lem:indices225}. With these
choices fixed, every flag on side \(i\) lifts \(P_i\) to
its integral perpendicular. Theorem~\ref{thm:frcount225}
and injectivity in Lemma~\ref{lem:indices225} bound the
possible pair classes by
\(n_T^{4t}(a^2b_1b_2)^t A_s(Z_1)A_s(Z_2)\).
Each individual flag retains the lifting property under
its changes of representative, since those act trivially
on the initial column and preserve the row lattice.

Set \(m=a b_1 b_2\le Y\), using \(\delta(S)\ge1\).
The mean number of primitive \(S\) with \(\delta(S)\le Y/m\)
is bounded by summing \eqref{eq:termupper221} over the
positive integers in this range. This gives
\[
 \frac{\mathfrak m_j\mathfrak m_s}{\mathfrak m_k}
          (Y/n_T)^k n_T^s Y^{\nu+1}m^{-k-\nu-1}.
\]
Use \((a^2b_1b_2)^t\le Y^{2t}\). After inserting the
index count, the remaining powers are \(m^{-s-5}\) since
\(\nu=4+\lambda\). The sum over the three positive indices
is at most \(\zeta(s+5)^3<8\). This proves
\eqref{eq:pairmean225}. For \(s=k\) forget the comparison
and use Theorem~\ref{thm:frcount225} with \(P=0\) twice.
The zero assertion follows from integrality of \(M\).

The pointwise estimate before averaging used only the
number of primitive sublattices in \(L\), uniformly in the
row above \(L\). At the original level the stabilizer of
the initial flag is the semidirect product of its unipotent
and column levels by Theorem~\ref{thm:midlevel131}. The
unipotent covolume over each initial column is the same by
the argument of Lemma~\ref{lem:covol127}: conjugation by
the column group has modulus one on the radical and the
rational unipotent is dense in its finite adeles. The
remaining normalized level weights push onto the genus
measure by Lemma~\ref{lem:termparent221}. This gives the
asserted initial-row average.
\end{proof}

\begin{corollary}[Partners of a fixed generic flag in codimension one]
\label{cor:partner225}
Fix one terminal flag with \(I_1=1\) and integral cross
class \(Z_1\). In the same row the number of possible
partners with \(I_2=1\), integral class \(Z_2\), an
admissible comparison and \(s=1\) is at most
\[
             (n_T^2 G_{0,1})^t A_1(Z_2),
             \qquad G_{0,1}=G_{Z_1}/Q^{2df}.
\]
\end{corollary}
\begin{proof}
Both rational graphs have dimension \(k-1\), so an
intersection of dimension \(k-1\) makes them equal. Their
integral intersections with the compatibility lattice are
equal by Lemma~\ref{lem:graphs223}, whose intersection
argument holds also for \(s=1\). Saturation gives
\(\mathcal L_i=\Lambda\cap C_i\) by
Proposition~\ref{prop:index219}. Thus \(P=p_1(\Lambda\cap C_1)\)
is already fixed by the first flag, has determinant \(G_{0,1}\),
and lifts to the integral perpendicular of every partner.
Its saturation in \(L\) has determinant at least one, so
\(a_P\le G_{0,1}\). Apply Theorem~\ref{thm:frcount225}.
\end{proof}

Equations \eqref{eq:cutoff223} and \eqref{eq:pairmean225}
give a mean for actual pairs separately for each pair of
integral cross classes. The codimension-one bound keeps its
first flag fixed. These are the counts entering
\srcloc{Formalism/SharpRestriction/CuspSpanNormBounds/CongruenceAtSplitFrobeniusHolds}{556}.
We next determine which codimension-one partners can occur at the original terminal prime.

\subsection{A congruence between the normal lines of two terminal flags}\label{nav:IX-16}
\label{sec:normal227}

Work at the original level on the fields of
Section~\ref{sec:saturation219}. The common norm of the three
terminal vectors is \(-T\), and \(T\) is prime to \(Q=g_1g_2\).
Fix an admissible comparison \(\beta\) for which both frame sums
\(S_i=W_0\oplus W_i\) are nondegenerate and \(s=1\).
In the two positive middle spaces \(P_0,P_2\) the graphs of
\(C_i=S_i^\perp\) coincide, after using \(\beta\) on the
second coordinate of the first graph. Indeed each graph and
their intersection have dimension \(k-1\). Write \(D,D'\)
for the two images of this common graph and \(N,N'\) for
their perpendicular lines.

We use the Hermitian pairing conjugate-linear in the first
variable. A vector of \(N\) has a row representative
\(l\in S_i\cap W_0^\perp\). To see this, lift to \(W_0^\perp\)
and project off \(C_i\). Its \(C_i\)-component is already
zero since its class is perpendicular to \(D\). In the
same way a vector of \(N'\) has a representative
\(l'\in S_i\cap W_i^\perp\), undoing \(\beta\) for \(i=1\).
In the direct sum \(S_i\), write
\[
 l-l'=x+\alpha z_0+x_i+\alpha' z_i,\qquad
                   x\in X_0,\quad x_i\in X_i.
\]
Changing either lift adds only an isotropic term. Let
\(A_i:N\oplus N'\to E^2\) send the two classes to
\((\alpha,\alpha')\). For lifts with a common \(C_i\)-component
the coefficients are still those of their normal components,
by subtracting that common vector. More generally, for lifts
in \(X_0^\perp,X_i^\perp\), one can first remove their respective
terminal-line components in the signed middle spaces
\(Ez_a\oplus P_a\). Their \(C_i\)-components coincide exactly
when their difference lies in \(S_i\). This follows by pairing
the difference with every vector of \(C_i\). Equivalently, the
projections on \(D,D'\) of the two positive coordinates lie on
the common graph. This last description is independent of \(i\).

The norm identity in the proof of Proposition~\ref{prop:radical223}
gives, at each embedding,
\[
          \|y\|^2-\|y'\|^2
       =T\bigl(|A_i(y,y')_1|^2-|A_i(y,y')_2|^2\bigr).
\]
Polarization identifies \(A_i\) as an isomorphism. Put
\(\gamma=A_2 A_1^{-1}\). It is an \(E\)-linear isometry for
\(\operatorname{diag}(1,-1)\). Equality \(\gamma=1\) says
precisely that the two pairs have the same coefficients for
all matched normal vectors.

For \(i=1,2\) retain \(B_i=\Lambda\cap S_i\), its sublattice
\(B_i^+\) from Section~\ref{sec:saturation219}, and
\(I_i=[B_i:B_i^+]\). Set
\[
 \mathfrak h_i=\operatorname{Ann}_{\mathcal O_E}(B_i/B_i^+),
 \qquad \mathfrak h=\overline{\mathfrak h_1}\mathfrak h_2,
 \qquad
 \mathcal Z=\{(u,v)\in\mathcal O_E^2:u+v\in Q\mathcal O_E\}.
\]
The norm \(N\mathfrak h_i\) is at most \(I_i\). At a prime
the exponent of the annihilator is the largest elementary
divisor exponent of the quotient, bounded by their sum.

\begin{theorem}[Integral congruence for the normal comparison]
\label{thm:congruence227}
For every such ordered comparison,
\begin{equation}
       \mathfrak h\,(\gamma-1)\mathcal Z\ \subset\ T\mathcal Z .
       \label{eq:congruence227}
\end{equation}
\end{theorem}
\begin{proof}
For \(b\in S_i\) let \(\kappa_i(b)\) denote its two terminal
coefficients modulo \(X_0+X_i\). We first compute
\(\kappa_i(QB_i^+)=\mathcal Z\). Away from \(Q\) the two
individual primitive frames generate \(B_i^+\) and \(Q\)
is a unit. At \(v\mid Q\) choose representatives with equal
reductions. The image of \(QB_i^+\) is generated by \(Q(1,0)\)
and \((Q/\pi_v)(-1,1)\). All the other generators have zero
coefficients. These vectors span exactly the local module
of pairs of integral coordinates whose sum is divisible
by \(Q\). Localization proves the assertion.

Fix \(p=(u,v)\in\mathcal Z\) and \(a\in\mathfrak h_1\).
On \(S_1\) take the conjugate-linear functional \(\ell\)
which vanishes on \(X_0+X_1\) and has values \(u,-v\) on
\(z_0,z_1\), respectively. It is integral on \(B_1^+\).
For the terminal divided difference at a marked place one
uses \(u+v\in Q\mathcal O_E\), and the other divided
differences give zero. Representatives of \(z_a\) differing
by a vector of \(X_a\) give the same values. Thus
\(\bar a\ell\) is integral on \(B_1\), since \(aB_1\subset B_1^+\).
The quotient of \(\Lambda\) by \(B_1\) is torsion-free and
finitely generated over the Dedekind ring \(\mathcal O_E\),
hence projective. Extend the integral functional across a
module complement. By self-duality it is represented by
\(w\in\Lambda\), so
\[
 w\perp X_0+X_1,\qquad
 \langle z_0,w\rangle=\bar a u,\qquad
 \langle z_1,w\rangle=-\bar a v .
\]
Put \(u'=\bar a u/T,\ v'=\bar a v/T\). The positive
coordinates of \(w\) in the two middle spaces have respective
representatives \(w+u'z_0,\ w-v'z_1\). Their difference is
\(u'z_0+v'z_1\). Thus their \(C_1\)-components coincide
and their normal coordinates have coefficient pair
\((u',v')\) under \(A_1\).

Use \(w,\beta w\) for the signed lifts in the second pair.
The initial positive coordinate is the same, and the terminal
one is its previous terminal coordinate transported by \(\beta\).
The common-graph description therefore gives equal
\(C_2\)-components. Also \(\beta z_1=z_2\bmod X_2\). It follows
that \(w-\beta w\in S_2\) and
\[
              \kappa_2(w-\beta w)=(\gamma-1)(u',v').
\]
Both signed lifts are integral. At each marked place the
residual middle quotients identify with
\(\bar X^\perp/\bar X\), and the comparison is the identity
there by Theorem~\ref{thm:unique221}. Their difference reduces
into the common \(\bar X\). Consequently
\[
 w-\beta w\in QB_2+(\Lambda\cap X_0)+(\Lambda\cap X_2).
\]
Indeed one may subtract local integral isotropic lifts of the
residues. The remainder lies in \(Q\Lambda\cap S_2=QB_2\).
Away from \(Q\) this holds since \(Q\) is a unit, so local
membership gives the asserted global membership. For any
\(b\in\mathfrak h_2\) the pair \(b\kappa_2(w-\beta w)\)
therefore belongs to \(\kappa_2(QB_2^+)=\mathcal Z\).
Multiplication by \(T\) and variation of \(a,b,p\) prove
\eqref{eq:congruence227}.
\end{proof}

Write \(D_i=\det Z_i\), \(R_i\) for its leading minor of
size \(d-1\), and \(\mathcal G_i=D_i\bar D_i-T^2R_i\bar R_i\).
It is totally positive, so \(D_i\ne0\). For one embedding
\(\tau\) over each real place put
\[
 U_{i,\tau}=\frac{|\tau D_i|}{\sqrt{\tau\mathcal G_i}},
 \qquad V_{i,\tau}=\frac{T|\tau R_i|}{\sqrt{\tau\mathcal G_i}},
 \qquad U_{i,\tau}^2-V_{i,\tau}^2=1 .
\]

\begin{theorem}[Separation of two comparisons in codimension one]
\label{thm:separation227}
Write \(c=\gamma_{21}\). If \(c\ne0\), then
\begin{equation}
 \sqrt{I_1I_2}\prod_{\tau}
       (U_{1,\tau}V_{2,\tau}+V_{1,\tau}U_{2,\tau})
                    \ \ge\ (T/Q)^f .                  \label{eq:angles227}
\end{equation}
If \(c=0\), either \(\gamma=1\) or
\(\sqrt{I_1I_2}\ge (T/2)^f\). In the first case one has
\begin{equation}
                        R_1/D_1=R_2/D_2\quad\text{in }E.
                        \label{eq:ratio227}
\end{equation}
Thus for \(I_1=I_2=1\) every pair has identical coefficient
maps or satisfies \eqref{eq:angles227}.
\end{theorem}
\begin{proof}
The vector \((Q,0)\) belongs to \(\mathcal Z\). Theorem~\ref{thm:congruence227}
gives \((Qc/T)\mathfrak h\subset\mathcal O_E\). When \(c\ne0\)
the norm of this nonzero integral ideal is at least one, whence
\(\prod_\tau|\tau c|\ge (T/Q)^f/\sqrt{N\mathfrak h}\).
Choose unit bases of \(N,N'\) at \(\tau\) and multiply both
output coordinates of \(A_i\) by \(\sqrt T\). The resulting
matrix \(M_i\) preserves \(\operatorname{diag}(1,-1)\).
Its first column has entry magnitudes \(U_{i,\tau},V_{i,\tau}\).
Indeed, if \(v_i\) is the signed cofactor column of the top
\(d-1\) rows of \(Z_i\), expansion along the last row gives
\((v_i)_d=R_i\) and \(\langle z_0,Y_i v_i\rangle=D_i\).
Here \(Y_i\) is the terminal frame. The vector
\[
       y_i=Y_iv_i+(D_i/T)z_0\in S_i\cap W_0^\perp
\]
has norm squared \(\mathcal G_i/T\), since
\(\langle Y_iv_i,Y_iv_i\rangle=-T R_i\bar R_i\).
Its nonzero class \(n_i\) generates \(N\), with
\(A_i(n_i,0)=(D_i/T,R_i)\). This calculation includes
\(R_i=0\). The column and row identities for a matrix
preserving \(\operatorname{diag}(1,-1)\) force its second
column to have magnitudes \(V_{i,\tau},U_{i,\tau}\).
The inverse is obtained by conjugate transposition and the
two signature matrices. Taking the lower left entry of
\(M_2M_1^{-1}\), in the same bases for both maps, now gives
\(|\tau c|\le U_{1,\tau}V_{2,\tau}+V_{1,\tau}U_{2,\tau}\).
Use \(N\mathfrak h\le I_1I_2\).

When \(c=0\), orthogonality of the two columns makes
\(\gamma\) diagonal, with diagonal entries \(a,e\) of
absolute value one at every embedding. Apply
\eqref{eq:congruence227} to \((1,-1)\). For \(z=a,e\)
one gets \((z-1)\mathfrak h/T\subset\mathcal O_E\).
If \(z\ne1\), its norm and \(|\tau(z-1)|\le2\) give
\(N\mathfrak h\ge(T/2)^{2f}\). Otherwise both entries
are one. In that case \(A_1=A_2\). Write \(n_2=h n_1\)
in the common rational line \(N\). The cofactor identities
then give \((D_2/T,R_2)=h(D_1/T,R_1)\), proving
\eqref{eq:ratio227}.
\end{proof}

\begin{corollary}[A restriction on the Gram data]
\label{cor:gap227}
For a pair with \(I_1=I_2=1\) and distinct ratios
\(R_i/D_i\), some \(i,\tau\) have \(R_i\ne0\) and
\[
 \frac{\tau\mathcal G_i}{T^2|\tau R_i|^2}
          \le \frac{2}{\sqrt{1+(T/Q)^2}-1}.
\]
\end{corollary}
\begin{proof}
Distinct ratios exclude \(\gamma=1\). The terminal prime
exceeds two, so the second alternative for \(c=0\) is
impossible. Put \(H=\max_{i,\tau}V_{i,\tau}\). Every factor
in \eqref{eq:angles227} is at most \(2H\sqrt{1+H^2}\).
Solving \(2H\sqrt{1+H^2}\ge T/Q\) yields the assertion
at a place where the maximum is attained.
\end{proof}

This restricts the two actual integral Gram classes in
Corollary~\ref{cor:partner225} before forming the period sum.
The ratios and the displayed normalized gaps are invariant
under the two unitriangular changes of frame. The normalized
gaps are also invariant under Lemma~\ref{lem:match209},
since \(D_i/R_i\) and the terminal norm parameter both
scale in magnitude by \(Q^{2d-3}\). The separation in
\eqref{eq:angles227} uses the original value of \(T\).
The scalar condition and the comparison of its local and
archimedean restrictions appear in
\srcloc{Formalism/SharpRestriction/CuspSpanNormBounds/IsPrimitiveSub}{252}
and \srcloc{Formalism/SharpRestriction/CuspSpanNormBounds/LatticeGluingRadicals}{3177}.

\subsubsection{Counting the partners of a rational label}\label{nav:IX-17}
\label{sub:partners230}

For a nondegenerate label \(\eta\), write \(D_\eta=D_d\),
\(R_\eta=D_{d-1}\), \(z_\eta=TR_\eta/D_\eta\), and
\(G_\eta=|N_{F/\QQ}\mathcal G_\eta|\). The tuple of slopes
\((\tau z_\eta)_\tau\) belongs to the product of unit disks. Put
\[
 \Phi_\eta=\prod_\tau
    \frac{|\tau D_\eta|+T|\tau R_\eta|}{\sqrt{\tau\mathcal G_\eta}}
          =\prod_\tau(U_{\eta,\tau}+V_{\eta,\tau}).
\]
This is \(\Phi((\tau z_\eta)_\tau)\) from
Section
\ref{nav:VII-02}.
It is also the product of the real angular factors in the proposed
local amplitude bound. Indeed for \(R_\eta\ne0\) that factor is
\((1+c)/\sqrt{c^2-1}\) with \(c=|\tau D_\eta|/(T|\tau R_\eta|)\),
and for \(R_\eta=0\) it is one. These formulas specify the
weight independently of any estimate on the amplitude.

\begin{proposition}[The common determinant and the slope]
\label{prop:slope230}
Let two flags admit an original-level comparison of codimension
one. At every finite place \(v\) of \(F\), write \(q_v\) for
the residue cardinality, \(v\) also for the normalized
valuation, and \(I_{a,v},i_{a,v}\) for the local indices of
\(B_a^+\subset B_a\) and \(\Lambda\cap C_a\subset\mathcal L_a\)
on side \(a\), respectively. Then
\[
 \frac{q_v^{v(\mathcal G_1/Q^{2d})}}{I_{1,v}i_{1,v}}
       =\frac{q_v^{v(\mathcal G_2/Q^{2d})}}{I_{2,v}i_{2,v}}.
\]
At the levels under consideration \(Q\ge4\). Distinct slopes
imply \(\sqrt{I_1 I_2}\,\Phi_1\Phi_2\ge(2T/Q)^f\).
For equal slopes, either \(A_1=A_2\) or
\(\sqrt{I_1 I_2}\ge(T/2)^f\).
In particular for saturated pairs the two determinant ideals
in \(\mathcal O_F\) agree, and the two coefficient maps agree
exactly when the two slopes agree.
\end{proposition}
\begin{proof}
The two rational graphs coincide in codimension one.
Lemma~\ref{lem:graphs223} therefore identifies the entire
lattices \(\mathcal L_1\) and \(\mathcal L_2\) isometrically.
For a local lattice \(L\) use
\(\delta_v(L)=[L^\vee:L]^{1/2}\) when it is integral.
Increasing a lattice by index \(a\) divides its determinant
\(\delta_v\) by \(a\), since the dual decreases by the same
index. Self-duality and primitivity give
\(\delta_v(B_a)=\delta_v(\Lambda\cap C_a)\). Indeed the
common quotient of \(\Lambda_v\) by the orthogonal direct
sum of those two primitive lattices identifies with the
dual quotient of each, by restricting integral functionals.
Such functionals extend from a primitive summand. In the
local frame basis the determinant ideal of \(B_a^+\) is
generated by \(\mathcal G_a/Q^{2d}\). At a marked place
the \(d\) divided differences remove two factors of the
uniformizer apiece from the Gram determinant, and at any
other place \(Q\) is a unit. The local index formula of
Proposition~\ref{prop:index219} gives
\(\delta_v(B_a^+)=q_v^{v(\mathcal G_a/Q^{2d})}\).
Equating the two values of \(\delta_v(\mathcal L_a)\) proves
the first assertion.

Write \(g=A_2 A_1^{-1}\) and
\(\mathfrak h=\overline{\mathfrak h_1}\mathfrak h_2\) as in
Theorem~\ref{thm:congruence227}. The cofactor vectors give
positive lines of slopes \(z_1,z_2\) as the respective
images of \(N\oplus0\) under \(A_1,A_2\). Thus \(gz_1=z_2\).
When \(g_{21}\ne0\), use \eqref{eq:angles227} and the inequality
\[
 2(U_1 V_2+V_1 U_2)\le(U_1+V_1)(U_2+V_2)
\]
at each embedding, since \((U_1-V_1)(U_2-V_2)\ge0\).
When \(g_{21}=0\) and the slopes differ, the alternative in
Theorem~\ref{thm:separation227} gives
\(\sqrt{I_1 I_2}\ge(T/2)^f\).
Now \(Q\ge4\) and \(\Phi_a\ge1\), giving the stated bound
also in this case.

If \(z_1=z_2\), at each embedding \(g\) preserves a positive
line and its negative perpendicular, with unit-modulus
eigenvalues. In the lattice basis \((1,-1),(0,Q)\) every
entry of \(g-1\) belongs to \(T\mathfrak h^{-1}\) by
Theorem~\ref{thm:congruence227}. If \(g-1\) has rank two,
its determinant belongs to \(T^2\mathfrak h^{-2}\), is nonzero
and has absolute value at most four everywhere. Taking the
norm of the integral ideal
\(\det(g-1)\mathfrak h^2/T^2\) yields
\(N\mathfrak h\ge(T/2)^{2f}\). In rank one use
\(\operatorname{tr}(g-1)\in T\mathfrak h^{-1}\). It is
nonzero because the two invariant lines diagonalize \(g\)
and precisely one eigenvalue differs from one. Its absolute
value is at most two everywhere, giving the same norm bound.
Since \(N\mathfrak h\le I_1I_2\), only rank zero is possible
below this threshold. For saturated flags both global indices
are one and \(i_a\mid I_a\) by Proposition~\ref{prop:index219}.
The local equality gives the determinant ideals, and \(A_1=A_2\)
gives equal slopes by \eqref{eq:ratio227}.
\end{proof}

\begin{theorem}[A weighted sum in one original row]
\label{thm:partners230}
Fix a saturated first flag \(*\) in that row and let its
terminal norm be \(-T\). Sum over its saturated, relevant,
nondegenerate partner classes admitting a codimension-one
comparison. For a represented rational label \(\eta\), use the
monomial matrix \(M_\eta\), the index bounds \(h_t\) and the
class count \(C(M_\eta)\) of Theorem
\ref{thm:crossclasses230}.
Put \(t=d-1\), \(n_T=T^f\), \(G_{0,*}=G_*/Q^{2df}\), and
\[
 B_*=(n_T^2G_{0,*})^t,\qquad
 K_\eta=C(M_\eta)\,9^f\max(1,G_*/n_T)
                 \min\{G_*,Q^{2ft}h_t\}^{t+1}.
\]
Let \(\mathcal A\) be the represented relevant labels of these
partners with \(z_\eta\ne z_*\), so
\((\mathcal G_\eta)=(\mathcal G_*)\) in \(\mathcal O_F\).
For any nonnegative coefficients \(a_\eta\) on \(\mathcal A\)
and real \(p\ge3\), define
\[
 S=\sum_{\substack{\text{partner classes }F\\z_{\eta(F)}\ne z_*}}
     \frac{a_{\eta(F)}}{(\Phi_*\Phi_{\eta(F)})^p},
 \qquad
 L(z)=\sum_{\substack{\eta\in\mathcal A\\z_\eta=z}}K_\eta a_\eta.
\]
One has
\begin{align*}
 S&\le B_*\,10^f\Phi_*^{-3}
           \left(\frac{Q}{2T}\right)^{f(p-3)}
              \sup_{z\in E} L(z),\\
 S&\le B_*\left(\frac{Q}{2T}\right)^{fp}
       \left(\frac{128\log(2Q)}{Q^4}\right)^{fd/2}
       \left(\sum_{\eta\in\mathcal A}w(D(\eta))
                         K_\eta^2 a_\eta^2\right)^{1/2}.
\end{align*}
For the partners with \(z_\eta=z_*\) the corresponding bound is
\(B_*\Phi_*^{-2p}\sum_{z_\eta=z_*}K_\eta a_\eta\), with the
sum over their labels.
\end{theorem}
\begin{proof}
Aligning the two residue frames at the marked primes gives
there the cross matrix \(-T E_{dd}\) modulo each marked prime
by the self pairing of the initial frame. The two unitriangular frame changes
preserve this congruence. Thus the number of integral classes
over \(\eta\) is at most \(C(M_\eta)\). In the notation of
Theorem~\ref{thm:frcount225}, their indices \(J_Z\) are bounded
by \(Q^{2ft}h_t\) from the indicated minor, and by \(G_Z\)
from the full Gram determinant. Proposition~\ref{prop:slope230}
gives \(G_Z=G_*\). Corollary~\ref{cor:partner225} now bounds
the number of actual partner classes above \(\eta\) by
\(B_*K_\eta\). The uniqueness in Lemma~\ref{lem:indices225}
accounts for the comparison attached to each partner.

For saturated pairs, Theorem~\ref{thm:congruence227} puts
\(A_2A_1^{-1}\) in the group \(\Gamma\) of Lemma \ref{lem:binary230}.
The slope of every represented label is therefore \(gz_*\)
for some \(g\in\Gamma\). The stabilizer assertion of that
lemma makes \(g\) unique for each slope. Grouping the label
bound by slope and applying Theorem
\ref{thm:disk230}
gives the first estimate. For the second use
\(\Phi_*\Phi_\eta\ge(2T/Q)^f\) on every distinct slope.
Cauchy--Schwarz bounds \(\sum_{\eta\in\mathcal A}K_\eta a_\eta\)
by the square root of the weighted sum displayed above times
\((\sum_{\eta\in\mathcal A}w(D(\eta))^{-1})^{1/2}\).
Theorem~\ref{thm:levelsum203} bounds the latter factor for the
full set of relevant supported labels. Equal slopes give
\(\Phi_\eta=\Phi_*\) directly. All estimates apply first to
finite sets, and then to their increasing union by positivity.
\end{proof}

The power \(p=2k\) is available in this arithmetic sum. It is
the exponent of the angular factor in
\srcloc{Formalism/SharpRestriction/CuspSpanNormBounds/CongruenceAtSplitFrobeniusHolds}{469}.
Application to the projected period still requires an estimate
for the remaining weighted amplitudes of its specified sections
and the contribution of equal slopes and nonsaturated pairs.

\subsection{Summation over the nondegenerate pairs}\label{nav:IX-24}
Keep the unit column tuples and the transported real kernels
from Corollary~\ref{cor:square216}. Write \(G^\#\) for the
absolute norm over \(F\) of the transported Gram invariant and
\(A_j^\#\) for the absolute norm over \(E\) of the transported
nonzero endpoint minor of size \(j\). Put
\(B_\gamma^\#=G^\#\prod_{j}A_j^\#\), the product over these
endpoints. For a finite set of supported nondegenerate labels
let
\[
 H_\theta(\mathcal A)=\sum_{\gamma\in\mathcal A}
       w(\beta_\gamma)(B_\gamma^\#)^\theta
                   P_\infty^\#\|\widehat K_{\gamma^\#}\|_{\rm HS}^2 .
\]

\begin{theorem}[A weighted real estimate sufficient for absolute summation]
\label{thm:moment218}
With \(h=k-1\), \(S=d(d-1)\) and \(m=h-2d+1\), the original
nondegenerate summands on the common cohomology intervals satisfy
\begin{align*}
 \sum_{\gamma\in\mathcal A,c}|J_{\gamma,c}|
 \le {}&
 \frac{A_*\,\mathfrak m_h\,8^{fd/2}}{\sqrt{c_{\rho_1}^*c_{\rho_2}^*}}\,
  e^{O(f2^{-m})}R^{-fS}
  Q^{f(S/2+d-2-d^2h)} H_{h/2}(\mathcal A)^{1/2}.
\end{align*}
The constant is uniform in \(\mathcal A\).
\end{theorem}
\begin{proof}
Combine the finite mass bound in Theorem~\ref{thm:finite216}
with the exact transport of the real kernel. At a proper block
endpoint \(j_\nu\) the exponent in \(P_\infty^{1/2}W_\gamma\)
is \(h_{\nu+1}\), while at the last endpoint it is zero. The
remaining Gram factor is \(|N_{F/\QQ}\mathcal G|^{h/2}\).
After transport \(G^\#\le A_d^\#\), by the real Gram identity
at every embedding. Each positive proper endpoint exponent
is at most \(d-1\le h/4\) at large \(x\). Thus the transported
product of norm factors is bounded by \((B_\gamma^\#)^{h/4}\).
The minor rescalings give the displayed power of \(Q\).
Cauchy--Schwarz and the bound \(8^{fd}\) on the sum of the
inverse endpoint weights give the assertion. The exact rescaling
and all finite measure constants are proved in
Theorem \ref{thm:moment217}.
\end{proof}

\begin{proposition}[Divergence with a smaller separated weight]
\label{prop:necessary218}
In the same Cauchy--Schwarz separation, putting \(H_\theta\)
in the real factor leaves an infinite arithmetic squared sum
when \(\theta\le(h-2)/2+1/(4f)\).
\end{proposition}
\begin{proof}
Fix the field and interval parameters. The rational unitary
transvections of
Proposition \ref{prop:escape217}
give a diagonal cross matrix whose first \(d-1\) entries are
fixed and whose last entry is \(T(-1+t\lambda)\),
\(t=1,2,\ldots\), with fixed nonzero purely imaginary \(\lambda\).
All are supported at the original finite level. The compatible
mass and the real density both grow as \(t^{2fh}\) up to fixed
positive factors. The two weights \(w(\beta_\gamma)\) and
\(B_\gamma^\#\) each grow as \(t^{4f}\). The arithmetic squared
summand on this family is
\((\mathcal M_\gamma^{=})^2/
 [P_\infty w(\beta_\gamma)(B_\gamma^\#)^\theta]\), since
\(a=0\). Its sum dominates a positive constant times
\(\sum_{t\ge1}t^{2f(h-2-2\theta)}\), which diverges in the
stated range.
\end{proof}

An explicit upper bound of the required size for \(H_{h/2}\),
or an estimate keeping the kernel and the masses together, is
still needed for this period argument. Proposition~\ref{prop:necessary218}
concerns the separated arithmetic sum. The exceptional strict
cohomology bound and the designated integral lower bound are
also open.

\subsection{Saturation at the original integral level}\label{nav:IX-25}
Fix one nondegenerate supported label and a compatible row
lattice \(\Lambda\) in its finite orbital mass. Write \(W_i\)
for the spans of its two decorated frames, \(X_i\) for their
isotropic parts, \(W=W_0\oplus W_1\) and \(C=W^\perp\).
Set \(U=(\Lambda\cap W_0)\oplus(\Lambda\cap W_1)\) and
\(B=\Lambda\cap W\). Enlarge \(U\) at each \(v\mid Q\) by
adjoining the differences of two integral frame representatives
of equal reduction, divided by a uniformizer. The lattice
\(B^+\) so obtained is independent of the representatives. In
fact further unitriangular changes on the two lists must agree
modulo the uniformizer, by primitivity of their reduction.
Write \(I=[B:B^+]\), \(h=k-1\) and \(m=h-2d+1\).

\begin{theorem}[The saturation index in the original period]
\label{thm:saturation220}
Give the compatible finite cells their normalized orbital
measure. On the common fields and intervals, uniformly in
the nondegenerate label,
\[
   \Pr(I>1)\ll f\,2^{-m},\qquad
   \mathbb E[I^t]\le \exp\bigl(O(f2^{-(m-t)})\bigr)
                                  \quad(0\le t\le m-2).
\]
Let \(\mathcal L\subset C\) consist of vectors \(c\) admitting
two integral lifts \(c+x_i\in\Lambda\), \(x_i\in X_i\), whose
difference belongs to
\(Q\Lambda+(\Lambda\cap X_0)+(\Lambda\cap X_1)\). Then
\([\mathcal L:\Lambda\cap C]\mid I\), and
\[
 \delta(\Lambda\cap C)=
             |N_{F/\QQ}\mathcal G|/(Q^{2df} I),
 \qquad \delta(M)=[M^\vee:M]^{1/2}.
\]
\end{theorem}
\begin{proof}
The ordered-frame calculation behind Theorem~\ref{thm:finite216}
sums over saturations \(M'\supset M_0\) with factor
\([M':M_0]^{-h}\). For the actual test this index is \(I_v\),
including at the marked places after changing to divided
differences. Removing the primitive term and inserting
\(I_v^t\) leaves a proportion \(O(q_v^{-(m-t)})\) by summation
over triangular lattice bases. Multiplying and using
\(\sum_v q_v^{-b}\ll f2^{-b}\) for \(b\ge2\) proves the
claims on probability. Sending \(c\) to \(x_0-x_1\bmod U\)
embeds \(\mathcal L/(\Lambda\cap C)\) into \(Q(B/U)\), a
quotient of \(B/B^+\). Finally \(B\) and its perpendicular
in a self-dual lattice have discriminant modules of equal
order, and \([B^+:U]=Q^{2df}\). The local calculations and the
bound on the original period summands over the exceptional
cells are proved in
Section \ref{nav:IX-10}.
\end{proof}

On the rows with \(I=1\) the compatibility module is the
integral perpendicular. The comparison of two terminal flags
in the same row still depends on how many pairs give a given
subspace in the positive definite column.

\subsection{Primitive subspaces and comparisons in a single row}\label{nav:IX-26}
Let \(T\) be the common terminal prime and \(F=E^+\) the
field already chosen on the interval. The isotropic subspace
\(X\) and the terminal vector \(z\) of the initial flag give
the positive rank-\(k\) column \(P=(X+Ez)^\perp/X\).
Its integral lattice \(L\) is the image of the integral row
vectors in \((X+Ez)^\perp\). For \(j+s=k\),
\(N_j(L,D)\) counts sublattices \(M=L\cap EM\) of rank \(j\)
and determinant \(\delta(M)=|M^\vee/M|^{1/2}=D\).
Give the genus of the initial column its normalized
reciprocal-automorphism measure. Write \(\mathfrak m_i\) for
the mass of the standard self-dual positive genus in rank
\(i\), \(R=|\operatorname{disc}F|^{1/f}\) and
\(\nu=4+\log_2(2f)\).

\begin{theorem}[A genus bound on the common fields and intervals]
\label{thm:genus222}
For \(1\le s\le k-1\), \(j=k-s\) and integral \(D\ge1\),
\[
   \mathbb E_T N_j(L,D)\le
    \frac{\mathfrak m_j\mathfrak m_s}{\mathfrak m_k}
              (D/T^f)^k(T^f)^s D^\nu .
\]
For \(2\le s\le k-2\) the genus factor is at most
\(\mathfrak m_2(C/(kR))^{s(k-s)f}\), with
\(\mathfrak m_2\le b^{Cf}\) for large \(x\), \(b=\exp(x^\eps)\).
The bounds use absolute constants.
\end{theorem}
\begin{proof}
Over a self-dual parent a primitive lattice and its
perpendicular have opposite finite discriminant pairings.
The self-dual parent is recovered as the inverse image of the
graph of this identification. Counting graphs with their
stabilizers and applying the definite Hermitian mass formula
gives the bound
\((\mathfrak m_j\mathfrak m_s/\mathfrak m_k)D^{k+\nu}\)
for the parent genus. The full argument, including the
stabilizer volumes for nonmaximal lattices and primes over
two, is in
Section \ref{nav:IX-11}.

In the first split factor at each \(v\mid T\) the terminal
perpendicular cuts the parent by a single residual hyperplane.
The other factor stays dual to the full parent factor. The
hyperplanes are uniform and independent conditional on the
parent, by orbit--stabilizer. Each cuts a fixed primitive
\(j\)-plane with index \(T\) or \(1\), and the probability
of index \(1\) is \(p=(T^s-1)/(T^k-1)\). For index \(T^e\)
use the parent bound at determinant \(D/T^e\). Averaging
the factor \(T^{-ke}\) gives
\((p+(1-p)T^{-k})^f=T^{-jf}\), proving the estimate.
For the genus ratio, each successive mass is given by the
Dirichlet functional equations on the abelian field in terms
of a factorial, a power of \(R\) and an Euler product at an
integer at least two. After retaining \(\mathfrak m_2\),
divide these factors in ranks \(\min(j,s)\) and
\(\max(j,s)\) as in
Proposition \ref{prop:massdecay221}.
\end{proof}

\begin{proposition}[The role of a primitive subspace in the pair sum]
\label{prop:comparison222}
For two fixed admissible terminal flags in the same row there
is at most one comparison on their middle quotients satisfying
the original congruence. For comparisons whose rational
common subspace has dimension \(j\) and whose determinant
cutoff is \(D_{\max}\), the number of distinct such subspaces
in the initial column, with compatibility-index product
\(\ge b_0\ge1\), has mean at most
\[
    \frac{\mathfrak m_j\mathfrak m_s}{\mathfrak m_k}
       (D_1/T^f)^k (T^f)^s D_1^{\nu+1},
                    \qquad D_1=D_{\max}/b_0
\]
for \(1\le j\le k-1\), and is zero for \(D_1<1\).
\end{proposition}
\begin{proof}
Each terminal flag and the initial flag have the same
reduced isotropic subspace \(U\) at \(v\mid Q\). By
self-duality their integral middle quotients all reduce to
\(U^\perp/U\). The congruence for a comparison forces the
identity on this common quotient. Two comparisons then
differ on the middle quotient by an integral isometry which
fixes the terminal vector and acts on its positive definite
perpendicular with finite order. At an odd completely split
prime dividing \(Q\) the principal congruence kernel in
\(\operatorname{GL}_{k+1}(\ZZ_p)\) is torsion-free, so the two maps agree.

The common subspace has a full comparison lattice \(M_\beta\)
whose image in the initial column is integral by the lifting
condition defining the two compatibility modules. The
cutoff is \(\delta(M_\beta)i_1i_2\le D_{\max}\) for the two
compatibility indices \(i_1,i_2\). Passing to the primitive
saturation divides the determinant by the index and hence
gives determinant at most \(D_1\). Distinct subspaces have
distinct saturations. Summing Theorem~\ref{thm:genus222}
over the positive integers \(D\le D_1\) proves the bound.
The two integral identifications and the uniqueness on the
entire middle quotient are proved in
Corollary \ref{cor:subspaces221} and Theorem \ref{thm:unique221}.
\end{proof}

Proposition~\ref{prop:comparison222} counts distinct common subspaces in the initial column. We now bound their determinants in terms of the two Gram matrices and then count the pairs over each subspace.

\subsection{The cutoff as a consequence of the two Gram matrices}\label{nav:IX-27}
For the two nondegenerate initial--terminal pairs write
\(C_i=(W_0+W_i)^\perp\) and set
\(G_{0,i}=|N_{F/\QQ}\mathcal G_i|/Q^{2df}\). The actual
projections of \(C_i\) into the initial and terminal positive
columns are isometries. Identify the two terminal columns
by their comparison. Write \(B\) for the difference between
the initial and terminal pairings of a vector from each
\(C_i\), after removing their common comparison subspace
of dimension \(k-s\). At an archimedean place \(\tau\) put
\(\kappa_{i,\tau}=1\) if the leading isotropic cross minor
\(D_{i,d-1}\) vanishes. Otherwise put
\[
 c_{i,\tau}=\frac{|\tau(D_{i,d}/D_{i,d-1})|}{T},\qquad
 \kappa_{i,\tau}=\frac{1+c_{i,\tau}}{\sqrt{c_{i,\tau}^2-1}},
                   \qquad K_i=\prod_\tau\kappa_{i,\tau},
\]
where \(D_{i,d}\) is the full cross determinant and the
product defining \(K_i\) uses either case at each place.

\begin{theorem}[The comparison determinant on the original interval]
\label{thm:cutoff224}
For \(2\le s\le k\) the rank of \(B\) is \(s-1\) or \(s-2\).
The comparison lattice \(M_\beta\) in Proposition~\ref{prop:comparison222}
and its two compatibility indices satisfy
\[
 \delta(M_\beta)i_1i_2\le \sqrt{G_{0,1}G_{0,2}}\,
 \begin{cases}
    (2/Q)^{(s-1)f}, & \operatorname{rank}B=s-1,\\
    (2/Q)^{(s-2)f}K_1K_2/T^f, &\operatorname{rank}B=s-2 .
 \end{cases}
\]
These give cutoffs in the genus bound for the distinct
common subspaces when \(2\le s\le k-1\).
\end{theorem}
\begin{proof}
In the product of the two positive columns the two
compatibility lattices are intersections of the two rational
graphs with the same integral lattice of pairs whose
residues agree at \(Q\). Their intersection on the common
graph is precisely \(M_\beta\). Orthogonal projection away
from it gives rank-\((s-1)\) lattices \(H_i\) with
\(\delta(\mathcal L_i)=\delta(M_\beta)\delta(H_i)\).
The difference of the two forms has norm at most two and
takes values in \(Q\mathcal O_E\). After removing the common
graph the two graphs are disjoint isotropic \((s-1)\)-planes in a space
of signature \((s,s)\). The radical on either side injects
into a space of signature \((1,1)\) and has dimension at most
one. On the quotient of rank \(n=\operatorname{rank}B\), the
determinant of the pairing takes values in \(Q^n\mathcal O_E\)
on exterior lattices. The archimedean norm bound thus gives
a lower bound \((Q/2)^{2nf}\) on the product of the two
projected lattice determinants.

If a radical line occurs, write \(R_i\) for its lattice on
side \(i\), \(I_i\) for the saturation index of
Theorem~\ref{thm:saturation220}. Against the opposite pair
\(a\), a radical vector gives two integral lifts with equal
perpendicular projection. The coefficient of the initial
terminal vector in their difference belongs to a fractional
ideal \(J\supset\mathcal O_E\) with
\([J:\mathcal O_E]\le I_a/i_a\). On the projected radical
its absolute value is
bounded by \(\kappa_{a,\tau}/\sqrt T\) times the local norm.
It is nonzero for every nonzero radical vector, as the two
normal components have equal length and the map to their
two terminal coefficients is an isometry of indefinite
planes. Multiplication gives
\(\delta(R_i)\ge T^f i_a/(I_aK_a^2)\). This cancels the
saturation factors in
\(\delta(\mathcal L_i)i_i^2=G_{0,i} i_i/I_i\).
Combining the determinant bounds proves the formula. Both
normal-line calculations, including the zero-minor case,
and the ideal and index computations are in
Section \ref{nav:IX-13}.
\end{proof}

This determines the cutoff to use in the genus bound. Theorem~\ref{thm:pairs226} below incorporates the number of ordered pairs over each primitive subspace.

\subsection{Counting the ordered flags above each subspace}\label{nav:IX-28}
Put \(t=d-1\), \(n_T=T^f\), and let \(Z\) be an integral
cross Gram matrix at the original level, modulo the two
integral unitriangular frame changes. Write \(C_Z\) for its
first \(t\) rows, \(J_Z=[\mathcal O_E^t:C_Z\mathcal O_E^d]\)
and \(G_Z=|N_{F/\QQ}\mathcal G|\). For the nondegenerate
frame sums considered here \(1\le J_Z\le G_Z\). Set
\[
     A_s(Z)=9^{sf}\max(1,G_Z/n_T)^s J_Z^{s+t}.
\]

\begin{theorem}[Ordered flag pairs with prescribed integral Gram matrices]
\label{thm:pairs226}
Fix the two integral cross classes \(Z_1,Z_2\) and impose
\(\delta(M_\beta)i_1i_2\le Y\), \(Y\ge1\). For \(1\le s\le k-1\),
\(j=k-s\), the normalized mean number of their ordered
comparison classes of common dimension \(j\) in the initial
row measure is at most
\[
   8\,\frac{\mathfrak m_j\mathfrak m_s}{\mathfrak m_k}
      (Y/n_T)^k\,n_T^{s+4t}Y^{\nu+1+2t} A_s(Z_1)A_s(Z_2),
          \qquad \nu=4+\log_2(2f).
\]
For \(2\le s\le k-1\) one may insert the corresponding
rank-dependent cutoff of Theorem~\ref{thm:cutoff224}. In
codimension one, fix a first flag with \(I_1=1\). The number
of partners of class \(Z_2\) with \(I_2=1\) is bounded per
row by \((n_T^2G_{Z_1}/Q^{2df})^t A_1(Z_2)\).
\end{theorem}
\begin{proof}
In an integral Witt splitting of the initial row, write
the columns of a terminal frame in upper, middle and lower
blocks. The lower block is \(C_Z\). If a prescribed
rank-\(j\) sublattice \(P\subset L\) lifts to the integral
frame perpendicular, its lift differs from the initial
column embedding by \(t\) integral functionals on \(P\).
Integral unipotent shears leave at most
\((n_T^2[L\cap EP:P])^t\) inequivalent choices of these
functionals. For fixed choices, contract the middle block
with the cofactor vector of \(C_Z\), which spans its rational
kernel. Projecting off the terminal vector gives a vector
in an \(s\)-dimensional positive space, with prescribed
squared norm \(\mathcal G/T\). Differences between possible
values belong to an integral lattice. Comparing disjoint
balls in the product over the embeddings gives at most
\(9^{sf}\max(1,G_Z/n_T)^s\) values. After this choice, the
middle-block differences factor through \(C_Z\), giving at
most \(J_Z^s\) classes by restriction of integral
homomorphisms. The remaining upper-block differences have
skew-adjoint coefficient and give at most \(J_Z^t\) classes.
This proves the single-flag bound with factor \(A_s(Z)\).

In the common subspace put \(S=L\cap EM_\beta\), identifying
the comparison lattice with its initial image. Let
\(a=[S:M_\beta]\), and let \(b_i\) be the index in \(M_\beta\)
of the sublattice lifting to the \(i\)-th integral
perpendicular. Then \(b_i\le i_i\) and
\(\delta(S) a b_1 b_2\le Y\). The number of sublattices of
index \(n\) in a projective rank-\(j\) lattice is at most
\(n^{j+\log_2(2f)}\), by successively taking simple
quotients at a prime dividing the remaining index. Counting
the three sublattices this way, using the single-flag bound
twice and Theorem~\ref{thm:genus222} for \(S\), leaves the
summable series \(\sum_{a,b_1,b_2}(ab_1b_2)^{-s-5}<8\).
Uniqueness of the comparison persists after independent
changes of terminal frame because both changes act trivially
on the marked middle quotient. For \(s=1\) the two rational
graphs coincide. When both flags satisfy \(I_i=1\), the
first integral perpendicular projects to a fixed sublattice
of determinant \(G_{Z_1}/Q^{2df}\) lifting to the second,
so one application of the single-flag count suffices. The
coordinate and measure arguments are proved in
Sections \ref{nav:IX-14} and \ref{nav:IX-15}.
\end{proof}

For codimension one the terminal prime also puts a restriction on the two Gram classes that can occur together.

\subsection{Which codimension-one pairs have to be summed}\label{nav:IX-29}
For the two Gram classes in Theorem~\ref{thm:pairs226}, put
\(D_i=\det Z_i\), let \(R_i\) be the leading minor of size
\(d-1\), and write \(\mathcal G_i=D_i\bar D_i-T^2R_i\bar R_i\).
Assume the two frame sums are nondegenerate, as in that
theorem. Then \(\mathcal G_i\) is totally positive.

\begin{theorem}[Separation at the original terminal prime]
\label{thm:normal228}
Suppose a comparison of common dimension \(k-1\) occurs
on a row satisfying \(I_1=I_2=1\). Either \(R_1/D_1=R_2/D_2\)
in \(E\), or at one real place and for one of the two flags,
\[
     R_i\ne0,\qquad
     \frac{\tau\mathcal G_i}{T^2|\tau R_i|^2}
                  \le \frac{2}{\sqrt{1+(T/Q)^2}-1}.
\]
All quantities refer to the field \(E_x\) and the primes
selected on the interval in \eqref{eq:region198}.
\end{theorem}
\begin{proof}
The two positive comparison graphs have dimension \(k-1\)
and coincide. Their perpendiculars in the initial and terminal
columns are lines \(N,N'\). For either flag pair, lift a
vector from each line to the corresponding frame perpendicular.
The coefficients of the difference along the two terminal
vectors give an isomorphism \(A_i:N\oplus N'\to E^2\). The
difference of the two squared lengths is \(T\) times the
difference of the squared coefficient magnitudes. Thus
\(\gamma=A_2A_1^{-1}\) is unitary for \(\operatorname{diag}(1,-1)\).
The integral congruence in
Theorem \ref{thm:congruence227}
gives \((\gamma-1)\mathcal Z\subset T\mathcal Z\), where
\(\mathcal Z=\{(u,v)\in\mathcal O_E^2:u+v\in Q\mathcal O_E\}\).
Its proof extends a functional with values \(u,-v\) on the
two terminal vectors to the same self-dual row and compares
its two integral lifts at the marked primes.

If \(c=\gamma_{21}\) vanishes, \(\gamma\) is diagonal.
Applying the congruence to \((1,-1)\) makes each diagonal
entry congruent to \(1\) modulo \(T\). Its absolute value
is one at every embedding, so the norm of its difference
from one forces equality as \(T>2\). In the common initial
normal line the signed cofactor vectors of the two frames
have coefficient pairs \((D_i/T,R_i)\). Since \(A_1=A_2\),
these pairs are proportional over \(E\).

If \(c\ne0\), apply the congruence to \((Q,0)\). The nonzero
integer \(Qc/T\) of \(E\) gives
\(\prod_\tau|\tau c|\ge(T/Q)^f\). In unit bases, after
scaling the output by \(\sqrt T\), the diagonal entry
magnitudes of \(A_i\) are \(\sqrt{1+V_{i,\tau}^2}\) and
the off-diagonal magnitudes are
\(V_{i,\tau}=T|\tau R_i|/\sqrt{\tau\mathcal G_i}\).
These follow by dividing the cofactor vector by its length
\(\sqrt{\tau\mathcal G_i/T}\). The product \(A_2 A_1^{-1}\)
therefore gives \(|\tau c|\le2H\sqrt{1+H^2}\) with
\(H=\max_{i,\tau}V_{i,\tau}\). Taking the \(f\)-th root
and solving for \(H^2\) proves the estimate. The same
argument, with the annihilators of the two saturation
quotients retained, is given in
Theorem \ref{thm:separation227}.
\end{proof}

The fixed-class partner bound can consequently be set to
zero whenever both alternatives fail. The sum over the
surviving classes, the designated integral Bessel pairing
in Corollary~\ref{cor:original200}, and the remaining
cohomology orders in \eqref{eq:region198} still enter the
proof of the main theorem.

\subsection{Primitive sublattices in the initial column}\label{nav:IX-30}
Work on the common intervals with polynomial conductor \(N_x\).
Keep the terminal prime \(T\) and the two level primes whose
product is \(Q\). Write \(f=[F:\QQ]\), \(n=T^f\) and
\(R=|\operatorname{disc} E|^{1/(2f)}\). For an integral
Hermitian lattice \(B\) write \(\delta(B)=[B^\vee:B]^{1/2}\).
The initial column of a row with first flag \(X+Ez\) contains
the lattice \(L\) of Section~\ref{sec:termmean221}. We select a
set of initial columns once, before specifying the second flag.

\begin{theorem}[Simultaneous lower bounds for primitive determinants]
\leavevmode\label{thm:columns235}
Fix \(A>0\). There is a constant \(C_A\) such that, for
sufficiently large \(x\) on the common intervals, a fraction at
least \(1-\exp(-Af(k-1))\) of the original refined class
coordinates has the following property. For every \(1\le j<k\)
and every primitive rank-\(j\) sublattice \(B\subset L\),
\[
 \delta(B)>Y_j,\qquad
 Y_j=n^{j/k}(kR/C_A)^{fj(k-j)/k}.
\]
On the same coordinates every primitive hyperplane also satisfies
\[
        \delta(B)>Y,\qquad
        Y=n\bigl(kR/(2\pi\exp(A+4))\bigr)^f .
\]
These coordinates form a union of full terminal scalar orbits.
The retained fraction is the same for every word, for each
specified terminal character and for the space with zero sums in
the individual scalar coordinates.
\end{theorem}
\begin{proof}
Put \(s=k-j\) and
\(\Xi_j=\mathfrak m_j\mathfrak m_s/\mathfrak m_k\).
Sum \eqref{eq:termupper221} over the positive integers \(D\le y\).
The probability of even one primitive sublattice of this rank
and determinant at most \(y\ge1\) is at most
\[
           \Xi_j(y/n)^k n^s y^{\nu+1},
                  \qquad \nu=4+\log_2(2f).
\]
For \(2\le j,s\), Proposition~\ref{prop:massdecay221} bounds
\(\Xi_j\) by \(\mathfrak m_2(C/(kR))^{fjs}\), with
\(\log^+\mathfrak m_2=O(f\log(2N_x))\). For the two end ranks,
\(\Xi_1=\Xi_{k-1}=\mathfrak m_1/a_k\). The mass formula in the
same proof gives \(\mathfrak m_1\le 2^{1-f}R^f\) and
\[
 a_k\ge R^{f(k-1/2)}\bigl((k-1)!/(2\pi)^k\bigr)^f\zeta(k)^{-f}.
\]
Stirling's formula gives
\(\log\Xi_1\le-fk\log(kR/(2\pi e))+O(f\log(2kR))\).
Thus for all ranks \(\log\Xi_j\le-fjs\log(kR/C)
+O(f\log(2kN_x))\) after enlarging \(C\).

At \(y=Y_j\) the factor \(n\) cancels. The remaining error is
uniformly \(o(fjs)\). Indeed \(js\ge k-1\), \(R\le N_x\),
and
\[
 \frac{\log Y_j}{fjs}
       =\frac{\log T}{ks}+\frac{\log(kR/C_A)}{k}.
\]
Here \(\nu=O(\log x)\), \(\log N_x=O(\log x)\),
\(k\asymp x^{9/2+10\eps}\), and
\(\log T=(17+o(1))m x^\eps\) with
\(m\asymp x^{1/2-\eps/16}\) by Theorem~\ref{thm:prime209}.
Choose \(\log(C_A/C)\ge A+4\). For the additional hyperplane
test substitute \(y=Y\) and \(s=1\) in the same probability
bound. Its logarithm is at most
\[
 -(A+3)fk+O\bigl((\nu+2)f\{\log T+\log(2kR)\}\bigr)
                 =-(A+3)fk+o(fk).
\]
Adding the \(k\) failure probabilities now gives the assertion.

For the counting measure, Lemma~\ref{lem:termparent221} shows
that forgetting all markings pushes the normalized mass to
the terminal genus measure. At the two auxiliary levels the
class stabilizers are trivial by Corollary~\ref{cor:rank111};
each refined class therefore carries the same mass, and its
space of real covectors has the same dimension for a fixed word.
Multiplication of a terminal covector by a nonzero residue scalar
fixes its kernel, hence fixes \(L\). Each retained scalar orbit
contributes one copy of every character and the same number of
vectors to the zero-sum space. This proves the fraction assertions.
\end{proof}

\begin{corollary}[A necessary condition on an original orbital cell]
\leavevmode\label{cor:gram235}
Take any retained initial column. For a relevant original finite
cell with nonzero Gram determinant put \(G=|N_{F/\QQ}\mathcal G|\)
and let \(I=[B:B^+]\) be its congruence index in
Proposition~\ref{prop:index219}. Then
\[
                       G>Q^{2df} IY,\qquad d=m/2 .
\]
The exclusion is simultaneous for every second flag above this
column, including those with \(I>1\).
\end{corollary}
\begin{proof}
The common positive perpendicular \(C\) has rank \(k-1\).
For these labels the two frames have disjoint isotropic parts:
their full cross minor is nonzero by \(\mathcal G>0\).
Any common isotropic vector would give a nontrivial left kernel
of the cross matrix. The projection to the first definite
column is therefore isometric on \(C\). It sends
\(\Lambda\cap C\) into \(L\) by the definition of the column
lattice. Its saturation in \(L\) is a primitive hyperplane,
and passing to this saturation divides its determinant by the
positive integral lattice index. Theorem~\ref{thm:saturation220}
gives \(\delta(\Lambda\cap C)=G/(Q^{2df}I)\).
The second inequality of Theorem~\ref{thm:columns235} now proves
the claim.
\end{proof}

\subsection{Full-rank cells with transverse isotropic subspaces}\label{nav:IX-31}
In the first moment for an identical word,
Theorem~\ref{thm:common133} permits a full-rank common definite
space. Take \(s=d-1\ge1\) and let \(\mathcal S\) be the
relevant, finitely supported labels of this kind for which
\(X_0\cap X_1=0\). We count each rational orbit once.

\begin{proposition}[The relative negative line and the isotropic blocks]
\leavevmode\label{prop:class236}
A label in \(\mathcal S\) is determined by a composition of
\(s\), its nonzero reversed-block scalars, and an element
\(\lambda\in E^\times\) with \(\lambda\bar\lambda=1\).
If \(a=D_s/Q^s\), then \(a\in\mathcal O_E\setminus\{0\}\) and
\[
         D_d=-T\lambda D_s,\qquad
         Ta(1-\lambda)\in Q\mathcal O_E .
\]
For fixed block scalars, the number of possible \(\lambda\) is
at most \(3^f|\mu(E)|\,\tau_E((a))\), where
\(\tau_E\) counts the integral ideal divisors.
\end{proposition}
\begin{proof}
By Theorem~\ref{thm:cross132}, \(Y_0\cap Y_1\) is a negative
line over \(E\). Projection identifies it with \(Y_0/X_0\).
Choose its unique vector \(z_*\) inducing the original final
vector there. The analogous vector for \(Y_1\) is
\(\lambda z_*\); both have norm \(-T\). These replacements
of the terminal representatives add vectors of the corresponding
isotropic flags and preserve their decorations. Each \(X_i\)
is perpendicular to \(z_*\). The cross pairing on
\(X_0\times X_1\) is nonsingular by the kernel argument in
Proposition~\ref{prop:denom235}. The full cross matrix now has
this \(s\)-square in its upper left and the scalar \(-T\lambda\)
in its lower right, with both rectangular corners zero.

Independent upper unitriangular changes of isotropic
representatives reduce the \(s\)-square to a monomial matrix
by successive pivot elimination as in Lemma~\ref{lem:forms203}.
For an inversion \(i<j\), \(w(i)>w(j)\), write \(b_i,b_j\)
for its pivot entries. At a real place the two elementary
changes with respective entries \(u\) in position \(ij\) and
\(-\bar u b_i/b_j\) in position \(w(j),w(i)\) preserve the
cross matrix. They give an isometry of \(X_0+X_1\), which is
nondegenerate. Extend it by the identity on the negative line
and on the definite perpendicular. This element belongs to both
unipotent radicals: it acts unitriangularly on the two flags
and trivially on their definite quotients. Equality of the
characters for every complex \(u\) implies that every adjacent
descent decreases by one and that the two adjacent pivots there
are negatives of each other. Thus the descending runs are
disjoint consecutive intervals in increasing order, with the
successive entries in each run alternating in sign. Their
nonzero leading minors are exactly the endpoints, and a run
of length \(h\) multiplies the preceding endpoint by \(c^h\)
in the alternating convention.

These data determine the Gram matrix on a basis consisting of
both isotropic lists and \(z_*\). Their span is nondegenerate.
If two pairs give the same data, the isometry of these spans
carrying one basis to the other extends to the ambient Hermitian
space by Witt's theorem. It carries both decorated frames to the
specified frames because their terminal representatives are
\(z_*,\lambda z_*\). This proves the required injectivity on
rational orbits. The finite congruences of
Proposition~\ref{prop:shift203} apply to every supported pair
and give both assertions about \(a\).

Now \(b=Ta\lambda\) is integral and \(b\bar b=T^2a\bar a\).
At a prime of \(E\) fixed by conjugation the exponent of
\((b)\) is determined. At a split pair the two exponents are
nonnegative and have a fixed sum \(e+2v(T)\), where
\(e=v_{\mathfrak p}(a)+v_{\bar{\mathfrak p}}(a)\).
There are \(e+1\) choices off \(T\) and at most \(3(e+1)\)
over \(T\), since \(T\) splits completely. Each \(e+1\) is
at most the product of the two individual divisor counts of
\((a)\). Thus there are at most \(3^f\tau_E((a))\) possibilities
for \((b)\). Two generators satisfying the norm equation for
the same ideal have ratio an integral unit of absolute value
one at every complex embedding. Such a unit is a root of
unity: the monic polynomials of its powers have bounded
integral coefficients and bounded degrees, so the powers take
only finitely many values. Counting these units gives the claim.
\end{proof}

\begin{theorem}[Weighted count of the transverse full-rank cells at the original level]
\leavevmode\label{thm:trans236}
Use \(w_\gamma=\prod_{j:D_j\ne0}\prod_\tau
(1+|\tau D_j|^2)^2\) and put
\(H_a=|N_{E/\QQ}a|^{1/(2f)}\). For \(H_0\ge1\) and \(0\le u<2\),
\[
 \sum_{\substack{\gamma\in\mathcal S\\H_a\ge H_0}}w_\gamma^{-1}
 \le 3^f|\mu(E)|\,H_0^{-2f}(9M_u+2)^{fs}
                          T^{-4f} Q^{-(6+u)sf}.
\]
At \(H_0=1\) the bound is
\(\exp(-(264-o(1))fd\,x^\eps)\). It holds on the entire
original class space, in particular for every retained
terminal character space.
\end{theorem}
\begin{proof}
All \(|\tau\lambda|=1\), so for fixed isotropic blocks the
product of terminal factors is at most
\[
 \prod_\tau(T|\tau D_s|)^{-4}
             =T^{-4f}Q^{-4sf}|N_{E/\QQ}a|^{-2}.
\]
For an integral ideal, its divisor count is at most its norm,
since \(e+1\le2^e\le(N\mathfrak p)^e\) at every prime of
exponent \(e\). Proposition~\ref{prop:class236} therefore
bounds the sum of the terminal factors by
\(3^f|\mu(E)|T^{-4f}Q^{-4sf}H_0^{-2f}\) in the indicated
range. For the isotropic endpoints \(j\le s\), the minors
are nonzero integers divisible by \(Q^j\) in \(\mathcal O_E\).
The product-distance packing and radial integration of
Theorem~\ref{thm:levelsum203} at each such endpoint yield
\((9M_u)^fQ^{-(2+u)jf}\). A prescribed endpoint tuple has
at most the product of its block lengths as its scalar
multiplicity by the equations of Proposition~\ref{prop:class236}.
Summing over compositions as in that theorem and retaining the
power at the last isotropic endpoint \(s\) gives at most
\((9M_u+2)^{fs}Q^{-(2+u)sf}\). Multiplication proves the bound.

If \(w=|\mu(E)|\), then \(\varphi(w)\le2f\) and
\(w/\varphi(w)^2\le2\), by multiplying
\(p^{2-e}/(p-1)^2\) over prime powers \(p^e\Vert w\).
Hence \(w\le8f^2\). Choose \(u=2-1/\log(2Q)\).
The factors preceding \(T^{-4f}Q^{-(6+u)sf}\) at \(H_0=1\)
have logarithm \(o(fd\,x^\eps)\). Using \(s=d-1\),
\(\log Q=(16+o(1))x^\eps\) and
\(\log T=(34+o(1))d x^\eps\) gives \(136+128=264\).
Restricting to coordinates or to labels meeting a character
space can only decrease the positive label sum.
\end{proof}

\begin{corollary}[A restriction on the terminal poles of the relative scalar]
\leavevmode\label{cor:poles236}
For a label as above with \(\lambda\ne1\), let \(l\) count the
primes \(\mathfrak t\mid T\) of \(E\) at which
\(a(1-\lambda)\) is nonintegral. Then
\[
           0\le l\le f,\qquad
           T^l(2H_a)^{2f}\ge Q^{2f}.
\]
For every fixed \(B\), \(H_a\le x^B\) forces
\(l\ge(32/17-o(1))f/m\) on the common intervals. If \(\lambda\)
is integral and different from one, then \(a\in Q\mathcal O_E\)
for all sufficiently large \(x\).
\end{corollary}
\begin{proof}
The congruence \(Ta(1-\lambda)\in Q\mathcal O_E\) implies
integrality off \(T\), divisibility by \(Q\) there, and poles
of order at most one over \(T\). Since \(\lambda\bar\lambda=1\)
and \(a\) is integral, a pole can occur on at most one member
of each split pair. Multiplying the principal fractional ideal
\((a(1-\lambda))\) by those \(l\) primes gives a nonzero integral
ideal divisible by \(Q\mathcal O_E\). Its norm is at most
\(T^l 2^{2f}|N_{E/\QQ}a|\) by the archimedean triangle
inequality. This proves the first assertion. Take logarithms
and use \(\log H_a=O(\log x)\) for the second.
An integral norm-one \(\lambda\) is a root of unity by the
argument in Proposition~\ref{prop:class236}. The primes
dividing \(Q\) exceed \(8f^2\) eventually, and hence do not
divide its order. Reduction preserves the order of a root
of unity of prime-to-residue-characteristic order: the
principal units have no torsion of that order by the binomial
expansion. Thus \(1-\lambda\) is a unit above \(Q\), proving
the last assertion.
\end{proof}

In Proposition~\ref{prop:denom235} the denominator of the
rational action on a transverse cell has root norm at most
\(H_a\). Small \(H_a\) thus permits the normalized character
estimate, while the complementary range has the factor
\(H_0^{-2f}\) above. The full signed first moment still
contains the orbital amplitudes and cells with nonzero
isotropic intersection.

\subsection{The weighted sum over all relevant rational labels}\label{nav:IX-32}
Let \(\mathcal S_t\) consist of the relevant, finitely supported
labels for which \(\dim(X_0\cap X_1)=s-t>0\), \(1\le t<s\).
Let \(\mathcal R_*\) consist of all relevant rational labels
outside the stabilizer of \(X_0\) that have an original finite
cell on a retained initial column of Theorem~\ref{thm:columns235}.
These sets count each rational double coset once. For every label
use the same weight
\(w_\gamma=\prod_{1\le j\le d:D_j\ne0}\prod_\tau
(1+|\tau D_j|^2)^2\).

\begin{theorem}[The complete arithmetic weight and its leading scale]
\leavevmode\label{thm:sum237}
For \(1\le u<2\),
\[
 \sum_{\gamma\in\mathcal S_t}w_\gamma^{-1}
       \le |\mu(E)|16^f(9M_u+2)^{ft}Q^{-(2+u)ft}.
\]
On restricting to \(\lambda\ne1\) on these labels the right hand
side can also be multiplied by \((2/Q)^{(u-1)f}\). On the common
fields and intervals the sum over \(\mathcal R_*\) satisfies
\[
       \log\left(\sum_{\gamma\in\mathcal R_*} w_\gamma^{-1}\right)
                      =-(64+o(1))f x^\eps .
\]
\end{theorem}
\begin{proof}
For a fixed set of prefix blocks, put \(a=D_t/Q^t\) and
\(b=a\lambda\). By Proposition~\ref{prop:finite237} the element
\(b\) is integral and \(b\bar b=a\bar a\). The split-prime
divisor count of Proposition~\ref{prop:class236}, now without
the terminal factor \(T\), gives at most \(\tau_E((a))\)
possible ideals \((b)\). Each admits at most \(|\mu(E)|\)
such generators. By Theorem~\ref{thm:suffix237} is
at most one rational label for given block scalars and
\(\lambda\), and all nonzero endpoint weights are in the prefix.

For any integral ideal \(\mathfrak a\) one has
\(\tau_E(\mathfrak a)\le16^f(N\mathfrak a)^{1/2}\).
Indeed \(e+1\le(N\mathfrak p)^{e/2}\) for prime ideals of norm
at least four. For norm two or three one has
\(e+1\le2(N\mathfrak p)^{e/2}\); there are at most \(4f\)
such primes over the two rational primes. Multiply these
inequalities for the ideal factorization.

Count prefix endpoint values in boxes of geometric mean radii
\(R_j\) as in Theorem~\ref{thm:levelsum203}. At endpoint \(j\)
the count conditional on earlier values is bounded by
\(9^f(R_j/Q^j)^{2f}\mathbf1_{R_j\ge Q^j}\).
At the last endpoint \(t\), the preceding divisor estimate
adds at most \(16^f(R_t/Q^t)^f\), since
\(|N_{E/\QQ}a|^{1/2}\le(R_t/Q^t)^f\).
Use exponent \(u-1\) on the last support indicator and \(u\)
on the others. The total power of every \(R_j/Q^j\) is then
\((2+u)f\). When \(\lambda\ne1\) the final box must satisfy
\(R_t/Q^t\ge Q/2\) by Proposition~\ref{prop:finite237},
which supplies the extra factor in the statement.
Integrating against the radial densities of
Theorem~\ref{thm:levelsum203} gives \((9M_u)^f Q^{-(2+u)jf}\)
per endpoint. The equations for the prefix scalars give
multiplicity at most the product of the block lengths.
The binomial sum over compositions, keeping the power of
\(Q^{-1}\) from the last endpoint \(t\), proves both bounds.

For the total sum take \(u=2-1/\log(2Q)\). Write
\(D=(9M_u+2)^fQ^{-(2+u)f}\). Then
\(\log D=-(64+o(1))f x^\eps\), and
\(\sum_{t=1}^{s-1}D^t\le D/(1-D)\) for large \(x\).
The factor \(|\mu(E)|16^f\) has logarithm \(o(fx^\eps)\)
by \(|\mu(E)|\le8f^2\). This covers every nontransverse label
outside the stabilizer. For transverse labels of zero Gram
determinant, Theorem~\ref{thm:cross132} gives the full-rank
case of Theorem~\ref{thm:trans236}, whose bound is
\(\exp(-(264-o(1))fd\,x^\eps)\). For all nonzero Gram
determinants on the retained columns, Theorem~\ref{thm:powers236}
gives \(\exp(-(132-o(1))fd\,x^\eps)\). As \(d\to\infty\),
these two contributions preserve the upper bound.

For the lower bound, Proposition~\ref{prop:sharp237} puts
the same label of weight \((1+16Q^2q)^{-2f}\) on every
refined column, hence on a retained one. Here \(q\le N_x\)
and \(N_x\) is polynomial on the interval. Therefore its
logarithm is \(-4f\log Q+O(f\log(2q))=-(64+o(1))f x^\eps\).
This proves the two-sided assertion.
\end{proof}

The rank-one example also gives equal oriented coordinates
in the individual-character trace and the identity on the
definite real space. Consequently a bound for the signed
period needs information on the orbital integrals already at
this label. The weighted count controls the arithmetic
multiplicity of labels, while the finite cell measures and
the prescribed real functions determine their contributions.

\Needspace{7\baselineskip}

\section{Terminal frequencies and stabilizers}\label{nav:g-XI-1}

\subsection{Orthogonality for all flags in a fixed isotropic subspace}\label{nav:VI-28}
Let \(F_+=E^+\), let \(L=\langle u_1,\ldots,u_s\rangle_E\), and let
\(Q\subset G\) be the parabolic subgroup stabilizing \(L\). Its
rational elements may change all the smaller spaces of the original
flag. Write \(\mathbb A=\mathbb A_{F_+}\). The seed \(\Phi\) and
its norm on \(B(F_+)\backslash G(\mathbb A)\) are those of
Sections \ref{nav:VI-15} and \ref{nav:VI-16}.
The row levels at \(g_1,g_2\) have unipotent reduction by
Proposition~\ref{prop:index90}.

\begin{theorem}[Summation over the stabilizer of the isotropic subspace]
\label{thm:maximal130}
For the fields, words, frequencies and terminal levels of
Theorem~\ref{thm:commonnorm128}, summation of the seeds over
\(B(F_+)\backslash Q(F_+)\) defines an isometry from their span
in \(L^2(B(F_+)\backslash G(\mathbb A))\) to
\(L^2(Q(F_+)\backslash G(\mathbb A))\), onto its image.
For seeds of compact support modulo \(B(F_+)\) the sum is the
ordinary locally finite series. For the prescribed square-integrable
seeds it is the limit in \(L^2\) of this operation on compact
truncations. Thus all selected columns have the same positive
factor \(C_K\) in their probability norms after this summation,
and the trace estimate \eqref{eq:tr128} holds for its full Gram
operator on the refined zero-sum space.
\end{theorem}
\begin{proof}
Let \(N_*\) be the unipotent radical of \(Q\).
The Levi quotient acts by
\(\operatorname{GL}_E(L)\times U(W)\), using the chosen Witt splitting. The
unipotent group for the full flag is \(N_*\rtimes U_s\), where
\(U_s\) consists of upper unitriangular transformations on \(L\).
Matrices on the dual isotropic space are determined by these
actions. The subgroup \(B\) adds the compact-at-infinity group
\(H=U(W_0)\) fixing \(e\in W\).

First suppose the real-place functions have compact support
modulo \(B_\infty\). Then the seeds have compact support modulo
\(B(F_+)\). Indeed the finite support reduces to finitely many
classes \(h_iK\), with compact finite part. Above each class the
remaining fibre from the real period group is
\(\Lambda_i\backslash B_\infty\). Its unipotent factor is the
quotient of the compact adelic unipotent quotient by the compact
open subgroup in Lemma \ref{lem:covol127}. Its other factor is the
compact \(H_\infty\).
A compact part of \(B_\infty\backslash G_\infty\) can be covered
by finitely many relatively compact local sections. This gives
the assertion. For any intermediate rational subgroup the ordinary
periodization is now locally finite. To see this, lift the compact
support to a compact set of representatives in \(G(\mathbb A)\),
using local sections for its discrete rational group. On a compact
test set in \(G(\mathbb A)\), contributing rational cosets have a
representative in the product of this compact lift and the inverse
test set. A closed discrete rational group meets that product in
only finitely many elements. The resulting function has compact
support on the intermediate quotient as the image of a compact set.

We use the following integration argument at each step. Suppose
\(J'\subset J\) are consecutive groups and \(D_0\subset J'\) is
unipotent and normal in \(J\). A summand at index \(\gamma\)
and frequency \(\beta\) transforms along
\([D_0]=D_0(F_+)\backslash D_0(\mathbb A)\) by the character of
the input conjugated by \(\gamma\). If two simultaneously supported
terms of different index or frequency have different characters,
their products integrate to zero on this compact quotient by
right translation. The squared sum then integrates to the sum of
the diagonal squares, after collecting the inputs of each frequency.
Integration over \(J(F_+)D_0(\mathbb A)\backslash G(\mathbb A)\)
and unfolding the diagonal over \(J'(F_+)\backslash J(F_+)\)
gives the sum of the input squared norms. For the groups below
\(J(F_+)D_0(\mathbb A)\) is closed as the inverse image in
\(J(\mathbb A)\) of the rational points of its displayed Levi
quotient. It is unimodular by the product formula for the rational
determinant on \(\operatorname{Lie}(D_0)\). These facts give the
invariant quotient integration formula with fibre \([D_0]\).
The sums are locally finite on the compact fibre. The nonnegative
diagonal integrals are finite by unfolding, so this calculation
also proves square integrability of the output. Inputs at the same
frequency may have any column types, since their sum transforms by
that one unipotent character.

For \(0\le t\le s-1\), put
\[
 M_t=\left\{\begin{pmatrix}A&Y\\0&u\end{pmatrix}
    :A\in\operatorname{GL}_t(E),\ u\in U_{s-t}(E)\right\}\subset\operatorname{GL}_s(E)
\]
in rational coordinates, with \(M_0=U_s(E)\). Use the corresponding
algebraic subgroups to set \(J_t=N_*\rtimes(M_t\times H)\), so
\(J_0=B\). At the step \(J_{t-1}\subset J_t\), \(1\le t<s\),
take \(D_t\) to be \(N_*\) times the unipotent block part with
\(A=1\), with identity \(H\)-component. It is normal in \(J_t\)
and contained in \(J_{t-1}\). Its character is the restriction of
\(\chi_\beta\). Every element of \(J_{t-1}\) preserves this
restriction. The upper unitriangular part acts by inner conjugation
in the full flag unipotent group, \(H\) fixes the character, and
the block \(\operatorname{GL}_{t-1}\) fixes the last coordinate of a column of
length \(t\). The surviving simple coordinates of \(D_t\) in the
upper matrix start at position \((t,t+1)\).
The terminal character on \(N_*\) is also fixed, since \(M_{t-1}\)
fixes the last coordinate covector on \(L\). In particular the
sum formed up to \(J_{t-1}\) still transforms by
\(\chi_\beta|_{D_t}\), as is seen term by term.

Now an element of \(J_t\) with top block \(A\) can preserve that
character only if the last row of \(A\) is the last coordinate
covector. Test the unipotents with arbitrary column in rows
\(1,\ldots,t\) of column \(t+1\). The character reads its
last entry, and conjugation by the diagonal block replaces the
column by its product with \(A\) on the left. The unipotent
factor and \(H\) already preserve the character. Nondegeneracy
of the trace pairing gives the asserted row equality. It says
exactly that the element belongs to \(J_{t-1}\).
Also \(J_t\) fixes the terminal character on \(N_*\), so two
different frequencies remain distinct there. The integration
argument proves isometry at each of these steps.

For the last inclusion \(J_{s-1}\subset Q\) use \(D_0=N_*\).
This includes \(s=1\), when there were no preceding steps.
Let \(b_s\) be the last coordinate covector on \(L\). On the
block \(X:W\to L\) the character is the trace character
with coefficient \(\beta b_s X e\).
All linear maps \(X\) can be tested at a complex place. Indeed
the unipotent isometry in Lemma \ref{lem:transport129} realizes any
one row by using its corresponding hyperbolic plane. Products
add their \(W\)-to-\(L\) blocks. Other blocks introduced by
multiplication have zero character. Every element of \(J_{s-1}\)
fixes \(b_s\) and \(e\) in the Levi, so the intermediate input
has the stated \(N_*\)-character.

Suppose two terms at indices \(\gamma,\gamma'\) have the same
character and meet at the same point. Write \((A,M)\) for the
Levi part of \(\lambda=\gamma\gamma'^{-1}\).
Test all \(X\) in the equality of characters. It gives the equality
of nonzero simple tensors
\[
       \beta(b_s A)\otimes M^{-1}e=\beta' b_s\otimes e .
\]
Indeed conjugation sends \(X\) to \(AXM^{-1}\), and any
unipotent component acts trivially on the character. Hence
\(b_s A=a b_s\), \(\beta a M^{-1}e=\beta'e\), for some
\(a\in E^\times\). These identities hold over \(E\) by injectivity
of its complex embedding. In particular
\(a\bar a=\beta'\bar\beta'/(\beta\bar\beta)\).

The finite support of the intermediate sum is contained in
\(J_{s-1}(\mathbb A_f)K\). Consequently
\(\lambda_f=b_1 k b_2\) with \(b_i\in J_{s-1}(\mathbb A_f)\),
\(k\in K\). All these elements preserve \(L\). The \(b_i\)
preserve \(\ker b_s\subset L\) and act trivially on its line
quotient. The equation for \(b_s A\) shows that \(k\) also
preserves that hyperplane and acts on the quotient by \(a\).
Compactness gives that \(a\) is a unit everywhere at finite
places. The norms \(1,T\) give \(\beta=\beta'\). As in Lemma \ref{lem:transport129}, \(a\bar a=1\) makes this unit a root of unity,
and the unipotent reductions at the two distinct auxiliary
primes give \(a=1\). Thus \(M\in H(F_+)\) and \(A\in M_{s-1}\).
We have proved \(\lambda\in J_{s-1}(F_+)\), which gives equal
rational indices. This proves the required orthogonality at the
last step. The successive locally finite sums are the sum over
\(B(F_+)\backslash Q(F_+)\), by partition into cosets.

For general \(W\) truncate on an increasing compact exhaustion
of \(B_\infty\backslash G_\infty\), using scalar cutoffs bounded
by one and equal to one on increasing compact sets. The input
seeds converge in \(L^2\) by the quotient norm calculation of
Theorem \ref{thm:identity127} and dominated convergence. Differences
of compact truncations satisfy the same isometry, for which
the real vectors need no unit normalization. Hence the output
is Cauchy in \(L^2\) and its limit is independent of the cutoffs.
The single-frequency norm equalities and mixed orthogonality,
followed by polarization within a frequency, give the isometry
on the full span. The final assertions follow from
Theorem~\ref{thm:commonnorm128} and Corollary~\ref{cor:cost128}.
\end{proof}

The remaining periodization is over \(Q(F_+)\backslash G(F_+)\).
Distinct indices here correspond to distinct rational isotropic
\(s\)-dimensional subspaces. The preceding result evaluates all
the contributions with the same such subspace. It leaves open
the cancellation between distinct subspaces for the prescribed
theta vectors and their comparison on the integral lattice.

\Needspace{14\baselineskip}

\subsection{The global Gram matrix at the terminal prime}\label{nav:VII-29}
Use the fields \(E=B_J(\sqrt{-q},\sqrt{-d_0})\) and intervals of
Theorem~\ref{thm:invar148}. In this paragraph \(r\) is the row rank
and \(f=[E^+:\QQ]\). In \eqref{eq:letters111} choose
\begin{gather*}
 L=8(f+1)(r+1)^3,\qquad M=L+8r^2+4,\qquad \Delta=M+1,\\
 H_0=\max\{M+1,\lceil k+m(t+3)+r^{19}f^5x^\eps\rceil\}.
\end{gather*}
These spacings are polynomial in \(x\). The smallest \(B_\tau\)
in the formula for \(b_{j,\tau}\) exceeds every fixed power of
\(x\) eventually. The row therefore has gaps greater than \(M\)
at every real place. The prime and word construction of
Theorem~\ref{thm:common111} applies. The choice changes the row
and column weights within that construction. It uses the same
fields and the same \(x\) as the cohomology estimate.

\begin{theorem}[Absolute convergence at the prescribed finite level]
\label{thm:ordinary164}
For these parameters, the untruncated ordinary Bessel
periodizations on the normalized inputs \(J_U(U\otimes D)\)
converge absolutely everywhere at \(K^0\). The resulting
\(L^2\) vectors already lie in the prescribed real row and its
lowest compact type. The Gram matrix on \(D\) equals
\(G_D=AA^*\) of Theorem~\ref{thm:global158}. Its difference
from the identity is the absolutely convergent sum of the
orbital integrals outside the rational flag parabolic. The
zero-determinant cells in the blocks between distinct words
contribute zero. In addition the height moment
\(\|H^{2D_0}J_c v\|_2\) is finite for every \(v\in U\).
\end{theorem}
\begin{proof}
For a compact-finite matrix coefficient of the row, tensoring
with \(M\) defining representations bounds the coefficient
times the \(M\)-th power of the operator norm. The gaps just
verified ensure that the shifted discrete series all stay in
one regular chamber. Integrating the resulting bound over
\(B_\tau\) gives decay of the local Bessel sections with power
\(L\) in \(\nu_\tau(g)=\inf_{h\in B_\tau}\|hg\|_{\rm op}\).
Lemma \ref{lem:translate163} and Proposition \ref{prop:weight163}
prove both estimates and the stated height moment. Their
constants may depend on \(x\). In particular finiteness is
compatible with the lower bound of Theorem~\ref{thm:localmass156}
and Corollary~\ref{cor:mom156}.

With \(R=\prod_\tau\nu_\tau(g_\tau)\) as height, the number
of rational positions in a ball of radius \(R\ge1\) is
\(O_x(R^{4fr^2})\) by Lemma \ref{lem:cosets163}. Since
\(L>4fr^2\), the rational series and the unfolded pairings
are absolutely convergent as proved in Theorem~\ref{thm:ordinary163}.
The integrable formal-degree coefficient fixes each summed
vector at its input lowest-type vector. Thus projection
preserves these sums. The identity and parabolic terms give
the seed inner product, and the compact stabilizer kills the
indicated mixed cells. This proves the formula for the same
\(G_D\) as in \eqref{eq:gram158}.
\end{proof}

\begin{theorem}[Nonvanishing after refinement at the terminal prime]
\label{thm:depth164}
Fix any one place \(v\mid T\) of \(E^+\), with residue size
\(Q\). For a \(K_v^0\)-stable lattice let \(J_j\) be the
subgroup acting as the identity modulo \(\mathfrak p_v^j\),
with \(J_0=K_v^0\). Multiply each original seed at this place
by \(\mathbf1_{B_vJ_j}\). Let \(S_j\) be its global
periodization from \(U\otimes D\), and put
\[
     \alpha_j=[K_v^0:(B_v\cap K_v^0)J_j]^{-1},\qquad
     \eta=L-4fr^2 .
\]
At each fixed \(x\) and translate there is a finite \(C_x\)
such that for every \(j\ge0\)
\[
       \|\alpha_j^{-1}S_j^*S_j-1\|_{\rm op}
                      \le C_x Q^{-j\eta},
       \qquad e_{K_v^0}S_j=\alpha_j S_0 .
\]
Here the average uses the prescribed level character. For
all sufficiently large \(j\) every nonzero vector of the
original finite domain, including each refined zero-sum column,
has a nonzero projected automorphic image.
\end{theorem}
\begin{proof}
The period group \(B\) is cut out inside the unitary group by
affine linear equations fixing the graded isotropic frame and
the terminal vector modulo that frame. If a nonidentity
rational coset meets both finite supports, one of those
equations takes a nonzero value divisible at \(v\) to depth
\(j\), up to fixed denominators. Its field norm gives a lower
bound \(Q^j/C_x\) for the product of the two real heights.
Lemma \ref{lem:deep163} proves this after compact adelic
reduction, uniformly in the two points at fixed field and
level. The rational tail of the second seed is bounded by
\(C_x Q^{-j\eta}\) times the first height to the power
\(\eta\). That factor is integrable against the first seed
by Proposition \ref{prop:weight163}. The identity term unfolds
with volume ratio \(\alpha_j\). Summing the matrix entries
over a basis gives the displayed estimate, as in Theorem \ref{thm:deep163}. The local support has fraction \(\alpha_j\)
in each original right level coset. Averaging it gives the
second identity by Proposition~\ref{prop:average163}. The first estimate
gives injectivity as soon as \(C_x Q^{-j\eta}<1\).
\end{proof}

The prescribed integral theta period uses the component at
\(K_v^0\). Its complex Gram matrix at that level is \(S_0^*S_0\).
The last identity gives the exact projection onto it from a
finer level. The estimate for large \(j\) gives no positive
lower bound for this projected component at \(j=0\).
Uniform control of its nonidentity orbital terms, the comparison
with the arithmetically normalized theta rows along the full
chain, and the remaining strict orders in \eqref{eq:tail114}
are still needed for \eqref{eq:window}.

\subsection{A first moment for a prescribed terminal character}\label{nav:VII-32}
Fix a word, the original level \(K^0\), frequency \(\varpi\), and
a unit vector of the lowest real type as in
Proposition~\ref{prop:word166}. Write \(S:D\to\mathcal A\) for
its normalized periodization on the full finite class space, with
the class inner product of Theorem~\ref{thm:adjoint157}. Thus
\(S^*S=1+\sum_\gamma O_\gamma\), where the sum is over the
nonparabolic rational double cosets and is absolutely convergent
in operator norm. Let \(Z=\prod_{v\mid T}Z_v\), with
\(Z_v=\mathbb F_T^\times\). Use the unitary translations
\(U_v(t)\) on the class coordinates, with the level character
included. The refined fibers are free \(Z\)-sets by Corollary~\ref{cor:rank111}. Choose unitary coordinates on one
fiber of values over each \(Z\)-orbit and transport them around
the orbit. This gives an orthonormal basis \(\{e_i\}\) permuted
by \(Z\), including multiplicities of the archimedean covectors.

\begin{theorem}[Signed class vectors testing the original Bessel map]
\label{thm:dyad234}
Let \(J\subset D\) be the space of a specified character
\(\lambda=\prod_v\lambda_v\) of \(Z\), with each \(\lambda_v\)
nontrivial. Alternatively let \(J\) be the whole space of vectors
of zero sum on each local scalar orbit in each coordinate.
Write \(P_J\) for orthogonal projection and \(d_J=\dim J>0\).
In the first case take \(i\) uniform in the coordinate basis and put
\[
 b=\frac{1}{\sqrt{|Z|}}\sum_{t\in Z}\overline{\lambda(t)}\,U(t)e_i .
\]
In the second, choose \(i\) in the same way, choose the \(t_v\)
independently and uniformly in \(Z_v\setminus\{1\}\), and put
\[
             b=2^{-f/2}\prod_{v\mid T}(1-U_v(t_v))e_i .
\]
These are unit vectors in \(J\). For any orthogonal projection
\(P\) on \(\mathcal A\), their mean squared projected period is
\[
 \mathbb E\|PSb\|^2 =
   1+\frac{1}{d_J}\sum_\gamma
                 \operatorname{Re}\operatorname{tr}(P_J O_\gamma)
       -\frac{\|(1-P)SP_J\|_{\rm HS}^2}{d_J}.
\]
When \(P=1\), positivity is decided by this single sum of traces.
For a specified \(\lambda\), its individual trace weights are
\(|Z|^{-1}\sum_{t\in Z}\overline{\lambda(t)}
 \operatorname{tr}(U(t)O_\gamma)\).
Whenever the displayed expectation is positive, one of the
displayed vectors attains at least that squared norm at the
original level and the stipulated character or zero-sum subspace.
\end{theorem}
\begin{proof}
If \(N=\dim D\), the mean of the rank-one operators \(e_ie_i^*\)
is \(1/N\) times the identity. In the first case \(b=\sqrt{|Z|}
P_J e_i\). Freeness shows that its \(|Z|\) summands have disjoint
support and that \(d_J=N/|Z|\). Consequently
\(\mathbb E[bb^*]=P_J/d_J\).

In the second calculation first choose each \(t_v\) uniformly
in all of \(Z_v\). The mean of
\((1-U_v(t_v))(1-U_v(t_v))^*\) is \(2(1-P_{1,v})\), where
\(P_{1,v}\) averages over \(Z_v\). The factors at different
places commute and their choices are independent. Averaging the
displayed vector's rank-one operator therefore gives \(P_J/N\).
The vector is zero whenever a \(t_v=1\). Otherwise its \(2^f\)
terms have disjoint supports by freeness and its norm is one.
The probability of the latter event is
\(((T-2)/(T-1))^f=d_J/N\). Conditioning gives precisely \(P_J/d_J\).

Take the trace against \(S^*PS\). Pythagoras gives the last
term in the assertion, and absolute convergence permits taking
the trace of \(S^*S\) term by term. Expanding the character
idempotent gives the individual trace weights. The maximum of
finitely many nonnegative numbers is at least their mean.

In particular the orbital sum can be taken over fewer matrix
entries at the original level. Write \(i=(a,z_i)\), where \(a\)
labels the orbit and the chosen base covector, and \(z_i\in Z\).
Put \(O^\gamma_{ij}=\langle e_i,O_\gamma e_j\rangle\), with
inner product linear in the second variable. Terms with different
\(a\)'s vanish in the trace. For equal \(a\) let \(\nu(i,j)\)
be the number of places at which \(z_i\) and \(z_j\) differ.
For the two choices of \(J\), respectively, we have the exact
normalized formulas
\[
 \frac{\operatorname{tr}(P_JO_\gamma)}{d_J}
 =\frac1N\sum_{i,j:\,a_i=a_j} O^\gamma_{ij}
 \begin{cases}
    \overline{\lambda(z_jz_i^{-1})}, & J=P_\lambda D,\\
    \bigl(-1/(T-2)\bigr)^{\nu(i,j)}, & J=\prod_v(1-P_{1,v})D .
 \end{cases}
\]
Indeed the only scalar sending \(i\) to \(j\) is \(z_jz_i^{-1}\).
For the second formula the local projector has diagonal entry
\(1-1/(T-1)\) and off-diagonal entry \(-1/(T-1)\) within the
scalar orbit. Dividing by \(d_J=N(1-1/(T-1))^f\) gives the
displayed power. Thus a mismatch at each further place introduces
an additional factor \(1/(T-2)\) for the full zero-sum trace.

Before unit normalization each mixed difference has \(2^f\) entries
with coefficients \(1\) or \(-1\). These arrays generate the
integral zero-sum module on each product orbit. Indeed the
augmentation kernel in \(\ZZ[Z_v]\) has basis
\(\delta_t-\delta_1,\ t\ne1\). Splitting off \(\ZZ\delta_1\)
at each coordinate identifies the simultaneous kernels with the
tensor product of these augmentation kernels. In the first case
the unnormalized array generates its character line on an orbit
over the ring containing the character values, since its identity
coefficient is one. These integrality statements concern finite
coefficients and hold with any given integral base covector.
\end{proof}

\begin{proposition}[Testing the trace with the adjoint theta lift]
\label{prop:adjtheta234}
Let \(P\) project onto the cuspidal row space \(\mathcal H\) of
Theorem~\ref{thm:norm177}, in the stipulated real
line and with the original first-edge test at \(T\). Use the common
row Hilbert space to take this projection. Let
\(\Theta:\mathcal H\to\mathcal K\) be that complex theta iteration
to its definite endpoint, and set
\(a=C_xP(\mathbf a)e^{-\Delta_x}>0\) at depth zero.
Every vector of \(\mathcal H\) is the adjoint lift of an endpoint
vector. For each of the unit columns \(b\) above one has
\[
  \sup_{\|y\|\le1}|\langle Sb,\Theta^* y\rangle|^2
         =\|\Theta P S b\|^2\ge a\|P S b\|^2 .
\]
The pairing on the left is the normalized prescribed Bessel
functional on the row \(\Theta^*y\). In particular a positive
first moment in Theorem~\ref{thm:dyad234} gives a nonzero pairing
on a reverse theta lift. This uses only the original finite level.
\end{proposition}
\begin{proof}
The norm theorem on the original row gives
\(a\,1\le\Theta^*\Theta\le C_xP(\mathbf a)e^{\Delta_x}\,1\).
Thus \(W=\Theta(\Theta^*\Theta)^{-1}\) satisfies
\(\Theta^*W=1\) and
\(\|Wv\|^2=\langle v,(\Theta^*\Theta)^{-1}v\rangle
\le a^{-1}\|v\|^2\). The adjoint here is the Hilbert-space
adjoint on the stipulated row. Its pairing with a row vector
is the ordinary Petersson pairing against its forward theta
lift. Pairing the forward lift of each test row with the endpoint
vector defines the reverse pairing weakly on the row.
Projection onto the stipulated row is part of this definition.
The bounded-extension and unfolding formula of
Theorem~\ref{thm:adjoint157} identifies the pairing against
\(Sb\) with the original Bessel functional. Cauchy--Schwarz,
with equality for the normalized vector \(\Theta P S b\) when
nonzero, proves the formula.

There is also a simultaneous choice when the words have a common
endpoint Hilbert space. Suppose one has found a column \(b_c\)
with \(y_c=\Theta_c P_c S_c b_c\ne0\) for each of \(N_0\) words.
In any finite-dimensional endpoint space containing the \(y_c\),
fix a basis \(v_1,\ldots,v_M\). One vector \(y=\sum_j n_jv_j\),
\(0\le n_j\le N_0\) integers, satisfies
\[
    |\langle y_c,y\rangle|\ge
                \tfrac12\max_j|\langle y_c,v_j\rangle|>0
                          \qquad\text{for every }c .
\]
For a fixed \(c\) leave a maximizing coordinate free and fix all
the others. Two successive possible pairings are spaced by
the positive maximum. An open disk of half that radius contains
at most one of them. In the full integer box the fraction failing
for this word is at most \(1/(N_0+1)\). The union over the words
has fraction less than one. When the chosen basis belongs to a
common endpoint lattice this gives a vector in that very lattice.
Its norm is at most \(N_0\sum_j\|v_j\|\).
\end{proof}

The trace formula places the remaining nonvanishing problem on
the actual scalar character if one is prescribed. Its signed
orbital trace and, when present, the component outside the
specified cuspidal row still require estimates. The bounded
right inverse concerns the full complex endpoint. Controlling
the arithmetic normalization of the prescribed pairing on the
common fields is a further requirement.

\subsection{The trace on each full-rank cell}\label{nav:VII-33}
Let \(P\) be the definite \(k\)-dimensional column space over
\(E/F\), with \(f=[F:\QQ]\), and \(\rho\) a word representation
of \(U(P)_\infty\). Its local highest weights are
\(\lambda_{\tau,1}\ge\cdots\ge\lambda_{\tau,k}\). Put
\(n_{\tau,i}=\lambda_{\tau,i}+k-i\) and
\(L=\min_{\tau,i<k}(n_{\tau,i}-n_{\tau,i+1})\).

\begin{proposition}[A character estimate in terms of a rational denominator]
\leavevmode\label{prop:weyl235}
Let \(g\in U(P)(F)\) and let \(\mathfrak a\subset\mathcal O_E\)
be a nonzero integral ideal such that
\(\mathfrak a\,gM\subset M\) for an \(\mathcal O_E\)-lattice
\(M\) spanning \(P\). Write \(m_a\) for the multiplicities of
the distinct eigenvalues and \(q=\sum_{a<b}m_am_b\). Then
\[
 \frac{|\operatorname{tr}\rho(g)|}{\dim\rho}
 \le\min\!\left(1,
   \left[\frac{2k}{L}(N_{E/\QQ}\mathfrak a)^{1/(2f)}\right]^{fq}\right).
\]
For a noncentral \(g\), \(q\ge k-1\).
\end{proposition}
\begin{proof}
At one real place let the distinct eigenvalues be \(z_a\).
They all have modulus one, and \(g\) is diagonalizable by
unitarity. Temporarily omit the place from \(n_{\tau,i}\).
In the Weyl character formula for \(U(k)\) take the limit of
\(\det(z_j^{n_i})/\det(z_j^{k-i})\) as \(m_a\) of the column
variables approach \(z_a\). Divide both determinants by the
Vandermonde products within those groups. The repeated columns
are now derivatives of orders \(0,\ldots,m_a-1\), divided by
their factorials, by Taylor expansion. The denominator has
absolute value
\(D_\tau=\prod_{a<b}|z_a-z_b|^{m_am_b}\).

Expand the remaining numerator in minors on these column groups.
Assigning a set of \(m_a\) exponents to the \(a\)-th group
gives a minor of absolute value
\(\prod_{i<j\ {\rm in\ that\ set}}|n_i-n_j|/
\prod_{t=1}^{m_a-1}t!\).
Indeed unit powers of \(z_a\) factor from the rows and columns,
and the falling factorials in \(n_i\) are monic polynomials of
the successive degrees. Divide by the Weyl dimension
\(\prod_{i<j}|n_i-n_j|/\prod_{t=1}^{k-1}t!\).
The \(q\) uncancelled differences are each at least \(L\).
The quotient of the factorial products is at most \(k^q\):
write the indices \(1,\ldots,k\) in contiguous groups of the
given sizes and cancel the within-group factors in
\(\prod_{i<j}(j-i)\). There are at most \(2^q\) assignments.
Indeed each interleaving of the ordered lists is determined by
the order of each of the \(q\) pairs from different lists.
The normalized character at this place is at most
\((2k/L)^q/D_\tau\).

Consider the endomorphism \(g\otimes1-1\otimes g\) on
\(P\otimes_E P\). Its \(2q\) nonzero eigenvalues have product
\(D(g)\in E^\times\), up to sign the lowest nonzero coefficient
of its characteristic polynomial. The rank and eigenvalue
multiplicities persist at all embeddings. At every finite prime
multiply the endomorphism by a local generator of \(\mathfrak a\).
It now preserves \(M\otimes M\) locally. Its characteristic
polynomial is integral, hence \(\mathfrak a^{2q}D(g)\subset
\mathcal O_E\). At one complex embedding \(|\tau D(g)|=D_\tau^2\),
so the absolute norm over both embeddings at each place gives
\(\prod_\tau D_\tau=|N_{E/\QQ}D(g)|^{1/4}
\ge(N_{E/\QQ}\mathfrak a)^{-q/2}\). The tensor product character
and dimension multiply over the places, proving the second bound.
The first follows from unitarity. A noncentral unitary has at
least two distinct eigenvalues, and the smallest possible number
of cross pairs then is \(1\cdot(k-1)\).
\end{proof}

\begin{corollary}[The original character trace after selecting columns]
\leavevmode\label{cor:trace235}
Intersect either terminal space \(J\) of Theorem~\ref{thm:dyad234}
with the coordinate subspace of Theorem~\ref{thm:columns235},
and denote the result by \(J'\). The first-moment formula of
Theorem~\ref{thm:dyad234} holds with \(J'\) and its dimension
when the initial coordinate is uniform on the retained set.
Every relevant nondegenerate term meets the cutoff of
Corollary~\ref{cor:gram235}. For a full-rank cell in this trace,
the sum over an orthonormal basis of real covectors, divided by
\(\dim\rho\), has absolute value
\[
              |a|\ \frac{|\operatorname{tr}\rho(g)|}{\dim\rho}
\]
for a rational isometry \(g\in U(P)(F)\). Here \(a\) is the
scalar of its orbital intertwiner, including the common finite
cell coefficient and the weight from the terminal projection.
For different words the full-rank orbital intertwiner is zero.
\end{corollary}
\begin{proof}
The coordinate projection onto the retained classes commutes
with all terminal translations because it retains full scalar
orbits. The proof by averaging rank-one operators in
Theorem~\ref{thm:dyad234} gives covariance
\(P_{J'}/\dim J'\) inside this subspace. Taking its trace against
the same projected Gram operator gives the stated first moment,
including its term for projection to the chosen cuspidal row.
In the basis expansion only two classes in the same terminal
scalar orbit can contribute. Each such pair has retained initial
column. The cutoff applies before any averaging over second flags.

For a full-rank pair let \(C\) be a rational positive complement
to the common radical as in Theorem~\ref{thm:cross132}.
Projection gives two rational isometries
\(\pi_0,\pi_1:C\longrightarrow P\), transporting the second
column to the first model by \(\gamma^{-1}\). Thus
\(\iota=\pi_1\pi_0^{-1}\) identifies the two actions of
\(U(C)\). By Theorem~\ref{thm:common133} the real cell
operator for an identical word is a scalar times \(\rho(\iota)\);
the scalar absorbs the phase of the unitary identification.
For different words Schur's lemma gives zero.

To take the trace, choose one finite class representative \(h\)
at the base of the common scalar orbit. Choose its other
representatives by right multiplication by scalar unit lifts
at the terminal places. These act only on the finite coordinate.
For each side of the cell the finite Levi coordinate can
therefore be written \(q_{i,f}h t_i k_i\), with \(q_i\in
U(P)(F)\), \(t_i\) such a lift and \(k_i\) in the original
column level. This is exactly equality in the corresponding
double coset. The rational factor is unique for fixed
representative by Lemma~\ref{lem:covol127}.
The class covector is transported by \(\rho(q_i)^{-1}\), as
in the construction preceding \eqref{eq:real127}. The finite
phases are common to the whole real basis. In these coordinates
the remaining unitary map is \(\rho(q_1^{-1}\iota q_0)\), or
its adjoint for the opposite convention. Summing its diagonal
entries proves the absolute-value formula with
\(g=q_1^{-1}\iota q_0\).
\end{proof}

In the words of \eqref{eq:letters111} the shifted weight gaps
are at least \(3\Delta\), since interior entries descend by
\(4\Delta\) before a change of either bit and the endpoint gaps
satisfy the same bound. For fixed \(A,B>0\), one may choose
polynomial spacing \(\Delta\ge 2e^A kx^B\) in
Theorem~\ref{thm:common111} before selecting \(T\). Every
noncentral full-rank cell admitting a denominator with
\((N_{E/\QQ}\mathfrak a)^{1/(2f)}\le x^B\) then contributes at
most \(e^{-Af(k-1)}|a|\) to this normalized real trace.
This leaves the central terms, large denominators, and the total
size of the scalar amplitudes to be estimated for the
prescribed period.

\begin{proposition}[An integral denominator from the two isotropic flags]
\leavevmode\label{prop:denom235}
On any original full-rank cell in the preceding trace let
\(X_0,X_1\) be the isotropic subspaces, \(I=X_0\cap X_1\),
and \(\Lambda\) the common row lattice. Put
\(U_i=(\Lambda\cap X_i)/(\Lambda\cap I)\). With the convention
that \(h\) is linear in its second slot, consider the map
\[
 D:U_0\longrightarrow U_1^\sharp,\qquad
   x\longmapsto \bigl(y\longmapsto h(y,x)/Q\bigr),
 \qquad
 \mathfrak d=\operatorname{Ann}_{\mathcal O_E}
                      (U_1^\sharp/D U_0).
\]
Here \(U_1^\sharp\) is the module of conjugate-linear integral
functionals. This is a map of lattices of equal rank with finite
cokernel. Its ideal \(\mathfrak d\) is a denominator for \(g\)
in Proposition~\ref{prop:weyl235}. If \(I=0\) and
\(s=d-1\), then
\[
            N_{E/\QQ}\mathfrak d
                \ \le\ |N_{E/\QQ}(D_s/Q^s)|.
\]
Thus the determinant or, more sharply, the largest local
elementary divisors of this integral map test the denominator
in the original trace.
\end{proposition}
\begin{proof}
Write \(Y_i=X_i+Ez_i\) and \(A=(Y_0+Y_1)^\perp\).
By Theorem~\ref{thm:cross132} the full-rank case gives a common
negative vector \(z_*\in Y_0\cap Y_1\) and
\(Y_i=X_i+Ez_*\). The kernel of the pairing between \(X_0\)
and \(X_1\) is \(I\) on either side. Indeed, an \(x\in X_0\)
perpendicular to \(X_1\) is also perpendicular to \(z_*\).
Thus it lies in \(Y_1^\perp\), whose isotropic vectors all
belong to \(X_1\) by positivity of the column quotient.
Interchanging the indices gives the assertion on the other
side. The resulting pairing on the quotients is perfect over
\(E\).

At every place above \(Q\), the primitive integral representatives
of the two isotropic flags have the same reduction up to an
upper unitriangular change. In particular each vector of
\(\Lambda_v\cap X_1\) is congruent to one of
\(\Lambda_v\cap X_0\) modulo \(Q\Lambda_v\). The original
level imposes this congruence in both split factors, as in
Proposition~\ref{prop:shift203}. Hence \(h(y,v)\in Q\mathcal O_E\)
for \(y\in\Lambda\cap X_1\) and \(v\in\Lambda\cap Y_0^\perp\);
away from \(Q\) this just uses integrality. Since every such
\(v\) is orthogonal to \(I\), \(y\mapsto h(y,v)/Q\) defines
an element of \(U_1^\sharp\). This includes \(v\in\Lambda\cap
X_0\), proving that \(D\) is integral. The rational
nondegeneracy just proved makes its cokernel finite.

Fix \(\beta\in\mathfrak d\). For every such \(v\) there exists
\(u\in\Lambda\cap X_0\) with \(h(y,\beta v-u)=0\) for every
\(y\in X_1\), by the definition of the annihilator and by
clearing \(Q\). Now \(\beta v-u\) lies in \(Y_0^\perp\), is
perpendicular to \(X_1\), and is perpendicular to \(z_*\).
It belongs to \(\Lambda\cap A\). The projections of \(A/I\)
onto both rational column spaces are isometries, so its image
in the second column is \(\iota\) of its image in the first.
In the first the image is precisely \(\beta[v]\).
Consequently \(\mathfrak d\,\iota L_0\subset L_1\) for the
two projected integral column lattices. In the common class
coordinates of Corollary~\ref{cor:trace235}, these lattices
are \(q_0M,q_1M\) in their two models. Indeed the terminal
scalar changes preserve the defining hyperplanes, and hence
their lattices by Lemma~\ref{lem:termparent221}.
This gives \(\mathfrak d\,gM\subset M\).

Over a local discrete valuation ring the cokernel is a sum of
cyclic modules by diagonal reduction. Its annihilator has the
largest elementary-divisor exponent, bounded by their sum, the
valuation of the determinant in integral bases. If \(I=0\),
the primitive ordered flag representatives give such bases at
each finite prime. Their allowed upper unitriangular changes
leave the leading cross minor unchanged. The determinant of
\(D\) is therefore \(\overline{D_s}/Q^s\) in these local bases,
up to a unit. Multiply the local bounds to obtain the stated
norm inequality. The zero-rank case gives the unit ideal.
\end{proof}

\subsection{The relative scalar on the common isotropic part}\label{nav:VII-34}
Let \(X_b\) be the span of the \(s=d-1\) isotropic vectors of
the \(b\)-th decorated frame, for \(b=0,1\), and \(Y_b=X_b+Ez_b\)
with \(h(z_b,z_b)=-T\). Write \(X_{b,i}\) for its prefix of
length \(i\). Suppose \(I=X_0\cap X_1\) has dimension \(e>0\)
and put \(t=s-e\). The exceptional rank in a signed first
moment can range from \(t=0\) to \(t=s-1\).

\begin{theorem}[Classification data with a common isotropic subspace]
\leavevmode\label{thm:suffix237}
For any relevant pair there are allowed changes of the two
representatives such that their isotropic lists and final
vectors have the form
\[
\begin{array}{llll}
 u_{0,1},\ldots,u_{0,t},& v_{t+1},\ldots,v_s,& z_*,\\
 u_{1,1},\ldots,u_{1,t},& \lambda v_{t+1},\ldots,\lambda v_s,
       &\lambda z_* .
\end{array}
\]
Here the \(v_i\) form a basis of \(I\), \(\lambda\bar\lambda=1\),
and the cross matrix on the first \(t\) positions is a direct
sum of reversed blocks of the kind in Lemma~\ref{lem:forms203}.
The lengths, block scalars, and \(\lambda\) determine the
rational double coset. The only nonzero leading cross minors
are the block endpoints in the first \(t\) positions.
\end{theorem}
\begin{proof}
Theorem~\ref{thm:success135} gives the same filtration on
\(I\) from both flags, with one jump at each of
\(t+1,\ldots,s\). It also supplies a common negative vector.
The projections from \(Y_0\cap Y_1\) to \(Y_b/X_b\) have
the same kernel \(I\): a vector of the intersection lying in
one isotropic subspace has norm zero and hence lies in the
other, since \(Y_b/X_b\) is negative. Choose \(z_*\) inducing
the prescribed final vector for the first flag. The second
one is induced by \(\lambda z_*\) for a unique \(\lambda\).
Equality of the negative norms gives \(\lambda\bar\lambda=1\).

At a jump \(i\), take \(v\in I\) with leading coefficients
\(a_0,a_1\) in the two lists. In the proof of
\eqref{eq:success135}, if \(i<s\), the common successor
has precisely these same coefficients at position \(i+1\).
Iterate to position \(s\). At that last position
\eqref{eq:success135} says \(a_0 z_*-a_1\lambda z_*\in X_0+X_1\).
But \(z_*\) is perpendicular to this sum and has nonzero
norm. Hence \(a_0=\lambda a_1\) at every jump. At each \(i\)
choose \(v_i\) with first coefficient one. Changing each list
by adding earlier vectors then gives the two suffixes in the
statement. These vectors are rational, since the flags and
filtrations are defined over \(E\).

The kernel of the cross pairing on \(X_0,X_1\) is exactly \(I\)
on each side, again by Theorem~\ref{thm:success135}. Both
prefixes of length \(t\) complement \(I\), so their pairing
matrix is invertible. Reduce it to monomial form by upper
unitriangular changes of the prefixes. Their sum is a
nondegenerate subspace perpendicular to \(I\) and to \(z_*\).
Each inversion test of Proposition~\ref{prop:class236}
acts on this sum and can be extended by the identity on its
perpendicular. It then fixes the common suffix and the two
terminal representatives, and is in both unipotent radicals.
The two character functionals on this element are exactly the
ones for the prefixes, with the coefficients at all later
simple positions zero. The same adjacent-descent computation
therefore gives consecutive reversed blocks. For \(j>t\) the
first \(j\) rows include a nonzero vector of \(I\), which pairs
to zero with the entire second frame. This proves the minor
assertion, including at \(d\).

If the data agree for two pairs, map the vectors of the two
prefixes, the first suffix and \(z_*\) to their counterparts.
Together these vectors form a basis of their span. The Gram
matrices agree: the prefixes pair by the specified blocks,
the negative vector has norm \(-T\), and \(I\) is the radical
of this span. Witt's extension theorem extends the isometry of
the subspaces to the ambient space. Equivalently one can first
match the nondegenerate prefixes and negative line, and then
split off and match hyperbolic partners of the common isotropic
basis in their perpendiculars. The same map also carries the
second suffix and final vector because \(\lambda\) agrees.
The resulting pairs are in one rational orbit.
\end{proof}

\subsection{Counting rational orbits in the terminal period}\label{nav:VIII-14}
\label{sec:orbits203}

Let \(E/F\) be the CM extension of Section~\ref{sec:char94},
\(f=[F:\QQ]\), and choose one embedding of \(E\) above each real
place \(\tau\) of \(F\). We write \(\tau x\) for its value on \(x\).
Put \(d=s+1\ge2\). Write \(Y=(u_1,\ldots,u_s,z)\) for the ordered
frame, with \(z=\varpi e\) and \(h(z,z)=-T\). The two factors in the
unfolded period use the same \(T\). At the finite places we use
the level \(K^0\) of Theorem~\ref{thm:midlevel131}. Set
\(Q=g_1g_2\). The congruence at \(Q\) will enter the count separately
at every nonzero leading minor.

A decorated flag is the class of \(Y\) under right upper
unitriangular changes of basis. Thus it specifies each successive
flag vector modulo its predecessors, including the last vector.
Its stabilizer in \(G=U(h)\) is \(B=NH\). Indeed an element
stabilizing the class acts trivially on each one-dimensional
isotropic quotient. On the middle Hermitian quotient it fixes the
image of \(z\), so its action there belongs to \(H=U(W_0)\).
The kernel of these graded actions is \(N\).
For \(\gamma\in G(F)\), let \(Y'=\gamma Y\), and write
\[
 Z=h(Y,Y'),\qquad D_0=1,\qquad
 D_j=\det(Z_{ab})_{1\le a,b\le j},\qquad
 \mathcal G=D_d\overline{D_d}-T^2D_{d-1}\overline{D_{d-1}} .
\]
Independent changes of representative of the two flags send \(Z\)
to \(U^*ZV\), where \(U,V\) are upper unitriangular. Each \(D_j\)
is unchanged, since its leading square transforms by the corresponding
leading unitriangular matrices. A pair is \emph{relevant} if its two
characters agree on the common real unipotent group in
Lemma~\ref{lem:zero133}. A pair is \emph{supported at the finite
places} if \(\gamma_f\in B_fK^0B_f\). Both properties descend to the
rational double coset.

\begin{lemma}[Normal form of a relevant pair]\label{lem:forms203}
Suppose \(\mathcal G\ne0\). The normal form of \(Z\) under the two
unitriangular changes determines its rational double coset. For a
relevant pair the normal form is block diagonal. Each block of
length \(h\) has successive antidiagonal entries
\(c,-c,c,\ldots,(-1)^{h-1}c\), with \(c\in E^\times\), and all its
other entries vanish. The nonzero leading minors have indices
precisely the successive block endpoints. For a composition
\(h_1+\cdots+h_l=d\) with endpoints \(j_\nu=h_1+\cdots+h_\nu\),
a prescribed tuple of endpoint minors occurs for at most
\(\prod_{\nu=1}^l h_\nu\) rational double cosets.
\end{lemma}

\begin{proof}
The determinant of the concatenated Gram matrix is
\((-1)^d\mathcal G\) by \eqref{eq:cross132}. Its \(2d\) vectors
are therefore independent and have nondegenerate span. By
Theorem~\ref{thm:cross132} the orthogonal complement is positive
definite of rank \(k-1\), and \(\mathcal G\) is totally positive.
In particular \(D_d\ne0\).

For an invertible \(Z\) clear its first row using its leftmost
nonzero entry. Add multiples of this pivot column to later columns
and multiples of the first row to later rows to clear the pivot
column. Repeat on the remaining rows and columns in their original
order. This gives a monomial matrix \(M\) with entries
\(a_i=M_{i,w(i)}\ne0\). Such a form is unique. The ranks of all
northwest rectangles are invariant and determine the location of
each nonzero entry by taking successive differences. If
\(M'=LMV\) has the same permutation, with \(L\) lower and \(V\)
upper unitriangular, the \(i\)-th diagonal entry of
\(M'V^{-1}M^{-1}=L\) is \(a'_i/a_i=1\).

Two pairs with the same normal form admit representatives with the
same concatenated Gram matrix. Both diagonal Gram blocks are
\(-T E_{dd}\), where \(E_{ij}\) denotes a matrix unit.
The linear map carrying the first concatenated basis to the second
is an isometry on their spans. Witt's extension theorem for
Hermitian spaces extends it over \(E\) to \(G(F)\). In this
nondegenerate case one can see the extension directly. Diagonalize
the first span and match its orthogonal basis vectors in turn. In
the perpendicular to the vectors already matched take vectors
\(v,v'\) of the same nonzero norm \(a\). For some \(y\in\{v',-v'\}\)
the number \(b=h(v-y,v)\) is nonzero, since the two choices give
numbers with sum \(2a\). Put \(w=v-y\). The map
\(t\mapsto t-w h(w,t)/b\) sends \(v\) to \(y\) and preserves \(h\),
as \(b+\bar b=h(w,w)\). If \(y=-v'\), compose with minus the
identity on this perpendicular. Extend by the identity on the
vectors already matched. Induction gives the required isometry.
Thus the two ordered pairs of decorated flags are in the same
\(G(F)\)-orbit, which is equivalent to equality in
\(B(F)\backslash G(F)/B(F)\).

On a frame an element of \(N\) acts by an upper unitriangular
matrix \(U\), and its character at a real place is the additive
character of the trace of \(\sum_{i<d}U_{i,i+1}\). The coefficient
in the last position uses \(z=\varpi e\), as in Section~\ref{sec:char94}.
Conjugation by an upper unitriangular matrix preserves the first
superdiagonal, so the formula applies to the representatives used
to obtain \(M\).

Take an inversion \(i<j\), \(w(i)>w(j)\), and work at one real
place. For any \(x\in\CC\) put
\[
 U_x=1+xE_{ij},\qquad
 V_x=1-\bar x(a_i/a_j)E_{w(j),w(i)} .
\]
Here the \(a_i\) are evaluated at the chosen place.
Multiplication gives \(U_x^* M V_x=M\). On the two independent
spans these upper unitriangular maps preserve their self pairings.
Together they define a unipotent isometry, extended by the identity
on the orthogonal complement. Its projection to each of the two
compact column groups is unipotent and unitary, hence the identity.
It is therefore in both unipotent radicals. Equality of the
characters, applied to all real multiples of \(x\), gives
\[
 {\bf1}_{j=i+1}\Re(x)
  ={\bf1}_{w(i)=w(j)+1}\Re\bigl(-\bar x\,a_i/a_j\bigr).
\]
An adjacent descent must decrease by one. At each such descent
\(x=1\) and \(x=\sqrt{-1}\) give \(a_{i+1}=-a_i\). The equality
also holds in \(E\) because the embedding is injective. Each maximal
descending run of \(w\) is a consecutive interval of values in
reverse order. Since two successive runs are disjoint intervals
joined by an ascent, their value intervals are in increasing order.
They partition \(\{1,\ldots,d\}\), proving the asserted block form.

Within a partially included reversed block the first row of the
partial block is zero. A whole block has determinant \(c^h\),
since the reversal sign and the alternating row signs cancel.
Consequently \(D_{j_\nu}/D_{j_{\nu-1}}=c_\nu^{h_\nu}\), with
\(j_0=0\). Each equation has at most \(h_\nu\) roots in the field
\(E\). The uniqueness just proved gives the claimed multiplicity.
\end{proof}

\begin{proposition}[Congruences of the leading minors]\label{prop:shift203}
Every rational pair supported at the finite places satisfies
\[
  D_j\in Q^j\mathcal O_E\quad(1\le j<d),\qquad
  D_d+T D_{d-1}\in Q^d\mathcal O_E .
\]
All \(D_j\) are algebraic integers. If \(\mathcal G\ne0\), then
\(\lvert\tau D_d\rvert>T\lvert\tau D_{d-1}\rvert\) for every
\(\tau\).
\end{proposition}

\begin{proof}
Write the finite translate as \(b_1kb_2\) with \(k\in K^0\).
Each exterior product of the first \(j\) frame vectors is fixed
by \(B_f\). Pairing it with its translate gives the same minor
as for \(k\). The frame is integral in the local Hermitian
lattices and \(k\) preserves them, including the two mutually
dual factors at \(T\), by Theorem~\ref{thm:midlevel131}.
Thus every minor is integral at all finite places.

Fix either prime \(g\mid Q\). At each place above \(g\) the
reduction of \(kY\) is \(Y\) times an upper unitriangular matrix
in both split factors. The two actions are described in the proof
of Theorem~\ref{thm:gap132}. Since the self pairing of the frame
is \(-T E_{dd}\), it follows that every entry of \(h(Y,kY)+T E_{dd}\)
is in \(g\) times the local integer ring. A leading \(j\)-square
for \(j<d\) therefore has determinant in \(g^j\). The full
determinant after adding \(T E_{dd}\) is \(D_d+T D_{d-1}\) by
expansion in that entry. It lies in \(g^d\). The two rational
primes are distinct and split completely, giving the asserted
global inclusions. The real inequality follows from total positivity
of \(\mathcal G\).
\end{proof}

For a tuple of minors with nonzero endpoints indexed by \(I\), put
\[
      w(D)=\prod_{j\in I}\prod_\tau (1+|\tau D_j|^2)^2 .
\]
The source uses constant antidiagonal entries on a reversed block.
Multiplication of the cross matrix on the left by
\(\operatorname{diag}(1,-1,1,-1,\ldots)\) changes the alternating
convention of Lemma~\ref{lem:forms203} into this convention. It
multiplies \(D_j\) by \((-1)^{j(j-1)/2}\) and preserves \(w(D)\).
For constant blocks the successive endpoint ratios are
\((-1)^{h(h-1)/2}c^h\). Either convention therefore gives the
same bound for the number of choices of block scalars.

\begin{theorem}[Sum over all integral block data]\label{thm:allorbits203}
Over any CM field \(E\) of degree \(2f\), sum \(w(D)^{-1}\) over
all compositions of \(d\ge1\) and all nonzero block scalars in \(E\)
for which the endpoint minors belong to \(\mathcal O_E\). In either
sign convention this sum is bounded by
\[
             L^f(L^f+2)^{d-1}\le 8^{fd},\qquad L=81/16 .
\]
In particular the bound \(64^{fd}\) in the source holds for every
pair of positive terminal norms in its definition of integral
block points.
\end{theorem}

\begin{proof}
For \(\delta\in\mathcal O_E\) set
\(p(\delta)=\prod_\tau(1+|\tau\delta|^2)^{-2}\). The centers
\((\tau\delta)_\tau\) of distinct algebraic integers have sup
distance at least one. Indeed the square of the product of the
coordinate distances is the absolute norm of a nonzero algebraic
integer, hence at least one. The open polydiscs of radius \(1/2\)
around the centers are disjoint.

In one coordinate average \((1+|y|^2)^{-2}\) over the disc of
radius \(1/2\) about \(z\) with uniform area probability.
The average of \(|y|^2\) is \(|z|^2+1/8\), by rotational symmetry
and radial integration. Convexity of \(t\mapsto(1+t)^{-2}\) gives
a lower bound \((1+|z|^2+1/8)^{-2}\), which is at least
\((8/9)^2(1+|z|^2)^{-2}\). The disc area is \(\pi/4\). The
integral of the same density on \(\CC\) is \(\pi\), by polar
coordinates. Multiply over the \(f\) coordinates and use disjointness
on any finite set of centers. This gives
\(\sum_{\delta\in\mathcal O_E}p(\delta)\le[4(9/8)^2]^f=L^f\)
on passing to the supremum over finite sets.

For \(l\) endpoints the sum over independent integral values is at
most \(L^{fl}\). The successive scalar equations give multiplicity
at most \(\prod_\nu h_\nu\le 2^{d-l}\), since \(h\le2^{h-1}\)
for \(h\ge1\). There are \(\binom{d-1}{l-1}\) compositions of this
length. The binomial theorem gives \(L^f(L^f+2)^{d-1}\), and
\(L^f+2\le(L+2)^f<8^f\). An integral representative in the source's
cell definition has integral leading minors because the lower and
upper unitriangular factors preserve them. The further conditions
on its Gram determinant restrict the set being summed. This proves
the last assertion, uniformly in the two terminal norms.
\end{proof}

\begin{theorem}[The sum at the prescribed finite level]\label{thm:levelsum203}
Let \(\mathcal R\) be the set of relevant rational double cosets
supported at the finite places with \(\mathcal G\ne0\). For
\(0\le t<2\) set
\[
       M_t=\int_0^\infty \frac{4r^{3+t}}{(1+r^2)^3}\,dr .
\]
Then
\begin{equation}
    \sum_{\gamma\in\mathcal R} w(D(\gamma))^{-1}
      \le (9M_t+2)^{fd}Q^{-(2+t)fd}.                         \label{eq:levelsum203}
\end{equation}
In particular the sum is at most
\(\bigl(128\log(2Q)/Q^4\bigr)^{fd}\).
\end{theorem}

\begin{proof}
Fix a composition with endpoint set \(I\). For \(j\in I\) set
\(P_j=Q^j\), with one modification. When the terminal block has
length one put \(P_d=Q^{d-1}\min(Q,T)\). For every endpoint tuple
under consideration,
\[
      \prod_\tau|\tau D_j|\ge P_j^f .
\]
For \(j<d\) this follows by taking the norm of the nonzero
integer \(D_j/Q^j\). For a terminal block of length greater
than one we have \(D_{d-1}=0\) and the same argument applies at
\(d\). For a terminal singleton we have \(d-1\in I\), and
Proposition~\ref{prop:shift203} gives
\(\prod_\tau|\tau D_d|>T^fQ^{(d-1)f}\). Also, given all preceding
endpoint values, distinct choices of \(D_j\) differ by an element
of \(Q^j\mathcal O_E\). At \(d\) this uses the fixed residue
\(-T D_{d-1}\), with \(D_{d-1}=0\) when that index is absent.
In particular the product of the coordinate distances of two such
choices is at least \(P_j^f\).

Impose bounds \(|\tau D_j|\le r_{j,\tau}\), where all radii are
positive, and put \(R_j=(\prod_\tau r_{j,\tau})^{1/f}\).
There are no tuples unless \(R_j\ge P_j\) at every endpoint.
For one endpoint at a time divide each complex coordinate by its
radius. The centers now lie in the unit polydisc and have sup
distance at least \(P_j/R_j\) by the product bound on differences.
The disjoint open polydiscs of radius \(P_j/(2R_j)\) around them
fit in the polydisc of radius \(1+P_j/(2R_j)\). Comparison of
volumes gives at most
\((1+2R_j/P_j)^{2f}\le 9^f(R_j/P_j)^{2f}\) choices.
Iterating conditionally on earlier values and allowing the scalar
multiplicity bounds the number of rational labels in the box by
\[
 \Bigl(\prod_\nu h_\nu\Bigr)
 \prod_{j\in I}9^f (R_j/P_j)^{2f}{\bf1}_{R_j\ge P_j}.
\]

The identity
\((1+u^2)^{-2}=\int_u^\infty4r(1+r^2)^{-3}\,dr\)
expresses the required weighted count as the integral of these
box counts against the product of the radial densities.
Tonelli's theorem applies since the summands are nonnegative.
Use \({\bf1}_{R_j\ge P_j}\le(R_j/P_j)^{tf}\). For one endpoint
the resulting integral is at most
\((9M_t)^f P_j^{-(2+t)f}\). It converges for \(0\le t<2\).
For a terminal block of length greater than one,
\(\prod_{j\in I}P_j\ge Q^d\) comes from the last endpoint alone.
For a terminal singleton the last two endpoints contribute at least
\(Q^{2(d-1)}\ge Q^d\), since \(d\ge2\) and \(T\ge1\).
The binomial sum in Theorem~\ref{thm:allorbits203}, here with
\(L=9M_t\), now proves \eqref{eq:levelsum203}.

For the explicit form take \(\eta=1/\log(2Q)\), \(t=2-\eta\).
Since \(Q\ge2\), \(0<\eta<1\). Splitting the integral at one
and bounding its denominator below by \(1\) and \(r^6\),
respectively, gives \(M_t\le1+4/\eta\). Hence
\(9M_t+2\le47\log(2Q)\), while \(Q^\eta\le e\) and \(47e<128\).
\end{proof}

\begin{corollary}[The measure on the individual finite cells]
\label{cor:amplitude203}
Fix two different words and unit vectors in the corresponding input
spaces of the untruncated periodization at \(K^0\). Denote the
scalar contribution of a finite cell \(c\) above \(\gamma\) by
\(J_{\gamma,c}\), including the finite coefficient and its volume,
with the normalization of \eqref{eq:orb133}. If
\(\sum_c |J_{\gamma,c}|\le A/w(D(\gamma))\) for every
\(\gamma\in\mathcal R\), then the absolute value of this mixed
Gram entry is at most
\(A\bigl(128\log(2Q)/Q^4\bigr)^{fd}\).
\end{corollary}

\begin{proof}
The sum is absolutely convergent by Theorem~\ref{thm:ordinary163}.
Lemma~\ref{lem:zero133} removes the irrelevant terms and
Corollary~\ref{cor:zero135} removes the terms with zero determinant
for distinct words. Apply the triangle inequality and
Theorem~\ref{thm:levelsum203} to the remaining labels.
\end{proof}

The source specifies its block points, weight, and summation
in \srcloc{Formalism/SharpRestriction/CuspSpanNormBounds/Point}{11},
\srcloc{Formalism/SharpRestriction/CuspSpanNormBounds/LeadMinorK}{150},
and \srcloc{Formalism/SharpRestriction/CuspSpanNormBounds/OnePlaceEstimates}{35}.
Its bound \(64^{fd}\) occurs in
\srcloc{Formalism/SharpRestriction/CuspSpanNormBounds/BesselShellMomentBounds}{265}.
Theorem~\ref{thm:allorbits203} proves this estimate from the scalar
equations and the area integral. Theorem~\ref{thm:levelsum203} gives
the further saving from the prescribed congruences. In the mixed
Gram entry the constant \(A\) includes the sum of the absolute
orbital integrals over the individual finite cells of one rational
orbit. Controlling this quantity uniformly and controlling the
singular terms for a word paired with itself remain necessary for
an estimate of the whole Gram matrix by this argument.

\subsubsection{Integral cross matrices within a rational orbit}\label{nav:VIII-15}
\label{sub:crossclasses230}

Put \(O=\mathcal O_E\). Write \(U^-(S)\) and \(U^+(S)\)
for the lower and upper unitriangular groups over a ring
\(S\). A rational label specifies a double orbit under these
groups over \(E\). Counting integral flags requires the
smaller equivalence relation over \(O\). The following bound
accounts for all the integral classes above the label.

\begin{theorem}[Integral classes of a marked cross matrix]
\label{thm:crossclasses230}
Let \(M\) be the monomial normal form of a relevant
nondegenerate label in Lemma~\ref{lem:forms203}, and write
\(a_i=M_{i,w(i)}\) for its nonzero entries. Put
\[
 b_0=1,\qquad b_j=Q^{-j}\prod_{i=1}^j a_i\quad(1\le j<d),
 \qquad
 u_j=\#\{i>j:w(i)>w(j)\}.
\]
Consider the integral matrices in \(U^-(E)M U^+(E)\)
which satisfy \(Z\equiv-T E_{dd}\bmod QO\), where
\(E_{dd}\) has its only nonzero entry, one, at \((d,d)\).
If such a matrix exists, all \(b_j\) are algebraic integers.
With \(h_j=|N_{E/\QQ}b_j|\), the number of
\(U^-(O),U^+(O)\) double classes is at most
\[
             C(M)=\prod_{j=1}^{d-1} h_j^{\,d-j+u_j}.
\]
If \(t=d-1\), the absolute index of the image in \(O^t\)
of the top \(t\) rows of any such \(Z\) is at most
\(Q^{2ft}h_t\), where \([E:\QQ]=2f\).
\end{theorem}
\begin{proof}
Both unitriangular groups fix \(E_{dd}\) on their respective
sides. Thus \(Z\mapsto Z'=(Z+TE_{dd})/Q\) carries the
specified matrices bijectively to the integral matrices of
one rational orbit and respects integral equivalence in both
directions. Its monomial form \(M'\) has the same permutation.
If \(w(d)=d\), its last entry is \((a_d+T)/Q\), and the
others are \(a_i/Q\). The last entry is nonzero because
\(\mathcal G>0\) at every real place, so
\(|D_d|>T|D_{d-1}|\) there. If \(w(d)\ne d\), row
\(w^{-1}(d)<d\) of \(M\) has its sole nonzero entry in column
\(d\). A rational lower unitriangular row operation removes
the added entry at \((d,d)\), giving \(M'=M/Q\).

Permute the columns of \(Z'\) into the order
\(w(1),\ldots,w(d)\), and call the resulting matrix \(B\).
It has a factorization \(B=L\Delta U\), where \(L\) is lower
unitriangular, \(\Delta\) has diagonal \(M'_{i,w(i)}\), and
\(U\) is upper unitriangular with \(U_{ij}=0\) when
\(i<j\) and \(w(i)>w(j)\). To see this, write
\(B=L_0\Delta V\) using the given rational orbit. The
matrix \(V\) is unitriangular in the order defined by \(w\),
so an off-diagonal entry \(V_{ij}\) can be nonzero only for
\(w(i)<w(j)\). Eliminate its entries below the ordinary
diagonal by adding multiples of earlier rows, clearing each
pivot column before proceeding. A nonzero entry to be
cleared in position \((i,j)\) has \(w(i)<w(j)\). Every
nonzero entry of the pivot row is in a column \(k\) with
\(w(j)\le w(k)\), so the row operation preserves the allowed
positions and the diagonal ones. This gives \(V=L_1U\).
Then \(L=L_0\Delta L_1\Delta^{-1}\) has the required form.

For \(j<d\), the leading \(j\)-minor of \(B\) is \(b_j\).
It is nonzero and integral. For \(i>j\), the minor with
rows \(1,\ldots,j-1,i\) and columns \(1,\ldots,j\) is
\(b_jL_{ij}\). For \(i>j\), the minor with rows
\(1,\ldots,j\) and columns \(1,\ldots,j-1,i\) is
\(b_jU_{ji}\). These identities follow by taking determinants
in the displayed triangular factorization. Consequently
the entries below column pivot \(j\) of \(L\) and to the
right of row pivot \(j\) of \(U\) lie in \(b_j^{-1}O\).

Choose representatives of \(b_j^{-1}O/O\), with zero
representing zero. There are \(h_j\) choices. In each row
\(i\) of \(L\), reduce its entries in columns \(j<i\)
in decreasing order of \(j\) by adding integral multiples
of row \(j\). The later columns already fixed in that row
are unchanged. Every entry in column \(j\) stays in
\(b_j^{-1}O\), since the operations are integral.
In each column \(i\) of \(U\) use the corresponding column
operations with \(j<i\) in decreasing order. A nonzero
coefficient is used only for \(w(j)<w(i)\), so in the
original column order the operation is an allowed integral
upper unitriangular change. These operations preserve the
zero pattern just proved, as well as the denominators within
each row. The two reductions act on opposite sides of
\(L\Delta U\). They give a list of representatives of length
at most the displayed product, with \(d-j\) positions on
the left and \(u_j\) on the right for index \(j\).

For \(t>0\), the columns \(w(1),\ldots,w(t)\) of the first
\(t\) rows of \(Z\) have determinant \(Q^t b_t\) in that
order, since adding \(TE_{dd}\) leaves these rows unchanged.
Their image is contained in the image of all the columns.
Taking the index gives \(|N_{E/\QQ}(Q^t b_t)|=Q^{2ft}h_t\).
For \(t=0\) the index and bound are both one, with \(h_0=1\).
\end{proof}

The bound uses only the block scalars of the rational label.
It has a simpler expression in terms of the endpoints. For
each proper endpoint put \(H_{j_\nu}=|N_{E/\QQ}(D_{j_\nu}/Q^{j_\nu})|\),
and put \(H_0=1\) and
\(H_d=|N_{E/\QQ}((D_d+TD_{d-1})/Q^d)|\).
Every \(H_{j_\nu}\) is a positive integer for a represented
label. At successive endpoints the magnitudes of the prefix
pivot products of \(M'\) are precisely these norms.
Inside a block of length greater than one the pivot magnitudes
are constant. The logarithms of their prefix product norms
therefore interpolate linearly between the two endpoint
logarithms. It follows that each \(h_j\) is at most
\(\prod_\nu H_{j_\nu}\). Since
\(\sum_{j<d}(d-j+u_j)\le d(d-1)\), the number of classes is
at most \((\prod_\nu H_{j_\nu})^{d(d-1)}\).
The finer product in the theorem can be used separately for
each label.

\subsection{The rational orbits at the prescribed period level}\label{nav:VIII-29}
Fix \(x\) in one of the intervals of Theorem~\ref{thm:all202}, and
one of its translates. In the period over \(E_x/F\), \(F=E_x^+\),
let \(f=[F:\QQ]\), \(d=s+1=m_0/2\), and \(Q=g_1g_2\).
The Gram matrix in Theorem~\ref{thm:ordinary164} is the sum on the
original level \(K^0\) at the terminal frequency \(\varpi\).
For two distinct words, the nonzero terms in the off-diagonal block
are supported on relevant rational double cosets with nonzero Gram
determinant. Write \(\mathcal R\) for all such rational labels allowed
by the finite support, including any whose orbital integral is zero.
We count a label once even when it supports several finite cells.

\begin{theorem}[Weighted count on the full period fields]
\label{thm:orbits204}
For \(\gamma\in\mathcal R\) let \(I_\gamma\) index the nonzero
leading minors \(D_j\) of the cross matrix of the two ordered frames.
Put \(w_\gamma=\prod_{j\in I_\gamma,\tau}(1+|\tau D_j|^2)^2\),
where \(\tau\) runs through one complex embedding above each real
place of \(F\). At every \(x\) and every translate as above,
\begin{equation}
  \sum_{\gamma\in\mathcal R} w_\gamma^{-1}
   \le \left(\frac{128\log(2Q)}{Q^4}\right)^{fd}
   \le \exp\!\bigl(-(64-o(1))fd\, x^\eps\bigr).
   \label{eq:orbits204}
\end{equation}
The second estimate is uniform on the intervals \([V,V^{3/2}]\)
as \(V\) tends to infinity. The subset for which \(D_{d-1}\ne0\)
satisfies the stronger bound
\[
   \sum_{\substack{\gamma\in\mathcal R\\D_{d-1}\ne0}}
       w_\gamma^{-1}\ \le\
       \exp\!\bigl(-(192-o(1))fd\,x^\eps\bigr).
\]
\end{theorem}

\begin{proof}
For \(Y=(u_1,\ldots,u_s,\varpi e)\) write \(Z=h(Y,\gamma Y)\).
Changing the two ordered representatives multiplies \(Z\) by a lower
and an upper unitriangular matrix. Gaussian elimination leaves one
nonzero entry in each row and column. Agreement of the two additive
characters on each common elementary unipotent forces each descending
run of the permutation to reverse a consecutive block, with successive
entries \(c,-c,c,\ldots\). A block of length \(h\) multiplies the
preceding nonzero leading minor by \(c^h\). Each tuple of nonzero
minors therefore has at most the product of the block lengths as its
multiplicity. The two auxiliary level primes force \(Q^j\mid D_j\)
for \(j<d\) and \(Q^d\mid D_d+T D_{d-1}\). The standard separation
of distinct algebraic integers by the product of the distances at
their complex embeddings now bounds the number of tuples in any
product of discs. Integrating the resulting count against
\(4r(1+r^2)^{-3}\,dr\) in every coordinate gives the first inequality
in \eqref{eq:orbits204}. The character calculation, congruences, volume
bound and integral are proved in
Section \ref{nav:VIII-14}, Lemma \ref{lem:forms203} through Theorem \ref{thm:levelsum203}. Since \(b=\exp(x^\eps)\) and
\((kb)^8\le g_i\le2(kb)^8\) by Proposition~\ref{prop:index90},
with \(k\) of polynomial size, \(\log Q=(16+o(1))x^\eps\)
uniformly. This proves the second inequality.

Here is the additional saving when the last block has length one,
the case \(D_{d-1}\ne0\). Write \(R_j\) for the geometric mean
of the disc radii at endpoint \(j\). For \(j=d-1\), a box with a
point must have \(R_{d-1}\ge Q^{d-1}\). At \(d\), the positivity
of \(|\tau D_d|^2-T^2|\tau D_{d-1}|^2\) gives
\(R_d\ge T Q^{d-1}\), and distinct last minors with the preceding
data fixed differ by an element of \(Q^d\mathcal O_{E_x}\).
Here \(T\ge Q\) for large \(V\), since Theorem~\ref{thm:common111}
gives \(\log T\ge(16-o(1))m_0 x^\eps\) and \(m_0\to\infty\)
uniformly. The count for the last endpoint is at most
\[
       9^f(R_d/Q^d)^{2f}\,{\bf1}_{R_d\ge TQ^{d-1}} .
\]
For \(0<t<2\) bound the indicator by
\((R_d/(TQ^{d-1}))^{tf}\). Writing
\(M_t=\int_0^\infty4r^{3+t}(1+r^2)^{-3}\,dr\), the integral
at \(d\) is at most \((9M_t)^f Q^{-2df}(TQ^{d-1})^{-tf}\).
The endpoint \(d-1\) supplies
\((9M_t)^f Q^{-(2+t)(d-1)f}\). Use the same integral at the
remaining endpoints, discard their further powers of \(Q^{-1}\),
and sum the multiplicities over compositions as in Section~\ref{nav:VIII-14}.
This gives
\[
     (9M_t+2)^{fd}Q^{-(2+t)(2d-1)f}(Q/T)^{tf}.
\]
Taking \(t=2-1/\log(2Q)\) gives \(9M_t+2\le47\log(2Q)\).
The factor in \(Q^{-(2+t)(2d-1)f}\) saves
\((128-o(1))fd\,x^\eps\) in the logarithm. Since \(m_0=2d\) and
\(d\to\infty\), the last factor saves at least
\((64-o(1))fd\,x^\eps\). The positive logarithm from the first
factor is \(o(fd\,x^\eps)\), proving the assertion.
\end{proof}

\begin{corollary}[A sufficient bound for a mixed Gram entry]
\label{cor:amplitude204}
Consider two distinct words and unit input vectors at \(K^0\).
Let \(J_{\gamma,c}\) be the scalar orbital integral of finite
cell \(c\) above \(\gamma\), with its finite coefficient and volume.
Suppose for some \(\delta>0\) that at all sufficiently large
\(x\) their absolute sums satisfy
\[
          \sum_c |J_{\gamma,c}|
            \le w_\gamma^{-1}\exp\!\bigl((64-\delta)fd\,x^\eps\bigr)
                    \qquad(\gamma\in\mathcal R).
\]
Then the absolute value of this mixed entry is at most
\(\exp(-(\delta-o(1))fd\,x^\eps)\).
\end{corollary}

\begin{proof}
Theorem~\ref{thm:ordinary164} gives absolute convergence and
vanishing of the zero-determinant mixed cells. Averaging over the
common unipotent group removes the irrelevant cells. Sum the
displayed inequality over the remaining rational labels and use
\eqref{eq:orbits204}.
\end{proof}

Theorem~\ref{thm:orbits204} and the cohomology estimates in
Theorem~\ref{thm:all202} use exactly the same \(E_x\) at each
point of the interval. The hypothesis of
Corollary~\ref{cor:amplitude204} concerns orbital integrals with
their actual volumes and normalizations. Its verification, the
terms for equal words, and the integral lower bound for the projected
Bessel coefficient \(p_{\lambda,H}\mathcal F A\) of
Corollary~\ref{cor:original200} remain open in this argument. On the
cohomological side the total order of the surviving groups in
\eqref{eq:region198} also remains to be bounded.

\subsection{A common real scale and the change of finite Fourier measure}\label{nav:VIII-30}
Take the fields of Theorem~\ref{thm:conductor232},
with its polynomial modulus. Write \(F=E_x^+\), \(E=E_x\),
\(f=[F:\QQ]\), and denote the fixed conductor of \(K_\psi\) by
\(F_\psi\). Keep the two split prime ideals \(\mathfrak p,\mathfrak q\)
and positive generators \(\beta,\gamma\) of
Proposition~\ref{prop:cost177} and Theorem~\ref{thm:balanced179}.
The rational ramified prime \(q\) and the split ideal \(\mathfrak q\)
lie over different rational primes.

\begin{theorem}[An exact real scale with the original skew form]
\label{thm:scale232}
Let \(\mathfrak d_F\) be the absolute different and let
\(\mathfrak a_0\OO_E=(\sqrt{-d_0})\), as in Lemma~\ref{lem:norm173}.
The ideal \(\mathfrak b=\mathfrak d_F\mathfrak a_0^{-1}\) is integral
and supported on places split in \(E/F\). Set
\(\beta'=(d_0\gamma)^{-1}\). Then
\[
  (\beta')\mathfrak d_F=\mathfrak b\mathfrak q^{-1},\qquad
  (\gamma)\mathfrak a_0=\mathfrak q,\qquad \beta'\gamma=1/d_0.
\]
The trace character scaled by \(\beta'\) has conductor \(\OO_v\) at
every inert finite place. At every real place keep the vacuum, coordinates,
and homogeneous tests of Proposition~\ref{prop:dilation179}. For every
tuple of degrees covered by that proposition,
\[
   R_\infty(\beta',\gamma)
             =d_0^{\sum_\tau D_\tau}R_\infty(1,1).
\]
For the two original conductor ideals, every choice of their generators
has nonrational product \(\beta\gamma\).
\end{theorem}
\begin{proof}
At \(d_0\) the only ramified quadratic layer in
\(K_\psi C_J(\sqrt{-q},\sqrt{-d_0})\) comes from \(\sqrt{-d_0}\).
The other factors have conductor prime to \(d_0\), and adjoining this
square root to an unramified completion has ramification index two.
At \(q\) the same argument uses \(\sqrt{-q}\). Inertia at either
prime fixes the other imaginary quadratic field, on which complex
conjugation acts nontrivially. Its intersection with
\(\operatorname{Gal}(E/F)\) is trivial. Thus the absolute
ramification index of \(F\) is two at both primes. Its different
exponent is one there. Indeed, over the maximal unramified subfield
of the completion, a ramified quadratic extension of odd residue
characteristic is generated by a root of \(X^2-a\) for \(a\) of
valuation one. The powers of that uniformizer form an integral basis,
and the derivative has valuation one.

Squaring the defining ideal of \(\mathfrak a_0\) and descending gives
\(\mathfrak a_0^2=(d_0)\). Its valuation is therefore one above
\(d_0\) and zero elsewhere. This proves integrality of \(\mathfrak b\)
and removes all places above \(d_0\) from its support. Every remaining
rational prime in the support divides \(2F_\psi q\), by the conductor
of the compositum. At an odd divisor \(p\) of \(F_\psi q\),
\(-d_0\equiv1\pmod p\), by the congruences in Theorem~\ref{thm:long232},
on every sufficiently distant interval. Hensel's lemma gives a square
root in \(\QQ_p\). At two, \(-d_0\equiv1\pmod8\) likewise makes it a
square in \(\QQ_2\). Since \(E=F(\sqrt{-d_0})\), every place above
these primes splits in \(E/F\).

The displayed ideal identities follow by multiplication, using
\((d_0)=\mathfrak a_0^2\). The local annihilator of \(\OO_v\) is the
inverse of \((\beta')\mathfrak d_{F,v}\), hence equals \(\OO_v\) at
an inert place. The skew scalar is still \(\gamma\sqrt{-d_0}\) with
the prescribed sign. At a real embedding the positive factor of the
symplectic pairing is \(\tau(\beta'\gamma)=1/d_0\). The Fock isometry
of Proposition~\ref{prop:dilation179} multiplies a degree-\(d\) vector
by \(d_0^{d/2}\) and its diagonal coefficient by \(d_0^d\). Multiplying
over the edges and the real places gives the formula.

Finally take a place \(v\mid q\) and normalize \(v\) to have value
group \(\ZZ\). The original ideals \(\mathfrak p,\mathfrak q\) were
chosen away from \(q\). Their defining equalities give
\(v(\beta\gamma)=-v(\mathfrak d_F\mathfrak a_0)=-1\), whereas
\(v(t)=2v_q(t)\) for every rational \(t\ne0\).
\end{proof}

\begin{proposition}[Finite Fourier measures and dual lattice indices]
\label{prop:fourier232}
Put \(u=\beta'/\beta\). At a finite place \(v\) fix a nonsingular
bilinear pairing of an \(n\)-dimensional \(F_v\)-space with its dual.
Use self-dual measures in dual bases, keeping the same coordinate
conventions for both additive characters. For every Schwartz function
\(h\), and every full lattice \(L\) in the domain, one has
\[
   \widehat h_{\beta'}(y)=|u|_v^{n/2}\widehat h_\beta(uy),
   \qquad L^\perp_{\beta'}=u^{-1}L^\perp_\beta.
\]
The support of the valuations of \(u\) consists of split places. For
all edges of the theta chain together, the sum of the logarithms of
the mutual lattice indices and of the absolute logarithmic changes
of Haar measure is \(O_{\psi,\varepsilon}(x^C)\) for a fixed \(C\),
uniformly in the lattices and in \(x\) on the selected intervals.
\end{proposition}
\begin{proof}
With the original measure the new Fourier kernel is the old kernel at
\(uy\). Applying the Fourier transform twice with this substitution
introduces \(|u|_v^{-n}\). Fourier inversion thus gives the positive
measure multiplier \(|u|_v^{n/2}\) in dual coordinates. The two
annihilators follow by testing the character on \(L\). For a lattice
\(M=L^\perp_\beta\), one of \(M,u^{-1}M\) contains the other, and
the product of their indices over their intersection equals
\((Nv)^{n|v(u)|}\). Here \(Nv\) is the residue field cardinality.

The two conductor equalities give
\((u)=\mathfrak b/(\mathfrak p\mathfrak q)\). The numerator and
denominator have disjoint support, by the choice of the two primes.
Consequently
\[
 \sum_{v<\infty}|v(u)|\log Nv
 =\log\frac{|D_F|}{d_0^{f/2}}+
              \log N\mathfrak p+\log N\mathfrak q.
\]
Indeed \(N\mathfrak d_F=|D_F|\) and \(N\mathfrak a_0=d_0^{f/2}\).
The conductor--discriminant formula gives \(|D_F|\le N_x^f\), since
each of the \(f\) characters of the abelian field \(F\) has conductor
dividing \(N_x\). The polynomial bound on \(N_x\), the degree bound
in Theorem~\ref{thm:conductor232}, and Proposition~\ref{prop:cost177} give a
polynomial bound for the displayed sum. A Schr\"odinger coordinate
space for an edge between ranks \(n_1,n_2\) has dimension \(n_1n_2\),
half the \(F_v\)-dimension of the tensor space. By
Section~\ref{sec:realtower169} the largest rank is
\(n_m=k+m-1+mt_{\rm r}\), so the sum of the coordinate dimensions
over the \(m\) edges is at most \(mn_m^2\).
Theorem~\ref{thm:rangenew50} gives
\(t_{\rm r}=O(x^{400})\) and polynomial bounds for \(k,m\).
Multiplication by three halves times this bound accounts for both
the lattice indices and the Haar multipliers.
\end{proof}

The Fourier identity evaluates the original transform at \(uy\).
For a prescribed terminal frequency \(y\), the integral period still
requires a lower bound at that transported argument and control of
the signed projected sum. Theorem~\ref{thm:transport233} gives an exact comparison when the Schwartz functions and splittings are transported together.

\subsection{Rational orbits meeting the retained columns}\label{nav:VIII-31}
Use \(Y,Q,T,d,f\) from Corollary~\ref{cor:gram235}. For a
relevant rational label \(\gamma\), let \(D_j\) be its leading
cross minors. Its weight is
\(w_\gamma=\prod_{j:D_j\ne0}\prod_\tau(1+|\tau D_j|^2)^2\),
where one complex embedding above each real place is used. Let
\(\mathcal R_Y\) consist of the labels with \(\mathcal G\ne0\)
for which some original finite cell has retained initial column.

\begin{theorem}[Weighted count of the supported rational labels]
\leavevmode\label{thm:orbits235}
For \(0\le u<2\), with \(M_u=\int_0^\infty
4r^{3+u}(1+r^2)^{-3}\,dr\), one has
\[
 \sum_{\gamma\in\mathcal R_Y}w_\gamma^{-1}
 \le(9M_u+2)^{fd}Q^{-(2+u)fd}Y^{-u/2}.
\]
On the common intervals this is at most
\(\exp(-(98-o(1))fd\,x^\eps)\). On labels with a single
reversed block of length \(d\) it is at most
\(\exp(-(132-o(1))fd\,x^\eps)\).
\end{theorem}
\begin{proof}
The nonzero leading minors have as indices the endpoints \(J\)
of a composition of \(d\), including \(d\). At the nonterminal
endpoints \(D_j\in Q^j\mathcal O_E\). At \(d\) the inclusion is
\(D_d+TD_{d-1}\in Q^d\mathcal O_E\) by
Proposition~\ref{prop:shift203}. Thus, with the preceding
values fixed, the product of the complex distances between two
distinct values of \(D_j\) is at least \(Q^{jf}\). The product
of the magnitudes of each nonzero \(D_j\), \(j<d\), has the
same lower bound. At the terminal endpoint,
\[
 \prod_\tau|\tau D_d|\ \ge\
 |N_{F/\QQ}(\,D_d\bar D_d-T^2D_{d-1}\bar D_{d-1}\,)|^{1/2}
 >Q^{df}\sqrt Y .
\]
Here both the total positivity and the cutoff refer to the very
same cell. Since \(Y>1\) eventually, all endpoints satisfy the
ordinary lower bound as well.

In a product of discs of coordinate radii \(r_{j,\tau}\), put
\(R_j=(\prod_\tau r_{j,\tau})^{1/f}\). Dividing each coordinate
by its radius puts the centers in a unit polydisc, separated in
sup norm by at least \(Q^j/R_j\). Packing polydiscs of half
that radius gives at most \(9^f(R_j/Q^j)^{2f}\) centers whenever
the box is nonempty. Bound the support indicator by
\((R_j/Q^j)^{uf}\) at \(j<d\), and by
\((R_d/Q^d)^{uf}Y^{-u/2}\) at \(d\).
These inequalities apply successively with all preceding values fixed.
Integrate each radius against \(4r(1+r^2)^{-3}\,dr\). Its tail
from \(|\tau D_j|\) equals \((1+|\tau D_j|^2)^{-2}\), so
Tonelli gives a factor \((9M_u)^fQ^{-(2+u)jf}\) at each
endpoint and a single additional factor \(Y^{-u/2}\).
For block lengths \(h_i\), the number of scalar tuples with
the given minors is at most \(\prod_i h_i\le2^{d-|J|}\),
since each successive scalar satisfies an equation of degree
\(h_i\). Sum over the subsets of \(\{1,\ldots,d-1\}\) of
earlier endpoints, retaining at least the power of \(Q^{-1}\)
from \(d\). The binomial formula gives the first bound.

For one block write \(D_d=\pm c^d\). All earlier minors vanish
when \(d>1\); we are on intervals with \(d\to\infty\).
The terminal congruence implies \(c/Q\in\mathcal O_E\), by the
valuation at every finite prime. The nonzero centers \(c\) are
separated in the product of embedding distances by \(Q^f\).
In a box of geometric mean radius \(R_0\), the preceding
terminal inequality requires \(R_0>QY^{1/(2df)}\). Packing as
above, and using the indicator with exponent \(af\), bounds
this subsum by
\[
 (9M_{d,a})^f Q^{-(2+a)f}Y^{-a/(2d)},\qquad
 M_{d,a}=\int_0^\infty\frac{4d\,r^{2d+1+a}}{(1+r^{2d})^3}\,dr .
\]
Indeed the new radial weight is \((1+r^{2d})^{-2}\); its
negative derivative supplies the density. For \(0\le a<4d-2\),
splitting at one gives \(M_{d,a}\le2+4d/(4d-2-a)\).

Take \(u=2-1/\log(2Q)\) and \(a=4d-2-1/\log(2Q)\).
The bound \(M_u\le1+4/(2-u)\) follows by the same split.
We have \(\log Q=(16+o(1))x^\eps\) and
\(\log Y=(34+o(1))fd\,x^\eps\) by \(m=2d\) and
Theorem~\ref{thm:prime209}. All logarithms of the moment
prefactors are \(o(fd\,x^\eps)\) after their indicated powers.
The two exponents in the statement follow.
\end{proof}

For a matrix entry supported on these columns, an upper bound
of the form \(\sum_b|J_{\gamma,b}|\le
w_\gamma^{-1}\exp((98-\eta)fd\,x^\eps)\), with fixed
\(\eta>0\), would therefore bound the entire nondegenerate part
by \(\exp(-(\eta-o(1))fd\,x^\eps)\). Here \(b\) runs through
the actual finite cells and \(J_{\gamma,b}\) includes the
archimedean integral. The theorem controls the sum of label
weights. Bounding these integral amplitudes remains necessary.

\subsection{Counting the last block scalar}\label{nav:VIII-32}
Retain the columns of Theorem~\ref{thm:columns235}, with their
hyperplane threshold \(Y\). For a single reversed block,
Theorem~\ref{thm:orbits235} counted the block scalar before
raising it to the block length. The same calculation works with
preceding blocks once their contribution to the denominator of
the last scalar is included.

\begin{theorem}[A bound uniform in every terminal block length]
\leavevmode\label{thm:powers236}
On the same intervals, for all the supported nondegenerate
labels meeting the retained columns,
\[
             \sum_{\gamma\in\mathcal R_Y} w_\gamma^{-1}
                      \le \exp\bigl(-(132-o(1))fd\,x^\eps\bigr).
\]
\end{theorem}
\begin{proof}
Let \(h\ge2\) be the length of the last reversed block and set
\(j=d-h\). For fixed preceding data put \(A=D_j\), taking
\(D_0=1\). Then \(A=Q^j a\) for a nonzero
\(a\in\mathcal O_E\). The final minor is \(D_d=\pm A c^h\).
Since \(D_{d-1}=0\), Proposition~\ref{prop:shift203} gives
\(a(c/Q)^h\in\mathcal O_E\). At every prime \(\mathfrak p\)
this says \(v_{\mathfrak p}(c/Q)\ge
-\lfloor v_{\mathfrak p}(a)/h\rfloor\).
Consequently \(c/Q\) ranges in a fractional ideal \(\mathfrak l\)
with \(N\mathfrak l\ge |N_{E/\QQ}a|^{-1/h}\).

In the complex embedding \(\tau\) count the possible scalars in
the variable \(v_\tau=|\tau A|^{1/h}\tau c\). For distinct
choices of \(c\) the difference divided by \(Q\) belongs to
\(\mathfrak l\), and the norm of a nonzero element of this ideal
is at least \(N\mathfrak l\). Thus the product of the coordinate
distances is at least
\[
  |N_{E/\QQ}A|^{1/(2h)}Q^f(N\mathfrak l)^{1/2}
                  \ge Q^{df/h}.
\]
The dependence on the earlier value \(a\) cancels. Also
\(|v_\tau|^h=|\tau D_d|\), so that in any box of these variables
with geometric mean radius \(R_0\) a retained contribution
requires \(R_0>Q^{d/h}Y^{1/(2hf)}\) by
Corollary~\ref{cor:gram235}. The count in that box is at most
\(9^f(R_0/Q^{d/h})^{2f}\) when it is nonempty, by the polydisc
packing argument of Theorem~\ref{thm:levelsum203}.
For \(0\le b_0<4h-2\), bound the support indicator by
\((R_0/Q^{d/h})^{b_0 f}Y^{-b_0/(2h)}\). Integrating the
negative radial derivative of \((1+r^{2h})^{-2}\) gives the
following bound for the sum of the final weight at the fixed prefix:
\[
 (9M_{h,b_0})^f Q^{-(2+b_0)df/h}Y^{-b_0/(2h)},\qquad
 M_{h,b_0}\le 2+\frac{4h}{4h-2-b_0}.
\]
As in Theorem~\ref{thm:orbits235}, the moment is the integral
of \(4h\,r^{2h+1+b_0}(1+r^{2h})^{-3}\). Counting \(c\) itself
already includes all possible last-block roots.

Write \(\eta=1/\log(2Q)\). Sum the earlier block scalars with
their original weights. All their endpoint minors \(D_i\) are
nonzero multiples of \(Q^i\). Packing at each earlier endpoint,
inserting the exponent \(2-\eta\) and summing over compositions
as in Theorem~\ref{thm:levelsum203}, bounds that sum by
\[
       (9M_{2-\eta}+2)^{fj}Q^{-(4-\eta)jf}.
\]
For \(j=0\) the empty prefix contributes one. In this estimate
the power of \(Q^{-1}\) comes from the final prefix endpoint
\(j\); the root multiplicities \(\prod h_i\) of the prefix
are included in the binomial sum. It applies to any subset of
prefixes because all summands are nonnegative.

Choose \(b_0=4h-2-\eta\). Uniformly over \(2\le h\le d\),
the logarithms of the two moment prefactors together are
\(o(fd\,x^\eps)\): the first is \(O(f\log(d\log(2Q)))\)
and the second \(O(fd\log\log(2Q))\). The saving from the
last endpoint alone, divided by \(fd\,x^\eps\), is
\(132-34/h+o(1)\), since \(\log Q=(16+o(1))x^\eps\) and
\(\log Y=(34+o(1))fd\,x^\eps\). For \(h\ge\sqrt d\) this
suffices. For \(2\le h<\sqrt d\) the prefix contributes a
further \(64j/d+o(1)=64-o(1)\). The sum over the at most \(d\)
lengths only costs \(\log d\) in the logarithm. When \(h=1\)
the terminal congruence is shifted. The terminal-singleton
estimate of Theorem~\ref{thm:orbits204}
instead saves at least \(192-o(1)\) on the same scale.
This completes the proof.
\end{proof}

Accordingly the amplitude hypothesis following
Theorem~\ref{thm:orbits235} suffices with \(132-\eta_0\) in
place of \(98-\eta_0\), for any fixed \(\eta_0>0\).

\section{Rational orbits on the first plane}\label{nav:g-XI-2}

\subsection{Short rational translates at the initial period}\label{nav:IX-34}
Keep the field \(E/F\) and the bottom decomposition \(X=Y\perp Ez\)
of Proposition~\ref{prop:depth245}. Write \(h\) for its Hermitian form,
linear in the second variable, and \(\eta=h(z,z)\in\{1,-1\}\).
Write \(G=U(X)\), \(H=U(Y)\), and \(\iota:H\to G\) for extension
by the identity on \(z\). Fix \(g\in H(\mathbb A_F)\). The finite
lattice and the real norms at this point are
\[
 \Lambda_g=X(E)\cap\iota(g_f)\widehat\Lambda_X
           =\Lambda_{Y,g}\oplus\mathcal O_Ez,\qquad
 |y+bz|_{\tau,g}^2=|y|_{\tau,g}^2+|\tau b|^2 .
\]
Here \(X(E)\) denotes the rational vectors of the \(E\)-space \(X\).
At each real place of \(F\) choose a frame in which \(h\) has
diagonal entries \(\pm1\) and use the Euclidean norm after applying
\(\iota(g_\tau)^{-1}\) in the adapted frame. Thus
\(|\tau h(v,w)|\le |v|_{\tau,g}|w|_{\tau,g}\).
For a nonzero rational vector put
\[
 \mathfrak c_g(v)=
 \{\lambda(v):\lambda\in\operatorname{Hom}_{\mathcal O_E}
                    (\Lambda_g,\mathcal O_E)\},\qquad
 H_g(Ev)=\frac{\prod_\tau |v|_{\tau,g}}{\sqrt{N\mathfrak c_g(v)}}.
\]
Each integral linear form is extended \(E\)-linearly, and \(N\)
denotes absolute ideal norm. The product runs over one embedding
in each conjugate pair. Scaling \(v\) scales the numerator and
denominator by the same positive number. Let \(W_g\) be the
\(E\)-span of the lines of height strictly less than one.

For the original nonzero integral modulus \(\mathfrak m\), set
\(\mathfrak b=\mathfrak m\cap\bar{\mathfrak m}\). Write \(K(\mathfrak m)\)
for the unitary stabilizer of \(\widehat\Lambda_X\) acting trivially
modulo \(\mathfrak m\widehat\Lambda_X\). A class \(\gamma H(F)\)
in \(G(F)/H(F)\) meets the finite test if
\(\iota(g_f)^{-1}\gamma\iota(h_f)\in K(\mathfrak m)\) for some
\(h_f\in H(\mathbb A_{F,f})\).
In the grid, \(\mathfrak m=\prod_{v\in S_0}\mathfrak w_v^{e_v}\),
with one prime \(\mathfrak w_v\) above each chosen split place
\(v\) of \(F\). Hence \(\mathfrak m\) and \(\bar{\mathfrak m}\)
are coprime, \(\mathfrak b=\mathfrak m\bar{\mathfrak m}\), and
\(\sqrt{N\mathfrak b}=N\mathfrak m\). The choice of the primes
and modulus is recorded in
\srcloc{Development/AbelianVarietyTheory/Definitions/Divisors/BasePlaceNotMemChiGridBadPlaces}{43--82}.

\begin{theorem}[All short displacements at the original modulus]
\label{thm:short246}
\(W_g\) is a totally isotropic subspace of \(Y\). Every class
meeting the finite test has displacement \(w=\gamma z-z\) in
\(\mathfrak b\Lambda_g\). The nonidentity classes meeting the test
and satisfying \(\prod_\tau |w|_{\tau,g}<\sqrt{N\mathfrak b}\)
correspond bijectively, by their displacements, to the vectors
\[
 0\ne w\in \mathfrak b\Lambda_{Y,g}\cap W_g,\qquad
                  \prod_\tau |w|_{\tau,g}<\sqrt{N\mathfrak b}.
\]
These assertions hold at every \(g\), including on the dyadic
components of the stated fields.
\end{theorem}
\begin{proof}
In local lattice bases the content ideal is generated by the
coordinates of \(v\). Thus \(v\in\mathfrak a\Lambda_g\) exactly
when \(\mathfrak c_g(v)\subset\mathfrak a\), for any fractional
ideal \(\mathfrak a\). Integrality of the Hermitian form on this
lattice gives \(h(v,w)\in\overline{\mathfrak c_g(v)}\mathfrak c_g(w)\).
When \(h(v,w)\ne0\), taking norms of the principal ideal and using
the real Cauchy--Schwarz inequalities yields \(H_g(Ev)H_g(Ew)\ge1\).
Taking \(v=w\) and taking two distinct lines shows that all the
short lines are pairwise orthogonal and isotropic. If \(v=y+bz\)
with \(b\ne0\), the integral projection onto the coefficient of
\(z\) gives \(b\in\mathfrak c_g(v)\). Then
\(\prod_\tau|v|_{\tau,g}\ge |N_{E/\QQ}b|^{1/2}\ge
\sqrt{N\mathfrak c_g(v)}\). This places \(W_g\) in \(Y\).

The lattice \(\widehat\Lambda_X\) is Hermitian self-dual.
If \(k\in K(\mathfrak m)\), then \(k^{-1}-1\) also takes this
lattice into \(\mathfrak m\widehat\Lambda_X\). Taking its Hermitian
adjoint shows that \(k-1\) takes it into
\(\bar{\mathfrak m}\widehat\Lambda_X\): duality gives this directly
by pairing against each integral vector. Hence
\(K(\mathfrak m)=K(\mathfrak b)\). Since \(\iota(h_f)z=z\),
the finite support condition implies \(w\in\mathfrak b\Lambda_g\).
If \(w\ne0\) satisfies the strict product bound in the statement,
the content inclusion gives \(H_g(Ew)<1\). Distinct classes yield
distinct vectors because the rational stabilizer of \(z\) is
exactly \(H(F)\).

For the converse the field contains
\(t=(1+\sqrt{-q})/2\in\mathcal O_E\) with \(t+\bar t=1\),
since the selected prime is \(3\bmod 4\). Given such a \(w\), set
\[
      N_w(v)=\eta^{-1}\big(w h(z,v)-z h(w,v)\big),
                  \qquad T_w=1+N_w+tN_w^2 .
\]
Orthogonality and isotropy give \(N_w(z)=w\), \(N_w(w)=0\),
\(N_w^*=-N_w\), and \(N_w^3=0\), since its image lies in
\(Ew+Ez\). Consequently
\(T_w^*T_w=1+(t+\bar t-1)N_w^2=1\), and \(T_wz=z+w\).
Since \(w\in\mathfrak b\Lambda_g\) and
\(\bar{\mathfrak b}=\mathfrak b\), both \(N_w\) and \(T_w-1\)
carry \(\Lambda_g\) into \(\mathfrak b\Lambda_g\).
The same holds for \(T_w^*-1=-N_w+\bar tN_w^2\). Hence
\(\iota(g_f)^{-1}T_w\iota(g_f)\in K(\mathfrak b)\), using
\(\widehat{\Lambda_g}=\iota(g_f)\widehat\Lambda_X\) by the
local description of a lattice. This gives the required rational
class at exactly the original test. The integral element \(t\)
works at the primes above two as well.
\end{proof}

\subsection{Short vectors in every terminal Hermitian lattice}\label{nav:IX-38}
Let \(N=n_m\) be the rank of the terminal even core, \(f=[F:\QQ]\),
\(D=2fN\), and \(D_E\) the absolute field discriminant.
Use the definite Hermitian lattices \(\Lambda_c\) of
Proposition~\ref{prop:terminalcore255} and their given level
\(\mathfrak m\). Set
\[
 R=2D^{D/2}|D_E|^{N/2},\qquad B=(\sqrt N\,R)^{Nf},\qquad
 H_*=(1+2R^2)^{2fN^2} R^{4fND}
                        \operatorname{Nm}(\mathfrak m)^{N^2}.
\]

\begin{lemma}[Uniform bases over the field and a class count]\leavevmode
\label{lem:terminalclass259}
For each class \(c\) there are \(x_1,\ldots,x_N\in\Lambda_c\)
forming an \(E\)-basis. In the Hermitian metrics at every selected
complex embedding the column map from the standard orthonormal
basis to these vectors has all singular values in
\([B^{-1},B]\). Contraction with the same vectors gives integral
linear functionals \(\xi_1,\ldots,\xi_N\) in the dual lattice.
Their column map into the Hermitian dual has the same bounds.
The number of terminal level classes satisfies \(h_m\le H_*\).
\end{lemma}
\begin{proof}
Write \(\Lambda=\Lambda_c\) and \(h\) for its definite Hermitian
form. Give its underlying \(\ZZ\)-lattice the positive bilinear
form \(\operatorname{Tr}_{E/\QQ}h\). Its dual for this form is
\(\mathfrak D_{E/\QQ}^{-1}\Lambda\). Indeed trace integrality
against all multiples of an element of \(\Lambda\) is precisely
the condition that its Hermitian pairing take values in the
inverse different. Hermitian unimodularity gives the formula
after localization at each ideal. The index of \(\Lambda\)
in the trace dual is \(|D_E|^N\), since \(\Lambda\) has local
rank \(N\). The Euclidean covolume is therefore
\(|D_E|^{N/2}\). Each nonzero vector has squared trace length a
positive integer and hence length at least one.

The Euclidean unit ball contains the cube of side \(2/\sqrt D\).
Minkowski's second theorem bounds the product of the \(D\)
successive minima by \(D^{D/2}|D_E|^{N/2}\). The preceding
lower bound for the first minimum gives \(D\) independent
vectors of length at most \(R\). They span the rational space,
so we can choose an \(E\)-basis \(x_1,\ldots,x_N\) among them.
Write \(X_\tau\) for their column map at \(\tau\). Each column
has norm at most \(R\) since its squared norm contributes to
the trace norm. The Gram determinant
\(\Delta=\det(h(x_a,x_b))\) is a nonzero totally positive
integer of \(F\). At each real embedding its value is
\(|\det X_\tau|^2\le R^{2N}\) by Hadamard's inequality.
Its norm to \(\QQ\) is an integer at least one. Thus
\(|\det X_\tau|\ge R^{-N(f-1)}\) at each embedding. All
singular values are at most \(\sqrt N R\). Dividing the
determinant by the product of the largest \(N-1\) bounds the
least singular value below by \(B^{-1}\). The Hermitian form
identifies each \(x_a\) with a linear functional \(\xi_a\) by
placing \(x_a\) in the conjugate-linear slot. Integrality and
the equality of singular values follow from this anti-isometry.

For the count, every class at maximal level is determined by
the isometry type of its lattice in the rational Hermitian
space. Conversely a rational isometry carrying one lattice to
another relates their representatives by the adelic stabilizer.
In the basis just found every Gram entry belongs to
\(\mathcal O_E\) and has modulus at most \(R^2\) at all
embeddings, by Cauchy--Schwarz. Two distinct algebraic integers
differ by at least one in maximum modulus over the embeddings,
by the norm of their difference. In \(\CC^f\) with maximum
complex modulus, the open balls of radius \(1/2\) about these
integers are disjoint and lie in the ball of radius \(R^2+1/2\).
Comparing real volumes bounds their number by
\((1+2R^2)^{2f}\). Taking all \(N^2\) entries gives the first
factor in \(H_*\).

For a fixed Gram matrix use \(M=\sum_a\mathcal O_E x_a\) as
coordinate lattice. Then \(M\subseteq\Lambda\subseteq M^\#\),
where the sharp denotes its Hermitian dual. In this basis the
determinant index gives
\(|M^\#/M|=|\operatorname{Nm}_{E/\QQ}\Delta|\le R^{4fN}\).
The image \(\Lambda/M\) is generated as an abelian group by
at most \(D\) elements since \(\Lambda\) has \(\ZZ\)-rank \(D\).
There are at most \(|M^\#/M|^D\) such subgroups, by counting
ordered generating lists padded with zero. Each preimage
defines at most one lattice for this fixed matrix. Finally,
the number of level classes over one maximal-level class
is at most the level index. At a marked split place this index
is \(|\mathrm{GL}_N(\mathcal O_v/\varpi_v^{a_v})|
\le (Nv)^{a_vN^2}\), and elsewhere it is one. Multiplying
these bounds proves \(h_m\le H_*\).
\end{proof}

\subsection{Separation of the arithmetic returns at the chosen point}\label{nav:IX-39}
Let \(G=U(C_0)\), \(H=U(Y)\), \(K_f=K(\mathfrak m)\), and
\(g\in H(\mathbb A_F)\) be a point of
Theorem~\ref{thm:ratio248}. Put \(c=\iota(g)\). The real
subgroup associated with its level component, in coordinates
centered at \(c_\infty\), is
\[
 \Gamma_c=
 \{c_\infty^{-1}\gamma_\infty c_\infty:
             \gamma\in G(F),\ c_f^{-1}\gamma_fc_f\in K_f\}.
\]
Write \(p=4n^2+2,\ d_{\mathbb R}=2fn,\ w(h)=\prod_\tau W(h_\tau)\),
where \(W\) is the defining operator norm. With the line height
\(b_g\) of Theorem~\ref{thm:packing246}, define
\[
 \lambda=b_g^{1/f},\qquad
 R_g=2d_{\mathbb R}^{d_{\mathbb R}/2}|D_E|^{n/2}
                         \lambda^{-(d_{\mathbb R}-1)},\qquad
 \rho=\min\{1/4,\lambda/(16 R_g)\}.
\]

\begin{lemma}[A uniform sum over all arithmetic returns]\leavevmode
\label{lem:returns261}
Every nonidentity member of \(\Gamma_c\) satisfies
\(\max_\tau\|\gamma_\tau-I\|_{\rm op}\ge\lambda/R_g\).
For the product real group one has
\[
 \sum_{\gamma\in\Gamma_c} w(\gamma)^{-p}
 \le 2^{fp}\,T_0^{Cfn^3}(1+4/\rho)^{2fn^2}
\]
with \(T_0\) as in Theorem~\ref{thm:first250}. At every
point above half the supremum of its prescribed lowest-type
row, the logarithm of this bound is at most a polynomial in
\(x\), uniform in the row, on the same original intervals.
\end{lemma}
\begin{proof}
Use the underlying integral lattice
\(\Lambda= C_0\cap c_f\widehat\Lambda_{C_0}\), using the
diagonal embedding of the Hermitian \(E\)-space.
At infinity use the positive Euclidean metrics in the frames
translated by \(c_\infty\) and sum their squared norms. Every
nonzero \(v=y+t z\) in \(\Lambda\) has product of its local
norms at least \(b_g\). Indeed the lattice is
\(\Lambda_{Y,g}\oplus\mathcal O_Ez\), and if \(t\ne0\) that
product is at least \(|N_{E/\mathbb Q}t|^{1/2}\ge1\).
If \(t=0\) and \(h(y,y)\ne0\), the integral value \(h(y,y)\)
has nonzero integer norm in \(F\), and
\(|\tau h(y,y)|\le |y|_{\tau,g}^2\) everywhere. In the
isotropic case the product is \(H_g(Ey)\sqrt{N\mathfrak c_g(y)}\),
with integral content, and \(H_g(Ey)\ge b_g\). This uses the
original \(g\) and its level lattice. In particular every
nonzero Euclidean length is at least \(\lambda\).

The absolute determinant of the trace form on this
\(\mathbb Z\)-lattice is \(|D_E|^n\). Its trace dual is
\(\mathfrak D_E^{-1}\Lambda\) by Hermitian unimodularity,
with this index over \(\Lambda\). In the positive metric the
trace form is twice the Euclidean form composed with the
orthogonal signature involution. The Euclidean covolume is
therefore \(2^{-fn}|D_E|^{n/2}\). Applying Minkowski's second
theorem in dimension \(d_{\mathbb R}\), as in
Lemma~\ref{lem:terminalclass259}, gives \(d_{\mathbb R}\)
independent integral vectors of length at most \(R_g\), using
the lower bound for all preceding minima. A nonidentity
arithmetic return moves at least one of these vectors.
Its displacement is a nonzero integral vector, since the finite
return belongs to the stabilizer of the level lattice. Comparing
its length to that of the starting vector gives the stated
separation.

Let \(U_\rho\) be the product of the local operator-norm
balls of radius \(\rho\) about \(I\) in the groups. The
sets \(\gamma U_\rho\) are disjoint, since the distance to
\(I\) of a local member of \(U_\rho U_\rho^{-1}\) is at most
\(2\rho/(1-\rho)<\lambda/R_g\).
At one place cover the group ball
\(\|h-I\|_{\rm HS}\le1/4\) by translates of \(U_{\rho,\tau}\).
A maximal \(\rho/4\)-separated set in the Hilbert--Schmidt
norm does this, since the inverses of its centers have
operator norm at most \(4/3\). The disjoint Euclidean balls
of radius \(\rho/8\) about the centers lie inside the ball of
radius \(1/4+\rho/8\). Thus at most
\((1+4/\rho)^{2n^2}\) translates suffice. Write \(d_\pi\)
for the product formal degree of the core series. Applying
the lower volume estimate of Theorem~\ref{thm:tensor249} at
each place gives
\[
 d_\pi\operatorname{vol}(U_\rho)
       \ge T_0^{-Cfn^3}(1+4/\rho)^{-2fn^2}.
\]
For \(t\in\gamma U_\rho\) one has
\(w(t)^{-p}\ge2^{-fp}w(\gamma)^{-p}\). Integrating on the
disjoint translates and using the upper integral estimate of
the same theorem, \(d_\pi\int w(t)^{-p}\,dt\le T_0^{Cfn^2}\),
proves the sum bound after increasing \(C\).

At a point above half the supremum,
\(\log(1/b_g)\ll f^2n^5\log b\) by the proof of
Theorem~\ref{thm:ratio248}. Also \(\log T_0\ll\log b=x^\eps\),
\(f,n\) are polynomially bounded, and
\(\log|D_E|\le2f\log M\) for the cyclotomic modulus of
Lemma~\ref{lem:dimension239}. Its logarithm is \(o(\log b)\).
The definitions of \(\lambda,R_g,\rho\) now give the claimed
polynomial bound for the logarithm.
\end{proof}

\subsection{The finite mass of a nontransverse orbit on the common interval}\label{nav:IX-48}
For a first-plane orbit let \(m_f(\gamma)=\prod_{v<\infty}
m_v(\gamma)\) with the quotient measures of
Theorem~\ref{thm:finite281}. Equivalently this is the mass in
the finite adelic orbital quotient divided by its identity
counterpart. Indeed the groups and their test sets are
restricted products, the local stabilizers are compact open,
and after division all but finitely many factors are one.
This measure is taken before pairing the column coordinates,
inserting the unipotent characters, or multiplying by the
real orbital integral. Write \(f=[F:\QQ]\).

\begin{theorem}[A finite mass at the original cuspidal period level]\leavevmode
\label{thm:product281}
On each interval of Theorem~\ref{thm:conductor232} and for
every sufficiently large \(x\) in it, use its very primes
\(q,d_0,Q,T\) and set \(c=Q\sqrt{-q}\). For \(b_v=\mathrm Nv\)
the entire finite factor above \(\gamma_c\) is
\[
 m_f(\gamma_c)=\prod_{v\mid q} b_v^{\,2s-2}
  \left[b_v^{2k}+
      \bigl(b_v-3+2s(b_v-1)^2\bigr)b_v^k+1\right].
\]
The rational label has weight \(w_{\gamma_c}^{-1}
=(1+Q^2q)^{-2f}\) in Theorem~\ref{thm:sum237}. Uniformly
on these intervals,
\[
 \log m_f(\gamma_c)
  =(k+s-1)f\log q+O\bigl(fs\,q^{\,2-k}\bigr),
 \qquad \frac{\log(m_f(\gamma_c)/w_{\gamma_c})}
                {fk\log q}\ \longrightarrow\ 1 .
\]
More generally, for any \(c\) of Proposition~\ref{prop:plane281}
whose ideal beyond \(Q\mathcal O_E\) is the extension of a
squarefree product of split places outside \(QT\), the exact
product runs over those places in place of \(v\mid q\).
\end{theorem}
\begin{proof}
The two primes of \(Q\) split completely in \(E\) and differ
from \(q,d_0,T\). Thus \(c\) has valuation one there and is a
unit elsewhere except above \(q\). Theorem~\ref{thm:finite281}
gives local mass one at all these minimal valuations,
including at two and at the terminal graph. At \(q\),
\(-d_0\equiv1\bmod q\) and \(q\) is odd. The simple roots
of \(X^2+d_0\) lift to \(\QQ_q\). Since
\(E=F(\sqrt{-d_0})\), every place of \(F\) above \(q\)
splits in \(E\).

The fields \(K_\psi,C_J\) are unramified at \(q\) by their
conductors. Adjoining the square root of \(-q\) to their
local compositum has ramification index two by the Eisenstein
polynomial \(X^2+q\). The other quadratic field is already
split there. This proves that each completion of \(F\) at
\(q\) has ramification index two over \(\QQ_q\) and that \(\sqrt{-q}\)
has valuation one in either split component. The row level
there is the full lattice stabilizer and \(T\) is a unit.
Apply Theorem~\ref{thm:split281} at every such place.
The same two local theorems give the general squarefree
assertion. Proposition~\ref{prop:plane281} gives the single
minor; each of its archimedean squared moduli equals \(Q^2q\).

At a place over \(q\) the bracket divided by \(b_v^{2k}\)
is \(1+O(s b_v^{2-k})\), with an absolute constant, by
direct division. Its logarithm satisfies the same bound
for large \(k\), since \(s=O(\sqrt{x})\) and the original
column ranks satisfy \(m^5x^\varepsilon=o(k)\) as in the
dimension estimate for the chosen words. The local degree
formula gives \(\sum_{v\mid q}\log b_v=(f/2)\log q\).
These identities prove the asymptotic for \(\log m_f\).
Finally \(\log w_{\gamma_c}=2f\log(1+Q^2q)\) and
\(\log Q=(16+o(1))x^\varepsilon\) on the same intervals.
Since \(s=o(k)\) and \(x^\varepsilon=o(k\log q)\), division
by \(fk\log q\) gives the limit.
\end{proof}

The positive finite factor can therefore exceed the weight used
to count rational labels even for a single orbit. To estimate
its contribution to the signed trace of
Theorem~\ref{thm:range239}, the two column coordinates and
their phases must be contracted with the real kernel.
The signed trace and the normalized integral theta pairing
still require estimates for this contraction and for the other orbits.

\subsection{The additive phases and the operator on the column}\label{nav:IX-49}
Use \(c=Q\sqrt{-q}\) in Theorem~\ref{thm:product281} and
the finite Bessel seed of equation~\eqref{eq:seed127} with
the original character \(\chi_\varpi\). At \(v\mid q\)
let \(A_v\) denote the sum of the translations to all
neighbors of the first split column lattice. This is the
unnormalized spherical Hecke operator supported on the
double coset with exponents \((1,0,\ldots,0,-1)\).
Use the normalization in which the identity orbital factor
acts by the identity on the column.

\begin{theorem}[The finite operator with its Bessel phases]\leavevmode
\label{thm:phase282}
For the right-invariant seed at level \(K^0\) the finite
operator for \(\gamma_c\) is the product over \(v\mid q\)
of
\[
                     F_v=a_b\,1+u_b A_v,\qquad b=\mathrm Nv ,
\]
with \(a_b,u_b\) of Proposition~\ref{prop:lattice282}.
All the local additive phases in this formula are one.
\end{theorem}
\begin{proof}
At \(v\mid q\), \(F_v/\QQ_q\) has ramification index two,
and \(q\) is odd. Its different has exponent one. For example,
in either split completion the extension is obtained by
adjoining \(\sqrt{-q}\) to an unramified field. The Eisenstein
generator in this totally ramified quadratic extension gives
an integral basis \(1,\sqrt{-q}\): in \(a+b'\sqrt{-q}\)
the two terms have different valuation parities. Its different
is therefore generated by \(2\sqrt{-q}\), also by direct
trace duality in that basis. The absolute trace character
is trivial on \(c^{-1}O\), the inverse different at \(v\).
This holds in both split factors.

Take each first-plane representative \(n(l,H,z)\) from
Theorem~\ref{thm:split281}. Its simple flag entries have
denominators dividing \(c\) on both sides. The first finite
phase is therefore one. To compute the phase at
\(g^{-1}=\gamma_c^{-1}n(l,H,z)\), take the integral primitive
representatives constructed in that same proof. If \(i=\infty\)
the primal list needs no changes. Otherwise every change uses
the fixed combination \(Au+cv\) of its earlier integral
columns with \(A\) a unit. The coefficients added to a
later column have denominators dividing \(c\). The dual
changes likewise have denominators dividing \(c\), using
\(D\) a unit when \(j<\infty\). Integral lifts of the
coordinates of \(X,Y\) only multiply these coefficients by
integers. Denote these upper unitriangular changes by \(V\)
in the two split factors. By Lemma~\ref{lem:flag215} an
element \(k_0\in K_v^0\) carries the reference decorated
list to \(g\) times the reference list changed by \(V\).
Thus \(k_0^{-1}g\in B(F_v)\) acts on the list by \(V^{-1}\).
For a unitriangular matrix, inversion negates its first
superdiagonal. The simple entries for \((k_0^{-1}g)^{-1}\),
which is the left factor of \(g^{-1}\), are consequently
in \(c^{-1}O\) too. Passing between \(e\) and \(\varpi e\)
uses only units here. The trace character of their sum is
one. These computations use the conductor of precisely the
character occurring in the period.

At a minimal place Theorem~\ref{thm:finite281} has the
representative \(1\) and \(\gamma_c^{-1}\in K_v^0\).
At \(v\mid q\) the first column is unchanged and its second
column is the lattice just calculated. In integrating a pair
of column vectors over their common column argument, a left
and a right integral change of the relative position leave
the pairing unchanged by Haar invariance and the level
condition. The stabilizer weights count precisely the
right \(U_L\)-cosets by Theorem~\ref{thm:finite281}.
Proposition~\ref{prop:lattice282} now gives the identity
and neighbor coefficients with their stated multiplicities.
Different places act by commuting translations. The
restricted product has just the nonidentity factors displayed.
\end{proof}

For a determinant character of the row level that is trivial
on all its finite units the same formula applies directly.
At a split full level, any unit character through the
determinant can also be removed locally by twisting the
translation action by its inverse, extended to give a
uniformizer value one. Relative column representatives here
can be chosen with determinant one, since both column
lattices have equal determinant valuation. These observations
leave the formula for the ordinary seed independent of
choices of bases or of uniformizer.

\subsection{An explicit tail for rational first-plane orbits}\label{nav:IX-50}
Keep the original fields, intervals, word, finite level and
terminal frequency. Let \(\mathcal T\) consist of the positive
squarefree integers coprime to \(2q d_0QT\) all of whose prime
divisors split completely and are unramified in \(E\).
There are infinitely many such primes by Chebotarev.
Put \(c_0=Q\sqrt{-q}\), \(c_j=j c_0\), and let \(O_j\)
be the orbital Gram operator for the ordinary seed at
\(\gamma_{c_j}\), normalized by its diagonal term as in
Theorem~\ref{thm:real282}. Define the positive real number
\[
 B_x=m_f(\gamma_{c_0})\, C_E\,M_D(W)
                  \prod_\tau (2/|\tau c_0|)^D C_{n,h_\tau}.
\]

\begin{theorem}[A summable tail on the prescribed column]\leavevmode
\label{thm:tail283}
For every integer \(Y\ge1\), every sufficiently distant
common interval, and \(A=(2s+2)f\), one has
\[
   \sum_{\substack{j\in\mathcal T\\j>Y}}\|O_j\|_{\rm op}
                  \le 2B_x\,\frac{Y^{1-A}}{A-1}.
\]
The bound applies also after compression to a terminal
character or to the retained coordinates. For the actual
real-frequency sections, with \(m=2(s+1)\), its positive
prefactor satisfies
\[
       \liminf_{x\to\infty}
       \frac{\log B_x}{(2r-2)f m x^\varepsilon}\ \ge\ 8
\]
along these intervals, uniformly in the prescribed words.
\end{theorem}
\begin{proof}
The distinct \(j\) give distinct labels since their first
minors are distinct. At the places above \(q\), multiplying
\(c_0\) by a unit leaves the unsigned finite mass unchanged
by Theorem~\ref{thm:split281}. Above a prime \(p\mid j\)
every local factor has valuation one, residue size \(p\),
and is at full split level. There are \(f\) such places
in \(F\). At the other places use the minimal-valuation
calculation in Theorem~\ref{thm:finite281}.
Thus the extra unsigned mass is
\[
 \prod_{p\mid j}
  \left(p^{\,2k+2s-2}
            [1+O(s p^{2-k})]\right)^{\! f}.
\]
The constant is absolute by the polynomial in
Theorem~\ref{thm:split281}. For \(k\ge5\),
\(\sum_{a\ge2}a^{2-k}\le 2^{2-k}+2^{3-k}/(k-3)\)
by integral comparison. Since \(fs\) is polynomial in \(x\)
and \(k\) grows faster than \(m^5\), the product of all the
bracketed errors, even over all rational primes, is at most
two on sufficiently distant intervals.

Disintegrate the orbital pairing as in
Theorem~\ref{thm:real282}. At an extra prime the local
phase can have modulus one without being trivial. On a
fixed representative both finite factors are unitary
translations of column arguments, with scalar phases.
The integral of their pairing over the common \([H]\)
is at most the product of the column norms, by
Cauchy--Schwarz and Haar invariance. This remains true
if a representative is not spherical before averaging.
Summing the finite weights therefore costs at most the
unsigned mass. The bound for the real integral of norms
in Theorem~\ref{thm:moment283} gains \(j^{-Df}\).
Combining the estimates gives
\(\|O_j\|_{\rm op}\le2B_x j^{-A}\), since
\(D=2r-2=2k+4s\). Summing by integral comparison proves
the first assertion. Orthogonal projections have norm at
most one.

To measure the prefactor use the simultaneous concentration
of Theorem~\ref{thm:localmass156}, writing \(b=\exp(x^\varepsilon)\).
On a set of relative \(\|W\|^2\)-mass at least
\(1-\exp(-c_1br^2)\), with \(c_1>0\) absolute, every
\(R(g_\tau)=a_{1,\tau}/\sqrt2\) is at least
\(\sqrt{T/2}(br^4)^{-s}\). Thus \(\log M_D(W)/(Df)\)
is at least
\(\tfrac12\log(T/2)-s\log(br^4)+o(1)\).
By Theorem~\ref{thm:prime210},
\(\log T=(17+o(1))m x^\varepsilon\) and
\(\log Q=(16+o(1))x^\varepsilon\). Also \(\log q=O(\log x)\)
by Theorem~\ref{thm:conductor232}. The lower bound of
Theorem~\ref{thm:fiber283} gives
\(\log C_{n,h}\ge-O(n\log(2n))\).
The discriminants have logarithms \(O(f\log x)\) by the
polynomial cyclotomic conductor and the conductor-discriminant
formula. Consequently the loss from \(C_E\) and the local
constants is \(O(Df\log x)\). Finally the unsigned mass is
at least one, \(\log|\tau c_0|=\log Q+\tfrac12\log q\),
and \(s/m\to1/2\). Division by \(Df m x^\varepsilon\)
gives the claimed lower limit.
\end{proof}

The lower limit concerns the positive prefactor in the tail
bound. The signed contribution at \(j=1\) is the real
integral of Theorem~\ref{thm:real282} times its finite
operator. Its possible cancellation and that of the other
rational orbits still enter the period estimate in Section~\ref{sec:missing}.

\subsection{The first orbit at the original terminal prime}\label{nav:IX-51}
Keep the original fields, power intervals, level, word
and column of Theorem~\ref{thm:real282}.
Use \(c_0=Q\sqrt{-q}\), \(U_0=Q\sqrt q\), and write \(O_1\)
for the orbital Gram operator at \(\gamma_{c_0}\), normalized
by the diagonal. Let \(M_*={\cal M}_{2K_*+1}(W_1)\), the
full-flag height moment of the unit-frequency real section
from Theorem~\ref{thm:norm284}. Put
\[
 H_x=m_f(\gamma_{c_0})\,C_E A_*^f
        (2/U_0)^{(r-1)f}
           (1+2\log^+(2T/U_0))^{(s-1)f}.
\]

\begin{theorem}[The critical power of the terminal prime]\leavevmode
\label{thm:first284}
On these same intervals, including after projection to the
prescribed terminal character and retained coordinates,
\[
 \|O_1\|_{\rm op}\le H_x M_*,
 \qquad
 \limsup_{x\to\infty}
        \frac{\log H_x}{(r-1)f x^\varepsilon}\le -16 .
\]
Both inequalities are uniform in the words in question.
The remaining moment \(M_*\) uses the actual word at that
point of the interval.
\end{theorem}
\begin{proof}
By Theorem~\ref{thm:real282} the norm is at most the
absolute real factor times the norm of the finite operator.
The latter is bounded by the unsigned mass, since every
translation entering the finite pairing has norm one before
spherical averaging. Orthogonal compression preserves this
bound. Apply Proposition~\ref{prop:scale284} and then
Theorem~\ref{thm:norm284} at
\(\delta_\tau=U_0/(2T)\). The resulting product is
\[
       T^{-(r-1)f}\bigl(2T/U_0\bigr)^{(r-1)f}
                  =(2/U_0)^{(r-1)f}.
\]
This leaves precisely \(H_xM_*\).

For the limiting upper bound,
\(\log m_f=(k+s-1)f\log q+O(fs q^{2-k})\) by
Theorem~\ref{thm:product281}, and
\(\log q=O(\log x)\) by Theorem~\ref{thm:conductor232}.
The discriminant formula of Proposition~\ref{prop:measure283}
gives \(C_E\le1\). Since \(d!\ge1\),
\(\log A_*\le O(r+\log(2K_*))\).
The terminal-prime estimate of
Theorem~\ref{thm:prime210} gives
\(\log(1+2\log^+(2T/U_0))=O(\log x)\).
All these logarithmic contributions are
\(o((r-1)f x^\varepsilon)\). Finally
\(\log Q=(16+o(1))x^\varepsilon\), as follows from the
two level primes between \((k\exp(x^\varepsilon))^8\)
and \(2(k\exp(x^\varepsilon))^8\), with \(k\) polynomial
in \(x\). Substitution proves the upper limit.
\end{proof}

The estimate controls the first term by a positive majorant
whose explicit factor decays on the original intervals.
A uniform upper bound for \(M_*\) at the required scale
would turn it into a numerical bound for this term.
The signed sum over all terms and the comparison with the
integral theta lattice in Section~\ref{sec:missing} are still open.

\subsection{The sum over all split squarefree first-plane multiples}\label{nav:IX-52}
Fix a field and a point of the intervals in Section~\ref{sec:missing}, one
of the prescribed words and a unit vector of the lowest row
type. Put \(c_0=g_1g_2\sqrt{-q}\), \(r=k+2s+1\),
\(R=r-1\) and \(b=\exp(x^\varepsilon)\). Let \(\mathcal T\)
be the positive squarefree integers prime to \(2qd_0g_1g_2T\)
whose prime factors split completely in \(E\), including \(1\).
For \(j\in\mathcal T\), take the ordinary orbital Gram operator
\(O_j\) at \(\gamma_{j c_0}\) with the diagonal normalization
of Theorem~\ref{thm:real282}.

\begin{theorem}[Exponential decay for the entire split squarefree family]\leavevmode
\label{thm:sum287}
On the original fields, intervals and finite level,
\[
       \sum_{j\in\mathcal T}\|O_j\|_{\rm op}
          \le \exp\bigl(-(16-o(1))(r-1)f x^\varepsilon\bigr),
                \qquad f=[E^+:\QQ] .
\]
The estimate is uniform in the word and the unit vector.
It holds also after compression to the original terminal
character or the retained coordinates. In particular the
normalized signed trace of this whole family on any such
nonzero space has absolute value at most the same bound.
\end{theorem}
\begin{proof}
The relative element \(\gamma_{j c_0}\) commutes with the
same column group and the same stabilizer \(J\) as
\(\gamma_{c_0}\), by Proposition~\ref{prop:plane281}.
Thus the unfolding of Theorem~\ref{thm:real282} applies
with each of the finite multipliers in
Proposition~\ref{prop:split287}. More explicitly, for
each finite unipotent coset tuple one integrates over
the common column quotient. Averaging over its right compact
real column group contracts the real pair by Schur orthogonality,
while the finite covectors are paired after their relative
Hecke translation. The factor \(1/\dim\rho\) cancels against
the diagonal. The finite additive phases depend only on the
finite tuple and remain in the finite operator. The modulus
of the common real factor is bounded by \(I_j^+\).
All disintegrations are legitimate for each \(j\) by the
absolute convergence of Theorem~\ref{thm:ordinary163} or,
for the positive real integral, by Lemma~\ref{lem:real287}.

Multiplying Proposition~\ref{prop:split287} with
Lemma~\ref{lem:real287} and cancelling \(m_f(c_0)\) gives
\[
       \|O_j\|_{\rm op}\le H_x^{f\omega(j)}
                    \min\{A_0,A_1j^{-Rf}\}.
\]
Put \(L_x=\log H_x/\log 2\), so \(H_x^{\omega(j)}\le j^{L_x}\)
by \(2^{\omega(j)}\le j\). Since \(L_x=O(\log r)\) and
\(R=r-1\to\infty\), the number \(\theta=(L_x+2)/R\)
lies in \((0,1)\) for large \(x\). Bound the minimum by
the product of its first and second entries raised to
\(1-\theta\) and \(\theta\), respectively. This yields
\[
        \|O_j\|_{\rm op}
           \le A_0^{1-\theta}A_1^\theta j^{-2f}.
\]
Summing over every positive integer on the right costs at
most \(1+\int_1^\infty t^{-2f}\,dt\le2\).
By Lemma~\ref{lem:real287} the logarithm of the product
of constants is bounded by
\[
 -(1-\theta)(16-o(1))Rf\log b + C\theta Rfs\log b.
\]
Here \(\theta=o(1)\) and \(\theta s=o(1)\) because
\(s\log r=o(r)\). Indeed the common construction has
\(m=2s+2\), \(m^5x^\varepsilon=o(k)\), and \(r\) polynomial
in \(x\). This proves the asserted exponent. Compression
by an orthogonal projection decreases each operator norm.
The absolute value of a normalized trace is at most the
operator norm, so absolute convergence and the same bound
also hold for its signed sum.
\end{proof}

This sum includes arbitrarily large rational multipliers.
The projected period in Section~\ref{sec:missing} also contains other
first-plane parameters and other relative flag positions.
The estimates for the full period and for the strict cohomology
on these intervals remain open.

\subsection{All first-plane terms prime to the marked level}\label{nav:IX-53}
For each \(x\) in the power intervals of Section~\ref{sec:missing}
use the skew coefficients \(\xi\in\mathcal O_E^-\) and
the descended ideal \(\mathfrak b(\xi)\) defined before
Theorem~\ref{thm:skew289}. Denote by \(\mathcal S_{\rm full}\)
all nonzero \(\xi\) whose ideal is supported away from \(QT\)
at places of full row and column lattice level. The prime
exponents of this ideal are arbitrary. As before \(O_\xi\)
is the normalized orbital operator at \(\gamma_{Q\xi}\).

\begin{theorem}[The full-level first-plane sum with unrestricted valuations]\leavevmode
\label{thm:good290}
For \(f=[E^+:\QQ]\), \(R=r-1\), on the original fields,
intervals and words one has
\[
   \sum_{\xi\in\mathcal S_{\rm full}}\|O_\xi\|_{\rm op}
       \le \exp\bigl(-(16-o(1))Rf x^\varepsilon\bigr).
\]
The estimate holds on the terminal-character projection
and the retained coordinates at the same level.
\end{theorem}
\begin{proof}
Let \(C_k=4(k+1)^2\), \(h=s-1\),
and let \(\Omega(\mathfrak b)\) count prime factors
with multiplicity. Proposition~\ref{prop:powers290}
at each divisor, and the minimal-valuation result
of Theorem~\ref{thm:finite281} at all the remaining
places, bound the finite operator by
\[
       C_k^{\Omega(\mathfrak b)}
                     (N\mathfrak b)^{R+2s+1}.
\]
These are the local factors of the original rational
orbit by the unfolding in Theorem~\ref{thm:skew289}.
For \(t_\tau=|\tau\xi|\) let
\(H(\xi)=\prod_\tau\max(1,t_\tau)\). Then
\(\prod_\tau t_\tau=N\mathfrak b\), and
\(2^{\Omega(\mathfrak b)}\le N\mathfrak b\le H(\xi)\).

Use the two signed real bounds with constants \(a'_0,a'_1\)
in the proof of Theorem~\ref{thm:skew289}. Their common
power \(\prod_\tau t_\tau^{-R}\) cancels
\((N\mathfrak b)^R\), leaving at most \(H(\xi)^{2s+1}\)
from the finite norm power and \(H(\xi)^{\log C_k/\log2}\)
from the multiplicities. The product of logarithmic factors
is at most \(H(\xi)^{2h}\). For \(t_\tau\ge1\) interpolate
with the higher real bound with exponent
\[
         \theta=\frac{\log C_k/\log2+2s+1+2h+3}{R}.
\]
At \(t_\tau<1\) use the first bound alone. Since
\(a'_1\ge a'_0\), this gives
\[
 \|O_\xi\|_{\rm op}
 \le C_E(a'_0)^{(1-\theta)f}(a'_1)^{\theta f}H(\xi)^{-3}.
\]
The sum of the last height power over all nonzero skew
integers is at most \(8^f\) by the skew-lattice count in
Theorem~\ref{thm:skew289}. Here
\(\theta s=O((s\log k+s^2)/R)=o(1)\). In particular
\(0<\theta<1\) on the sufficiently distant intervals.
Using \(C_E\le1\),
\(\log a'_0\le-(16-o(1))R x^\varepsilon\), and
\(\log a'_1\le C R s x^\varepsilon\) proves the stated
exponent. The norm bound passes to orthogonal compressions.
\end{proof}

The remaining first-plane coefficients can have extra
valuation at the marked level. The full projected period
also contains other relative flag positions. The integral
lower comparison and the uniform strict cohomology bound
on these fields and intervals still require proofs.

\subsection{A finite window for the marked first-plane contribution}\label{nav:IX-54}
At each point of the common intervals consider every nonzero
skew integer \(\xi\) and its orbit \(\gamma_{Q\xi}\).
Let \(\mathfrak m\) be the part of the descended ideal of
\(\xi\) above \(g_1g_2T\). Let \(H=H(\xi)\) be the product
of the larger of \(1\) and \(|\tau\xi|\) over the real
places of \(F\), using one extension of each place to \(E\).

\begin{theorem}[The entire first-plane sum and its marked range]\leavevmode
\label{thm:window291}
Uniformly in the prescribed words on the original intervals,
\[
 \sum_{0\ne\xi\in\mathcal O_E^-}\|O_\xi\|_{\rm op}
                      \le\exp(C R f s\log b).
\]
For every fixed \(0<\delta<1\),
\[
 \sum_{N\mathfrak m\le (Q k^{1-\delta})^f}\|O_\xi\|_{\rm op}
                 \le\exp\bigl(-(\delta-o(1))k f\log k\bigr).
\]
For every \(A>0\) there is a constant \(M_A\) for which
\[
  \sum_{\log H\ge M_A f s\log b}\|O_\xi\|_{\rm op}
                         \le \exp(-A R f s\log b).
\]
All three conclusions hold after restriction to the original
terminal-character projection and retained coordinates.
\end{theorem}
\begin{proof}
For all \(y\ge1\) the height summation of
Theorem~\ref{thm:skew289} and Lemma~\ref{lem:factorial291}
give
\[
 \sum_{N\mathfrak m\le y}\|O_\xi\|_{\rm op}
       \le y^{R-1}Q^{-Rf}\{(k+1)!\}^{-f}\exp(C R f).
\]
The \(8^f\) from summing \(H^{-3}\) is included in the
exponential. Insert \(y=(Q k^{1-\delta})^f\).
The factors involving \(Q\) give \(Q^{-f}\).
The integral of \(\log u\) from \(1\) to \(k+1\) gives
\(\log((k+1)!)\ge k\log k-O(k)\). Since \(R=k+2s\)
and \(s=o(k)\), the second assertion follows.

For an unrestricted parameter use the higher bound of
Lemma~\ref{lem:power291} at every place with \(t_\tau\ge1\),
and the first real bound of Lemma~\ref{lem:factorial291}
at \(t_\tau<1\), where \(t_\tau=|\tau\xi|\).
The per-place constants are at most \(\exp(C R s\log b)\).
After removing the common product of \(t_\tau^{-R}\)
the large places give \(H^{-2R}\) and the small-place
logarithms cost at most \(H^{2h}\), \(h=s-1\).
For the remaining finite factors note that
\(2s+1+\log C_k/\log 2\le R-1\) for large \(x\).
Consequently Proposition~\ref{prop:marked291} gives
\[
   C_k^{\Omega(\mathfrak g)}
    (N\mathfrak g)^{2s+1}(N\mathfrak m)^{R-1}
       \le (N\mathfrak g\,N\mathfrak m)^{R-1}\le H^{R-1}.
\]
The field factor satisfies \(C_E\le1\). The unfolding gives
\[
     \|O_\xi\|_{\rm op}\le\exp(C R f s\log b)H^{-R-1+2h}.
\]
Since \(R=k+2s\), the last power is exactly \(H^{-k-3}\).
Summing \(H^{-3}\) costs at most \(8^f\), proving the
unrestricted bound. Above a height cutoff \(H_0\), the
remaining factor is at most \(H_0^{-k}\). Taking
\(\log H_0=M_A f s\log b\), using \(R/k\to1\), and
choosing \(M_A>2(C+A+1)\) proves the third assertion.
An orthogonal compression only decreases the norms.
\end{proof}

For example, \(\xi=Q\sqrt{-q}\) has
\(N\mathfrak m=Q^f\). The second assertion includes all
its multiples by \(a\in\mathcal O_F\) which are units above
\(QT\), with no size restriction at good primes.
The part of the first-plane sum outside the two decay
estimates is finite for each \(x\), since \(H<H_0\) bounds
every embedding and the skew integers form a discrete
lattice there. For instance \(\xi=Q^2\sqrt{-q}\) belongs to
this range. Its marked norm is \(Q^{2f}\) and its logarithmic height
is \(f(2\log Q+\tfrac12\log q)\), which is below
\(M_A f s\log b\) for large \(x\). Theorem~\ref{thm:completeplane293} bounds the entire skew family at these marked valuations.

\subsection{Arbitrary powers at both unitriangular primes}\label{nav:IX-55}
For a nonzero skew integer \(\xi\in\mathcal O_E^-\) use
the descended ideal \(\mathfrak b\) and the normalized
operator \(O_\xi\) at \(c=Q\xi\) of
Theorem~\ref{thm:skew289}. Let \(\mathfrak t\) be the
part of \(\mathfrak b\) supported over \(T\). Define
\[
 D_Q(\xi)=\prod_{\substack{v\mid Q\\v(\xi)>0}}Nv .
\]
In particular \(D_Q(\xi)\le Q^f\), where \(f=[F:\mathbb Q]\).

\begin{theorem}[Decay with unrestricted valuations above \(g_1g_2\)]\leavevmode
\label{thm:twolevels292}
For \(y\ge1\), uniformly on the fields, words and intervals
of Section~\ref{sec:missing},
\[
 \sum_{N\mathfrak t\le y}\|O_\xi\|_{\rm op}
    \le y^{R-1} Q^{-(R-k)f}\{(k+1)!\}^{-f}\exp(CRf),
               \qquad R=k+2s .
\]
In particular for fixed \(0<\delta<1\), the entire sum
with \(N\mathfrak t\le k^{(1-\delta)f}\) is at most
\(\exp(-(\delta-o(1))k f\log k)\) as \(x\to\infty\).
There is no restriction on the valuations above either
unitriangular prime. Both inequalities hold on the
original terminal projection and on the retained columns.
\end{theorem}
\begin{proof}
Write \(\mathfrak a=\mathfrak b/\mathfrak t\) and
\(C'_k=4(k+1)^4\). Multiplication of
Propositions~\ref{prop:powers290}, \ref{prop:marked291}
and \ref{prop:markedpowers292} bounds the finite factor by
\[
 D_Q(\xi)^k (C'_k)^{\Omega(\mathfrak a)}
       (N\mathfrak a)^{R+2s+1}(N\mathfrak t)^{2R-1}.
\]
Each active place above \(Q\) contributes its extra
\((Nv)^k\) exactly once, independently of valuation.
Remove \(\prod_\tau|\tau\xi|^{-R}=(N\mathfrak b)^{-R}\)
from the real estimates. The remaining good and
unitriangular powers, apart from \(D_Q(\xi)^k\), cost at
most \(H^{2s+1+\log C'_k/\log 2}\), with
\(H=\prod_\tau\max(1,|\tau\xi|)\). The terminal factor
left over is \((N\mathfrak t)^{R-1}\).

At a place where \(|\tau\xi|\ge1\), interpolate the two
real estimates in Lemma~\ref{lem:factorial291} with weight
\(\theta=(2s+1+\log C'_k/\log 2+2h+3)/(2R)\) on the
higher exponent, where \(h=s-1\). At the other places
use the lower exponent. Their logarithmic terms cost at
most \(H^{2h}\) by \(\prod_\tau|\tau\xi|\ge1\).
This leaves \(H^{-3}\) in the bound. Since
\(\log C'_k=O(\log k)\), the estimates for
\(\theta\log((k+1)!)\) and \(\theta Rs\log b\) in that
lemma still give \(o(R)\). With the same diagonal and
field normalizations we obtain
\[
 \|O_\xi\|_{\rm op}\le D_Q(\xi)^k(N\mathfrak t)^{R-1}
 H^{-3}Q^{-Rf}\{(k+1)!\}^{-f}\exp(CRf).
\]
Use \(D_Q(\xi)\le Q^f\), since the two rational primes
split completely. The sum of \(H^{-3}\) is at most \(8^f\)
by Theorem~\ref{thm:skew289}. This proves the first bound.
Taking \(y=k^{(1-\delta)f}\), use
\(\log((k+1)!)\ge k\log k-O(k)\), \(s=o(k)\) and \(R\ge k\)
to obtain the stated decay. Orthogonal projections do
not increase the operator norm.
\end{proof}

For example every \(Q^j\sqrt{-q}\), \(j\ge0\), has
\(N\mathfrak t=1\), as does its product with any
nonzero \(F\)-integer prime to \(T\). All these terms
are included in one convergent sum. A large terminal
part can still occur below the height cutoff of
Theorem~\ref{thm:window291}. The next theorem bounds its contribution at every terminal valuation.

\subsection{Exponential decay for the whole first-plane family}\label{nav:IX-56}
Let \(F=E^+\), \(f=[F:\mathbb Q]\), \(R=k+2s\) and
\(Q=g_1g_2\). Use the original finite row level, terminal
character and real frequencies of Section~\ref{sec:missing}. For every
nonzero skew integer \(\xi\) let \(O_\xi\) be the orbital
Gram operator of \(\gamma_{Q\xi}\) divided by the
diagonal normalization of Theorem~\ref{thm:skew289}.

\begin{theorem}[The full first-plane contribution on the common intervals]\leavevmode
\label{thm:completeplane293}
Uniformly in the prescribed words and interval points,
\[
  \sum_{0\ne\xi\in\mathcal O_E^-}\|O_\xi\|_{\rm op}
    \le Q^{-(R-k)f}\{(k+1)!\}^{-f}\exp(CRf).
\]
In particular the entire sum is at most
\(\exp(-(1-o(1))kf\log k)\). The same estimates hold for
the sum after compression to the specified terminal
character and retained coordinates. On any nonzero such
space they also bound the absolute normalized trace of
the sum.
\end{theorem}
\begin{proof}
Let \(\mathfrak b\) be the descended ideal of \(\xi\),
\(C'_k=4(k+1)^4\) and
\(D_Q=\prod_{v\mid Q,\ v(\xi)>0}Nv\). Away from \(QT\)
use Proposition~\ref{prop:powers290}. Above \(Q\) use
Proposition~\ref{prop:markedpowers292}, and above \(T\)
use Proposition~\ref{prop:terminal293}. Since
\(k+4s\le R+2s+1\), their product in the common-column
unfolding is bounded by
\[
        D_Q^k(C'_k)^{\Omega(\mathfrak b)}
                      (N\mathfrak b)^{R+2s+1}.
\]
This estimate includes every valuation at every finite place.
The two real estimates at the physical frequency, the field constant,
and the height sum in the proof of Theorem~\ref{thm:twolevels292}
apply to the same parameters. Explicitly, the factor
\(\prod_\tau|\tau\xi|^{-R}\) cancels \((N\mathfrak b)^R\).
With \(H=\prod_\tau\max(1,|\tau\xi|)\), the remaining
finite powers other than \(D_Q^k\) cost at most
\(H^{2s+1+\log C'_k/\log 2}\). Take the interpolation
weight \((2s+1+\log C'_k/\log 2+2(s-1)+3)/(2R)\)
at the places \(|\tau\xi|\ge1\). The other places cost
at most \(H^{2(s-1)}\) in logarithmic factors. The
result of Lemma~\ref{lem:factorial291} with these powers is
\[
  \|O_\xi\|_{\rm op}\le
       D_Q^k H^{-3}Q^{-Rf}\{(k+1)!\}^{-f}\exp(CRf).
\]
The two unitriangular primes split completely, so
\(D_Q\le Q^f\). Theorem~\ref{thm:skew289} bounds the
sum of the height weights by \(8^f\). Absorb this in the
exponential and use
\(\log((k+1)!)\ge k\log k-O(k)\) and \(R/k\to1\).
The stated estimate also implies absolute convergence
in operator norm. Orthogonal compression is contractive,
and normalized trace on a finite-dimensional space has
norm one. This proves the last assertions.
\end{proof}

This exhausts the parameters of Theorem~\ref{thm:skew289}, including simultaneous divisibility by both unitriangular primes and the terminal prime. The next theorem allows arbitrary relative scalars in the same intersection dimension.

\subsection{All pairs whose isotropic spaces meet in codimension one}\label{nav:IX-57}
Let \(\mathcal S_1\) denote the relevant, finitely
supported rational labels in the ordinary one-word
expansion for which \(\dim(X_0\cap X_1)=s-1\), counted
once each as in Theorem~\ref{thm:sum237}. Write \(O_\gamma\)
for an individual orbital Gram operator in the identity
normalization on the original column.

\begin{theorem}[The entire codimension-one contribution]\leavevmode
\label{thm:codimension294}
On the common fields and intervals of Section~\ref{sec:missing},
uniformly in the prescribed word,
\[
       \sum_{\gamma\in\mathcal S_1}\|O_\gamma\|_{\rm op}
       \le Q^{-(R-k)f}\{(k+1)!\}^{-f}\exp(CRf),
             \qquad R=k+2s .
\]
The sum includes every allowable relative scalar. In
particular it is at most \(\exp(-(1-o(1))kf\log k)\).
The same bounds apply after the terminal-character or
retained-column projections and to the absolute normalized
trace on each nonzero projected space.
\end{theorem}
\begin{proof}
For \(\gamma=w_{Q\xi,\lambda}\) use
\(H(\xi)=\prod_\tau\max(1,|\tau\xi|)\), using one
embedding of \(E\) above each embedding of \(F\).
The product of the real absolute values is \(A_\xi\)
in Proposition~\ref{prop:finiteall294}. In particular
\(1\le A_\xi\le H(\xi)\). Factoring
\(\prod_\tau|\tau\xi|^{-R}=A_\xi^{-R}\) out of
Proposition~\ref{prop:representative294} cancels the first
\(R\) powers of \(A_\xi\) in the finite estimate. Besides
\(D_Q(\xi)^k\), the remaining finite cost is at most
\(H^{4s+1+2\log C_k/\log 2}\).

At a place with \(|\tau\xi|\ge1\) interpolate the two
real estimates with weight \(\theta\) on the second,
where
\[
    2R\theta=4s+1+2\log C_k/\log2+2(s-1)+7 .
\]
At smaller places use the first estimate. The product
of their logarithmic terms costs at most \(H^{2(s-1)}\),
because the sum of the negative logarithms of
\(|\tau\xi|\) is bounded by \(\log H\). Thus \(H^{-7}\)
remains. As in Lemma~\ref{lem:factorial291},
\(\theta=O((s+\log k)/R)\),
\(\theta\log((k+1)!)=o(R)\) and \(\theta Rs\log b=o(R)\)
on the prescribed intervals. The field factor is at most
one. We have obtained the individual bound
\[
  \|O_\gamma\|_{\rm op}\le D_Q(\xi)^k H(\xi)^{-7}
          Q^{-Rf}\{(k+1)!\}^{-f}\exp(CRf).
\]

For a fixed \(\xi\), put \(\eta=\xi\lambda\).
It is integral and \(\eta\bar\eta=\xi\bar\xi\). Hence
\((\eta)\) divides \((\xi\bar\xi)\). The number of
divisors of an integral ideal is at most its absolute
norm, because \(e+1\le2^e\) for each prime exponent \(e\).
Here the norm is \(A_\xi^4\). Two possible generators
of the same ideal differ by a unit of absolute value one
at all embeddings. Such a unit is a root of unity by
Kronecker's theorem. There are at most \(8f^2\) roots
of unity in \(E\), since a primitive root of order \(N\)
has degree \(\varphi(N)\ge\sqrt{N/2}\) and
\([E:\mathbb Q]=2f\). The inequality follows by
multiplying \(\varphi(p^e)\ge p^{e/2}\) for odd prime
powers and \(\varphi(2^e)\ge2^{(e-1)/2}\) for \(e\ge1\).
Thus there are at most
\(8f^2 H(\xi)^4\) possible \(\lambda\)'s. For each pair
\((\xi,\lambda)\) there is at most one label by
Theorem~\ref{thm:suffix237}.

For completeness, the sum of \(H^{-3}\) over all
\(\xi\in\mathcal O_E\setminus\{0\}\) is at most \(27^f\).
For coordinate radii \(r_\tau\ge1\), center a product
of open complex disks of radius \(1/2\) at each integer
in the box \(|\tau\xi|\le r_\tau\). Distinct products
are disjoint, since otherwise the difference of the
integers would have absolute norm strictly less than one.
Comparison with the product of disks of radii
\(r_\tau+1/2\) bounds the count by
\(\prod_\tau(1+2r_\tau)^2\le9^f\prod_\tau r_\tau^2\).
Integrate this count against
\(\prod_\tau 3r_\tau^{-4}\,dr_\tau\) on \([1,\infty)^f\).
The contribution of each point is \(H^{-3}\), and the
majorant integrates to \(27^f\).

Now \(D_Q(\xi)\le Q^f\). Sum the individual bound using
the two counts and absorb \(8f^2 27^f\) into \(\exp(CRf)\).
The factorial estimate and the assertions concerning
projections and traces follow as in
Theorem~\ref{thm:completeplane293}.
\end{proof}

In this ordinary one-word expansion, the case \(X_0=X_1\)
has only the identity label by Proposition~\ref{prop:finite237}.
Theorem~\ref{thm:allflags298} below treats every other proper nonzero intersection dimension and the transverse pairs with full common column.

\subsection{The full sum for all common isotropic flags}\label{nav:IX-58}
For \(1\le j<s\) let \(\mathcal F_j\) be the relevant,
supported single-word rational labels with
\(\dim(X_0\cap X_1)=s-j\). Let \(\mathcal F_s\) denote
the transverse labels with full common column. Each rational
label is counted once in its ordinary diagonal normalization.

\begin{theorem}[All proper intersections and the transverse full column]\leavevmode
\label{thm:allflags298}
On the fields and intervals of Section~\ref{sec:missing}, for every
prescribed word and every \(1\le j\le s\),
\[
            \sum_{w\in\mathcal F_j}\|O_w\|_{\rm op}
                 \le ((k+1)!)^{-jf}\exp(Crjf).
\]
The sum over all \(j\) is at most
\(((k+1)!)^{-f}\exp(Crf)\). There is no restriction on
the prefix composition, coefficients or individual
heights. These estimates hold at the original levels and
after the prescribed orthogonal projections.
\end{theorem}
\begin{proof}
For \(j=1\) use Theorem~\ref{thm:codimension294}.
For \(j\ge2\) take a composition as in
Proposition~\ref{prop:composition298} and let
\(H_i=\prod_\tau\max(1,|\tau a_i|)\). Set \(\eta_i=3\)
for \(i<t\) and \(\eta_t=\delta_k+7\), with \(\delta_k\)
as in Proposition~\ref{prop:compositionfinite298}.
The power constraint in Theorem~\ref{thm:compositionreal298}
holds because \(\sum_i J_i\le j^2\) and \(j\le s\).
Multiply its bounds over the real places. The factors
from the intermediate products cancel their entire finite
cost in Proposition~\ref{prop:compositionfinite298}, since
\[
        \prod_\tau|\tau G_p|^{-2u_p}
                          = |N_{E/\mathbb Q}G_p|^{-u_p}.
\]
The remaining rational measure factor is at most
\(2^{f\sum_p u_p}\) by Proposition~\ref{prop:composition298}
and the whole-plane measure calculation of
Lemma~\ref{lem:blockphysical295}. The common compact
column contributes the same Schur dimension factor as
on the diagonal, since \(w\) centralizes its group and
the characters agree on \(J\) for each relevant pair.
The lowest-type product and polarization argument of
Theorem~\ref{thm:blocktail295} therefore bounds the
ordinary operator by these products.

Now put \(A=|N_{E/\mathbb Q}G_j|^{1/2}
=Q^{jf}|N_{E/\mathbb Q}a_t|^{1/2}\). We have
\(1\le A\le Q^{jf}H_t\), and \(E_j=d-j+4hj+3j^2\ge k\).
Consequently the remaining norm power
\(A^{k+\delta_k-E_j}\) costs at most \(Q^{jf\delta_k}H_t^{\delta_k}\).
The contribution of \(Q^{jf\delta_k}\) and the measure cost
is \(\exp(O(rjf))\), since \(\log Q=O(\log b)\) and
\((s+\log(k+1))\log b=o(r)\). Thus each label is bounded by
\[
      ((k+1)!)^{-jf}\exp(Crjf)\,
               H_t^{-7}\prod_{i<t}H_i^{-3}.
\]

For given nonzero integers \(a_i\), the successive block
coefficients solve
\((c_i/Q)^{j_i}=a_i/a_{i-1}\), with \(a_0=1\) and with
the signs fixed by our definition of \(G_p\). Each equation
has at most \(j_i\) roots in \(E\), giving at most \(2^j\)
choices altogether. The integer \(Ta_t\lambda\) has
prescribed relative norm \(T^2 a_t\bar a_t\), so the ideal
divisor and root-of-unity count in
Theorem~\ref{thm:blocktail295} gives at most
\(8f^2 T^{4f}H_t^4\) possible scalars. The parameters
determine a rational label by Proposition~\ref{prop:composition298}.
It remains to sum \(\prod_i H_i^{-3}\) over the tuples.
The disk count of Theorem~\ref{thm:codimension294} gives
at most \(27^{ft}\). There are \(2^{j-1}\) compositions,
corresponding to subsets of the \(j-1\) possible cuts.
Since \(t\le j\) and \(\log T=O(s\log b)=o(r)\), all
these counting factors are absorbed by \(\exp(Crjf)\).
This proves the first assertion. Summing the geometric
majorant, as in Theorem~\ref{thm:allblocks297}, proves
the second, and orthogonal projections are contractions.
\end{proof}

When \(X_0=X_1\), Proposition~\ref{prop:finite237} leaves
only the identity at the given level. Thus the remaining
one-word contribution outside the cases above comes from
transverse pairs with smaller common column. For pairs of
distinct words the same positions have zero real operator
by Theorem~\ref{thm:common133}. Indeed the representative
in Proposition~\ref{prop:composition298} centralizes the
whole definite column of rank \(k\), and the restrictions
of the two real types to that column are inequivalent.
The mixed real integrals converge by the same weighted
Cauchy--Schwarz estimate as above, applying the diagonal
moment bound to each word separately.
The complete
projected period and the strict cohomology estimate of
Section~\ref{sec:missing} remain open.

\subsection{An integral plane for the first relative orbit}\label{nav:V-45}
Take the field \(E=E_x\) of Theorem~\ref{thm:conductor232},
put \(F=E^+\), and keep the row \(G\), the group \(B=N\rtimes H\)
and the original level \(K^0\) of Theorem~\ref{thm:midlevel131}.
Write \(Q=g_1g_2\) for the product of the two split level primes.
Assume \(s=m/2-1\ge2\). Put \(u=e_++e_-\) in the first signed
coordinate plane, with \(h(e_\pm,e_\pm)=\pm1\).
Let \(M\) be the perpendicular of that plane. The subsequent
flag vectors, the terminal vector \(\varpi e\), and the entire
definite column \(W_0\) belong to \(M\). Write \(J\) for their
Bessel group in \(U(M)\), using the flag with \(u\) deleted.
Embed \(J\) by the identity on the first plane.

\begin{proposition}[A first-plane orbit integral also at two]\leavevmode
\label{prop:plane281}
There is an integral isotropic partner \(v\) of \(u\) over \(E\)
such that \(u,v\) are a basis of the signed plane lattice.
For every \(0\ne c\in Q\mathcal O_E\) with \(\bar c=-c\), put
\[
               \gamma_c y=y+c\,v h(v,y).
\]
This is a rational isometry in \(K^0\) at every finite place.
It represents a relevant orbit with \(t=1\), relative scalar
one, and \(B\cap\gamma_c B\gamma_c^{-1}=J\).
Its sole nonzero leading minor is \(D_1=c\). On every oriented
column one finite representative gives the same column on
the two sides.
\end{proposition}
\begin{proof}
Let \(\alpha=(1+\sqrt{-d_0})/2\). It satisfies the integral
monic equation \(X^2-X+(1+d_0)/4=0\), and
\(\alpha+\bar\alpha=1\). Take \(v=e_+-\alpha u\).
Then \(h(u,v)=1\) and \(h(v,v)=1-\alpha-\bar\alpha=0\).
The determinant of \(u,v\) in the signed basis is \(-1\).
The identity \(\bar c=-c\) cancels the two linear terms in
the pairing of two images under \(\gamma_c\). Its quadratic
term vanishes by isotropy of \(v\). Its inverse is
\(\gamma_{-c}\). The difference from the identity and the
corresponding difference for the inverse have entries in
\(Q\mathcal O_E\). This gives the prescribed reduction at
\(g_1,g_2\). The graph defining the level above \(T\) changes
only the middle lattice in \(M\). Elsewhere the integral
basis computation applies to the full stabilizer.

The two first vectors are \(u\) and \(u+cv\).
An element of the common stabilizer fixes both, hence fixes
their plane pointwise. On \(M\) the two remaining decorations
are identical and their stabilizer is precisely \(J\).
Their simple characters coincide there. The coordinate of
\(u\) on the image of the second flag vector is zero there
for both flags. This proves relevance. The common isotropic
space is spanned by \(u_2,\ldots,u_s\), on which the two flags
have the same decoration. The cross matrix pairs the first
two first vectors to \(c\), all its other isotropic
coordinates to zero, and the two terminal vectors to \(-T\).
Since \(s\ge2\), this proves the minor assertion.
Finally \(\gamma_c\) commutes with \(H\) and lies in the finite
level. At a representative \(h_f\in H(\mathbb A_{F,f})\)
of the column, \(h_f\) and \(\gamma_c^{-1}h_f=h_f\gamma_c^{-1}\)
therefore give the same finite right coset.
\end{proof}

\subsection{A representative for every pair with one moving line}\label{nav:V-51}
Use the ordinary orbital expansion for one prescribed
word on the fields of Section~\ref{sec:missing}. Write \(X_0,X_1\)
for the two \(s\)-dimensional isotropic spaces. Suppose
\(\dim(X_0\cap X_1)=s-1>0\). Let \(u,v_1\) be the first
integral hyperbolic pair in Proposition~\ref{prop:plane281},
and \(M\) its perpendicular. For \(c\in E^\times\) and
\(\lambda\bar\lambda=1\) define a rational isometry by
\[
    w_{c,\lambda}u=c v_1,\qquad
    w_{c,\lambda}v_1=u/\bar c,\qquad
    w_{c,\lambda}|_M=\lambda I_M .
\]
The first two images form a hyperbolic pair and the last
map preserves the form on \(M\).

\begin{proposition}[Rational parameters and their archimedean estimates]\leavevmode
\label{prop:representative294}
Every relevant supported label with this intersection has
a representative \(w_{c,\lambda}\) with \(c=Q\xi\),
\(0\ne\xi\in\mathcal O_E\) and \(\xi\lambda\in\mathcal O_E\).
Its common stabilizer is the group \(J\) in
Proposition~\ref{prop:plane281}. For \(t=|\tau\xi|>0\)
the two real majorants used in Lemma~\ref{lem:factorial291}
apply to its local normalized coefficient. With \(h=s-1\),
these are
\[
 \begin{split}
 Q^{-R}\frac{\exp(CR)}{(k+1)!}\,t^{-R}
                      (1+2\log^+(1/t))^h,\qquad
 Q^{-R}\exp(C'Rs\log b)\,t^{-3R}.
 \end{split}
\]
The field factor for their product is the same \(C_E\le1\).
\end{proposition}
\begin{proof}
Theorem~\ref{thm:suffix237} gives a length-one prefix
with scalar \(c=h(u,u_1)\) and the scalar \(\lambda\) on
the common suffix and terminal vector. These data determine
the rational double coset. The displayed isometry has
those data. Proposition~\ref{prop:finite237} gives the
two assertions of integrality. The stabilizer of both
decorations fixes \(u\) and \(c v_1\), so it fixes their
nondegenerate plane pointwise. Its restriction to \(M\)
must preserve the remaining decoration, whose stabilizer
is \(J\). Conversely the scalar \(\lambda I_M\) commutes
with \(J\), proving equality and commutation with the
relative isometry. The common-column unfolding of
Theorem~\ref{thm:sum287} therefore has the same rational
quotient and the same Haar constants.

At a real place use the Euclidean unit Witt pair \(u',v'\)
of Theorem~\ref{thm:fiber283}. Passing from the original
opposite vector to \(v'/\sqrt2\) conjugates by a central
first unipotent. Such elements commute with \(J\) and
their Bessel character is trivial since \(s\ge2\).
In this unit pair the exchange on the plane has lower
left entry \(C=\tau c/2\) and upper right entry
\(1/\bar C\). Write \(\delta=|C|\). Multiplication by
the unit scalar \(C/(i\delta)\) on the plane changes the
exchange with lower entry \(i\delta\) into the exchange
with entry \(C\). On \(M\) the additional scalar
\(\tau\lambda\) also has absolute value one. Their block
diagonal product is Euclidean unitary, normalizes \(B\),
commutes with \(J\), and preserves each of the functions
\(l,\nu,\mathcal A\) on the left. For \(l\) and \(\mathcal A\)
this follows from the unit scalar on each vector of their
defining exterior products. For \(\nu\) it follows by
taking the infimum over \(B\) after using the normalizer
property. This block diagonal product commutes with the
plane exchange.

Write \(n(z)\) for a central first unipotent with purely
imaginary upper plane entry \(z\). Multiplying the exchange
with lower entry \(i\delta\) on the left and right by
\(n(1/(i\delta))\) gives the transvection with lower
entry \(i\delta\) and diagonal entries one. In a first
unipotent fiber integral, such changes are absorbed by
left invariance of the heights and translation of the
unipotent variable. Consequently the fiber bounds for
\(l^{-3R}\) in Lemma~\ref{lem:power291} and for the
\((\mathcal A,\nu)\) weight in Theorem~\ref{thm:low285}
hold with this \(\delta\). The reciprocal Cauchy--Schwarz
integrals use the inverse exchange, which has the same
absolute parameter and the inverse unit scalar on \(M\).
Only the heights in these fiber bounds are transported.
The original squared section norms in the base integrals
are bounded by the moments of Theorem~\ref{thm:moment286}.

The change from the terminal frequency \(\varpi_\tau\) to
unit frequency conjugates the exchange to one with lower
entry \(C/T\) and the reciprocal conjugate upper entry.
Its modulus on the quotient is \(T^{-R}\), by the
computation in Proposition~\ref{prop:scale284} for
\(J\subset B\). The resulting estimates, including the
factorial in the low estimate, are exactly the two
bounds displayed. Schur orthogonality and product Haar
measure give the field factor as in
Proposition~\ref{prop:places288}.
\end{proof}

\subsection{A uniform sum over the transverse rational classes}\label{nav:VII-38}
At \(g\) let \(\mathcal D_g\) be the set of displacements meeting
the finite test in Theorem~\ref{thm:short246}. It includes zero.
Group these vectors by their images in \(X/W_g\). Set \(n=\dim_E X\),
\(f=[E^+:\QQ]\), and
\[
 b_g=\min\bigl(1,\inf_{\ell}H_g(\ell)\bigr),\qquad
 D_g(w)=\prod_\tau |z+w|_{\tau,g}^{-(2n+5)}.
\]
The infimum can be taken over the isotropic rational lines of
\(Y\), with value \(+\infty\) for an empty set.
For every fixed \(g\), one has \(b_g>0\). To see this, rescale a
generator of each line to put its content in a fixed finite list
of ideals representing the ideal classes. Multiplication by a unit
keeps this content. The logarithms of the absolute values of the
units form a lattice in the hyperplane of coordinate sum zero,
by Dirichlet's unit theorem. Hence when \(H_g(\ell)\le1\) we can
rescale the generator further by a unit to give uniformly bounded
norm in each real block, with bounds depending on \(E,g\).
The rational vectors of each of the indicated contents lie in
the corresponding fractional ideal times \(\Lambda_g\).
Their real images form a subset of a lattice, so only finitely
many lines occur. All their heights are positive.

\begin{theorem}[Absolute mass for distinct classes modulo the short lines]
\leavevmode\label{thm:packing246}
Write \(C_n=3^{4n+5}2^{2n+5}\), so that \(\log C_n<6n+9\).
On the actual fields and at every point \(g\),
\[
 \begin{split}
 \sum_{\substack{C\in\mathcal D_g\bmod W_g\\C\ne0}}
          \ \sup_{w\in C\cap\mathcal D_g} D_g(w)
   &\le C_n^f\,(\sqrt{N\mathfrak b})^{-(2n+4)},\\
 \sum_{0\ne w\in\mathcal D_g}D_g(w)
   &\le C_n^f\,(b_g\sqrt{N\mathfrak b})^{-(2n+4)} .
 \end{split}
\]
In the first sum \(C\) denotes a coset in \(X\) represented in
\(\mathcal D_g\). The first bound is independent of the lengths
of the short isotropic vectors at \(g\).
\end{theorem}
\begin{proof}
We first give the packing estimate. Let \(\mathcal V\) be a set
in \(\prod_{\tau=1}^{f}\CC^n\) containing zero, with
\(\prod_\tau\|v_\tau-v'_\tau\|\ge R>0\) for distinct vectors.
For each nonzero \(v\) assign the integers \(k_\tau\ge0\) given by
\(2^{k_\tau}\le1+\|v_\tau\|<2^{k_\tau+1}\), and put
\(K=\sum_\tau k_\tau,\ B_\tau=2^{k_\tau+1}\). In an occupied
box \(R\le\prod_\tau B_\tau=2^{K+f}\) by separation from zero.
Set \(p=(R/\prod_\tau B_\tau)^{1/f}\le1\).
Products of open balls of radii \(pB_\tau/2\) around the points
in this box are disjoint, since an intersection would violate
the product separation. They are contained in the product of the
balls of radii \((1+p/2)B_\tau\) about zero. Comparing volumes in
real dimension \(2n\) per block bounds the number of points by
\[
       (1+2/p)^{2nf}\le 3^{2nf}\,2^{2n(K+f)}R^{-2n}.
\]
Each point has weight
\(\prod_\tau(3/(1+\|v_\tau\|))^{2n+5}
 \le 3^{(2n+5)f}2^{-(2n+5)K}\).
Multiplying and using \(R^4\le2^{4(K+f)}\) bounds the weight in
the box by
\(3^{(4n+5)f}2^{(2n+4)f}R^{-(2n+4)}2^{-K}\).
Summing \(2^{-K}\) over the \(f\) nonnegative indices gives
\(2^f\). This proves the stated constant for any finite set of
points and then for the full sum of nonnegative terms. Finally
\(3^4 2^2<e^6\) and \(3^5 2^5<e^9\).

For an admissible displacement \(w\) the vector \(z+w\) has
Hermitian length \(\eta=\pm1\). In each adapted Euclidean norm
\(|z+w|_{\tau,g}\ge1\) and \(|z|_{\tau,g}=1\). The triangle
inequality gives \(1+|w|_{\tau,g}\le3|z+w|_{\tau,g}\), so \(D_g(w)\)
is at most the weight just summed. Every difference of two
displacements lies in \(\mathfrak b\Lambda_g\). If the difference
is outside \(W_g\), its line has height at least one, and the product
of its real norms is at least \(\sqrt{N\mathfrak b}\). Choose one
point from each of finitely many nonzero cosets and adjoin zero.
The packing estimate applies independently of these choices.
Taking suprema over the choices and then over finite sets of cosets
proves the first inequality. For the complete set every nonzero
difference has height at least \(b_g\), by Theorem~\ref{thm:short246}
and the definition of \(b_g\). Applying the same packing estimate
with this smaller separation proves the second inequality.
\end{proof}

\subsection{The integral orbit and its measure}\label{nav:VII-42}
Take the notation of Proposition~\ref{prop:shell263}. Write
\(G=U(S)\), \(H=U(C)\), \(u=\dim_E S\), and let \(K_f\)
be the given level on \(G\). At an unmarked place it is the
stabilizer of \(\Lambda_S\), and at a marked place it is its
principal congruence subgroup of depth \(a_v\).

\begin{lemma}[The finite orbit cut out by the position test]\leavevmode
\label{lem:orbvol263}
For \(y\in S(\mathbb A_{F,f})\) of Hermitian length \(-1\),
\[
      y\in z+\mathfrak j\widehat L
             \quad\Longleftrightarrow\quad y\in K_f z .
\]
Consequently the support on \(H(\mathbb A_{F,f})\backslash
G(\mathbb A_{F,f})\), under \(g\mapsto g^{-1}z\), is the
image of \(K_f\). Give both lattice stabilizers at each place
volume one, before imposing any congruence, and take the quotient
measure. The measure of this support is
\[
       c_f=\prod_{v\in S_0}
             \frac{(Nv)^{-a_v(2u-1)}}{1-(Nv)^{-u}}
             \ \ge\ (N\mathfrak m)^{-(2u-1)} .
\]
\end{lemma}
\begin{proof}
The lattices \(L=\mathfrak l^{-1}\Lambda_S\) and \(\Lambda_S\)
have the same stabilizer, since locally their ratio is a scalar
on the entire space. The same holds for congruences, which
are conditions on the endomorphism \(g-1\). A compact translate
of \(z\) lies in the required coset because \(z\in L\)
and the action reduces to the identity at every marked place.

At an unmarked nonsplit place the residue characteristic is odd.
If \(y\) in the self-dual lattice has length \(-1\), its line
splits off integrally by orthogonal projection. Its orthogonal
complement is unimodular and has the same rank as the complement
of \(z\). Every unimodular Hermitian lattice over this unramified
quadratic extension admits an orthonormal basis. Indeed the
residue form has a vector of nonzero norm by polarization and
nondegeneracy. Split off a lift of unit norm value, rescaling
that value to one by surjectivity of the norm on units. The
last surjectivity follows from the residue norm and Hensel
lifting on principal units. Induction gives the basis.
An integral isometry of the two complements, extended by
\(z\mapsto y\), now gives the requisite compact element.

At a split \(v\), the two halves of \(L_v\) are an integral
module and its dual. Write \(z=(t,\lambda)\) and
\(y=(t',\lambda')\), so that \(\lambda(t)=\lambda'(t')=-1\).
The action on pairs is
\((t,\lambda)\mapsto(At,\lambda A^{-1})\).
Take an integral basis \(t,e_2,\ldots,e_u\) with the last
vectors in the kernel of \(\lambda\); such a basis exists since
\(\lambda(t)\) is a unit. At an unmarked place one can take
an analogous basis for \(y\) and map the former basis to it.
At a marked place the coset condition gives
\(t'\equiv t,\ \lambda'\equiv\lambda\bmod\pi^{a_v}\).
Define
\[
    At=t',\qquad Ae_j=e_j+\lambda'(e_j)t'\quad (2\le j\le u).
\]
Then \(A\equiv I\bmod\pi^{a_v}\), so it is invertible, and
\(\lambda' A=\lambda\) on the chosen basis. This proves the
claim at every split place, including the unmarked dyadic ones.
Taking the product of the local compact elements proves the
finite adelic assertion. The stabilizer in \(G\) of \(z\) is
\(H\) because the line is nondegenerate. This identifies the
stated support on the quotient.

Integration on the image of \(K_f\) gives
\(\operatorname{vol}_G(K_f)/\operatorname{vol}_H(K_f\cap H)\).
The intersection is the same level in dimension \(u-1\), by
the orthogonal splitting of \(L\). The unmarked quotient
volumes are one. At a marked place, reduction of
\(\mathrm{GL}_d(\mathcal O_v)\) modulo \(\pi^{a_v}\) gives
index \((Nv)^{a_vd^2}\prod_{k=1}^d(1-(Nv)^{-k})\).
Taking the ratio for \(u-1,u\) gives the claimed formula.
\end{proof}

\section{The mean terminal trace}\label{nav:g-XI-3}


\subsection{Counting the flags fixed by a residue comparison}\label{nav:III-65}
The upper unitriangular group \(U_k(\mathbb F_q)\) stabilizes
a complete flag with a specified nonzero vector on each
one-dimensional quotient. We compute the probability
\(P_k(A;q)\) of Proposition~\ref{prop:loops302} by forgetting
these quotient vectors.

\begin{lemma}[Fixed complete flags of a unipotent matrix]\leavevmode
\label{lem:fixedflags302}
If \(A\) is not unipotent then \(P_k(A;q)=0\).
Otherwise let \(\lambda\) be the partition of \(k\) given
by the Jordan blocks of \(A-1\). Write \(b_i=\lambda'_i\)
for the number of blocks of length at least \(i\), put
\(\ell=k-b_1\), and write \(P(\lambda)=P_k(A;q)\).
With \(P(\varnothing)=1\) one has
\[
 P(\lambda)=\sum_{i:b_i>b_{i+1}}
         \frac{q^{b_i}-q^{b_{i+1}}}{q^k-1}
         P(\lambda-\square_i),\qquad
 P(\lambda)\le k^\ell q^{-(k^2-\sum_i b_i^2)/2}.
\]
Here \(\lambda-\square_i\) shortens a length-\(i\) block
by one. Uniformly for \(q\ge2\) and sufficiently large \(k\),
\(P(\lambda)\le q^{-k\ell/3}\).
\end{lemma}
\begin{proof}
Every matrix conjugate into \(U_k\) is unipotent. Conversely
if \(A\) is unipotent and preserves a complete flag, it acts
as one on every successive quotient. All decorations of that
flag are then fixed. Each ordinary complete flag has the
same number of decorations, so it suffices to choose a
uniform ordinary flag and test its invariance.

Put \(N=A-1\). The first line \(L\) must lie in \(\ker N\).
For \(i\ge1\) the nested spaces \(\ker N\cap\operatorname{im}
N^{i-1}\) have dimensions \(b_i\), as is seen on each Jordan
block. The fraction of all first lines in the \(i\)-th space
and outside the next is the coefficient in the formula.
For such a line the image rank of \(N^a\) on the quotient
by \(L\) is \(\operatorname{rank}N^a-\mathbf{1}_{a<i}\) for
every \(a\ge0\). The list of ranks determines the Jordan
partition and gives precisely \(\lambda-\square_i\). The
rest of a uniform flag with first line \(L\) is uniform
on the quotient. This proves the recursion over every
finite field.

Write \(c(\lambda)=(k^2-\sum_i b_i^2)/2\).
Removing the indicated box decreases \(c\) by \(k-b_i\).
Each coefficient in the recursion is at most \(q^{-(k-b_i)}\),
because \(1-q^{-(b_i-b_{i+1})}\le 1-q^{-k}\).
On any sequence of removals the product of these upper
bounds telescopes to \(q^{-c(\lambda)}\). To bound the
number of sequences, fill each removed box with the number
of boxes present just before that removal. These fillings
increase along rows and columns and recover the entire
sequence. There are \(\ell\) boxes outside the first column.
Recording their entries in fixed spatial order gives at
most \(k^\ell\) possibilities. The remaining entries in
the first column must increase, so the record determines
the whole filling.

Finally \(\sum_i b_i^2\le b_1\sum_i b_i=b_1 k\), whence
\(c(\lambda)\ge k\ell/2\). For all sufficiently large \(k\),
\(k\le 2^{k/6}\le q^{k/6}\). This proves the last bound,
including the identity, for which \(\ell=0\) and \(P=1\).
\end{proof}

\subsection{The trace when the two maps agree on the whole perpendicular}\label{nav:III-66}
Keep the word, prescribed terminal space and retained fraction
\(\vartheta\ge1/2\) of Proposition~\ref{prop:tracemass301}.
The rank in the next result is the rank of
\(p_1-\beta p_0\) on the common rational perpendicular.

\begin{theorem}[The full rank-zero comparison contribution]\leavevmode
\label{thm:zerotrace303}
The sum of the absolute normalized cell traces over all
matching, relevant smaller-column transverse positions
with \(\mathcal G>0\) and comparison rank zero is at most
\[
                         (k!)^{-df}\exp(Cfkd)
\]
on the common fields and intervals. All endpoint heights
and all values of the saturation index occur in this estimate.
\end{theorem}
\begin{proof}
Fix a label, set \(R=A_-/Q^{df}\), and put \(t_0=k-2d-2\).
This is nonnegative on the common intervals. On every cell
being counted \(I\ge R\) by Theorem~\ref{thm:agreement303}.
The index is constant on all refinements of a coarse row
cell. Theorem~\ref{thm:moment219} with \(m=k-2d\) bounds
its \(t_0\)-th moment in the normalized compatible finite
measure by \(\exp(Cf)\). Hence the fraction with \(I\ge R\)
is at most \(\exp(Cf)\min(1,R^{-t_0})\), by Markov's
inequality and the trivial bound one.

In Proposition~\ref{prop:tracemass301} the estimate before
summing the finite covolumes has a uniform factor depending
only on the label. We may therefore multiply the bound of
Theorem~\ref{thm:product301} by this fraction when restricting
to the present cells. Give just the terminal endpoint power
\(\theta_t=k-t_0=2d+2\). Its weighted cost
\(L=d(2d+2)\) satisfies \(L\le r^2/2\), and the extra
exponent \(CfdL\log b\) is \(O(fkd)\) on the original intervals.
In the product bound the remaining factors obey
\[
     Q^{-dfk}A_-^k\min(1,R^{-t_0})H_t^{-(k-t_0)}
                    \le1 .
\]
For \(R\le1\) use \(H_t\ge1\). For \(R>1\) use
\(R\le A_-\le H_t\). The bound for this label is now
\((k!)^{-df}\exp(Cfkd)\prod_i H_i^{-3}\).
Summing by the integral endpoint count in
Theorem~\ref{thm:ranges301} costs at most \(4^d27^{fd}\),
which is absorbed in the stated exponential.
\end{proof}

\subsection{An integral column for the ordinary series}\label{nav:III-67}
On a good orbit of Theorem~\ref{thm:commonorbit305} fix a
representative with positive column lattice \(L\) and a
prescribed terminal character \(\lambda\). For an archimedean
covector \(\ell\), let \(b(\ell)\) be the column supported
on the orbit with coefficients obtained by transporting
\(\ell\) and multiplying by \(\overline{\lambda(z)}\) at
the scalar \(z\in Z\), in the convention of
Theorem~\ref{thm:dyad234}. Its squared class norm is
\(c_{\rho,K^0}|Z|\|\ell\|^2\).

Use the integral covector Schur model of
Proposition~\ref{prop:schurframe259}. If the polynomial
degree after subtracting the determinant power \(j_\tau\)
is \(e_\tau\) at \(\tau\), put
\(e_s=\sum_\tau e_\tau,\ a_s=\sum_\tau(e_\tau+k|j_\tau|)\).
Here the highest weights refer to the representation whose
dual supplies the covectors. Write \(d_\rho=\dim\rho\).
For \(D=2fk\) define
\[
 R_s=2D^{D/2}|\Delta_E|^{k/2}T^f,\qquad
                     B_s=(\sqrt{k}R_s)^{kf}.
\]

\begin{proposition}[A literal integral character column on a good orbit]\leavevmode
\label{prop:source305}
There are \(k^{e_s}\) integral covectors on the chosen class
whose frame operator on the complex covector space lies
between \(B_s^{-2a_s}\) and \(B_s^{2a_s}\) times the identity.
For at least one of these covectors,
\[
 \|\Theta S b(\ell)\|^2\ge
  c_{\rho,K^0}|Z|\,\alpha e^{-\Delta_x}
    \frac{d_\rho B_s^{-2a_s}}{k^{e_s}}M_{J,{\cal O}} .
\]
Each coefficient of \(b(\ell)\) belongs to the original
integral covector model after adjoining the character
values.
\end{proposition}
\begin{proof}
The lattice \(L\) has Hermitian determinant \(T^f\) by
Lemma~\ref{lem:termparent221}. Its trace lattice has
rank \(D\) and Euclidean covolume
\(|\Delta_E|^{k/2}T^f\). This follows from a free sublattice
by the trace discriminant identity and passage through the
index, as in Lemma~\ref{lem:terminalclass259}. Every
nonzero vector has trace norm at least one, by positivity
and integrality. The unit ball contains the cube of side
\(2/\sqrt D\). Minkowski's second theorem therefore gives
\(D\) rationally independent lattice vectors of length at
most \(R_s\). They span the space over \(E\), so choose
\(k\) that form an \(E\)-basis. At every complex embedding
their column matrix has largest singular value at most
\(\sqrt{k}R_s\). The product of the squared moduli of its
determinant over the \(f\) embeddings is the norm of a
nonzero integral Gram determinant and is at least one.
Hadamard's inequality at the other embeddings, and the
upper bound for the other \(k-1\) singular values at the
chosen one, give least singular value at least \(B_s^{-1}\).
The same vectors yield integral functionals by Hermitian
contraction, with the same singular values.

Take all ordered tensors of these functionals and restrict
them to the integral Young image. Include the determinant
factor by the exterior product of the chosen vectors or
functionals according to sign, exactly as on the dual side
of Proposition~\ref{prop:schurframe259}. The tensor and
exterior singular-value inequalities give the two claimed
frame bounds. In representatives obtained by terminal scalar
translation, the lattice \(L\) is the same because the
covector kernels are fixed. The transport of the coefficients
introduces only the indicated unitary level characters.
These have root-of-unity values on the finite markings.
Thus the orbit sum is integral in its stated model.

The map \(i:\ell\mapsto b(\ell)/(c_{\rho,K^0}|Z|)^{1/2}\)
is an isometry onto the orbit character space. The first
moment is \(d_\rho^{-1}\|S i\|_{\rm HS}^2\) by its trace
identity. The matching complex norm theorem
\ref{thm:range239} implies
\(\|\Theta S i\|_{\rm HS}^2\ge
\alpha e^{-\Delta_x}d_\rho M_{J,{\cal O}}\)
by summing over an orthonormal basis. Taking the trace
of \((\Theta S i)^*\Theta S i\) against the lower frame
bound gives a lower bound \(B_s^{-2a_s}\|\Theta S i\|_{\rm HS}^2\)
for the sum over the integral list. Some member attains
at least its average. Undoing the normalization of \(i\)
proves the assertion.
\end{proof}

\subsection{The finite mass in the normalized terminal trace}\label{nav:VI-57}
Work on a field and interval of Section~\ref{sec:missing} with \(d=s+1\),
\(f=[F:\QQ]\), \(h=k-1\), and \(Q=g_1g_2\). Fix one
prescribed word. Let \(J_0\) be a prescribed nontrivial
terminal character space or the simultaneous zero-sum space
of Theorem~\ref{thm:dyad234}. One may retain all its finite
class coordinates or a union of complete scalar orbits of
fraction \(\vartheta\ge1/2\) in the full set of finite classes.
Theorem~\ref{thm:columns235} gives such a set on the common
intervals. In this section \(P_0\) denotes the projection
onto this possibly restricted space. Write each ordinary
orbital operator \(O_\gamma\) as the sum of its finite-cell
terms \(O_{\gamma,c}\). For a supported relevant transverse
pair with \(\mathcal G>0\) set
\[
     B_\gamma=
       \frac{\sum_c|\operatorname{tr}(P_0 O_{\gamma,c})|}
            {\operatorname{tr} P_0}.
\]
Take the original level and the real group measures of
Theorem~\ref{thm:realtrans300}. The real sections have diagonal
norm one. Write \(a=\sum_\nu\binom{h_\nu}{2}\) for the block
unipotent dimension over \(E\), and \({\cal P}_\tau\) for the
positive, locally normalized orbital integral there.

\begin{proposition}[Cancellation of the fine level and of the word dimension]\leavevmode
\label{prop:tracemass301}
With the finite mass \(\mathcal M_\gamma^{=}\) of
Theorem~\ref{thm:allfinite215}, the projected trace satisfies
\[
 B_\gamma\ \le\
 \frac{1}{\vartheta}\,
 \frac{\mathfrak m_{k-1}}{\mathfrak m_k}\,
 \frac{\mathcal M_\gamma^{=}/\mathfrak m_{k-1}}{(T^k-1)^f}\,
       \prod_\tau{\cal P}_\tau .
\]
One can use the local factors for any unit pure tensor in
the prescribed lowest type.
\end{proposition}
\begin{proof}
Give every compact real column group probability and each
full finite row lattice stabilizer \(K_v^*\) volume one.
At \(T\) this is the modified row lattice of
Theorem~\ref{thm:midlevel131}. At \(Q\) write
\(I_\sigma=[K^*:K^0]\) and let \(I_\pi\) be the corresponding
index on the definite column. There are
\[
             n_{\rm cl}=\mathfrak m_k I_\pi(T^k-1)^f
\]
fine column classes. Indeed each terminal covector has
stabilizer of index \(T^k-1\) in its parent compact by
transitivity on nonzero covectors. Multiplying the genus
mass by the local indices gives the refined mass by
orbit--stabilizer, and the fine rational stabilizers are
trivial by Lemma~\ref{lem:covol127}. Each retained class
carries the same \(d_\rho=\dim\rho\) real coordinates.

In the coordinate formula of Theorem~\ref{thm:dyad234} two
classes can contribute only on the same scalar orbit. For
each such ordered pair the coefficient in the normalized
trace has modulus at most \(1/(\vartheta n_{\rm cl}d_\rho)\),
with the basis elements on the two classes matched by
unitary transport. On a cell, sum over the orthonormal
real covectors before taking the absolute value. At each
real integration point this sum is the inner product of
the transported two real vectors, times a phase of modulus
at most one. This follows from the completeness relation
of an orthonormal basis. Its modulus is bounded by the
product of the two norms, including when the common
compact group is only \(U(k-1)\).

The unit class-coordinate inner product includes
\(c_{\rho,K^0}\) from Theorem~\ref{thm:identity127}. The
sum of the covolume weights of all fine cells is
\(c_E^a(I_\pi/I_\sigma)\mathcal M_\gamma^{=}\), by the
refinement and mass computations of Theorem~\ref{thm:norm212},
with \(c_E=2^{-f}\sqrt{|\Delta_E|}\). The integral of the
two real norms over the common unipotent quotient equals
the integral over the quotient by the whole common group,
since the norms are invariant under the compact factor, which
has probability one. On a pure tensor this factors over
the real places. The trace of each signed cell, polarized
in its lowest-type row vectors, is invariant under simultaneous
right translation by the product maximal compact. Haar
invariance gives this directly in the convergent integral.
Schur's lemma on the irreducible product type therefore
gives the same signed trace on every unit row vector.

The identity constant has
\(d_\rho c_{\rho,K^0}=v_N/(I_\sigma\operatorname{vol}_f K_N^*)\),
where \(K_N^*=N_f\cap K^*\). Indeed the integral unipotent
levels agree at \(Q\) by Lemma~\ref{lem:flags211}. Substitution
leaves exactly the proposed right side times
\(c_E^a\operatorname{vol}_f(K_N^*)/v_N\).

This last scalar is at most one. For detail, separate \(N\)
into the radical of the full \(s\)-plane and the upper
unitriangular group on its \(s\) coordinates. With ordinary
Euclidean Witt coordinates at infinity, the covolume of
the radical is the one in Lemma~\ref{lem:blockphysical295}
at \(j=s,\ n=k+1\). It has \(s(k+1)\) complex rectangular
variables, \(\binom{s}{2}\) complex skew entries and \(s\)
trace-zero diagonal entries. The additional upper group
has \(\binom{s}{2}\) integral \(E\)-coordinates, with
additive covolume \(c_E\) each. The integral upper matrices
lift on the two isotropic strings and have determinant one
on the radical, so the covolumes multiply. The middle graph
at \(T\) preserves the unipotent volume by
Theorem~\ref{thm:midnorm131}. This gives
\[
   \frac{v_N}{\operatorname{vol}_f K_N^*}
     =2^{fs(k+d)}c_E^{\,s(k+s)}|\Delta_F|^{s/2}.
\]
Here the power of two comes from rescaling the integral
Witt pairs to Euclidean length one, as in
Lemma~\ref{lem:blockphysical295}. Since \(c_E\ge2^{-f}\),
\(|\Delta_F|\ge1\), and \(a\le s(k+s)\), the quotient
of \(c_E^a\) by this covolume is at most \(2^{-f(s+a)}\).
This proves the proposition.
\end{proof}

\subsection{Counting single flags with a prescribed comparison rank}\label{nav:VI-58}
Fix one integral cross class \(Z\) with \(\mathcal G>0\)
on the original row. Equivalence of terminal flags uses
integral upper changes of frame and the row group
\(N_\Lambda\) of Section~\ref{sec:frcount225}. For a fixed
row and initial flag count each terminal flag class once
if it has a representative admitting a scalar-orbit
comparison of rank \(h\) from Theorem~\ref{thm:rationalrank302}.
Use the normalized average on the initial terminal column
genus. The bound below also applies when an arbitrary row
over each column is chosen, or is averaged with a probability
measure.

\begin{theorem}[A mean bound for one integral cross class]\leavevmode
\label{thm:singlemean303}
Put \(c=h+1,\ j=k-c\), and suppose \(1\le j<k\). With
\(n_T=T^f,\ \nu=4+\log_2(2f)\), set
\[
    Y=2^{fh}\prod_\tau (|\tau D_d|+T|\tau D_s|),\qquad
    A_c(Z)=9^{cf}\max(1,G/n_T)^c J_Z^{c+s},
\]
where \(J_Z\) is the index of the maximal-minor ideal
in Theorem~\ref{thm:frcount225}. For \(Y\ge1\) the mean
number of the stated flag classes is at most
\[
       2\,\frac{\mathfrak m_j\mathfrak m_c}{\mathfrak m_k}\,
       (Y/n_T)^k n_T^{c+2s}\,Y^{\nu+1+s} A_c(Z).
\]
For \(Y<1\) the count is zero.
\end{theorem}
\begin{proof}
If a representative admits the comparison, let
\(M=p_0(\Lambda\cap\ker(p_1-\beta p_0))\) in the initial
positive column lattice \(L\). This has rank \(j\), lifts
integrally to \(\Lambda\cap C\), and has determinant at
most \(Y\) by Theorem~\ref{thm:agreement303}.
The existence of a sublattice with these properties suffices
for the upper bound. In particular the determinant of an
integral Hermitian lattice is at least one, which gives
the zero assertion.

Let \(M_*=(EM)\cap L\), \(D=\delta(M_*)\) and
\(a=[M_*:M]\). Thus \(Da\le Y\). The determinant \(D\)
is an integer by the local discriminant formula used in
Theorem~\ref{thm:termmean221}. For fixed \(M_*\) there
are at most \(a^{j+\lambda}\) submodules of index \(a\),
where \(\lambda=\log_2(2f)\), by Lemma~\ref{lem:indices225}.
For each one Theorem~\ref{thm:frcount225} bounds the
number of flags lifting it by \((n_T^2 a)^s A_c(Z)\).
These two bounds hold in every row over the given column.
The resulting upper bound is a function of \(L\) alone.
By Theorem~\ref{thm:termmean221} the mean number of
possible \(M_*\) at determinant \(D\) is at most
\[
       \frac{\mathfrak m_j\mathfrak m_c}{\mathfrak m_k}
                   n_T^{c-k}D^{k+\nu}.
\]
The normalized measure on any of the original full
markings gives the same mean for a function of \(L\)
by Lemma~\ref{lem:termparent221}. The formula concerns
the full mean, so restricting the count to specified
initial coordinates can only reduce its unnormalized sum.

After removing the constant factors we must bound
\(\sum_{a\le Y}a^{j+\lambda+s}
\sum_{D\le Y/a}D^{k+\nu}\). The inner sum is at most
\((Y/a)^{k+\nu+1}\). Since \(j=k-c\) and
\(\nu=4+\lambda\), the remaining sum of powers of \(a\)
is bounded by
\[
       \sum_{a\le Y} a^{s-c-5}
              \le Y^s\sum_{a\ge1}a^{-c-5}\le 2Y^s .
\]
Restoring the factors proves the result.
\end{proof}

\subsection{Two ranges in the original terminal trace}\label{nav:VII-59}
Keep the original word, terminal space, fields and intervals
of Theorem~\ref{thm:product301}. For \(\mathcal G>0\) write
\(H_*=\max_i H_i^{1/j_i}\), using the heights of its
normalized integral endpoints. Recall that \(A_-\) measures
the product of the differences
\(|\tau D_d|-T|\tau D_s|\), divided by \(Q^{df}\) and
truncated below at one.

\begin{theorem}[The small terminal differences and the large endpoint heights]\leavevmode
\label{thm:ranges301}
There is a constant \(M\) such that, for every fixed
\(0<\epsilon_0<1\), the sum of \(B_\gamma\) over all relevant
supported transverse pairs of smaller common column satisfying
at least one of
\[
             A_-\le (Q k^{\,1-\epsilon_0})^{df},
             \qquad H_*\ge b^{Mfd}
\]
is at most \(\exp(-(\epsilon_0-o(1))fkd\log k)\).
Here \(o(1)\to0\) on the common intervals, and the estimate
is uniform in the prescribed character and word. On each
field at each point of an interval the complement is a
finite set of tuples of algebraic integers.
\end{theorem}
\begin{proof}
For any fixed composition the original block data have
multiplicity at most \(\prod_i h_i\le2^d\) over the tuple of
the \(a_i\), by Lemma~\ref{lem:forms203}. Indeed the proper
minors are \(Q^{j_i}a_i\), and the last original minor is
\(Q^d a_t-TD_s\), so every original endpoint is recovered.
There are at most \(2^{d-1}\) compositions. For
\(H(z)=\prod_\tau\max(1,|\tau z|)\), the sum of \(H(z)^{-3}\)
over the nonzero integers of \(E\) is at most \(27^f\), as
in Theorem~\ref{thm:codimension294}. One can see this by
packing disjoint polydiscs of radius \(1/2\) at the integer
points in the complex embedding. Distinct centers have sup
distance at least one by the norm inequality. A box with
radii \(r_\tau\ge1\) thus contains at most
\(9^f\prod r_\tau^2\) points. Integrating against
\(\prod_\tau 3r_\tau^{-4}dr_\tau\) gives the stated bound.
Therefore the sum of \(\prod_i H_i^{-3}\) over all labels
is at most \(4^d 27^{fd}\).

For the first range use Theorem~\ref{thm:product301} with
all \(\theta_i=0\). The factor \(Q^{-dfk} A_-^k\) is at
most \(k^{(1-\epsilon_0)dfk}\). Since
\(\log(k!)\ge k\log k-k\), the factorial and the preceding
count give the required exponent.

For the second range take an index \(i\) attaining \(H_*\).
In the same theorem put power \(k\) at the terminal endpoint
and additional power \(kd/j_i\) at endpoint \(i\). Add the
two powers if the indices agree. Their weighted sum is
\(2kd\le r^2/2\) on sufficiently distant intervals.
The first power absorbs \(A_-^k\) by \(A_-\le H_t\).
The other contributes \(H_*^{-kd}\). Take \(M\) larger than
twice the constant of the term \(CfdL\log b\) in the product
theorem. The resulting negative power of \(b\) absorbs
that entire term. Summation as above gives at most
\((k!)^{-df}\exp(Cfkd)Q^{-dfk}\) on this range. The sum of
the two upper bounds proves the assertion for their union.

Outside the two ranges, each \(a_i\) is integral and every
one of its complex embeddings has modulus at most
\(H_i<b^{Mfdj_i}\). The integral embedding is a lattice
and hence meets the resulting compact box in finitely many
points. The block multiplicity already proved gives the
last assertion.
\end{proof}

Combining this result with Theorem~\ref{thm:allflags298}
in the first moment of Theorem~\ref{thm:dyad234} leaves only
the smaller-column transverse terms with both a larger terminal
difference and bounded endpoint heights. The theorem gives
a uniform contribution outside that range. Controlling the
signed trace inside it and comparing the resulting period
with the designated integral pairing are still required.

\subsection{A rank estimate in the ordinary terminal trace}\label{nav:VII-60}
Use the normalized terminal trace and the retained coordinate
space of Proposition~\ref{prop:tracemass301}. For a fine
transverse cell \(c\) with its two columns in one scalar orbit
take the unique comparison of Proposition~\ref{prop:loops302}.
Define
\[
 \ell_v(c)=\operatorname{rank}_{\mathbb F_{q_v}}
                      (A_{\alpha,v}-1),\qquad
 S(c)=\sum_{v\mid Q}\ell_v(c)\log q_v .
\]
These numbers are independent of the coordinates on the
coarse rational cell. The other cells have zero trace by the
character projector formula in Theorem~\ref{thm:dyad234}.

\begin{theorem}[The terminal difference and the comparison ranks]\leavevmode
\label{thm:ranktrace302}
Fix \(0<\epsilon_0<1\). On the fields and intervals of
Section~\ref{sec:missing}, sum the absolute normalized cell traces over
all relevant smaller-column transverse labels with
\(\mathcal G>0\). Restrict the sum to matching cells satisfying
\[
    A_-\le (Q k^{\,1-\epsilon_0})^{df}\exp(S(c)/3),
       \quad\text{or}\quad H_*\ge b^{Mfd}.
\]
Here \(A_-,H_*,M\) are as in Theorem~\ref{thm:ranges301}.
The sum is at most \(\exp(-(\epsilon_0-o(1))fkd\log k)\).
In the complementary range the height tuples form a finite
set on each field at each parameter value and
\[
 S(c)< 3\log\!\left(\frac{A_-}{(Qk^{\,1-\epsilon_0})^{df}}\right)
        \le 3Mfd^2\log b .
\]
\end{theorem}
\begin{proof}
First fix a label and impose \(S(c)\ge S_0\) with \(S_0\ge0\).
For a contributing comparison each marked reduction is
unipotent. Lemma~\ref{lem:fixedflags302} bounds its proportion
of all common refinements by \(\exp(-kS_0/3)\).
The number of coarse rational comparisons above one cell
is at most \(\exp(Cfk\log^2(2fk))\).
Proposition~\ref{prop:loops302} therefore gives the product
of these two bounds for the proportion of the finite measure
above each compatible coarse cell. The same upper bound
holds if we retain only the chosen class coordinates.

The proof of Proposition~\ref{prop:tracemass301} bounds
the absolute normalized trace of each retained cell by its
finite covolume times a bound independent of the cell for
this label. We may thus insert the preceding factor before
summing the covolumes. Every subsequent estimate in
Theorem~\ref{thm:product301} concerns the label alone.
Putting all its free endpoint powers equal to zero gives
the bound
\[
 \frac{\exp(Cfkd-kS_0/3)}{(k!)^{df}}\,
       Q^{-dfk} A_-^k\prod_i H_i^{-3}.
\]
We used \(\log^2(2fk)=O(d)\), since \(f,k\) are polynomially
bounded in \(x\) and \(d\asymp x^{1/2-\eps/16}\).

For the first range put \(V=(Q k^{1-\epsilon_0})^{df}\) and
take \(S_0=3\max(0,\log(A_-/V))\). Thus
\[
       Q^{-dfk}A_-^k\exp(-kS_0/3)\le k^{(1-\epsilon_0)dfk}.
\]
The sum of \(\prod_i H_i^{-3}\) over the labels is at most
\(4^d27^{fd}\) by the lattice count in
Theorem~\ref{thm:ranges301}. The factorial bound
\(\log(k!)\ge k\log k-k\) proves the desired exponent.
On the second range Theorem~\ref{thm:ranges301} already
bounds the sum of the absolute traces of all cells. Adding
the two bounds proves the assertion for the union.
Outside the union the first inequality is reversed and
\(A_-\le H_t\le H_*^d<b^{Mfd^2}\). These inequalities
give the displayed bounds. Finiteness follows from the
integrality of the endpoints as in Theorem~\ref{thm:ranges301}.
\end{proof}

\subsection{The first moment for the prescribed ordinary Bessel series}\label{nav:VII-61}
Work at the original terminal frequency, field, interval
and fine level of Section~\ref{sec:missing}. For each prescribed word
take a unit pure real vector in its stipulated lowest
type. Let \(S\) be its normalized ordinary Bessel
periodization, \(J\) a specified character space with all
local terminal characters nontrivial, or the simultaneous
zero-sum space. Allow the retained classes of
Corollary~\ref{cor:trace235} of fraction \(\vartheta\ge1/2\).
Write \(\xi\) for a random unit signed column from the
first-moment formula on \(J\), and \(M_J=\mathbb E\|S\xi\|^2\).

\begin{theorem}[The ordinary first moment on each prescribed terminal space]\leavevmode
\label{thm:ordinarymoment304}
Uniformly in the prescribed word and terminal space,
\[
          |M_J-1|\le\exp\bigl(-(1-o(1))fk\log k\bigr).
\]
For the matching complex theta operator in
Theorem~\ref{thm:range239}, with its norm factor
\(\alpha=C_xP(\mathbf a)>0\), some unit signed column satisfies
\[
   \|\Theta S\xi\|^2\ge
          \alpha e^{-\Delta_x}
          \left(1-\exp\bigl(-(1-o(1))fk\log k\bigr)\right)>0
\]
on sufficiently distant common intervals.
\end{theorem}
\begin{proof}
The ordinary Gram formula of Theorem~\ref{thm:dyad234} and
Corollary~\ref{cor:trace235} expresses \(M_J-1\) as the
real part of the sum of the nonidentity orbital traces
divided by \(\dim J\). For coincident isotropic spaces
only the identity contributes at this level by
Proposition~\ref{prop:finite237}. Theorem~\ref{thm:allflags298}
bounds the sum of operator norms for proper nonzero
intersections and for transverse pairs with full common
column by \(((k+1)!)^{-f}\exp(Crf)\). This also bounds
their absolute normalized traces. The remaining labels
are transverse with smaller common column and
\(\mathcal G>0\), as in Theorem~\ref{thm:cross132}.

Use Theorem~\ref{thm:ranktrace302} with \(\epsilon_0=1/2\).
Its two ranges together cost at most
\(\exp(-(\tfrac12-o(1))fkd\log k)\). Outside them,
\(P_-=\prod_\tau (|\tau D_d|-T|\tau D_s|)>1\)
and \(H_*<b^{Mfd}\). Theorem~\ref{thm:zerotrace303} already
bounds all comparison ranks equal to zero. For every other
contributing cell Theorem~\ref{thm:rationalrank302} gives
\(1\le h\le C_0d^2\) with one constant. We can enlarge
the sum to all cells of these ranks at the same labels.

For such a label multiply Proposition~\ref{prop:phase301}
over the real places with all free endpoint powers zero.
With \(\kappa=k-d,\ E_0=\kappa+2ds+s^2\) it gives
\[
 \prod_\tau{\cal P}_\tau\le
   \frac{\exp(Cfkd)}{(k!)^{sf}}\,
                   G^{-\kappa}P_-^{\,2\kappa-E_0} n_T^k .
\]
Indeed the between-block partial pivot products at any
index with \(u_p>0\) are integral after multiplication by
\((2T)^p\) in modulus. More explicitly, their product over
all places, after this rescaling at each one, is the
product of the moduli of the nonzero integral intermediate
pivot product \(\omega_p\) of Lemma~\ref{lem:classes304}.
The final phase leaves those pivots unchanged. Its normalized
inverse powers cost at most \((2T)^{2f\sum_p p u_p}\).
The exponent of \(T\) at one place becomes
\(dE_0+2\sum_p p u_p-s(k+d)=k+O(d^3)\).
The extra factor costs \(\exp(O(fd^4\log b))\) because
\(\log T=O(d\log b)\). The powers of two cost at most
\(\exp(Cfkd)\). This proves the displayed estimate.

Now apply Proposition~\ref{prop:flagtrace304}.
Its \(n_T^{-k}\) cancels \(n_T^k\). The remaining Gram power
is \(G^{k-\kappa}=G^d\), whose logarithm is \(O(fd^3\log b)\).
The exponent of \(P_-\) is \(2\kappa-E_0-k<0\), so this
factor is at most one. The error \(\exp(Cfd^6\log b)\)
fits in \(\exp(Cfkd)\) since \(d^5\log b=o(k)\) by the
parameter sizes in Theorem~\ref{thm:rangenew50}.
It remains to bound the genus factor times \(2^{fhk}\).

Put \(c=h+1,\ j=k-c\), so \(c,j\ge2\) and \(c=O(d^2)=o(k)\).
Theorem~\ref{thm:genus222} gives
\(\mathfrak m_j\mathfrak m_c/\mathfrak m_k
\le \mathfrak m_2(C/(kR_F))^{fjc}\), with root
discriminant \(R_F\ge1\) and \(\mathfrak m_2\le b^{Cf}\).
Uniformly on this range \(jc=k(1+h)-c^2\ge k(1+h/2)\)
for large \(x\), because \(c^2/(kh)=o(1)\).
For sufficiently large \(k\), the decay
\((k/C)^{-fkh/2}\) absorbs \(2^{fhk}\).
Consequently
\[
       2^{fhk}\frac{\mathfrak m_j\mathfrak m_c}{\mathfrak m_k}
                           \le (k!)^{-f}\exp(Cfk).
\]
This leaves \((k!)^{-df}\exp(Cfkd)\) for each fixed
rank and label.

The number of labels in the remaining range is at most
\(4^d27^{fd}\exp(Cfd^3\log b)\). In fact
\(\sum_\gamma\prod_i H_i^{-3}\le4^d27^{fd}\) by
Theorem~\ref{thm:ranges301}, and \(H_i\le b^{Mfd^2}\)
for at most \(d\) endpoints. The number of ranks is
at most \(k\). These counts leave the same bound
\((k!)^{-df}\exp(Cfkd)\) for the entire intervening
sum of absolute cell traces. Add the two height ranges
and the full-column estimate above. As \(d\to\infty\),
\(r=k+2s+1\sim k\), and \(\log(k!)\ge k\log k-k\), the
stated error for \(M_J\) follows.

For the specified whole cuspidal row the ordinary image
already belongs to the projected space. The two-sided
norm estimate in Theorem~\ref{thm:range239} now gives
\(\mathbb E\|\Theta S\xi\|^2\ge\alpha e^{-\Delta_x} M_J\).
The signed column ensemble is finite, so a member
attains at least this mean. This proves the second assertion.
\end{proof}

The vector supplied by the first moment depends on the word
and on the terminal character. The arithmetic comparison
in Corollary~\ref{cor:original200} asks for the specified
integral row and its distinguished pairing with controlled
cost. The strict cohomology sum in Section~\ref{sec:missing} and this
integral comparison are still open.

\subsection{The determinant of the integral agreement lattice}\label{nav:VIII-58}
Use the spaces of Lemma~\ref{lem:archdet303}. In the original
row \(\Lambda\) let \(K=\Lambda\cap C_0\), the integral lattice
on which the two column maps agree. Put \(G=|N_{F/\QQ}\mathcal G|\).

\begin{theorem}[A determinant cutoff from an integral comparison]\leavevmode
\label{thm:agreement303}
For every matching transverse cell of positive Gram determinant,
with comparison of rank \(h\) as in Theorem~\ref{thm:rationalrank302},
\[
 \delta(K)\le
 \frac{2^{fh}G}{\prod_\tau (|\tau D_d|-T|\tau D_s|)}
     = 2^{fh}\prod_\tau (|\tau D_d|+T|\tau D_s|).
\]
In particular, if \(h=0\), then \(A_-\le Q^{df} I\), where
\(I\) is the index of Theorem~\ref{thm:saturation220}.
\end{theorem}
\begin{proof}
The original row is self-dual. Put \(B'=\Lambda\cap W'\).
Orthogonal projection of \(\Lambda\) onto \(W'\) gives
\((B')^\vee\). Locally every integral functional on \(B'\)
extends to the row since \(B'\) is primitive, and self-duality
represents the extended functional. The same reasoning
applies to \(K\) in \(C_0\). The quotient
\(\Lambda/(B'\oplus K)\) identifies by projection with each
discriminant module. Indeed subtraction of a projected
element of \(B'\), or of \(K\) on the other side, gives the
kernel assertion. Thus \(\delta(B')=\delta(K)\), using the
positive square root of the discriminant index also on the
nondegenerate indefinite lattice \(B'\).

Take primitive integral frame bases for
\(U_i=\Lambda\cap Y_i\) as in Lemma~\ref{lem:coordinates225}.
Changing either frame by an upper unitriangular matrix
preserves its full cross determinant and its leading
\(s\)-minor. Choose a full integral lattice \(M_1\) in
\(C_1\) contained in \(\Lambda\), and set
\(U'=U_0\oplus U_1\oplus M_1\), \(b_0=[B':U']\).
The Gram determinant of the first two frames has absolute
norm from \(F\) equal to \(G\). The third summand is
orthogonal to them. Passing to the over-lattice divides
the determinant by its additive index. It follows that
\[
                    b_0=\frac{G\,\delta(M_1)}{\delta(K)}.
\]

The orthogonal projections of \(\Lambda\cap Y_i^\perp\)
onto \(W'\) are exactly \(D_i=(B')^\vee\cap S_i\).
In fact every lift of a vector of \(D_i\) by the surjection
just proved differs from it by a vector of \(C_0\), itself
orthogonal to \(Y_i\). For integral lifts \(v,w\) on the
two sides, the row pairing is integral and so is the
column pairing since \(\beta L_0=L_1\).
The difference \({\cal H}(v,w)\) depends only on the
projections to \(W'\). Hence \({\cal H}(D_0,D_1)\subset\OO_E\).

In the coordinates of Lemma~\ref{lem:archdet303} the lattice
\((U')^\vee\cap S_i\) is \(\OO_E^d\oplus M_1^\vee\).
The index of \(D_i\) in it equals \(b_0\). To see this,
restrict the functionals on \(B'\) and on \(U'\) to \(U_i\).
Both maps onto its module dual are surjective because \(U_i\)
is primitive in each lattice and the quotient is projective
over the Dedekind ring. Their kernels are the two lattices
in \(S_i\). Consequently their index is the full dual index
\([(U')^\vee:(B')^\vee]=b_0\).
With the standard positive metric on the \(d\) coordinates
and the given one on \(C_1\), each \(D_i\) therefore has
normalized Euclidean covolume \(b_0/\delta(M_1)=G/\delta(K)\).
Here the normalized covolume of a module of rank \(a\)
uses the factor \(c_E^{-a}\) of Theorem~\ref{thm:rationalrank302}.
In this convention a Hermitian dual has reciprocal covolume.

The nondegenerate integral pairing maps \(D_1\) into the
Hermitian dual of \(D_0\) by the matrices \(N_\tau\).
Its image has positive integral index, equal by covolumes to
\[
          \left(\frac{G}{\delta(K)}\right)^2
                              \prod_\tau|\det N_\tau|^2 .
\]
Taking square roots and applying Lemma~\ref{lem:archdet303}
gives the first inequality. The identity follows by factoring
\(\tau\mathcal G\) at each place. If \(h=0\), then
\(K=\Lambda\cap C\), whose determinant is
\(G/(Q^{2df}I)\) by Theorem~\ref{thm:saturation220}. Substitution
gives \(\prod_\tau(|\tau D_d|-T|\tau D_s|)\le Q^{2df}I\).
Since \(Q^{df}I\ge1\), this proves the assertion for the
truncated \(A_-\) as well.
\end{proof}

\subsection{From the mean flag count to the terminal trace}\label{nav:VIII-59}
For a supported smaller-column transverse label put
\[
  P_-=\prod_\tau(|\tau D_d|-T|\tau D_s|),\qquad n_T=T^f .
\]
For one word and one retained terminal space as in
Proposition~\ref{prop:tracemass301}, let \(B_{\gamma,h}\)
be the sum of absolute normalized traces of its matching
fine cells of comparison rank \(h\). The normalization
divides by the dimension of the retained space.

\begin{proposition}[The finite count with the original trace weights]\leavevmode
\label{prop:flagtrace304}
Fix \(C_0\). When \(\mathcal G>0,\ H_*<b^{Mfd}\) and
\(0\le h\le C_0d^2\), put \(c=h+1,\ j=k-c\). On
sufficiently distant common intervals,
\[
 B_{\gamma,h}\le
   \exp(Cfd^6\log b)\,
   \frac{\mathfrak m_j\mathfrak m_c}{\mathfrak m_k}
   \left(\frac{2^{fh}G}{P_-n_T}\right)^k
   \prod_\tau{\cal P}_\tau .
\]
The \({\cal P}_\tau\) are the real factors of
Proposition~\ref{prop:phase301}. The constant is uniform
in the label, the word and the retained fraction
\(\vartheta\ge1/2\).
\end{proposition}
\begin{proof}
We first count fine rational cells above one initial
column coordinate \(y\). The initial finite support is
\(B_f K^0/K^0\). Its quotient by \(B(F)\) is the fine
column class space. Indeed \(B=N\rtimes H\), the
induced level projects onto \(K_H\) by
Theorem~\ref{thm:midlevel131}, and \(N(F)\) is dense
in \(N_f\) by successive additive approximation in its
unipotent coordinates. We may take a finite representative
\(t_y\) in \(H_f\subset B_f\). In the row lattice
\(\Lambda\) it specifies, the rational stabilizer of
the coset \(t_y K^0\) in \(B(F)\) is exactly
\(N_\Lambda\) of Section~\ref{sec:frcount225}.
In fact its column component lies in the trivial fine
rational column stabilizer. Within \(N(F)\), integrality
is sufficient. At the marked primes the whole integral
unipotent radical reduces into the row level by
Lemma~\ref{lem:flags211}, and \(t_y\) normalizes the
rational radical. Elsewhere the row level is the lattice
stabilizer, using the original graph at \(T\). This proves
the asserted stabilizer equality.

Write a row coset above \(y\) as \(b t_y K^0\) with
\(b\in B(F)\). If it supports the second decoration in
the relative position \(\gamma\), transport this decoration
back to \(t_y\) by \(b^{-1}\). This maps the fine rational
cells above \(y\) injectively to the terminal flag classes
used in Theorem~\ref{thm:singlemean303}. Here is the
injectivity with the stabilizers included. If two transported
flags agree after a change \(n\in N_\Lambda\), then
\(b_2 n b_1^{-1}\) belongs to both \(B(F)\) and
\(\gamma B(F)\gamma^{-1}\), since it stabilizes the
second decoration as a rational flag. It identifies the
original row cosets because \(n\) fixes \(t_y K^0\).
Integral upper changes of frame represent exactly the
same decorated rational flag among its primitive integral
representatives, as in Lemma~\ref{lem:coordinates225}.
The image of a rank-\(h\) matching cell has a representative
with the integral lift of determinant at most
\(Y=2^{fh}G/P_-\) used in Theorem~\ref{thm:singlemean303}.

In the trace estimate before summation in
Proposition~\ref{prop:tracemass301}, the bound for one
fine cell is
\[
 \frac{\operatorname{vol}_f K^0}
      {\vartheta n_{\rm cl}\,d_\rho c_{\rho,K^0}}\,
 \operatorname{vol}(\Gamma_D\backslash D_\infty)
                         \prod_\tau {\cal P}_\tau ,
\]
where \(n_{\rm cl}\) is the full number of fine initial
column classes. This follows by contracting the real
covectors in the normalized projector and integrating
over the common unipotent quotient before taking the
product of real norms. With the coarse row compact of
volume one, the factor
\(\operatorname{vol}_f K^0/(d_\rho c_{\rho,K^0})\)
equals \(\operatorname{vol}_f K_N^*/v_N\) by the same
proposition. There it was proved that this times
\(c_E^{a_0}\) is at most one. Proposition~\ref{prop:covol304}
therefore bounds the weight of a cell by
\(\exp(Cfd^6\log b)\prod_\tau {\cal P}_\tau/
(\vartheta n_{\rm cl})\).

Summing over \(y\) now gives the uniform average of the
flag counts on all fine initial coordinates. The pointwise
bound by sublattice counts in Theorem~\ref{thm:singlemean303}
depends only on the column lattice. Its average over these
coordinates is its genus mean, by Lemma~\ref{lem:termparent221}.
Discarded coordinates can be included in an upper estimate.
There are at most \(\exp(Cfd^6\log b)\) integral cross
classes by Lemma~\ref{lem:classes304}. For each such class
apply Theorem~\ref{thm:singlemean303}, using
\(J_Z\le G\) from Section~\ref{sec:frcount225}. If \(Y<1\)
the count vanishes. Otherwise
\[
       \log Y=O(fd^2\log b),\quad \log G=O(fd^2\log b),
       \quad \log n_T=O(fd\log b).
\]
For the first inequality use
\(Y=2^{fh}\prod_\tau(|\tau D_d|+T|\tau D_s|)\)
and the endpoint height bounds of Lemma~\ref{lem:classes304}.
Since \(c=O(d^2)\) and \(\nu=O(\log x)=O(d)\), the logarithm
of \(n_T^{c+2s}Y^{\nu+1+s}A_c(Z)\) is \(O(fd^4\log b)\).
The factors for the cross classes, the covolume and
\(\vartheta^{-1}\le2\) all fit in the displayed exponential.
The remaining factors are precisely those in the proposition.
\end{proof}

\section{Ordinary traces at the two marked primes}\label{nav:chapter-marked-ordinary-traces}

\subsection{The determinant on the common positive hyperplane}
\label{sec:corankone534}

Fix the two words and lattice rows of Section~\ref{sec:partnerlifts531}.
Assume that their terminal vectors have the same negative norm
\(-T_0\), as happens in every diagonal pairing. Retain \(Q_0=g_1g_2\)
and \(k\ge2\). On a cell with \(\dim(A/I)=k-1\), the maps
\(\pi_l:A/I\to P_l\) of Section~\ref{sec:fullrankresidue533}
are isometric embeddings onto hyperplanes. We first extend the
resulting partial isometry without choosing a vector on either
normal line.

\begin{lemma}[The determinant supplied by the ordered frames]
\leavevmode\label{lem:flagdeterminant534}
The ordered frames and the rational isometry \(\gamma\) define
an isometry of determinant lines
\(d_\gamma:\det P_s\longrightarrow\det P_t\). At each marked
place it identifies the local determinant lattices and reduces
in the marked bases to \(\det i(\gamma)_v\).
There is a unique rational isometry \(j_\gamma:P_s\to P_t\) of
determinant \(d_\gamma\) extending
\(\pi_0\pi_1^{-1}\bar\gamma\) on
\(H_s=\bar\gamma^{-1}\pi_1(A/I)\). Write \(H_t=\pi_0(A/I)\).
\end{lemma}
\begin{proof}
For a frame \((e_1,\ldots,e_s,z)\) in the ambient space \(V\)
with isotropic span \(X\) and positive quotient \(P\), send
\(w_1\wedge\cdots\wedge w_k\) to the top exterior product of
an ordered list in \(V\) given by
\[
     e_1,\ldots,e_s,\quad\widetilde w_1,\ldots,\widetilde w_k,
             \quad z,\quad f_1,\ldots,f_s .
\]
Take lifts in \((X+Ez)^\perp\) and choose dual partners satisfying
\(h(e_i,f_j)=\delta_{ij}\), with \(h\) linear in its second
argument. Changing a lift adds a vector of \(X\). Changing a
partner adds a vector of \(X^\perp\), already spanned by the
preceding vectors when the \(w_i\) form a basis. This defines
a canonical linear isomorphism \(\kappa_Y:\det P\to\det V\).
An upper unitriangular change in the ordered frame preserves it,
since the dual partners change by the inverse adjoint on the
opposite graded space. In calculating the Hermitian determinant
one can take the dual partners hyperbolic and perpendicular to
the middle lifts. The ratio of the two determinant forms is then
\((-1)^s h(z,z)\).

Put \(d_\gamma=\kappa_{Y_t}^{-1}(\det\gamma)\kappa_{Y_s}\).
The two ratios just computed agree since the terminal norms
agree, so \(d_\gamma\) is an isometry. At a marked place
the isotropic vectors form a basis of a primitive summand and
the last vector has unit norm in the middle quotient. Complete
the indicated integral bases by dual partners. Reduction of
\(\gamma\) respects the isotropic filtration, has determinant
one on the ordered isotropic block and its dual, and takes the
source last vector to the target last vector modulo the isotropic
part. On the remaining block it acts by \(i(\gamma)_v\).
Taking the top exterior product proves both local claims.

The map on \(H_s\) is an isometry onto \(H_t\).
Use \(\det P_s=(\det H_s)\otimes H_s^\perp\) and the analogous
identity on the target. Dividing \(d_\gamma\) by the determinant
of the hyperplane map defines an isometry between the normal
lines. Its orthogonal direct sum with the hyperplane map is
the asserted \(j_\gamma\). Every isometric extension carries
normal line to normal line, where its determinant fixes its
only scalar. This proves uniqueness. The construction respects
the unipotent changes defining the lifted double quotient.
\end{proof}

Put \(R=(\Lambda\cap A)/(\Lambda\cap I)\), and denote its images
in \(L_s,L_t\) by \(R_s,R_t\) under the preceding hyperplane
identifications. At a finite prime \(\mathfrak p\) of \(E\) let
\(e_s,e_t\) be the local lengths of the respective quotients
\[
  L_s/(R_s\oplus(L_s\cap H_s^\perp)),\qquad
  L_t/(R_t\oplus(L_t\cap H_t^\perp)).
\]
Both are finite since their denominators span the whole spaces
over \(E\). Write \(m_{\mathfrak p}\) for the valuation of the
multiplier of \(d_\gamma:\det L_s\to\det P_t\) in a local basis
of \(\det L_t\). Define integral ideals and their normalized size by
\[
  v_{\mathfrak p}(\mathfrak c)=e_s+e_t+|m_{\mathfrak p}|,
       \qquad \mathfrak b_\gamma=N_{E/F}(\mathfrak c^2),
       \qquad h_\gamma=(N_{F/\QQ}\mathfrak b_\gamma)^{1/(2f)} .
\]
All these numbers are obtained from the original pair of
flags by taking indices and determinants of lattices. For
full-rank cells retain the smaller ideal and \(h_\gamma\)
of Lemma~\ref{lem:markedtransport533}.

\begin{lemma}[The hyperplane map in the integral marking]
\leavevmode\label{lem:hyperplanemark534}
The ideal \(\mathfrak c\) clears the denominators of \(j_\gamma\)
and \(j_\gamma^{-1}\) between the two column lattices.
For each \(v\mid Q_0\) outside \(\mathfrak b_\gamma\), its
reduction equals the actual period index \(i(\gamma)_v\).
\end{lemma}
\begin{proof}
Work first over the local discrete valuation ring at
\(\mathfrak p\). A basis of \(R\) gives bases of both \(R_s,R_t\).
Complete each on its own side by a generator of its integral
normal line. The determinants of the resulting column lists in
bases of \(L_s,L_t\) have valuations \(e_s,e_t\).
Thus the multiplier of \(j_\gamma\) on these two normal generators
has valuation \(m_{\mathfrak p}+e_s-e_t\). A uniformizer to
the power \(e_s\) sends all of \(L_s\) into the lattice spanned
by the source list, by the elementary divisor theorem. The
denominator exponent for \(j_\gamma\) is therefore at most
\[
        e_s+\max(0,e_t-m_{\mathfrak p}-e_s)
                          =\max(e_s,e_t-m_{\mathfrak p}) .
\]
For the inverse the bound is \(\max(e_t,e_s+m_{\mathfrak p})\).
Both are at most \(v_{\mathfrak p}(\mathfrak c)\).

If \(v\nmid\mathfrak b_\gamma\), the two orthogonal lists give
direct sums of the local column lattices at both conjugate primes.
At a marked \(v\) these column lattices are self-dual, since
the outer hyperbolic pairs and the unit-norm terminal line
split off from the self-dual row. The pairing on the \(R\)
summand after reduction is consequently perfect between the
two split factors. A member of \(R\) has a common representative
in \(\Lambda\cap A\). Its images in the two transported positive
quotients thus have the same class modulo the common reduced
isotropic space. After composing with \(\bar\gamma\), this
means that \(i(\gamma)_v\) agrees with \(j_\gamma\) on the
reduced hyperplane. The same argument in the dual split factor
shows that the map carries the orthogonal normal line to the
target normal line. By Lemma~\ref{lem:flagdeterminant534}
the two determinants agree, so their actions on that last
one-dimensional summand coincide as well.
\end{proof}

\begin{theorem}[The finite period comparison in both common ranks]
\leavevmode\label{thm:bothranks534}
In the equal-terminal-norm pairing the separation inequalities
of Theorem~\ref{thm:fullrankseparation533} hold for any two cells,
including a pair of different common ranks. Throughout the common
interval, for all cells with \(h_\gamma\le b_x^2\), each projected
array \(HZ^{(j)}\) of a single rational transport can be selected
integrally from their projected sum, for any number of cells.
If \(\mathfrak b_\gamma\) is prime to \(Q_0\) on this collection,
each finite row family recovers the one signed operator at each
distinct transport, with equality of normalized norms and of
minor ideals.
\end{theorem}
\begin{proof}
The two possible dimensions of \(A/I\) are \(k,k-1\) by
Theorem~\ref{thm:cross132}. At either rank \(j_\gamma\) is a
rational isometry of the same two totally positive spaces.
Its marked reduction outside \(\mathfrak b_\gamma\) is its
actual index by the two reduction lemmas. For
\(u=j_{\gamma_2}^{-1}j_{\gamma_1}\), the extended ideal of
\(\mathfrak b_{\gamma_1}\mathfrak b_{\gamma_2}\) clears both
\(u\) and \(u^{-1}\). For a hyperplane contribution this follows
from \(\mathfrak b_\gamma\mathcal O_E=(\mathfrak c\bar{\mathfrak c})^2\)
and Lemma~\ref{lem:hyperplanemark534}. For a full-rank contribution
it follows from Lemma~\ref{lem:markedtransport533}.
These are exactly the ingredients in the trace estimate of
Theorem~\ref{thm:fullrankseparation533}. It gives
\(Q_0\le4k(h_{\gamma_1}h_{\gamma_2})^4\) for a collision
of different maps and \(Q_0^2\le4k(h_{\gamma_1}h_{\gamma_2})^6\)
for equality of their indices. The stated cutoff excludes a
collision. Selecting the disjoint sets of rows and, in the
prime-to-\(Q_0\) case, undoing the character units proves all
remaining assertions as in Corollary~\ref{cor:transportarray533}.
\end{proof}

\begin{proposition}[The signed trace retained in a hyperplane orbit]
\leavevmode\label{prop:hyperplanetrace534}
Take an absolutely convergent cell of identical words and
common rank \(k-1\). Let \(U_j\) be a unitary identification
of the two column types induced by \(j=j_\gamma\). If
\(Q_0^2>4k h_\gamma^6\), then
\[
       \operatorname{Tr}(U_j^*\Pi_{M_i}B_\gamma)
          =\operatorname{Tr}(U_j^*B_\gamma)
          =\sum_{\mu}d_\mu c_\mu .
\]
Here \(c_\mu\) is the scalar of \(U_j^*B_\gamma\) on the
irreducible branch \(\mu\) of the common compact hyperplane
group, and \(d_\mu\) is its dimension. Thus any nonzero value
of this signed scalar survives the entire arithmetic orbit
and the finite projection. Its squared modulus divided by
\(\dim V_s\) bounds \(\|\Pi_{M_i}B_\gamma\|_{\rm HS}^2\) from
below. This holds uniformly for \(h_\gamma\le b_x^5\).
\end{proposition}
\begin{proof}
For \((a,b)\in M_i\) put \(u=j^{-1}ajb^{-1}\). The extended
ideal of \(\mathfrak b_\gamma\) clears both \(u,u^{-1}\), since
it is contained in \(\mathfrak c^2\) and \(a,b\) are integral.
At the good marked places \(u\equiv1\). The squared congruence
trace argument of Theorem~\ref{thm:orbitnoncancellation533}
gives \(u=1\) under the stated inequality. Consequently \(U_j\)
is fixed by the left and right action of \(M_i\). Taking its
inner product with the orthogonal average gives the first
identity and the squared-norm inequality. By
Theorem~\ref{thm:common133} the original orbital operator
intertwines the two actions on \(A/I\). On identifying the
endpoints by \(U_j\), the restricted representation has
multiplicity one at each branch by the compact unitary
branching rule at each archimedean place. Schur's lemma gives
the indicated scalars and their trace formula.
Theorems~\ref{thm:orbitaverage532} and~\ref{thm:doubleorbit532}
give the positive multiplicity at the base residue and integral
recovery of that operator on the orbit. They apply at the same
terminal primes as the strict vanishing in
Corollary~\ref{cor:whittakerterminal527}.
\end{proof}

\begin{proposition}[Distinct transverse partners with the same transport]
\leavevmode\label{prop:repeatedtransport534}
Over each of the period fields there are standard integral
Hermitian rows with infinitely many transverse partners satisfying
the marked congruences, all with \(j_\gamma=1\), \(i(\gamma)=1\)
and \(h_\gamma=1\). Their Gram invariants are nonzero and
pairwise distinct. The two rational unipotent groups have trivial
intersection in each pair. The ambient isometries can also be
taken congruent to the identity to any fixed precision at the
same terminal prime \(T\).
\end{proposition}
\begin{proof}
Use the orthogonal sum of \(s\ge1\) standard hyperbolic planes
with bases \(e_a,f_a\) and a middle space with orthogonal basis
\(b_1,\ldots,b_k,z\) of norms \(1,\ldots,1,-1\). Take the
standard integer lattice and use the same ordered frame at both
ends. Choose \(0\ne t\in Q_0\mathcal O_E\) with \(\bar t=-t\),
for example \(Q_0\) times the difference of an algebraic integer
and its distinct conjugate. Replace \(t\) by successive positive
integer multiples. On each outer plane send \(e_a\) to
\(e_a+t f_a\) and fix \(f_a\). On the middle space put \(v=z+b_1\)
and send \(w\) to \(w+t v h(v,w)\). The vector \(v\) is isotropic.
Expanding the Hermitian form gives an isometry on each block
since \(t+\bar t=0\). The inverse uses \(-t\), so these maps are
integral row isometries, all congruent to the identity at the
two marked primes and all of determinant one. Requiring
\(t\in Q_0 T^e\mathcal O_E\) gives any prescribed terminal
precision \(e\) as well.

The two outer spans are transverse with cross pairing \(tI_s\).
Their sum contains all the outer planes. The two last vectors
are \(z,z-t(z+b_1)\), so their constraints leave precisely
\(\langle b_2,\ldots,b_k\rangle\) in the common perpendicular.
The map on this hyperplane is the identity. Since the endpoint
frames coincide and the ambient determinant is one, its prescribed
determinant on the full column is one. Thus \(j_\gamma=1\).
The integral hyperplane and its normal line split the columns
in their given bases. The two lattice indices and the determinant
valuations are zero, giving \(h_\gamma=1\).

The leading and full cross determinants are \(t^s\) and
\(t^s(-1+t)\). Formula~\eqref{eq:cross132} gives the invariant
\((t\bar t)^{s+1}\). Its distinct values for the chosen multiples
separate the rational partners, since the unipotent changes
preserve both exterior pairings. Finally an element in both
unipotent groups acts by upper unitriangular matrices on the two
outer spans. Since their cross matrix is \(tI_s\), these two
matrices are inverse adjoints. They must both be the identity,
being simultaneously lower and upper unitriangular in the
respective dual bases. Their fixed outer sum is nondegenerate.
On its orthogonal complement the action equals its action on
the first middle quotient, which is trivial by definition of
the unipotent group. This proves the intersection claim.
\end{proof}

Proposition~\ref{prop:repeatedtransport534} distinguishes the
size of a rational transport from the number of partner flags
above it, even when every determinant and lattice index used
in its height is a unit. In a prescribed period the remaining
sum at such an index contains the real orbital integrals with
their phases. The definition of the arithmetic height gives
no bound for these scalar sums.


\subsection{The two terminal frequencies in the Gram bound}
\label{sec:mixedgram535}

Retain the actual integral rows and the two terminal frequencies
\(\beta\in\{1,\varpi\}\), where
\(\varpi\bar\varpi=T\) and \(\varpi\equiv1\bmod (g_1g_2)^2\)
by Theorem~\ref{thm:common111}. On a row with frequency \(\beta\)
write \(z\) for the terminal vector of norm \(-\beta\bar\beta\),
and \(X\) for the initial isotropic space. Its middle lattice
\[
    J=(\Lambda\cap X^\perp)/(\Lambda\cap X)
\]
is self-dual and \(\bar z\in J\) is primitive, by the integral
frame conditions of Lemma~\ref{lem:coordinates225}. In this
section \(\bar z\) denotes the vector in the quotient. Bars on
elements of \(E\) denote complex conjugation. Write
\(L=J\cap (E\bar z)^\perp\), and let \(L^\sharp\) be its
Hermitian dual in its positive rational space \(P\).

\begin{lemma}[A common integral determinant at the terminal prime]
\leavevmode\label{lem:completedcolumn535}
There is a self-dual lattice \(\widehat L\) in each \(P\) given
by the following local prescriptions. Use \(L\) everywhere for
\(\beta=1\). For \(\beta=\varpi\) use \(L^\sharp\) at the primes
of \(E\) dividing \(\varpi\), and \(L\) at all other primes.
It contains \(L\) and is preserved by the integral frame
stabilizer. Use \(z/\bar\beta\) as last vector solely in the
determinant-line map of Lemma~\ref{lem:flagdeterminant534}.
For every cell with source and target frequencies independently
in \(\{1,\varpi\}\), the resulting map
\[
      d_\gamma^\circ:\det P_s\longrightarrow\det P_t
\]
is an isometry carrying \(\det\widehat L_s\) onto
\(\det\widehat L_t\). At both auxiliary primes its matrix
is still \(\det i(\gamma)_v\) after reduction.
\end{lemma}
\begin{proof}
Self-duality of \(J\) and primitivity of \(\bar z\) make
\(y\mapsto h(\bar z,y)\) surjective onto the local integer
ring at every place. This also follows by extending an integral
functional on the primitive line and using the perfect pairing
in the other variable. The kernel is \(L\). In a basis of this
kernel completed by a vector with pairing one, the last
coordinate of \(\bar z\) is \(-\beta\bar\beta\).
Hence the determinant of \(L\oplus\mathcal O_E\bar z\) in \(J\)
generates the ideal \((\beta\bar\beta)\) at each prime. At
primes away from this ideal the two summands are orthogonal and
self-dual. For \(\beta=\varpi\) each \(v\mid T\) is split, and
the index exponent is one in each factor. The determinant of
the Hermitian pairing between the two sublattices has valuation
two, by taking one determinant from each side of the perfect
pairing of \(J\). The negative line supplies valuation one.
The pairing between the positive factors therefore has determinant
of valuation one. In particular \(L^\sharp/L\) has length one
in each split factor. Enlarging just the \(\varpi\)-factor
to the lattice dual to the opposite factor makes the pair
self-dual. Taking the intersection of these finitely modified
local lattices inside \(P\) gives \(\widehat L\) with the
stated localizations. This is precisely the parent of
Lemma~\ref{lem:termparent221}, where the terminal hyperplane is
imposed in the factor dividing \(\varpi\). Every frame
stabilizer preserves each prescribed lattice.

To check the determinant away from \(T\) as well, write
\(\Lambda\cap X=\bigoplus_a\mathfrak h_a e_a\) in the prescribed
ordered coordinates. The fractional ideals \(\mathfrak h_a\)
are fixed by the interval data and are the same in both rows,
as in
\srcloc{Formalism/SharpRestriction/CuspSpanNormBounds/FlagEEqHyperbolic}{93}.
Let \(\mathfrak h=\prod_a\mathfrak h_a\). The opposite graded
lattice is the conjugate dual of \(\Lambda\cap X\), since this
is a saturated summand locally in the self-dual row. Its
determinant in the paired dual frame is
\(\bar{\mathfrak h}^{-1}\). Consequently the exterior-product
map \(\kappa_{Y/\bar\beta}\) formed with \(z/\bar\beta\) satisfies
\[
   \kappa_{Y/\bar\beta}(\det L)
        =(\bar{\mathfrak h}/\mathfrak h)(\beta)\det\Lambda .
\]
Indeed the isotropic and dual basis factors contribute
\(\bar{\mathfrak h}/\mathfrak h\) relative to the ambient
lattice determinant, and the middle factor contributes
\((\beta\bar\beta)/(\bar\beta)=(\beta)\), by the kernel
calculation above. Enlargement at the primes dividing
\(\varpi\) cancels exactly the last ideal for \(\beta=\varpi\).
Thus \(\kappa_{Y/\bar\beta}(\det\widehat L)=
(\bar{\mathfrak h}/\mathfrak h)\det\Lambda\) for each frequency.
The rational isometry \(\gamma\) carries the ambient lattice
determinants to each other. Conjugating its determinant by
these two maps proves the integral assertion. Both rescaled
vectors have norm \(-1\), so the determinant map is an isometry
by the determinant-form calculation in
Lemma~\ref{lem:flagdeterminant534}. Since each \(\bar\beta\) is
one modulo \(Q_0=g_1g_2\), its marked reduction is the
prescribed one.
\end{proof}

Let \(R=(\Lambda\cap A)/(\Lambda\cap I)\) be the common positive
lattice of Section~\ref{sec:corankone534}, of rank \(k-1\).
Use the rational hyperplane isometry on its two images and
complete it by the determinant \(d_\gamma^\circ\) to an isometry
\(j_\gamma:P_s\to P_t\), as in Lemma~\ref{lem:flagdeterminant534}.
Let \(\mathfrak t_R\) be the integral ideal whose valuation at
\(\mathfrak p\) is the length of \((R^\sharp/R)_{\mathfrak p}\),
and set
\[
  \mathfrak b_R=N_{E/F}\mathfrak t_R,\qquad
  h_R=(N_{F/\QQ}\mathfrak b_R)^{1/(2f)}
         = [R^\sharp:R]^{1/(2f)} .
\]
The bracket is the order of the finite quotient.

\begin{theorem}[Comparison bounded by the discriminant of the common lattice]
\leavevmode\label{thm:mixedtransport535}
For either terminal frequency on each side, \(\mathfrak t_R\)
clears the denominators of \(j_\gamma\) and \(j_\gamma^{-1}\)
between the completed column lattices. At marked places outside
\(\mathfrak b_R\) its residue is \(i(\gamma)_v\).
In any fixed pair of rows, cells of common rank \(k-1\) with
\(h_R\le b_x^2\) and distinct \(j_\gamma\)'s occupy disjoint
entries of \(H\) in every row family. Their projected arrays
therefore have additive squared norms and are recoverable by
integral entry selection by transport, including at the common
terminal prime. If each \(\mathfrak b_R\) is prime to \(Q_0\),
the single signed raw operator for each transport is recoverable
in every row family.
\end{theorem}
\begin{proof}
At either endpoint the copy of \(R\) is contained in \(L\)
and hence in \(\widehat L\). Orthogonal projection from
\(\widehat L\) to the rational hyperplane of \(R\) takes
values in \(R^\sharp\), since it preserves all pairings against
\(R\) and \(\widehat L\) is integral. Its kernel is the
intersection of \(\widehat L\) with the normal line. Consequently
\[
    \widehat L\big/\big(R+(\widehat L\cap R_E^\perp)\big)
                    \lhook\joinrel\longrightarrow R^\sharp/R ,
\]
where \(R_E\) denotes the corresponding hyperplane. Thus both
local determinant exponents \(e_s,e_t\) in the proof of
Lemma~\ref{lem:hyperplanemark534}, calculated in the completed
lattices, are at most \(v_{\mathfrak p}(\mathfrak t_R)\).
Lemma~\ref{lem:completedcolumn535} makes the determinant
exponent \(m_{\mathfrak p}\) zero everywhere. The two denominator
bounds in that proof now reduce to \(\max(e_s,e_t)\) in each
direction. At a good marked place \(R\) splits as a self-dual
orthogonal summand at both ends, and the common representatives
in \(\Lambda\cap A\) match its two reductions. The determinant
agrees with the actual marked determinant. The hyperplane and
normal-line calculation in Lemma~\ref{lem:hyperplanemark534}
therefore proves the residue assertion.

The ideal \(\mathfrak t_R\) is invariant under conjugation.
Indeed the two conjugate local pairing matrices are conjugate
transposes in paired bases, and the determinant valuations give
the indicated lengths. In particular
\(\mathfrak b_R\mathcal O_E=\mathfrak t_R^2\). For a pair
of cells the product of their two ideals of \(F\) clears
the relative isometry and its inverse on the same completed
lattice. The trace and ideal-norm calculation of
Theorem~\ref{thm:fullrankseparation533} now applies with
\(h_{\gamma_l}=h_{R_l}\). A collision between unequal maps
would give \(Q_0\le4k(h_{R_1}h_{R_2})^4\), impossible at the
specified cutoff. The assertions on arbitrary arrays follow
by selecting the disjoint output supports, as in
Corollary~\ref{cor:transportarray533}. Only the lattices used
to bound the rational traces were enlarged. The matrix \(H\)
and all selected entries come from the original marking.
\end{proof}

\begin{corollary}[An estimate in terms of the actual cross Gram determinant]
\leavevmode\label{cor:actualgram535}
For a supported transverse pair with nonzero \(\mathcal G\),
write \(d=s+1\), and let \(I_\Lambda=[B:B^+]\) denote the
saturation index of Proposition~\ref{prop:index219}. With
\(\mathfrak h=\prod_a\mathfrak h_a\) from
Lemma~\ref{lem:completedcolumn535}, one has
\[
     h_R^f=\frac{(N_{E/\QQ}\mathfrak h)^2
                    |N_{F/\QQ}(\mathcal G)|}{Q_0^{2df} I_\Lambda}.
\]
In particular all such cells satisfying
\[
              |N_{F/\QQ}(\mathcal G/Q_0^{2d})|\le b_x^{2f}
\]
obey Theorem~\ref{thm:mixedtransport535} on the prescribed
fields and intervals. Neither the number of rational partners
nor the size of the terminal prime enters this test.
For an absolutely convergent identical-word cell at either
pair of frequencies, the signed trace identity of
Proposition~\ref{prop:hyperplanetrace534} holds whenever
\(|N_{F/\QQ}(\mathcal G/Q_0^{2d})|\le b_x^{5f}\).
\end{corollary}
\begin{proof}
The transverse spans give \(I=0\). The integral lattice \(R\)
is precisely \(C_\Lambda\) of Proposition~\ref{prop:index219}.
For a nondegenerate lattice \(M\) put
\(\delta(M)=[M^\sharp:M]^{1/2}\), so \(\delta(R)=h_R^f\).
Put \(U_l=\Lambda\cap Y_l\) and \(U=U_0\oplus U_1\).
Primitivity of the last vector in each middle quotient gives
\(U_l=(\Lambda\cap X_l)\oplus\mathcal O_Ez_l\), by projecting
a vector in \(U_l\) onto \(E\bar z_l\). The determinant module
of \(U\) in the two concatenated rational frames is therefore
\(\mathfrak h^2\). Their Hermitian Gram determinant is
\((-1)^d\mathcal G\) by \eqref{eq:cross132}, with the two
actual last-vector norms. Taking local bases in the determinant
module and multiplying the resulting local indices gives
\(\delta(U)=(N_{E/\QQ}\mathfrak h)^2|N_{F/\QQ}\mathcal G|\).
Here the norm of a fractional ideal is taken multiplicatively.
In primitive integral frame coordinates \(\mathfrak h=\mathcal O_E\).

At each marked place the two frames have the same reduction
up to an integral upper unitriangular change. Dividing each of
the \(d\) aligned differences by the local uniformizer enlarges
\(U\) into the row by index \(q_v^{2d}\), counting both split
factors. The columns are independent over the field, so this
index follows from the diagonal scaling in their basis.
Elsewhere take \(B^+=U\). The resulting intermediate lattice
has index \([B^+:U]=Q_0^{2df}\). Its saturation \(B\) in the
row has perpendicular \(R\). Projection of the self-dual row
onto the two rational summands has images \(B^\sharp,R^\sharp\)
because every integral functional extends locally from a saturated
summand to the row. Modulo \(B\oplus R\) either projection gives
an isomorphism onto the associated discriminant module, since an
integral projection can be subtracted to land in the other lattice.
Thus \(\delta(R)=\delta(B)=\delta(U)/[B:U]\), proving the formula.
The prescribed \(\mathfrak h_a\) are nonpositive powers of the
extended different of \(F/\QQ\), so \(N_{E/\QQ}\mathfrak h\le1\).
Since \(I_\Lambda\ge1\), the two displayed cutoffs give
\(h_R\le b_x^2\) and \(h_R\le b_x^5\), respectively.
For the trace assertion the arithmetic stabilizers preserve
the completed lattices by Lemma~\ref{lem:completedcolumn535}.
For \((a,b)\in M_i\), the denominators of \(j^{-1}ajb^{-1}\)
and its inverse are cleared by \(\mathfrak b_R\), since its
extension is \(\mathfrak t_R^2\). Reduction at a good marked
place is the identity. The squared congruence in
Theorem~\ref{thm:orbitnoncancellation533} forces \(aj=jb\)
because \(Q_0^2>4k h_R^6\). The rest of the proof of
Proposition~\ref{prop:hyperplanetrace534} applies to this
normalized \(j\). The original real functions and frequencies
enter its signed trace.
\end{proof}

\begin{proposition}[The retained columns and a test independent of the Gram size]
\leavevmode\label{prop:retainedtest535}
On the retained norm-\(T\) columns of Theorem~\ref{thm:columns235},
each supported relevant transverse cell of nonzero Gram determinant
satisfies \(h_R>b_x^5\) for all sufficiently large \(x\).

For any two hyperplane cells in fixed rows, let
\(\mathfrak a_l^+\) be the ideal of integral denominators of
\(j_l:\widehat L_s\to P_t\) relative to \(\widehat L_t\), and
\(\mathfrak a_l^-\) the analogous ideal for its inverse. Put
\(\mathfrak e_l=N_{E/F}(\mathfrak a_l^+\mathfrak a_l^-)\).
Let \(W\) be the product of \(N_{F/\QQ}v\) over the marked
places where both maps are integral isomorphisms with their
prescribed residue matrices. For \(j_1\ne j_2\), a collision
in a row of \(H\) requires
\[
       W\le(4k)^f N_{F/\QQ}(\mathfrak e_1\mathfrak e_2).
\]
Equality of the two residue indices requires the same bound
with \(W^2\) on the left. This test allows arbitrary values
of \([R^\sharp:R]\).
\end{proposition}
\begin{proof}
The image of \(R\) in the original positive lattice is contained
in its saturation, a primitive hyperplane. Passing to that
saturation divides \(\delta(R)\) by the finite index. With
\(D_E=|\operatorname{disc}E|^{1/(2f)}\), the hyperplane bound
in Theorem~\ref{thm:columns235} gives
\[
      h_R^f=\delta(R)>
        T^f\left(\frac{kD_E}{2\pi\exp(A_0+4)}\right)^f
\]
for fixed selection parameter \(A_0>0\).
Here \(D_E\ge1\), \(k\to\infty\), and
\(\log T\ge(16-o(1))m x^\varepsilon\) with \(m\to\infty\)
by Lemma~\ref{lem:size97}. This proves the first assertion.

For the second assertion, the extended ideal of
\(\mathfrak e_1\mathfrak e_2\) clears the relative map
\(u=j_2^{-1}j_1\) and its inverse from the definition of the
two denominator ideals. For unequal maps the trace
\(\operatorname{tr}((1-u)(1-u^{-1}))\) is totally positive
and has archimedean absolute values at most \(4k\). The norm of
the integral ideal obtained by multiplying its principal ideal
by \(\mathfrak e_1\mathfrak e_2\) is at most the right side.
At each place counted in \(W\) the trace vanishes after reduction
if the two indices collide. If the indices are equal, it vanishes
to order at least two by the product expression. Multiplication
by an integral ideal preserves each of these divisibilities.
This gives the two necessary inequalities.
\end{proof}

The Gram cutoff applies to the unrestricted rows and excludes the
nonzero-Gram cells on the retained norm-\(T\) columns just described.
On those columns the sharper test requires estimates for the
actual denominators and the agreement at the marked places.
Even when different transports separate, the signed sum over
partners inducing the same isometry still requires control.
Proposition~\ref{prop:repeatedtransport534} gives distinct
examples in a common fibre. These estimates and the residual
cohomology estimate of Section~\ref{sec:missing} are required
for an unconditional conclusion.


\subsection{A signed series on each retained diagonal row}
\label{sec:diagonalseries536}

Fix an original integral row \(\Lambda\) on the period field
and pair one word and its marking with themselves. Let \(s\ge1\)
be the length of its isotropic flag, \(d=s+1\), \(P\) its
positive column of dimension \(k\ge2\), and \(T_0\in\{1,T\}\)
the absolute value of the terminal norm. Let \(N_\Lambda\) be
the integral unipotent group of the flag. The following family
exists on every such diagonal row, including the retained
norm-\(T\) rows of Theorem~\ref{thm:columns235}.

\begin{theorem}[Infinitely many terms above a prescribed rational hyperplane]
\leavevmode\label{thm:everydiagonal536}
Choose any rational hyperplane \(H\subset P\), and fix a
precision at the terminal prime. There is a cyclic group of
integral row isometries \(\{\gamma^n:n\in\ZZ\}\) such that
all its nonidentity terms are supported, relevant, transverse,
and distinct in the lifted double quotient. At the chosen
precision they act trivially at \(T\). At both auxiliary primes
their actual period index is the identity. Each has common
positive hyperplane \(H\) and rational transport \(j_{\gamma^n}=1\).
In particular the two actual denominator ideals in
Proposition~\ref{prop:retainedtest535} are the unit ideal.
Their nonzero Gram invariants satisfy
\[
        \tau(\mathcal G_n)\ge |n|^{2d}\tau(\mathcal G_1)
                \qquad(\tau:F\hookrightarrow\RR,\ n\ne0).
\]
Their common positive lattice \(R\) is independent of \(n\).
On a retained norm-\(T\) row it has \(h_R>b_x^5\) at all
sufficiently large \(x\).
\end{theorem}
\begin{proof}
Choose rational hyperbolic partners \(f_a\) for the ordered
outer vectors \(e_a\), perpendicular to the terminal vector
and to the rational model of \(P\) in the middle. These exist
by the Hermitian Witt decomposition. Choose \(w\ne0\) in the
one-dimensional perpendicular of \(H\) in \(P\). By scaling
we may assume \(D=h(w,w)\in\mathcal O_F\), which is totally
positive. Write \(e=z/\beta\) with \(\beta\bar\beta=T_0\);
then \(h(e,e)=-1\). Use the identity on \(H\). On each outer
plane take \(e_a\mapsto e_a+t f_a,\ f_a\mapsto f_a\), where
\(t\ne0\) is an integral multiple of a nonzero \(t_0\in E\)
with \(\bar t_0=-t_0\). Such an element is the difference of
an integer of \(E\) and its distinct conjugate.

First suppose \(D\) is a square, say \(D=c^2\) in \(F\).
On the plane \(Ew+Ee\) put \(v=w+ce\) and use
\(y\mapsto y+t v h(v,y)\). This is an isometry of determinant
one since \(h(v,v)=0\) and \(t+\bar t=0\). Its \(n\)-th
power replaces \(t\) by \(nt\). If \(D\) is a nonsquare, put
\(K=F(\sqrt D)\). This is a totally real quadratic extension.
Modulo torsion the units of \(K\) have rank \(2f-1\) and
those of \(F\) have rank \(f-1\), by Dirichlet's unit theorem.
The norm image contains the squares of the units of \(F\).
Thus norm-one units have rank \(f\), in particular positive.
The order \(\mathcal O_F[\sqrt D]\) has finite additive index
in \(\mathcal O_K\), so for any rational integer \(L\ge1\)
one can choose an infinite-order norm-one unit
\(\epsilon=a+b\sqrt D\) with \(a-1,b\in L\mathcal O_F\).
Indeed choose an integer \(c_0\) with
\(c_0\mathcal O_K\subseteq\mathcal O_F[\sqrt D]\) and raise
an infinite-order norm-one unit to the order of its image in
\((\mathcal O_K/Lc_0\mathcal O_K)^\times\). Squaring also
makes the unit totally positive. In the basis \(w,e\) take
the matrix
\[
             \begin{pmatrix} a&b\\Db&a\end{pmatrix}.
\]
It preserves \(\operatorname{diag}(D,-1)\) and has determinant
\(a^2-Db^2=1\). Its \(n\)-th power has parameters \(a_n,b_n\)
defined by \(\epsilon^n=a_n+b_n\sqrt D\).

The rational basis used in these blocks spans a reference
integer lattice \(\Lambda_0\). There is a positive rational
integer \(h\) for which
\(h\Lambda_0\subseteq\Lambda\subseteq h^{-1}\Lambda_0\).
Choose \(L\) divisible by \(h^2Q_0T^l\), where \(l\) gives
the required precision, increasing \(l\) to include the original
row level if needed. Clear the finitely many denominators in the
square-case matrix and in the outer shears by multiplying \(t\).
This makes all blocks and their inverses the identity modulo
\(L\Lambda_0\). On \(\Lambda\) the difference from the identity
then has image in \(Q_0T^l\Lambda\). Any further prescribed
finite principal level can be imposed in the same choice of
\(L\). The product \(\gamma\) therefore preserves the actual
row and its marked reductions. The transported flag inherits
the primitive frame and fractional-ideal conditions by this
isometry. Its admissibility at the two auxiliary primes follows
from the congruence. All powers give supported terms there.

For \(n\ne0\) the matrix of cross pairings of the two outer
frames is \(ntI_s\), so their spans are transverse. On the
middle plane \(z,\gamma^nz\) are independent. In the square
case the transvection has a nonzero \(w\)-component on \(e\).
In the nonsquare case \(b_n\ne0\), since a norm-one element of
\(F\) in \(K\) is \(1\) or \(-1\). The common perpendicular
of both full frames is precisely the chosen \(H\). The action
on it is the identity and the ambient determinant is one, so
Lemma~\ref{lem:flagdeterminant534} gives \(j_{\gamma^n}=1\).
In particular \(R=\Lambda\cap H\) in this rational realization
is constant along the family.
The first cross pairing \(nt\) is invariant under \(N_\Lambda\)
on both sides, since \(N_\Lambda\) fixes \(e_1\). It distinguishes
all nonzero integers \(n\). A common unipotent acts on the two
outer spans by inverse adjoints in their displayed bases, since
their cross matrix is a scalar identity. Both matrices are upper
unitriangular, hence both are identities. The outer sum is
nondegenerate and the first middle quotient is fixed pointwise,
so the common unipotent is trivial. The same calculation for
strictly upper triangular Lie algebra elements gives zero.
The two defining linear characters agree on this intersection,
which proves relevance.

With \(D_s=(nt)^s\) the Gram invariant is
\(|D_s|^2\bigl(|h(z,\gamma^n z)|^2-T_0^2\bigr)\), where
\(|u|^2=u\bar u\) over \(F\). For the square case the second
factor is \(T_0^2 n^2(t\bar t)D^2\), by \(h(v,e)=-c\).
In the nonsquare case it is \(T_0^2(a_n^2-1)\).
At a real place \(a_n=\cosh(nu_\tau)\) with
\(u_\tau=\log|\tau'(\epsilon)|\ne0\) for either extension
\(\tau'\) to \(K\). Convexity on the positive axis gives
\(|\sinh(nu_\tau)|\ge |n|\,|\sinh u_\tau|\).
This proves positivity and the growth estimate. The retained
lower bound follows from Proposition~\ref{prop:retainedtest535}.
\end{proof}

\begin{corollary}[Saturation growth with constant transport]
\leavevmode\label{cor:saturationseries536}
The saturation indices of these terms satisfy
\[
        I_{\Lambda,n}/I_{\Lambda,1}
          =|N_{F/\QQ}(\mathcal G_n/\mathcal G_1)|
          \ge |n|^{2df}\qquad(n\ne0).
\]
\end{corollary}
\begin{proof}
In Corollary~\ref{cor:actualgram535} the ideals of the original
frames, the two marked primes, and the discriminant of the
common lattice are fixed along this family. Divide the two
equalities and multiply the bounds at the \(f\) real places.
\end{proof}

\begin{theorem}[A weighted tail for the ordinary signed operators]
\leavevmode\label{thm:diagonaltail536}
For the ordinary archimedean section \(W\) of
Proposition~\ref{prop:weight163} use the same test at both ends.
Write \(B_n\) for the corresponding orbital operator of the
term \(\gamma^n\), including its finite outer character and
the row normalization. Let \(r=\dim_E V\),
\(L_0=8(f+1)(r+1)^3\) and \(a=(L_0-4fr^2)/2>0\).
For this fixed field, row, test and cyclic family there is a
finite \(C_{\Lambda,W,\gamma}\) such that for \(M\ge1\)
\[
             \sum_{|n|>M}\|B_n\|_{\rm HS}
                   \le C_{\Lambda,W,\gamma} M^{-af}.
\]
Let \(P_\mu\) be the orthogonal projections onto the irreducible
branches of the compact group of \(H\) in the column type.
For \(n>0\) write \(B_n=\sum_\mu c_{\mu,n}P_\mu\).
The entire signed subseries on these terms at the identity
residue is the Hermitian operator
\[
             2\sum_\mu\left(\sum_{n\ge1}
                               \operatorname{Re}c_{\mu,n}\right)P_\mu .
\]
\end{theorem}
\begin{proof}
Let \(G_\infty\) be the ambient real group and \(B_\infty\)
the stabilizer of the decorated frame with unipotent radical
and compact positive quotient. In signed orthonormal coordinates
put \(w(g_\tau)=\|g_\tau\|_{\rm op}=\|g_\tau^{-1}\|_{\rm op}\).
Use the height
\(V(y)=\prod_\tau\inf_{b\in B_\tau}w(by_\tau)\ge1\) of
Proposition~\ref{prop:weight163}. Every member of \(B_\tau\)
fixes \(e_{1,\tau}\). Applying the operator norm bound to its
inverse after left translation therefore bounds
\(\|y_\tau^{-1}e_{1,\tau}\|\) by the local height times
\(\|e_{1,\tau}\|\). Apply this also to \(\gamma^{-n}y\).
The Hermitian pairing is bounded by the two Euclidean norms.
Since \(h(e_1,\gamma^ne_1)=nt\), multiplication gives
\[
      |n|^f\le D_0 V(y)V(\gamma^{-n}y),\qquad
      D_0=\prod_\tau\|e_{1,\tau}\|^2/|\widetilde\tau(t)|.
\]
Here either complex extension \(\widetilde\tau\) of \(\tau\)
gives the same absolute value.

Put \(F_a(y)=V(y)^a\|W(y)\|\). It is left invariant under
\(\Gamma_N=N_\Lambda\) by unitary equivariance. Its integral
on \(\Gamma_N\backslash G_\infty\) is finite by
Proposition~\ref{prop:weight163}. Indeed, the integral points
of the rational unipotent radical form a uniform lattice in
its real points (the unipotent lattice theorem applied after
restriction of scalars to \(\QQ\)). The row lattice and its
finite levels give a commensurable subgroup. The remaining
factor of \(B_\infty\) is compact. Let
\(\Gamma=\operatorname{Aut}(\Lambda,h)\).
The function \(S_a(y)=\sum_{\delta\in\Gamma_N\backslash\Gamma}
F_a(\delta y)\) is bounded uniformly in \(y\) for these fixed
data. Here are the counting details. For \(V(\delta y)\le R_0\)
choose a representative on the left by \(B_\infty\) with
norm at most \(2R_0\) in each factor. Compactness of
\(\Gamma_N\backslash B_\infty\) permits replacing this multiplier
by an element of \(\Gamma_N\) up to a fixed compact multiplier.
Relative to one chosen representative in the ball, the others
are distinct rational integral row isometries of norm
\(O_\Lambda(R_0^2)\) at each place. Their matrix entries lie
in a fixed fractional ideal in a rational basis, hence in a
discrete lattice of real dimension \(2f\). Counting its points
in the box gives \(O_\Lambda(R_0^{4fr^2})\) possibilities
for the \(r^2\) entries, uniformly in \(y\).
The ordinary section has \(\sup V^{L_0}\|W\|<\infty\).
Summing over dyadic shells converges since
\(L_0-a>4fr^2\).

Integrate \(F_a(y)S_a(y)\) on \(\Gamma_N\backslash G_\infty\).
Its integral is at most \(\|F_a\|_1\|S_a\|_\infty\).
Grouping the nonnegative summands into integral double cosets
unfolds the term of \(\delta\) to an integral of
\(F_a(y)F_a(\delta^{-1}y)\) over
\((\Gamma_N\cap\delta\Gamma_N\delta^{-1})\backslash G_\infty\).
Indeed the cosets in that double coset translate a fundamental
domain for \(\Gamma_N\) through one fundamental domain for the
intersection. Tonelli's theorem permits this regrouping.
For the distinct terms \(\delta=\gamma^n\), \(n\ne0\), the
intersection is trivial by Theorem~\ref{thm:everydiagonal536}.
The finite outer phase is one because the marking is identical
at both ends, the isometry is congruent to the identity at
both auxiliary primes, and the character is trivial on the
remaining integral differences. The volume scalar \(c_{\rm row}\)
is common. The rank-one operator \(W(y)W(\gamma^{-n}y)^*\)
has Hilbert--Schmidt norm the product of the two vector norms.
The height inequality therefore gives
\[
       \sum_{n\ne0}|n|^{af}\|B_n\|_{\rm HS}
          \le |c_{\rm row}|D_0^a\|F_a\|_1\|S_a\|_\infty .
\]
This proves the tail bound with the displayed finite constant.

Each orbital operator commutes with the compact group of \(H\)
by Theorem~\ref{thm:common133} and \(j_{\gamma^n}=1\). The
restriction to that group has multiplicity one, giving the
stated branch scalars. The integration domain is all of
\(G_\infty\) for \(n\ne0\). Taking the adjoint and changing
variable by \(\gamma^{-n}\) gives \(B_{-n}=B_n^*\), since the
row volume is real. Absolute convergence just proved permits
pairing inverse terms. All occur at the same identity residue,
which gives the signed formula.
\end{proof}

The bound controls the tail of this subseries also on the
retained rows of Proposition~\ref{prop:retainedtest535}. Its
constant includes the weighted norms of the prescribed ordinary
section and the lattice count for its row. An interval estimate
for the entire period still requires control of these constants
as the rows and fields vary and of the other terms at the same
residue. The signed branch sums above may include cancellation.


\subsection{Integral transport on a marked hyperplane}
\label{sec:parabolic537}

Fix two original period rows. For a cell with common positive
hyperplane, write \(j:P_s\to P_t\) for the isometry with prescribed
determinant in Theorem~\ref{thm:mixedtransport535}. Its two parent
lattices on the target space are \(M=j\widehat L_s\) and
\(N=\widehat L_t\). They are self-dual and have the same determinant
lattice. Let \(Y\subset P_t\) denote the image of the common
hyperplane and let \(J=Y^\perp\), a line. The two parents may induce
different extensions of their lattices in \(Y\) by lattices in \(J\).
These extensions determine the denominators of \(j\).

\begin{proposition}[Computing both denominator ideals]
\leavevmode\label{prop:gluing537}
Localize at a finite place of \(F\), keeping both factors at a split
place. For \(L=M,N\), put \(A_L=L\cap Y\), \(C_L=L\cap J\), and
\(X_L=\operatorname{pr}_Y L\). Then \(X_L=A_L^\sharp\), with
duality taken inside \(Y\). In each discrete valuation ring factor
choose a linear map \(\lambda_L:X_L\to J\) for which
\((x,\lambda_L x)\in L\), and extend it to the rational space.
The local factor of the forward denominator ideal is the set of
integral scalars \(u\) satisfying the three inclusions
\[
 \begin{split}
       uC_M&\subset C_N,\qquad uX_M\subset X_N,\\
       u(\lambda_M-\lambda_N)X_M&\subset C_N .
 \end{split}
\]
Interchanging \(M,N\) computes the inverse denominator.

Suppose \(A_M=A_N=A\) in all factors above this place of \(F\).
Then \(C_M=C_N=C\). If \(\mathcal O\) denotes the integer ring
in one factor and \(\pi\) its uniformizer, \(C^\sharp/C\)
there is cyclic of length \(d\). For \(d>0\) the two lattices
determine isomorphisms \(\phi_L:A^\sharp/A\to C^\sharp/C\).
Write multiplication by \(\xi\in(\mathcal O/\pi^d)^\times\)
for \(\phi_M\phi_N^{-1}\). The two denominator exponents are
both
\[
                d-\min\{d,v_\pi(\xi-1)\}.
\]
The capped valuation is taken on any lift. For \(d=0\) both
exponents vanish.
\end{proposition}
\begin{proof}
The module \(A_L\) is saturated in \(L\), so is a direct summand
locally. A functional on \(A_L\) extends to \(L\) and is represented
by an element of \(L\) by self-duality. Its representative in \(Y\)
is the orthogonal projection of that element. Conversely every
such projection pairs integrally with \(A_L\). This proves the
formula for \(X_L\), including the duality between the two factors
at a split place. The same argument gives
\(\operatorname{pr}_J L=C_L^\sharp\). Since \(X_L\) is free over
each discrete valuation ring we may lift a basis to \(L\) and
take \(\lambda_L\) from the second coordinates of the lifts. Thus \(L\)
consists of \((x,\lambda_L x+c)\) for \(x\in X_L,\ c\in C_L\).
Testing \(uM\subset N\) on the kernel and on these lifts gives
precisely the three inclusions, proving also their independence
of the lifts.

The determinant of \(L\) in the decomposition \(Y\oplus J\)
is \((\det X_L)\otimes C_L\), as is seen on the same lifts and a
generator of its kernel. Equal \(A_L\)'s give equal \(X_L\)'s.
Equality of the determinant lattices now gives equal \(C_L\)'s.
Modulo \(A\oplus C\), each \(L\) is a graph in
\((A^\sharp/A)\oplus(C^\sharp/C)\). Indeed both projections
are onto and both are injective on this graph, by the definitions
of \(A,C\) as intersections. The induced isomorphism is \(\phi_L\).
Since \(J\) is a line, the module \(C^\sharp/C\) is cyclic in
each factor. A scalar carries the first graph into the second
exactly when it annihilates \((\phi_M-\phi_N)(A^\sharp/A)\).
This is the image of multiplication by \(\xi-1\) on
\(\mathcal O/\pi^d\). The same image up to sign occurs in the
other direction. Its annihilator gives the asserted exponent.
\end{proof}

At each marked \(v\mid Q_0\) take the split factor used for the
finite general linear group. Write \(V_s,V_t\) for the two
reduced columns. The reductions of the integral common lattice
\(R\) have images in these spaces. For a collection of cells, say
explicitly that the images are fixed hyperplanes if there are
subspaces \(Y_s(v)\subset V_s,\ Y_t(v)\subset V_t\), each of
codimension one, that are these images for every member of the
collection. This asks for no nondegeneracy of their pairing
with the opposite split factor.

\begin{lemma}[The action on a reduced hyperplane and its quotient]
\leavevmode\label{lem:parabolic537}
Suppose \(j\) is an integral isomorphism at a marked \(v\) and
the images of \(R\) are hyperplanes \(Y_s(v),Y_t(v)\).
The maps \(i(\gamma)_v\) and \(\bar j\) agree on \(Y_s(v)\)
and on \(V_s/Y_s(v)\), carrying these spaces to their target
counterparts. Equivalently
\[
 i(\gamma)_v\bar j^{-1}\in
 D_v=\{1+A:\ A(V_t)\subset Y_t(v),\ A(Y_t(v))=0\}.
\]
For two such cells with fixed hyperplanes at \(v\), the relative
maps \(i_1 i_2^{-1}\) and \(\bar j_1\bar j_2^{-1}\) have the
same characteristic polynomial at \(v\).
\end{lemma}
\begin{proof}
A member of \(R\) is represented by a row vector in the common
orthogonal space. Its two reductions in the positive quotients
coincide under the actual index, because both quotient maps take
its class modulo the common reduced isotropic space. This is the
representative calculation in Lemma~\ref{lem:hyperplanemark534}.
It holds on the image of \(R\) without any assumption on its
pairing. The defining rational map \(j\) also identifies its
two images. They span the given hyperplanes after reduction.
The determinants of the two full maps agree by
Lemma~\ref{lem:completedcolumn535}. Taking the determinant on
the hyperplane shows that the induced maps on the
one-dimensional quotients agree. This proves the inclusion.
Every such relative map preserves \(Y_t(v)\), with the same
two diagonal blocks for the observed and rational relative maps.
Multiplying the characteristic polynomials of the blocks gives
the last assertion.
\end{proof}

\begin{theorem}[Separation and integral recovery on fixed hyperplanes]
\leavevmode\label{thm:hyperplanerecovery537}
Suppose \(j\widehat L_s=\widehat L_t\) for all cells under
consideration. Fixed hyperplane images at even one marked place
imply that distinct rational transports occupy disjoint entries
of the finite projection matrix \(H\) in every row family.

If the images are fixed hyperplanes at all marked places, set
\[
     \rho=\prod_{v\mid Q_0}\frac{q_v^k-q_v^{k-1}}{q_v^k-1},
     \qquad
     \eta=\prod_{v\mid Q_0}\frac{q_v^k}{q_v^k-1},
\]
where \(q_v\) is the residue cardinality. A proportion at least
\(\rho\) of the row families recovers every signed coefficient
\(Z_i\) by entry selection and multiplication by character units.
Here \(Z_i\) is the sum at the actual residue index \(i\). For
any Hilbert-space-valued coefficients on this support the
normalized pairing satisfies
\[
           \rho\sum_i\|Z_i\|^2
             \ \le\ Q(Z)\ \le\ \eta\sum_i\|Z_i\|^2 .
\]
The integral left inverse preserves all determinantal ideals
of a matrix of such columns, including at \(T\). These conclusions
hold for both terminal frequencies with arbitrary common-lattice
discriminant. Uniformly on the prescribed intervals the two
norm factors tend to one, with
\[
 0\le 1-\rho\le f(g_1^{-1}+g_2^{-1}),\qquad
 0\le \log\eta\le 2f(g_1^{-k}+g_2^{-k}).
\]
\end{theorem}
\begin{proof}
Consider \(u=j_1j_2^{-1}\) on the target. It preserves its positive
parent lattice, whose group of integral isometries is finite.
Indeed all matrices at the archimedean places lie in a compact
set and all their entries in a fixed rational basis lie in a
fractional ideal, a discrete lattice over all embeddings.
Write \(p\) for the residue characteristic of the chosen marked
factor. Here \(p\) splits completely in \(E\) and \(p-1>k\)
by Proposition~\ref{prop:ray50} and the choices
\(g_i\ge(kb_x)^8\). The order \(m\) of \(u\) is prime to \(p\).
Otherwise some power would have order \(p\) over \(\QQ_p\).
Its minimal polynomial would contain \(\Phi_p\), of degree
\(p-1\), since \(\Phi_p(1+X)\) is Eisenstein at \(p\).
A unipotent reduction of \(u\) therefore has order one, its
order being both a power of \(p\) and a divisor of \(m\).
Then \(1+u+\cdots+u^{m-1}\) reduces to the unit \(m\) times
the identity. Multiplying it by \(u-1\) proves \(u=1\).

An overlap of two indices in a single entry of \(H\) makes
their relative map conjugate into the upper unitriangular group
at every marked place, by Lemma~\ref{lem:whittakermatrix526}.
At the place with fixed hyperplanes Lemma~\ref{lem:parabolic537}
makes the reduction of \(u\) unipotent. The preceding argument
gives equality of the rational transports, proving separation.

Now fix one transport \(j\) with the hyperplane condition at
every marked place. Its possible indices belong to
\(\prod_v D_v\bar j_v\). In the row family \(aU\) the local
conjugate \(aUa^{-1}\) fixes its first flag line pointwise.
If this line lies outside \(Y_t(v)\), then its pointwise
stabilizer meets \(D_v\) only in the identity. Indeed an
element of \(D_v\) fixes \(Y_t(v)\) pointwise and the extra
line spans a complement. Thus distinct indices of this
transport occupy different entries if the first line is
outside at every place. Counting lines gives exactly the
proportion \(\rho\) of choices with this property. The same
choices work for all transports. Each selected coordinate
is one original coefficient times a root of unity. This
gives a left inverse over the ring of character values.
Cauchy--Binet applied to the matrix and to its left inverse
gives equality of the ideals of minors in both directions.

For the upper and lower norms take the full local coset
\(D_v\bar j_v\), putting zero on any missing entries.
It has \(q_v^{k-1}\) elements. Every nonidentity relative
element of \(D_v\) is a transvection with fixed space of
dimension \(k-1\). By Theorem~\ref{thm:char125} the normalized
character is \(-1/(q_v^k-1)\) on these elements and one at
the identity. The local Gram matrix therefore has eigenvalue
\((q_v^k-q_v^{k-1})/(q_v^k-1)\) on constants and
\(q_v^k/(q_v^k-1)\) on their orthogonal complement. The
Gram matrix for a full product coset is the tensor product of
these matrices, by the product character formula. The same
extreme eigenvalue bounds hold after zero extension for a
smaller support, and for vector coefficients by expansion in
an orthonormal basis of their span. Distinct transports have
disjoint projected supports, so their squared norms add.
There are \(f\) places of each residue size \(g_i\). The lower
factor has \(1-\rho_v=(q_v^{k-1}-1)/(q_v^k-1)<q_v^{-1}\),
and the upper has \(-\log(1-q_v^{-k})\le 2q_v^{-k}\).
The product inequality \(1-\prod_v\rho_v\le\sum_v(1-\rho_v)\)
and addition of logarithms give the final bounds. The degree
\(f\) grows at most polynomially by Lemma~\ref{lem:size97},
whereas \(b_x=\exp(x^\varepsilon)\).
\end{proof}

\begin{corollary}[A trace test allowing nontrivial reductions in the radical]
\leavevmode\label{cor:parabolictrace537}
For two distinct rational transports use the actual denominator
ideals \(\mathfrak e_l\) of Proposition~\ref{prop:retainedtest535},
as computed by Proposition~\ref{prop:gluing537}. Let \(S\)
be any set of marked places where the two maps are integral
isomorphisms and both cells have the same source and target
hyperplane images. If their actual indices share an entry of
\(H\), then
\[
       \prod_{v\in S}Nv\ \le\
                 (4k)^f N_{F/\QQ}(\mathfrak e_1\mathfrak e_2).
\]
An additional factor \(Nv\) is valid on the left at each
\(v\in S\) where the two reduced rational maps coincide.
\end{corollary}
\begin{proof}
The relative isometry \(u=j_1j_2^{-1}\) is integral at the
places in \(S\). There its reduction is unipotent by
Lemma~\ref{lem:parabolic537} and the row formula for \(H\).
Consequently \(t=2k-\operatorname{tr}u-\operatorname{tr}u^{-1}\)
vanishes in the residue field. If the reduced maps coincide,
the identity \(t=\operatorname{tr}((1-u)(1-u^{-1}))\) gives
two factors of the prime. At every real embedding \(t\) is
the squared Hilbert--Schmidt norm of \(1-u\), positive and at
most \(4k\). The product \(\mathfrak e_1\mathfrak e_2\)
clears both matrices as in Proposition~\ref{prop:retainedtest535}.
Taking the norm of the resulting integral ideal gives the bound.
\end{proof}

\begin{proposition}[Realization on the retained diagonal rows]
\leavevmode\label{prop:realizedhyperplanes537}
On every diagonal row of Theorem~\ref{thm:everydiagonal536},
prescribe a hyperplane in its positive reduction at each marked
place. There is a cyclic family as in that theorem for which
the images of \(R\) are exactly the prescribed hyperplanes
at both ends and all forward and inverse denominator ideals
are trivial.
\end{proposition}
\begin{proof}
At a marked prime the outer lattice is primitive with integral
ordered basis \(e_a\). Its prescribed fractional-ideal weights
are units there since the prime splits completely. The terminal
norm is a unit and its vector is perpendicular to the outer
basis. Split off this line from the self-dual local lattice.
Primitivity and perfect pairing give integral vectors
\(f_j^0\) perpendicular to the terminal line with
\(h(e_i,f_j^0)=\delta_{ij}\). Choose rational \(f_j^0\)
satisfying the same equalities and integral at all the marked
primes by weak approximation on this affine space of solutions.
If \(C=(h(f_i^0,f_j^0))\), subtract from the matrix of their
columns one half the matrix of the \(e_i\)'s times \(C\).
The resulting columns have zero mutual Gram matrix and keep
the prescribed dual pairings. The marked primes are odd, so
the columns remain integral there. The resulting outer planes
and the terminal line are unimodular at all these places.
Their orthogonal complement consequently realizes the positive
parent as the intersection with the row at the marked primes.

Choose a rational linear form on this positive space which
reduces to a defining nonzero form for each prescribed
hyperplane. This again follows from weak approximation in a
rational basis, since each residue condition specifies a
nonempty open subset. Its kernel \(Y\) is a rational hyperplane.
At the prescribed factors the form is primitive, so its integral
kernel reduces onto the kernel of the reduced form.
Use this Witt decomposition and \(Y\) in
Theorem~\ref{thm:everydiagonal536}. Its common lattice is
\(\Lambda\cap Y\), and its transports and residues are the
identity. This gives all the stated local conditions, also
on the retained rows with \(h_R>b_x^5\).
\end{proof}

The estimates here start with the sums \(Z_i\) at the actual
indices. Distinct rational partners can contribute with different
signs to the same \(Z_i\), as in Theorem~\ref{thm:diagonaltail536}.
The comparison in Section~\ref{sec:comparison} still requires
estimates for those signed sums and for cells outside the
hyperplane conditions, on the same interval as the residual
cohomology bound.


\subsection{The full ordinary diagonal at the two marked primes}
\label{sec:markeddiagonal540}

Use the same row, word and marking at both ends of the ordinary
period. Put \(Q=g_1g_2\), \(E=E_x\), \(F=E^+\),
\(f=[F:\mathbb Q]\), and let \(d_V\) be the dimension of the
ambient Hermitian space over \(E\). The lifted terms lie in
\(N_\Lambda\backslash\Gamma/N_\Lambda\), where
\(\Gamma=\operatorname{Aut}(\Lambda,h)\). Write the frame as
\((e_1,\ldots,e_s,z)\), put \(X_i=\sum_{j\le i}Ee_j\),
\(X_0=0\), and \(X=X_s\). We take integral frame coordinates as
in Lemma~\ref{lem:coordinates225}. Equivalently one may clear
the denominators of the frame by a common integer prime to \(Q\),
since it is integral at every place above \(Q\). In this
description the resulting factor belongs to the row constants below.

The rational stabilizer \(B\) of the decorated flag consists of
the isometries \(m\) satisfying
\[
       (m-1)e_i\in X_{i-1}\quad(1\le i\le s),\qquad
                         (m-1)z\in X .
\]
Its unipotent radical acts trivially on \(X^\perp/X\), and
\(N_\Lambda\) is the intersection of this radical with \(\Gamma\).
These are the defining equations in Lemma~\ref{lem:deep163}.
At the original marking the admissibility of a transported
frame means that it agrees with an upper unitriangular change
of the initial frame modulo \(Q\Lambda_{\mathfrak w}\) at
every \(\mathfrak w\mid Q\) of \(E\). In particular this
condition is imposed at both halves of each split place.
The definition is given in
\srcloc{Formalism/SharpRestriction/CuspSpanNormBounds/IsPrimitiveSub}{139}.
The finite index \(i(\gamma)\) records the action on the
remaining \(k\)-dimensional positive complement at the chosen
half of each marked place, as in Lemma~\ref{lem:middlelattice531}.

\begin{lemma}[The product of the marked primes in the real height]
\leavevmode\label{lem:frameheight540}
Write \(V(y)=\prod_\tau\inf_{b\in B_\tau}w(by_\tau)\) for the
height of Proposition~\ref{prop:weight163}. There is a constant
\(A_\Lambda\ge1\), determined by the integral frame, the norms
and a compact set of representatives for \(B_\infty/N_\Lambda\),
with the following property. For every admissible \(\gamma\in\Gamma\)
outside \(B(F)\) and every \(y\in G_\infty\),
\[
                   Q^f\le A_\Lambda V(y)V(\gamma^{-1}y).
\]
\end{lemma}
\begin{proof}
For each pair \((v,Y)=(e_i,X_{i-1})\) or \((z,X)\), choose
finitely many \(E\)-linear functionals \(\lambda\) annihilating
\(Y\), integral on \(\Lambda\), whose common kernel is \(Y\).
One obtains them by clearing the denominators of any basis of
the annihilator. Thus the numbers \(\lambda((m-1)v)\) are
algebraic integers for \(m\in\Gamma\), and their simultaneous
vanishing is exactly membership in \(B(F)\). At every
\(\mathfrak w\mid Q\), admissibility implies
\(\lambda((\gamma-1)v)\in Q\mathcal O_{E,\mathfrak w}\).
The same holds with \(\gamma\) replaced by
\(m=n_1\gamma n_2^{-1}\), \(n_1,n_2\in N_\Lambda\).
Indeed these changes preserve the lattice and are upper
unitriangular on the initial frame at the marked places.
Upper unitriangular changes remain so after multiplication
or inversion. Since \(N_\Lambda\subset B(F)\), such an
\(m\) is still outside \(B(F)\). Some tested value \(z_m\)
is therefore nonzero and belongs to \(Q\mathcal O_E\).
With one embedding \(\widetilde\tau:E\to\mathbb C\) above each
real \(\tau\) of \(F\), this implies
\[
        \prod_\tau|\widetilde\tau(z_m)|
          =|\operatorname{Norm}_{E/\mathbb Q}(z_m)|^{1/2}
          \ge Q^f .
\]

In signed orthonormal coordinates let \(w\) be the operator norm,
which equals the norm of the inverse for a unitary matrix of the
given signature. Write \(v_\tau(y)=\inf_{b\in B_\tau}w(by_\tau)\).
The integral unipotent subgroup is a uniform lattice in the real
unipotent radical, and its quotient Levi group in \(B_\infty\)
is compact. Thus choose a compact set \(C\) with
\(B_\infty=C N_\Lambda\), as in the proof of
Theorem~\ref{thm:diagonaltail536}. For any \(y\) choose
\(b=(b_\tau)\) with \(w(b_\tau y_\tau)\le2v_\tau(y)\).
Writing \(b=c n_y\) gives
\(w((n_yy)_\tau)\le K_\tau v_\tau(y)\), with
\(K_\tau=2\max(1,\sup_{c\in C}w(c_\tau^{-1}))\).
Apply this to \(y\) and to \(y'=\gamma^{-1}y\). Take
\[
      m=n_y\gamma n_{y'}^{-1}=(n_yy)(n_{y'}y')^{-1}.
\]
For this same \(m\) at all places we have
\(w(m_\tau)\le K_\tau^2v_\tau(y)v_\tau(y')\).
Because these heights are at least one, every tested value obeys
\[
  |\widetilde\tau(\lambda((m-1)v))|
    \le 2K_\tau^2\max(1,\|\lambda_\tau\|\|v_\tau\|)
                    v_\tau(y)v_\tau(y').
\]
Taking the maximum of the coefficient over the finite list at
each \(\tau\) and then their product defines \(A_\Lambda\).
Use the nonzero tested value in the preceding norm inequality.
\end{proof}

For a prescribed ordinary archimedean section \(W\), let
\[
 \begin{split}
  L_0&=8(f+1)(d_V+1)^3,\qquad a=(L_0-4fd_V^2)/2,\\
  F_a(y)&=V(y)^a\|W(y)\|,\qquad
  S_a(y)=\sum_{\delta\in N_\Lambda\backslash\Gamma}F_a(\delta y),\\
  M_{\Lambda,W}&=\|F_a\|_{L^1(N_\Lambda\backslash G_\infty)}
                          \sup_y S_a(y).
 \end{split}
\]
Here \(a>0\). Let \(c_\Lambda\) be the common volume factor of
Proposition~\ref{prop:partnerweights531} for the fixed row and
word. Put
\[
         \mathcal E_{\Lambda,W}
          =|c_\Lambda|M_{\Lambda,W}(A_\Lambda/Q^f)^a .
\]
All quantities in this expression use the actual marked row.

\begin{theorem}[All ordinary orbital terms off the decorated flag]
\leavevmode\label{thm:totaldiagonal540}
At the actual index \(i\) let \(Z_i\) be the sum of the signed
orbital operators of the diagonal row, and \(Z_i^B\) the sum
over those represented in \(B(F)\). Every admissible partner
and every lift is included. Both sums converge absolutely, and
\[
             \sum_i\|Z_i-Z_i^B\|_{\rm HS}\le\mathcal E_{\Lambda,W}.
\]
The same upper bound applies to the sum of the individual operator
norms over all lifts outside \(B(F)\), even before grouping by \(i\).
\end{theorem}
\begin{proof}
The estimates of Proposition~\ref{prop:weight163} and the
coset count in Theorem~\ref{thm:diagonaltail536} give
\(M_{\Lambda,W}<\infty\). More explicitly, at height at most
\(R\), representatives of the cosets can be taken with all
real norms after translation by \(y\) at most \(O_\Lambda(R)\),
by compact reduction as in the lemma. The ratios to one of
these representatives are integral matrices of archimedean
size \(O_\Lambda(R^2)\). Counting their \(d_V^2\) entries in a
fixed fractional ideal over \(E\) gives
\(O_\Lambda(R^{4fd_V^2})\) cosets, uniformly in \(y\).
Since \(W\) decays with exponent \(L_0\) and
\(L_0-a>4fd_V^2\), summation over shells bounds \(S_a\).
The weighted integral of \(W\) on \(B_\infty\backslash G_\infty\)
is finite by the same proposition. The compact fiber over that
quotient has finite volume, giving the asserted \(L^1\) norm.

Integrate \(F_a(y)S_a(y)\) over \(N_\Lambda\backslash G_\infty\).
Grouping the nonnegative summands by double cosets gives
\[
 \sum_{[\gamma]\in N_\Lambda\backslash\Gamma/N_\Lambda}
 \int_{(N_\Lambda\cap\gamma N_\Lambda\gamma^{-1})\backslash G_\infty}
             F_a(y)F_a(\gamma^{-1}y)\,dy
           \le M_{\Lambda,W}.
\]
Indeed the cosets in \(N_\Lambda\gamma^{-1}N_\Lambda\) translate
a fundamental domain of \(N_\Lambda\) through one for the
intersection in the display. Tonelli's theorem permits this
grouping. The absolute value of the finite outer phase is one
and the Hilbert--Schmidt norm of \(W(y)W(\gamma^{-1}y)^*\)
is the product of the two vector norms. These are the factors
of the lifted operator in Proposition~\ref{prop:partnerweights531}.
For each admissible \(\gamma\notin B(F)\), the lemma bounds this
product by \((A_\Lambda/Q^f)^a F_a(y)F_a(\gamma^{-1}y)\).
Sum, multiply by the common absolute volume, and then apply the
triangle inequality at each actual index. The nonnegative integral
bound also gives absolute convergence without the height gain,
since \(V\ge1\).
\end{proof}

\begin{corollary}[The projected diagonal and its error at the original level]
\leavevmode\label{cor:projecteddiagonal540}
Put \(d=[G:U]\) for the finite group and unipotent subgroup of
Lemma~\ref{lem:whittakermatrix526}, and use its matrix \(H\).
With the Hilbert direct-sum norm on arrays of operators,
\[
 \left| d^{-1/2}\|HZ\|_2-
             \left(\sum_i\|Z_i^B\|_{\rm HS}^2\right)^{1/2}\right|
                    \le\mathcal E_{\Lambda,W}.
\]
In every row family \(a_0U\) of \(H\), the entry containing the
identity index, after multiplication by the inverse of its known
character coefficient, differs from the identity orbital operator
by at most \(\mathcal E_{\Lambda,W}\) in Hilbert--Schmidt norm.
\end{corollary}
\begin{proof}
An element of \(\Gamma\cap B(F)\) acts on the positive
complement to \(E\bar z\) in \(X^\perp/X\) and preserves its
integral lattice. The kernel of this action is \(N_\Lambda\).
The image is finite, since all archimedean factors are compact
and the integral matrices are discrete. Its action on the
underlying rational space has dimension \(2fk\). A nonidentity
element of order \(g_j\) would have a primitive \(g_j\)-th root
as an eigenvalue, forcing \(g_j-1\le2fk\). Thus the order of
the finite image is prime to both \(g_j\) at the prescribed
marking. Reduction to a unipotent at even one marked place
forces the element to be the identity. Indeed its reduction has
both prime-to-\(g_j\) and \(g_j\)-power order. If \(m^h=1\),
\(g_j\nmid h\), and \(m\equiv1\), then the sum
\(1+m+\cdots+m^{h-1}\) is invertible over the local lattice,
and multiplication by \(m-1\) gives zero. This is the
faithful reduction argument of Lemma~\ref{lem:middlelattice531}.

Two terms from \(\Gamma\cap B(F)\) in a common entry would
have relative positive action reducing to a conjugate of an
element of \(U\) at every marked place. Its reduction at each
one is unipotent, so the relative action is trivial. The terms
are the same double coset since \(N_\Lambda\) is normal in this
stabilizer. Consequently in each row family the stabilizer terms
occupy disjoint entries and each has a coefficient of absolute
value one. This proves
\(d^{-1/2}\|HZ^B\|_2=(\sum_i\|Z_i^B\|_{\rm HS}^2)^{1/2}\),
and only the identity orbit contributes to its identity entry.
Every column of \(d^{-1/2}H\) has norm one, since it has a
single nonzero entry of modulus one in each of the \(d\) row
families. By the triangle inequality on Hilbert-valued columns,
the norm of its action on \(Z-Z^B\) is at most the sum in
Theorem~\ref{thm:totaldiagonal540}. The unnormalized error in
any one entry is bounded by the same sum. This proves both
assertions with the original finite projection.
\end{proof}

The estimate is at the two specified marked primes and includes
the different rational partners contributing to the same entry.
Its right hand side contains the real norms of the prescribed
section and the integral frame height. Their growth relative to
\(Q^f\) as the fields and rows vary is needed for an interval
bound. The mixed rows and their arithmetic normalization also
enter the comparison in Section~\ref{sec:missing}.


\subsection{Mixed ordinary periods at the original marked level}
\label{sec:mixedordinary541}

Fix \(x\) and the pair \(g_1,g_2\) on its original interval.
Use any finite collection of the source and target rows of the
same ambient rank, flag length and positive column dimension
in the lifted period. Write \(Q=g_1g_2\), \(E=E_x\) and
\(f=[E^+:\mathbb Q]\). A pair can contribute only if its
Hermitian lattices are integrally isometric. In each such
isometry class choose a representative \(\Lambda\) and transport
the rows to it, including their frames and their real sections.
Concretely, if the target and source identifications from the
model are \(u_t,u_s\), change a rational isometry \(\gamma\) to
\(u_t^{-1}\gamma u_s\) and the real integration variable \(y\)
to \(u_t^{-1}y\). This preserves Haar measure and the original
orbital integrals. Each row now uses the same arithmetic group
\(\Gamma=\operatorname{Aut}(\Lambda,h)\), with its own integral
unipotent subgroup \(N_j\), real stabilizer \(B_j\), height
\(V_j\), and section \(W_j\). Translating a section by a fixed
real isometry preserves its weighted decay with the transported
height, by submultiplicativity of the defining operator norm.

We first make the counting constant in the unfolding explicit.
Use one \(E\)-basis and write \(n=\dim_E V\). Choose an integer
\(h\ge1\) such that each matrix in \(h\operatorname{End}_{\mathcal
O_E}(\Lambda)\) has integral entries in this basis, viewing
endomorphisms after scalar extension to \(E\). At each \(\tau\)
let \(D_\tau\ge1\) bound the absolute value of each matrix
coordinate by \(D_\tau\) times the operator norm in signed
orthonormal coordinates. Such a constant is obtained from the
change of basis and its inverse. For row \(j\), take \(K_{j,\tau}\)
from the compact reduction in Lemma~\ref{lem:frameheight540}.
Put
\[
 B_j^*=\left(\prod_\tau 4hD_\tau K_{j,\tau}^2\right)^{1/f},
 \quad U_j(L)=\sup_y V_j(y)^L\|W_j(y)\|,\quad
 I_j(a)=\int_{N_j\backslash G_\infty}V_j(y)^a\|W_j(y)\|\,dy .
\]
In the following choose \(0<a<L-4n^2\) such that \(U_j(L)\)
and \(I_j(a)\) are finite. The values of \(L_0,a\) used in
Theorem~\ref{thm:totaldiagonal540} have this property for
the prescribed ordinary sections.

\begin{lemma}[An explicit lattice count for the weighted sum]
\leavevmode\label{lem:packing541}
Uniformly for \(y\in G_\infty\),
\[
 \sum_{\delta\in N_j\backslash\Gamma}
       V_j(\delta y)^a\|W_j(\delta y)\|
   \le U_j(L)\,(1+2B_j^*)^{2fn^2}
                  (1-2^{-(L-a-4n^2)})^{-f}=: C_j(a,L).
\]
\end{lemma}
\begin{proof}
For positive \(R_\tau\), the number of algebraic integers
\(v\in E\) with \(|\widetilde\tau v|\le R_\tau\) at all the
chosen embeddings is at most
\((1+2(\prod_\tau R_\tau)^{1/f})^{2f}\).
Indeed the nonzero difference of two such integers has product
of absolute values at least one. In the product of the scaled
complex planes with norm \(\max_\tau |v_\tau|/R_\tau\), distinct
points therefore have distance at least
\(b=(\prod_\tau R_\tau)^{-1/f}\).
The disjoint open balls of radius \(b/2\) about them fit inside
the ball of radius \(1+b/2\). Taking real \(2f\)-dimensional
volumes gives the bound.

Partition the cosets according to
\(2^{m_\tau}\le v_{j,\tau}(\delta y)<2^{m_\tau+1}\) at each
real place, with nonnegative integers \(m_\tau\).
In a nonempty box choose representatives by compact reduction
on the left, so their translates by \(y\) have norm at most
\(K_{j,\tau}2^{m_\tau+1}\). Divide each on the right by one
fixed representative in that box. The resulting distinct
integral matrices have local norms bounded by
\(4K_{j,\tau}^2 2^{2m_\tau}\), since the inverse has the same
norm. Their entries multiplied by \(h\) are integral. Applying
the preceding count independently to the \(n^2\) entries
bounds the number of cosets by
\[
           (1+2B_j^*)^{2fn^2}\,2^{4n^2\sum_\tau m_\tau}.
\]
Here one uses \(1+tu\le(1+t)u\) for \(u\ge1\).
The weighted size of one summand is at most
\(U_j(L)2^{-(L-a)\sum_\tau m_\tau}\).
Summing the resulting geometric series independently in each
\(m_\tau\) proves the lemma.
\end{proof}

Write the ordered frame of row \(j\) as
\((v_{j,1},\ldots,v_{j,s},v_{j,s+1})\), with \(v_{j,s+1}=z_j\).
Put \(X_{j,i}=\sum_{h\le i}Ev_{j,h}\) for \(0\le i\le s\).
Choose, for each \(X_{j,i-1}\), a basis of its annihilator
integral on \(\Lambda\) as in Lemma~\ref{lem:frameheight540}.
Here a basis means an \(E\)-basis taken in the integral dual.
For each row \(j\) let
\[
 \begin{split}
 V_{j,\tau}^*&=\max(1,\max_i\|v_{j,i,\tau}\|),\\
 D_{j,\tau}^*&=\max(1,\max_\lambda\|\lambda_\tau\|),\\
 A_j^*&=\prod_\tau 2K_{j,\tau}D_{j,\tau}^*V_{j,\tau}^* .
 \end{split}
\]
Frame vectors may be taken integral as before. If denominators
are cleared, use a single common multiplier prime to \(Q\)
on all frames in the isometry class, so that their relative
markings are preserved.

\begin{theorem}[The ordinary orbital operators between two rows]
\leavevmode\label{thm:mixedordinary541}
For target \(t\) and source \(s'\), let \(Z_i^{t,s'}\) be the
full signed operator at the marked index \(i\). Let
\(B_i^{t,s'}\) consist of its terms whose rational isometry
\(m\) satisfies
\[
        m v_{s',j}-v_{t,j}\in X_{t,j-1}
                          \qquad(1\le j\le s+1).
\]
For the last vector use \(X_{t,s}\). These conditions are
independent of the representative modulo the target and source
integral unipotent groups. With the volume \(c_{t,s'}\) of
the pair in Proposition~\ref{prop:partnerweights531}, one has
\[
 \begin{split}
 \sum_i\|Z_i^{t,s'}-B_i^{t,s'}\|_{\rm HS}
     &\le e_{t,s'},\\
 e_{t,s'}&=
 |c_{t,s'}|
  \left(\frac{A_t^* A_{s'}^*}{Q^f}\right)^a
  \min\{I_t(a)C_{s'}(a,L),\, I_{s'}(a)C_t(a,L)\}.
 \end{split}
\]
If the two terminal norms differ, every \(B_i^{t,s'}\) is zero.
For equal terminal norms, the distinct contributing rational
classes in \(B^{t,s'}\) occupy disjoint entries of \(H\) in
each row family. For a pair in different integral isometry
classes set its arrays and its error bound equal to zero.
\end{theorem}
\begin{proof}
An integral unipotent change at either end preserves the
successive graded frame vectors, so preserves the rational
conditions and their marked congruences. For an admissible
\(\gamma\) failing a rational condition, perform compact reduction
on \(y\) at the target and \(\gamma^{-1}y\) at the source. The
resulting matrix \(m=n_t\gamma n_{s'}^{-1}\) obeys
\[
  w(m_\tau)\le K_{t,\tau}K_{s',\tau}
                v_{t,\tau}(y)v_{s',\tau}(\gamma^{-1}y).
\]
Some \(\lambda(mv_{s',j}-v_{t,j})\) with \(\lambda\) from the
corresponding target annihilator is nonzero. It is integral and
is divisible by \(Q\) at every prime above \(Q\) by admissibility.
At each infinite place its absolute value is bounded by
\(\|\lambda_\tau\|(w(m_\tau)\|v_{s',j,\tau}\|+\|v_{t,j,\tau}\|)\).
Taking the product over \(\tau\), the definitions of \(A_t^*\)
and \(A_{s'}^*\) therefore give
\[
           Q^f\le A_t^* A_{s'}^*\, V_t(y)V_{s'}(\gamma^{-1}y).
\]
Integrate the product of the two weighted section norms over
\((N_t\cap\gamma N_{s'}\gamma^{-1})\backslash G_\infty\).
Summing over all double cosets unfolds to the integral on
\(N_t\backslash G_\infty\) of the target norm times the sum
of the source norm over \(N_{s'}\backslash\Gamma\).
Lemma~\ref{lem:packing541} bounds this by \(I_t(a)C_{s'}(a,L)\).
Inverting the cosets and changing variables
\(y\mapsto\gamma^{-1}y\) gives the other bound in the minimum.
The height inequality, the unitary finite phases and the
rank-one operator norm give the asserted estimate exactly
as in Theorem~\ref{thm:totaldiagonal540}.

A map satisfying the rational conditions sends the source
isotropic flag to the target flag. In the induced map on middle
quotients it sends the terminal vector class to the target
terminal vector class. Their Hermitian norms must agree.
Given two such integral maps, their quotient is in the rational
decorated stabilizer of the target and preserves its lattice.
Its action on the positive middle lattice has finite order prime
to both marked primes, as in Corollary~\ref{cor:projecteddiagonal540}.
In a common entry of \(H\) the reduction of this quotient on the
positive complement is unipotent. For these flag-preserving maps
their finite indices are precisely the reductions of their
middle actions, so the reduction of their quotient is the
quotient of the indices. The finite-order argument forces its
middle action to be trivial. It belongs to \(N_t\), hence the
two integral maps give one double coset. This proves the claim
about disjoint entries and shows that only finitely many rational
classes contribute to \(B^{t,s'}\).
\end{proof}

\begin{corollary}[A simultaneous error for the prescribed finite projection]
\leavevmode\label{cor:mixedmatrix541}
For these rows let \(\mathcal Z,\mathcal B\) be the block
operators with entries \(d^{-1/2}HZ^{t,s'}\) and
\(d^{-1/2}HB^{t,s'}\), respectively, with \(d=[G:U]\).
They act from the direct sum of the source coefficient spaces
to the direct sum over target rows and finite output entries.
Then
\[
 \|\mathcal Z-\mathcal B\|_{\rm op}
 \le
 \left(\max_t\sum_{s'}e_{t,s'}\ 
                  \max_{s'}\sum_t e_{t,s'}\right)^{1/2}.
\]
Every singular value changes by at most the right hand side.
On blocks between the two distinct terminal norms this is an
absolute bound for the ordinary projected operator itself.
\end{corollary}
\begin{proof}
Each normalized column of \(H\) has norm one. Applying the
triangle inequality to the Hilbert-valued columns bounds the
operator norm of one difference block by \(e_{t,s'}\).
For an input with component norms \(b_{s'}\), its output
difference on row \(t\) has norm at most
\(\sum_{s'}e_{t,s'}b_{s'}\). Cauchy--Schwarz with these
nonnegative weights, followed by summation in \(t\), gives the
display. The variational characterization of singular values
gives the second assertion. The same argument applies after
retaining any subset of rows and columns, in particular the
blocks between unequal terminal norms.
\end{proof}

This comparison uses the actual two terminal frequencies and the
same unrestricted signed orbital sums as the prescribed ordinary
period. To apply it throughout an interval one must bound the
displayed row and column sums, and compare the resulting ordinary
operators to the arithmetically normalized integral theta maps.


\subsection{Singular values of the flag-preserving terms in the projected period}
\label{sec:flagsingular544}

Fix one of the fields and pairs of marked primes of
Section~\ref{sec:mixedordinary541}. Use the same retained ordinary
rows as sources and targets. A row consists of its integral
Hermitian lattice, ordered frame, terminal vector, marking, and
nonzero real Bessel section \(W_j\). Its column space \(V_j\)
carries the prescribed irreducible unitary representation \(\rho_j\)
of the positive compact group at infinity. Write \(d_j=\dim V_j\).
Choose a collection \(\mathcal J\) in which, for distinct rows
\(s,t\), there is no admissible integral isometry carrying every
successive frame vector to its target modulo the preceding span,
including the terminal vector modulo the full isotropic span.
These are precisely the rational equations in
Theorem~\ref{thm:mixedordinary541}. For instance, rows in different
integral lattice classes may be included together. One can also
use a retained row for each of the two distinct supported terminal
norms at their common level. An isometry satisfying the equations
would preserve the norm in the middle Hermitian quotient.

For row \(j\), write \(\Gamma_j\) for its integral isometries,
\(B_j\) for the stabilizer of its ordered frame modulo earlier
vectors, and \(N_j\) for the integral unipotent subgroup in that
stabilizer. Let \(\mathcal D_j\) be the admissible double cosets
in \(N_j\backslash(\Gamma_j\cap B_j(F))/N_j\), where \(F=E^+\).
Set \(m_j=|\mathcal D_j|\). This number is finite by
Corollary~\ref{cor:projecteddiagonal540} and is positive because
the identity contributes. Define
\[
        \lambda_j=\frac{|c_{j,j}|}{d_j}
                \int_{N_j\backslash G_\infty}\|W_j(y)\|^2\,dy,
        \qquad u_j=\lambda_j\sqrt{m_j}.
\]
Here \(c_{j,j}\) is the volume factor of the lifted diagonal
orbital integral in Proposition~\ref{prop:partnerweights531}.
In particular \(\lambda_j>0\). For the unit sections of
\eqref{eq:real127} it depends only on the measures, the
unipotent lattice, and the column dimension. In fact, writing
\(N_{j,\infty}\) for the real unipotent radical and
\(H_{j,\infty}\) for the real compact factor, one has
\[
     \lambda_j=
       \frac{|c_{j,j}|}{d_j}
       \operatorname{vol}(N_j\backslash N_{j,\infty})
       \operatorname{vol}(H_{j,\infty})
\]
with the quotient measures of that equation. This follows by
integrating \(\|W_j\|^2\) along the two fibers, on which its value
is constant, and then using its prescribed quotient norm.
In the notation of
Corollary~\ref{cor:mixedmatrix541} the finite matrix is
\(d^{-1/2}H\), with \(d=[G:U]\). Denote its column at the
index of \(\gamma\) by \(h_{i(\gamma)}\). The target of the
period has a copy of \(V_j\) at each marked entry on row \(j\).

\begin{theorem}[The exact spectrum on the chosen rows]
\leavevmode\label{thm:flagisometry544}
Let \(\mathcal B\) denote the operator formed from the
flag-preserving terms with the original finite projection.
Then
\[
            \mathcal B^*\mathcal B=
            \bigoplus_{j\in\mathcal J} u_j^2\,1_{V_j}.
\]
More generally, keep any set \(\mathcal A_j\) of marked entries
on target row \(j\), and let \(P_{\mathcal A}\) be the
corresponding coordinate projection. Put
\[
       k_j=\sum_{\gamma\in\mathcal D_j}
         \big|\mathcal A_j\cap\operatorname{supp}h_{i(\gamma)}\big|.
\]
On \(V_j\) the Gram operator of \(P_{\mathcal A}\mathcal B\)
is exactly \(k_j\lambda_j^2/d\) times the identity, with zero
off-diagonal blocks. In particular, retaining just the support
of the identity column on each target row gives \(k_j=d\).
\end{theorem}
\begin{proof}
The chosen rational equations eliminate every off-diagonal
block of \(\mathcal B\). Consider a diagonal term represented
by \(\gamma\in\Gamma_j\cap B_j(F)\). It normalizes \(N_j\)
because it preserves both the lattice and the unipotent radical.
Thus the real integration domain in that term is
\(N_j\backslash G_\infty\). Write \(\gamma_\infty=nh\) in
the real unipotent radical times the positive compact group.
The real covariance of the Bessel section gives
\[
            W_j(\gamma^{-1}y)
                  =\zeta\,\rho_j(h)^{-1}W_j(y)
                 \qquad(|\zeta|=1).
\]
The phase is fixed by \(\gamma\). This follows directly from
the character on the real unipotent subgroup and its invariance
under the compact factor in \eqref{eq:real127}. The finite
outer character also has absolute value one. Consequently
the operator of this double coset is a unit scalar times
\[
          |c_{j,j}|
          \left(\int_{N_j\backslash G_\infty}
                    W_j(y)W_j(y)^*\,dy\right)\rho_j(h).
\]
For completeness, disintegrate the integral over
\(B_{j,\infty}\backslash G_\infty\). In each fiber integrate
first over the compact quotient of the real unipotent subgroup
by \(N_j\), then over the positive compact group. In the
rank-one operator the unipotent character cancels against
its conjugate. The second integral averages the operator
over conjugation by \(\rho_j\), using Haar measure on that
compact group. This is the fiber calculation of
Theorem~\ref{thm:identity127}. By irreducibility the resulting
operator is scalar. Its trace before multiplication by
\(|c_{j,j}|\) is \(\int\|W_j(y)\|^2\,dy\) on the stated
domain. It follows that the actual coefficient is
\(\lambda_j U_\gamma\) for a unitary operator \(U_\gamma\)
on \(V_j\).

The finite projected block sends \(v\in V_j\) to
\[
       \lambda_j\sum_{\gamma\in\mathcal D_j}
                  h_{i(\gamma)}\otimes U_\gamma v .
\]
Indeed each rational double coset has its original signed
coefficient at its original residue index. The columns in
this sum have pairwise disjoint support. To recall the
reason, sharing a marked entry makes the relative action
of the two isometries on the integral positive complement
unipotent after reduction. This action has finite order
prime to either marked prime, since a nontrivial element of
order that prime would have an eigenvalue of larger degree
than the dimension of the rational positive lattice. A
prime-to-residue-characteristic torsion element reducing to
a unipotent is the identity on that lattice, by faithful
reduction as in Corollary~\ref{cor:projecteddiagonal540}.
The kernel of the action is \(N_j\). Hence the two double
cosets would agree. Each column has one nonzero entry of
modulus \(d^{-1/2}\) in each of the \(d\) row families.
The squared norm after keeping the coordinates in
\(\mathcal A_j\) is therefore
\(\lambda_j^2 k_j\|v\|^2/d\). The different target rows
are orthogonal. Polarization proves the asserted Gram
identities, including the identity-column assertion.
\end{proof}

\begin{theorem}[Mean square error for the actual signed ordinary period]
\leavevmode\label{thm:ordinarysingular544}
Let \(\mathcal Z\) be the full ordinary operator on the same
rows, with all rational partners and the same finite projection.
Write \(D=\bigoplus_j u_j1_{V_j}\), \(\mathcal B_0=\mathcal B D^{-1}\),
and \(T=\mathcal ZD^{-1}\). For the explicit errors \(e_{t,s}\)
of Theorem~\ref{thm:mixedordinary541}, put
\[
   F_{t,s}=\frac{e_{t,s}}{u_s},\qquad
   E_2^2=\sum_{t,s\in\mathcal J} F_{t,s}^2,\qquad
   E_\infty=
     \left(\max_t\sum_s F_{t,s}\ \max_s\sum_t F_{t,s}\right)^{1/2}.
\]
For every orthogonal projection \(P\) on the output, let
\(b_P=\|(1-P)\mathcal B_0\|_{\rm HS}^2\). The singular values
\(\sigma_a(PT)\), counted with multiplicities up to the full
input dimension \(n_0=\sum_jd_j\) and including zeros, satisfy
\begin{align}
       \sum_{a=1}^{n_0}(\sigma_a(PT)-1)^2
         &\le E_2^2+b_P, \label{eq:svdisp544}\\
       \|(PT)^*PT-1\|_{\rm HS}^2
         &\le (2+E_\infty)^2(E_2^2+b_P). \label{eq:gramdisp544}
\end{align}
Always \(b_P\le\operatorname{codim}(\operatorname{im}P)\).
For the coordinate projection of the preceding theorem its
exact value is
\[
             b_{P_{\mathcal A}}
                 =\sum_j d_j\left(1-\frac{k_j}{d m_j}\right).
\]
\end{theorem}
\begin{proof}
For each pair \(t,s\) the sum of the Hilbert--Schmidt norms of
the unprojected difference coefficients is at most \(e_{t,s}\)
by Theorem~\ref{thm:mixedordinary541}. Tensoring each coefficient
with its finite column preserves its Hilbert--Schmidt norm
since that column has norm one. The triangle inequality in
the resulting operator space bounds the Hilbert--Schmidt norm
of the entire projected difference block, after division by
\(u_s\), by \(F_{t,s}\). Hilbert--Schmidt squares add over
both indices of a block matrix in orthogonal direct sums,
so \(\|T-\mathcal B_0\|_{\rm HS}^2\le E_2^2\). Applying
Cauchy--Schwarz with the nonnegative coefficients \(F_{t,s}\)
to the component norms, as in Corollary~\ref{cor:mixedmatrix541},
also gives \(\|T-\mathcal B_0\|_{\rm op}\le E_\infty\).

The preceding theorem makes \(\mathcal B_0\) an isometry.
The two terms of
\[
      PT-\mathcal B_0
         =P(T-\mathcal B_0)-(1-P)\mathcal B_0
\]
have orthogonal images. The squared Hilbert--Schmidt norm is
therefore at most \(E_2^2+b_P\). Choose an orthonormal basis
of right singular vectors \(v_a\) of \(PT\), also taking a
basis of its kernel. The reverse triangle inequality gives
\(
 |\sigma_a(PT)-1|\le\|(PT-\mathcal B_0)v_a\|
\)
because \(\|\mathcal B_0v_a\|=1\). Squaring and summing
gives \eqref{eq:svdisp544}. Moreover
\(\|PT\|_{\rm op}\le1+E_\infty\). On this same basis
the eigenvalues of the Gram difference are
\((\sigma_a(PT)-1)(\sigma_a(PT)+1)\), proving
\eqref{eq:gramdisp544}. Finally \(\mathcal B_0\mathcal B_0^*\)
is an orthogonal projection. Taking its trace on the range of
\(1-P\) proves the codimension bound. For \(P=P_{\mathcal A}\),
the preceding theorem computes \(\|P\mathcal B_0\|_{\rm HS}^2\)
as \(\sum_jd_j k_j/(dm_j)\). Subtract this from \(n_0\).
\end{proof}

In particular, for \(0<\delta<1\) the actual projected ordinary
map has at most \((E_2^2+b_P)/\delta^2\) singular values below
\(1-\delta\) in this normalization. The estimate applies at the
two original marked primes and permits both supported terminal
norms in the same calculation. To use \eqref{eq:gramdisp544}
in Proposition~\ref{prop:rank11} one still needs an interval
bound for the errors relative to \(n_0\), sufficient ranks
on the selected rows, and agreement of these row scalings with
the integral theta-period normalization in Section~\ref{sec:missing}.


\subsection{The identity coefficient at the two marked primes}
\label{sec:identitycoefficient551}

Use the original fields, frequencies and distinct decorated rows
of Section~\ref{sec:flagsingular544}. Put
\(G=\prod_v\operatorname{GL}_k(\mathbb F_{q_v})\), where \(q_v=g_1\)
or \(g_2\) at its \(f=[E_x^+:\QQ]\) places. Let
\(U=\prod_v U_k(\mathbb F_{q_v})\) be upper unitriangular and
\(d=[G:U]\). Denote the columns of \(H/\sqrt d\) in
Lemma~\ref{lem:whittakermatrix526} by \(h_i\), \(i\in G\).
For output row \(t\) let \(\iota_t\) be the marked index of its
identity isometry. For each pair \(t,s\) set
\[
 R_i^{ts}=Z_i^{ts}-B_i^{ts},\qquad
 \mathcal Z_{ts}=\sum_{i\in G}h_i\otimes Z_i^{ts}.
\]
Here \(Z_i^{ts}\) is the full signed ordinary coefficient and
\(B_i^{ts}\) is its flag-preserving contribution. By
Theorem~\ref{thm:mixedordinary541},
\(\sum_i\|R_i^{ts}\|_{\rm HS}\le e_{ts}\). Write \(\beta_t 1_{V_t}\)
for the identity orbital operator before the finite matrix. Its
modulus is \(|\beta_t|=\lambda_t>0\) by
Theorem~\ref{thm:flagisometry544}. In this subsection \(\beta_t\)
denotes this scalar, including its phase. Both terminal frequencies
in the earlier notation are allowed. Set \(D_0=\bigoplus_t
\beta_t1_{V_t}\) and \(q_*=\min(g_1,g_2)\).

\begin{theorem}[Contraction onto the marked identity column]
\leavevmode\label{thm:identitycoefficient551}
On each target row contract the finite coordinate against the unit
column \(h_{\iota_t}\), and denote the resulting operator by \(L\).
Then \(\|L\|_{\rm op}\le1\), and
\[
 (L\mathcal Z)_{ts}=\mathbf1_{t=s}\,\beta_t1_{V_t}
                +R_{\iota_t}^{ts}
                +\sum_{i\ne\iota_t}A_{\iota_t i}R_i^{ts}.
\]
The first term means the zero map for \(t\ne s\). Define
\[
 D_G=\prod_v\prod_{a=2}^{k}(q_v^a-1),\qquad
 1+\delta_G=\prod_v\prod_{a=2}^k(1-q_v^{-a})^{-1}.
\]
For the last sum alone the Hilbert--Schmidt norm is at most
\[
 \min\left\{\frac{\sum_{i\ne\iota_t}\|R_i^{ts}\|_{\rm HS}}{q_*^k-1},
       \ \sqrt{\delta_G}
       \left(\sum_{i\ne\iota_t}\|R_i^{ts}\|_{\rm HS}^2\right)^{1/2}
       \right\}=b_{ts}.
\]
In particular it is at most \(2e_{ts}(k b_x)^{-8k}\) on the
marked prime intervals, where \(b_x=\exp(x^\varepsilon)\).
For any lattice of coefficient tables the finite operation
\((Z_i)\mapsto\sum_i A_{\iota_t i}Z_i\) has denominator dividing
\(D_G\). On a free table supported on a specified subset \(S\subset G\)
its least positive integer denominator is the least common multiple
of
\[
     \prod_v\prod_{a=k-r_v+1}^k(q_v^a-1),\qquad
     r_v=\operatorname{rank}(\iota_{t,v}^{-1}i_v-1),
\]
over indices \(i\in S\) for which all the relative matrices are
unipotent. The empty least common multiple is one.
\end{theorem}
\begin{proof}
Contraction against a unit vector is a contraction on each
orthogonal summand. On a signed table it pairs \(h_i\) with
\(h_{\iota_t}\), giving \(A_{\iota_t i}\) by
Lemma~\ref{lem:whittakermatrix526}. The hypothesis on the rows
removes the flag-preserving terms for \(t\ne s\). On row \(t=s\)
the identity orbital integral is scalar by compact averaging
and Schur's lemma. The other admissible stabilizer terms have
supports disjoint from the identity column by
Theorem~\ref{thm:flagisometry544}. This proves the identity,
including its assertion about the full signed error at \(\iota_t\).

Corollary~\ref{cor:content125} gives zero for a nonunipotent
relative matrix at any one place. Otherwise it gives the
reciprocal of the displayed positive integer, with sign
\((-1)^{\sum_v r_v}\). The denominators on a free table are
tested on its individual basis vectors, proving the assertion
about the least common multiple. The absolute value at a distinct
index is at most \(1/(q_*^k-1)\). Alternatively the sum of the
squares of these coefficients over the distinct indices is
\(\delta_G\) by Theorem~\ref{thm:moment126}. Apply the triangle
inequality and then Cauchy--Schwarz to obtain the two bounds.
Finally \(q_*\ge(kb_x)^8\) on the chosen intervals. This proves
the claimed uniform factor. The same finite denominator at all
the rows has
\[
   \log D_G\le f\bigl(k(k+1)/2-1\bigr)\log(g_1g_2),
\]
by bounding each factor by \(q_v^a\).
\end{proof}

It follows, for example, that the squared Hilbert--Schmidt
distance of \(L\mathcal ZD_0^{-1}\) from the identity is at most
\(\sum_{t,s}(\|R_{\iota_t}^{ts}\|_{\rm HS}+b_{ts})^2/|\beta_s|^2\).
The Gram operator of \(\mathcal ZD_0^{-1}\) dominates that of
the contraction, since \(L^*L\le1\). The identity index can
contain other rational partners in addition to the identity
isometry. Their signed sum \(R_{\iota_t}^{ts}\) occurs with
coefficient one. The estimates at distinct indices alone cannot
bound this term.

\begin{theorem}[One integral coordinate choice for all the rows]
\leavevmode\label{thm:integralfamily551}
For \(i\in G\) put \(g_{ti}=i\iota_t^{-1}\), and let
\[
 \pi_t(i)=\prod_v
 \frac{|\operatorname{Cl}_{G_v}(g_{ti,v})\cap U_v|}
      {|\operatorname{Cl}_{G_v}(g_{ti,v})|},
\]
where \(\operatorname{Cl}_{G_v}\) denotes a conjugacy class.
Choose sets \(S_t\subset G\setminus\{\iota_t\}\) and suppose
\(\eta=\sum_t\sum_{i\in S_t}\pi_t(i)<1\).
There is a single row family \(aU\in G/U\) and a coordinate
selection \(L_a\) on the \emph{unnormalized} table \(HZ\) such
that the flag-preserving part gives \(D_0\), every index in
\(S_t\) is absent on row \(t\), and
\[
\begin{split}
 &\|(L_a HZ)D_0^{-1}-1\|_{\rm HS}^2\\
 &\quad\le\frac1{1-\eta}\sum_{t,s}\frac1{|\beta_s|^2}
 \left(\|R_{\iota_t}^{ts}\|_{\rm HS}
  +\sum_{i\notin S_t\cup\{\iota_t\}}\sqrt{\pi_t(i)}
                                         \|R_i^{ts}\|_{\rm HS}\right)^2 .
\end{split}
\]
The coordinate selection uses only roots of unity as coefficients
and has a right inverse over their integer ring. The same
selection exists with the indicated bound on any nonnegative
weighted sum of these errors for finitely many coefficient
tables with the same marked data, including tables at different
coefficient embeddings.
\end{theorem}
\begin{proof}
For a representative \(a\) and row \(t\) there is a unique
left coset \(Uh\) with \(\iota_t\in aUh\). In the \(aU\)
family keep its coordinate and multiply by the inverse of the
coefficient of \(\iota_t\) in that entry of \(H\). This
coefficient is a root of unity. Inserting a scalar in just
that coordinate with the reciprocal correction gives the
integral right inverse. Changing the representatives modifies
the coordinate and its correction by opposite phases. In
particular the identity coefficient is exactly one in every
family. The other stabilizer terms occupy different coordinates
there, as in the preceding theorem.

Index \(i\) occurs in the chosen entry precisely when
\(g_{ti}\in aUa^{-1}\); if present its corrected coefficient
has modulus one. For \(aU\) uniform the probability of this
event is \(\pi_t(i)\), by counting the fibers of the conjugation
map in each local group. If the relative element is unipotent
this is the probability that a uniform complete flag is fixed.
Indeed its triangular diagonal on a stable flag consists of
ones, and \(G_v/U_v\to G_v/B_v\) has constant fiber size.
At a nonidentity unipotent place one may therefore bound the
factor by \((k-1)/[k]_{q_v}\) using
Lemma~\ref{lem:flagcollision530}. A nonunipotent local factor
gives zero.

By the union bound, a proportion at least \(1-\eta\) of the
families avoid every prescribed index. Delete those indices
from the error arrays in all families. For a fixed block the
\(L^2(G/U)\) norm, with probability measure and Hilbert--Schmidt
coefficients, of the contribution of one surviving index is
\(\sqrt{\pi_t(i)}\|R_i^{ts}\|_{\rm HS}\). For the identity index
its norm is exactly \(\|R_{\iota_t}^{ts}\|_{\rm HS}\). The
triangle inequality in this Hilbert space gives the expression
inside parentheses. Square, divide by \(|\beta_s|^2\), and sum
over the blocks. Restricting this nonnegative sum to the good
families and dividing by their proportion gives the bound in
the assertion. Some good family attains at most the resulting
average. Multiplying by the given weights and summing over
all coefficient tables before making this last choice proves
the simultaneous form.
\end{proof}

The two tests use the ordinary orbital coefficients on the
original marked fields. The integer denominator and the unit
coordinate selection concern the finite operation on coefficient
tables. The passage to the integral theta-period lattice still
requires its normalization of the actual \(Z_i^{ts}\), together
with an interval estimate for the signed terms remaining at the
identity index.

\clearpage
\section{References}\label{nav:bibliography-XXI}
\begingroup\def\section*#1{}\begin{thebibliography}{99}\small
\bibitem{XXI-XVI-IV-ghy209} W. T. Gan, J. P. Hanke and J.-K. Yu, \emph{On an exact mass formula of Shimura}, Duke Math. J. \textbf{107} (2001), 101--133. \href{https://web.math.princeton.edu/~jonhanke/Web-02/Mass-Formula-of-Shimura/shimura-mass.pdf}{Section 4, formulas (4.2)--(4.3) and Propositions 4.4--4.5}.
\end{thebibliography}\endgroup
\end{document}
